% SPDX-License-Identifier: CC-BY-4.0
% Master file of the revised paper: Sections 1-7, the literature review,
% Section 8 (summary), the glossary, the acknowledgments and the
% bibliography (stages 0-8 of ../revision_log.md). Until stage 0 this file
% was called preview.tex. Stage 8 (November 12, 2026) brought it into line
% with ../../joyner-math-paper-style-guide.md; the former Appendix A
% (comparison with the 1995 notes) is now in ../revision_log.md.
% Build: latexmk -pdf sl2bar-trace-2026.tex
\documentclass[12pt,letterpaper]{article}
\usepackage[T1]{fontenc}
% Times for text and math. The style guide prefers newtx; it was not
% installed where this draft was built, so the fallback mathptmx is used.
% To switch, replace the next line by \usepackage{newtxtext,newtxmath}.
\usepackage{mathptmx}
\usepackage[letterpaper,margin=1in]{geometry}
\usepackage{amsmath,amssymb,amsthm}
\usepackage[inline]{enumitem}
\usepackage{longtable}
\usepackage{xcolor}
\usepackage[colorlinks=true,linkcolor=blue!50!black,citecolor=blue!50!black,urlcolor=blue!50!black]{hyperref}
\emergencystretch=3em

\numberwithin{equation}{section}
\theoremstyle{plain}
\newtheorem{lemma}{Lemma}[section]
\newtheorem{proposition}[lemma]{Proposition}
\newtheorem{theorem}[lemma]{Theorem}
\newtheorem{corollary}[lemma]{Corollary}
\theoremstyle{definition}
\newtheorem{example}{Example}[section]
\renewcommand{\theexample}{\thesection.\Alph{example}}
\theoremstyle{remark}
\newtheorem{remarkrev}{Remark}
% A lemma added after Proposition 7.5 without renumbering what follows.
\newcommand{\shellref}{\ref{prop:weil-char}$'$}
\theoremstyle{plain}
\newtheorem*{lemmaprime}{Lemma \shellref}
% A proposition added after Lemma 7.8 without renumbering what follows.
\newcommand{\normref}{\ref{lem:variance-gl2}$'$}
\newtheorem*{propprime}{Proposition \normref}
\renewcommand{\theremarkrev}{2026.\arabic{remarkrev}}

\newcommand{\Q}{\mathbb{Q}}
\newcommand{\R}{\mathbb{R}}
\newcommand{\C}{\mathbb{C}}
\newcommand{\Z}{\mathbb{Z}}
\newcommand{\A}{\mathbb{A}}
\newcommand{\OO}{\mathcal{O}}
\DeclareMathOperator{\Res}{Res}
\DeclareMathOperator{\ord}{ord}
\DeclareMathOperator{\diag}{diag}
\DeclareMathOperator{\Hom}{Hom}
\DeclareMathOperator{\Ad}{Ad}
\DeclareMathOperator{\vol}{vol}
\DeclareMathOperator{\sgn}{sgn}
\DeclareMathOperator{\Ind}{Ind}
\DeclareMathOperator{\Resid}{Res}

\begin{document}
\title{A hitch-hiker's guide to the metaplectic trace formula for $SL(2)$:
basic theory\thanks{This is a revision of the unpublished notes of the
same title by D.~Joyner (January 1995, slightly revised June 2002). The
revision was made in November 2026. It is restricted to the double cover
of $SL(2)$, while the 1995 notes treated the $n$-fold cover.}\\[6pt]
\large Draft 1. Corrected and expanded version}
\author{D. Joyner\thanks{\texttt{wdjoyner@gmail.com}} \and
Claude\thanks{Claude is an AI model made by Anthropic. See the
Acknowledgments for the division of work.}}
\date{Original draft: 1995. Corrected and expanded draft: 2026}
\maketitle

\begin{abstract}
This revision was written with the help of AI systems: Claude (Anthropic)
refereed the 1995 notes and drafted the revision as a co-author, and
ChatGPT (OpenAI) refereed the revised draft. See the Acknowledgments for
details. We write out the Arthur--Selberg trace formula, in its non-invariant and
in its invariant form, for the metaplectic double cover of
$G=\Res_{k/\Q}SL(2)$ over a number field $k$, and take the first steps
towards its stable form. The general theory is due to
Arthur and, for covering groups, to W.-W.~Li; here every step is
specialized to this group, and the terms are made explicit. The coarse
geometric expansion is proved directly from one inequality between
heights. On the spectral side we determine the classes fixed by the Weyl
group (Weil's characters $\chi_{\psi_b}$), the residual spectrum (the even
Weil representations $r^+_{\psi_{-b}}$), the normalizing factor
$L(2\lambda,\chi^2)/L(2\lambda+1,\chi^2)$ of the intertwining operator,
and the value $M(w,0)=-1$. The invariant distributions are constructed by
Fourier inversion on the abelian cover of the diagonal torus, without
appeal to a trace Paley--Wiener theorem, and the unipotent coefficients
of the fine geometric expansion are computed in terms of the values at
$1$ of quadratic Hecke $L$-functions. In the last section we determine
the action of $GL(2)$ on the cover: conjugation by $GL(2)$ can exchange
the two lifts of an element, and stable conjugacy in the sense of Adams
and Li is conjugation by $GL(2)$ twisted by the sign
$(\det h,\det(1+\gamma))$. The regular elliptic terms of the trace formula
then take the form $S(f)+S(f_-)$, with $S$ a stable distribution and
$f_-$ the translate of $f$ by the central element $-1$; the theta
functions contribute to the spectral side in the same form. The
stabilization of the remaining terms, which is due to Li, is not
reproduced. One ingredient is not computed here but taken from the work of
Szpruch: the explicit form of the local normalizing factors at the dyadic
and ramified places. This is a
revision, restricted to the double cover, of notes of the first author
from 1995; the explicit formulas are checked by accompanying scripts.
\end{abstract}

\tableofcontents

% SPDX-License-Identifier: CC-BY-4.0
% =====================================================================
% Section 1 of D. Joyner and Claude, "A hitch-hiker's guide to the
% metaplectic trace formula for SL(2)", revised for the 2-fold cover
% (n = 2), November 2026.
%
% New in the November 2026 draft, written in response to the second referee
% report (ChatGPT): overview, literature review, list of what is proved and
% what is cited, and the running example k = Q, S = {oo, 2}. It replaces the
% Introduction of ../original-1995/cctrace_new.tex for Sections 2-6; the
% stable form (the 1995 Section 7) is not yet revised.
% =====================================================================

\section{Introduction}\label{sec:intro}

Let $k$ be a number field, $G=\Res_{k/\Q}SL(2)$, and $\overline{G(\A)}$
the metaplectic double cover of $G(\A)=SL(2,\A_k)$. This paper writes out
the Arthur--Selberg trace formula for $\overline{G(\A)}$, first in its
non-invariant form and then in its invariant form, with every step
specialized to this group. The general theory exists: Arthur's trace
formula for connected reductive groups, and its extension to covering
groups by W.-W.~Li. The aim here is the opposite of generality. The group
has rank one, its only proper Levi subgroup is the diagonal torus $A$,
and the cover of $A$ is abelian. As a result most of the general
machinery reduces to statements that can be proved in a few lines and
checked numerically, and the terms of the formula can be written down.

\subsection{Overview}\label{sec:overview}

The paper has five parts, each resting on the ones before it.

\begin{description}
\item[The cover (Section~\ref{sec:background}).] The double cover is
  defined by Kubota's cocycle, place by place, and assembled into an
  adelic group that splits over $G(\Q)$. The cocycle is not continuous,
  so the topology of the cover and the spaces of test functions are
  defined intrinsically, and the coordinates $(g,\zeta)$ are used only as
  Borel coordinates. Test functions are \emph{anti-genuine}: they
  transform by $\zeta^{-1}$ under $\mu_2$.
\item[Rank $0$ (Section~\ref{sec:rank0}).] For the cover of a torus the
  trace formula is Poisson summation, twisted by a genuine automorphic
  character. For the split torus Weil's character $\chi_\psi$ is a
  canonical choice. This case is used twice later: to describe the
  cuspidal data of the Borel subgroup, and to match the correction terms
  on the two sides of the invariant formula.
\item[The geometric side (Section~\ref{sec:geometric}).] The kernel of
  $R(f)$ is truncated, and the integral of the truncated kernel is
  computed class by class: orbital integrals for the elliptic classes,
  weighted orbital integrals for the hyperbolic classes, and the finite
  part of a zeta integral for the two central classes.
\item[The spectral side (Section~\ref{sec:spectral}).] The same
  truncated integral is computed from the spectral decomposition. It
  gives the traces on the cuspidal spectrum and on the residual spectrum
  (even Weil representations), an integral of the logarithmic derivative
  of the intertwining operator $M(w,\lambda)$ along the unitary principal
  series, and one extra term $\frac14\operatorname{tr}(M(w,0)I(\mu_0,f))$
  for each character $\mu$ fixed by the Weyl group.
\item[The invariant form (Section~\ref{sec:invariant}).] The weighted
  orbital integrals and the term with $M^{-1}M'$ are not invariant under
  conjugation. Both fail in the same way, and subtracting a common
  correction, built from the normalized local intertwining operators by
  Fourier inversion on the cover of $A$, makes both sides invariant.
\end{description}

The result is Theorem~\ref{thm:itf} with Corollary~\ref{cor:itf-final}.
For $f=f_S\otimes e^S$, with $S$ a large enough finite set of places,
\begin{align*}
&\sum_{z=\pm1}\Bigl(\mathrm{vol}\cdot f_S(z)+\sum_{\xi}a^G(S,u_\xi)\,
J_G(zu_\xi,f_S)+\tfrac12c_k\,I_A(z,f_S)\Bigr)\\
&\qquad+\sum_{\mathcal O\ \mathrm{elliptic}}\mathrm{vol}(T_\mu)\,J_G(\mu_S,f_S)
+\sum_{\mathcal O\ \mathrm{hyperbolic}}c_k\,I_A(\gamma_S,f_S)\\
&=\sum_\pi m_{\mathrm{disc}}(\pi)\operatorname{tr}\pi(f)
-\frac14\sum_{b\in k^\times/k^{\times2}}\operatorname{tr}I(\chi_{\psi_b,0},f)\\
&\qquad-\frac12\sum_\mu\frac1{2\pi}\int_\R\frac{m'}m(\mu,it)\,\operatorname{tr}
I(\mu_{it},f)\,dt .
\end{align*}
Here $J_G$ are orbital integrals, $I_A$ the invariant distributions that
replace the weighted orbital integrals, $\pi$ runs over the discrete
spectrum (cuspidal and residual), $\mu$ over the genuine automorphic
characters of the cover of $A$, and $m$ is the global normalizing factor
of the intertwining operator, essentially
$L(2\lambda,\chi^2)/L(2\lambda+1,\chi^2)$. The notation is collected in
the Glossary at the end.

Three features distinguish the double cover from $SL(2)$ itself, and they
are the main mathematical content of Sections~\ref{sec:spectral}
and~\ref{sec:invariant}. First, the Weyl group does not act on the cover
of $A$ by inversion; the characters it fixes are the Weil characters
$\chi_{\psi_b}$, one for each square class $b$
(Lemma~\ref{lem:w-action}). Second, the residual spectrum is not the
trivial representation but the sum of the even Weil representations
$r^+_{\psi_{-b}}$, which come from the pole of $M(w,\lambda)$ at
$\lambda=\frac12$ (Theorem~\ref{thm:residual}). Third, $M(w,0)=-1$ for
every one of these characters (Proposition~\ref{prop:Mw0}), whereas for
$SL(2)$ and a nontrivial quadratic character the corresponding operator
is not a scalar and is the source of the multiplicity formula of Labesse
and Langlands~\cite{LL}.

\subsection{Literature}\label{sec:literature}

\paragraph{Covering groups and the Weil representation.} The double
cover of $SL(2)$ over a local field and over the adeles first appeared in
Weil's work on the representation that now carries his name~\cite{We64}.
Kubota gave the explicit cocycle used here, for the $n$-fold cover of
$SL(2)$ over a local field containing the $n$-th roots of unity, and the
splitting over the rational points~\cite{K1,K2}. Moore~\cite{M} and
Matsumoto~\cite{Mat} treated the central extensions of split groups
over local fields and adele rings in general. Kazhdan and
Patterson~\cite{KP} studied the covers of $GL(r)$ and their automorphic
forms, and Brylinski and Deligne~\cite{BD} placed all these covers in one
algebraic framework, as extensions by $K_2$. Gelbart's lecture
notes~\cite{Ge2} remain the standard account of the double cover of
$SL(2)$ and $GL(2)$, the Weil representation and theta functions.
Covering groups are not algebraic groups, and the Borel cocycles that
describe them are not continuous; the relevant general theorem on the
topology of such extensions is Mackey's~\cite{Ma57}. For the Langlands
program for covering groups, and its history, see Gan, Gao and
Weissman~\cite{GGW} and Weissman~\cite{We18}.

\paragraph{Half-integral weight and the Shimura correspondence.}
Shimura~\cite{Sh} attached to a cusp form of weight $k/2$ a form of even
weight $k-1$. Niwa~\cite{Ni} and Shintani~\cite{Shi} realized the
correspondence and its inverse by theta series. Gelbart and
Piatetski-Shapiro~\cite{GPS80,GPS83} and Waldspurger~\cite{Wa80,Wa81,Wa91}
recast it in terms of automorphic representations of the double cover.
Waldspurger's description of the discrete spectrum of the double cover of
$SL(2)$, through the theta correspondence with $PGL(2)$ and its inner
forms, is the definitive result for the group studied here. The
identification of the residual spectrum with Weil representations, used
in Theorem~\ref{thm:residual}, belongs to this circle of
results~\cite{Ge2,GPS80,Wa80}. Gan and Savin~\cite{GS12}
obtained the local Langlands correspondence for the double cover of
$Sp(2n)$ from the theta correspondence, and Gan and Ichino~\cite{GI}
extended Waldspurger's global description to the generic part of the
discrete spectrum of the double cover of $Sp(2n)$.

\paragraph{The trace formula.} The trace formula goes back to
Selberg~\cite{Sel}. For $GL(2)$ it was put in adelic form by Jacquet and
Langlands~\cite{JL}; the treatment of Gelbart and Jacquet~\cite{GJ} is the
closest in spirit to the present paper, and we compare signs and
constants with it in Section~\ref{sec:spectral}. The classical Selberg
trace formula with multiplier systems and real weight, as in
Hejhal~\cite{He}, covers forms of half-integral weight on congruence
subgroups of $SL(2,\Z)$. For a general connected reductive group the
formula is Arthur's: the coarse expansions~\cite{A2,A3}, the invariant
form~\cite{A6}, the fine expansions~\cite{A12,A14}, the local theory of
weighted orbital integrals and normalized intertwining
operators~\cite{A7,A9,A11}, and the invariant trace
formula~\cite{A10}; his Clay lectures~\cite{Ar05} are the best
introduction. The spectral decomposition on which the spectral side
rests is due to Langlands~\cite{L2}, and M\oe glin and
Waldspurger~\cite{MW2} wrote it in a form that includes finite central
covers. For $SL(2)$ the stable trace formula is the subject of Labesse
and Langlands~\cite{LL}.

\paragraph{The trace formula for covering groups.} Flicker~\cite{Fl} used
a trace formula for the $n$-fold covers of $GL(2)$ to prove the
metaplectic correspondence with $GL(2)$, and Flicker and
Kazhdan~\cite{FK} did the same for $GL(r)$; see also Kazhdan and Patterson~\cite{KP2}. Mezo~\cite{Me1,Me2} prepared
the comparison of the full invariant trace formulas of $GL(r)$ and its
covers. The first author studied a metaplectic analog of Kazhdan's
endoscopic lifting in~\cite{J88}. The 1995 version of the present notes
wrote out Arthur's formula for the $n$-fold cover of $SL(2)$, with
several steps conjectural. Those steps are now theorems of
Li, who extended Arthur's trace formula to covering groups of connected
reductive groups in the four papers~\cite{LiI,LiII,LiIII,LiIV}: the fine
geometric expansion, the local harmonic analysis including normalized
intertwining operators, the fine spectral expansion, and the invariant
distributions (the invariant formula under a hypothesis at the real
places, which he verifies for the double cover of $Sp(2n)$). For the double cover of $Sp(2n)$ Li went on to
develop endoscopy, building on Adams~\cite{Ad} and Renard~\cite{Re} at
the real place: the transfer of orbital integrals~\cite{Li11}, the
weighted fundamental lemma~\cite{Li12}, the stabilization of the elliptic
terms~\cite{Li15}, local character relations~\cite{Li19}, the
stabilization of the full trace formula~\cite{Li21}, and a multiplicity
formula for the discrete spectrum in terms of Arthur
parameters~\cite{Li24}. For $n=1$ these results apply to the group of
this paper (Li works with an eight-fold cover that contains the double
cover). A $p$-adic character lifting between
the double cover of $SL(2)$ and $SO(1,2)$ is in Schultz's
thesis~\cite{Sch}.

\paragraph{Local harmonic analysis on the double cover of $SL(2)$.} The
irreducibility of the unitary genuine principal series is due to
Moen~\cite{Mo2}. Szpruch~\cite{Sz09} computed the local coefficients of
the genuine principal series of the double cover of $SL(2,F)$, and
Goldberg and Szpruch~\cite{GS16} the Plancherel measures for the $n$-fold
covers. These results bear on the normalizing factors at the places where
our Section~\ref{sec:normalized} is not explicit. The unramified computation
of the intertwining operator (the Gindikin--Karpelevich formula) for
general covers is in~\cite{KP,McN,Gao}; see also Savin~\cite{Sav} for the
Hecke algebra correspondence.

\paragraph{The place of this paper.} Nothing in
Sections~\ref{sec:background}--\ref{sec:invariant} goes beyond what the
general theorems of Arthur and Li give. What the paper adds is an
account in which the group is fixed, the constants are determined, the
proofs that are short for this group are given in full, and the
remaining appeals to the literature are listed. We hope that it is
useful as a worked example of Li's theory and as a companion to the
classical treatments of half-integral weight.

\subsection{What is proved here and what is cited}\label{sec:proved-cited}

Each section ends with a subsection ``Status of the results'', which
lists the statements taken from the literature and the statements
checked numerically. In summary:

\begin{center}
\begin{tabular}{p{0.11\textwidth}p{0.40\textwidth}p{0.40\textwidth}}
Section & Proved in the text & Taken from the literature\\\hline
\ref{sec:background} & the splittings, the center, commutativity of
  preimages of tori, the local triviality of the cocycle, the adelic
  cover and its splitting over $G(\Q)$ &
  Kubota's cocycle and the existence of the topological
  cover~\cite{K1,M}; the splitting over $K_v$~\cite{K2}; strong
  approximation\\
\ref{sec:rank0} & the trace formula for covers of tori and its
  reduction to Poisson summation & the properties of the Weil
  index~\cite{We64}\\
\ref{sec:geometric} & the coarse geometric expansion and the
  polynomials $J^T_{\mathcal O}$; the local distribution $J_A(z,f)$ &
  the fine expansion at the central classes~\cite{A12,LiI}, whose
  coefficients $a^G(S,u_\xi)$ are not evaluated\\
\ref{sec:spectral} & cuspidal data and the Weyl action; the unramified
  intertwining operator; the terms $J^T_\chi$ from the Maass--Selberg
  relations; the matching $V_b\cong r^+_{\psi_{-b}}$ &
  the spectral decomposition and Eisenstein series~\cite{L2,MW2}; the
  convergence and polynomiality of the spectral
  side~\cite{A3,LiIII}; the residual spectrum as theta
  functions~\cite{Ge2,GPS80,Wa80}\\
\ref{sec:invariant} & the variance formulas; the real normalizing
  factors; the invariance of $I_A$; the invariant trace formula;
  $M(w,0)=-1$ & the existence of normalizing factors~\cite{A11,LiII},
  which are not explicit at the dyadic and ramified places
\end{tabular}
\end{center}

Two ingredients of the final formula are therefore defined but not
evaluated: the unipotent coefficients $a^G(S,u_\xi)$, and the local
normalizing factors $r_v$ at the places where the additive character or
the inducing character is ramified, and at the dyadic places. Remarks
numbered ``2026.$n$'' record these and other points that rest on
citation or on numerical checks only. The explicit formulas are tested
by the scripts in \texttt{checks/}; they are evidence, not proofs.

The stable form of the trace formula, Section~7 of the 1995 notes, is not
part of this draft. The 1995 notes define stable conjugacy on the cover
through conjugation in a cover of $GL(2)$, following~\cite{LL}; its
relation to Li's endoscopy~\cite{Li11,Li15} has to be worked out when
that section is revised. Section~\ref{sec:gl2} prepares for it.

\subsection{A running example}\label{sec:running}

The case $k=\Q$, $S=\{\infty,2\}$ is the smallest one in which every
feature of the formula appears, and it is followed through the paper. We
collect it here; $\psi$ is the character of $\A/\Q$ with
$\psi_\infty(x)=e^{-2\pi ix}$.

\begin{itemize}
\item \emph{The cover.} At the real place the cover is $Mp(2,\R)$; its
  center is cyclic of order $4$, and so is the center at the place $2$
  (Lemma~\ref{lem:local}(e), Example~\ref{ex:real}). At an odd prime $p$
  the cover splits over $SL(2,\Z_p)$ by $k\mapsto(k,\kappa_p(k))$, and on
  $SL(2,\Z)$ the product of the $\kappa_p$ is a Jacobi symbol
  (Example~\ref{ex:jacobi}). At $2$ we use a splitting over a congruence
  subgroup (Lemma~\ref{lem:local}(c), Remark~\ref{rem:dyadic}). Test functions are
  $f=f_\infty\otimes f_2\otimes\bigotimes_{p\ \mathrm{odd}}e_p$.
\item \emph{Rank $0$.} $\chi_{\psi,\infty}(a(-1),1)=i$ and
  $\chi_{\psi,p}(a(p),1)\in\{1,i\}$ for odd $p$
  (Examples~\ref{ex:order4} and~\ref{ex:weil-Q}): a genuine character of
  the cover of $A$ is not an ordinary character times a sign.
\item \emph{Geometric side.} $c_\Q=1$. Besides the elliptic classes
  there are the two central classes $\pm1$ and the hyperbolic classes of
  $a(t)$, $t=\pm2^j$, $j\ge1$. The unipotent terms are indexed by
  $\Q_S^\times/\Q_S^{\times2}$, which has $2\cdot8=16$ elements. The
  constant term of the unipotent contribution is computed for a Gaussian
  in Example~\ref{ex:unip-Q}.
\item \emph{Spectral side.} The characters fixed by the Weyl group that
  are unramified outside $S$ are the four $\chi_{\psi_b}$,
  $b\in\{\pm1,\pm2\}$ (Example~\ref{ex:ramified-Q}). Each carries a
  residual representation $V_b\cong r^+_{\psi_{-b}}$; $V_1$ contains the
  classical theta function $\sum_ne^{2\pi in^2z}$
  (Example~\ref{ex:theta}). The partial normalizing factor is
  $(1-2^{-2\lambda})\zeta(2\lambda)/((1-2^{-1-2\lambda})\zeta(2\lambda+1))$,
  with residue $2/\pi^2$ at $\lambda=\frac12$
  (Examples~\ref{ex:theta} and~\ref{ex:mprime-Q}).
\item \emph{Invariant form.} At the real place the normalized
  intertwining operator is $1$ on the $K$-type of weight $\pm\frac12$ and
  a product of factors $(c-\lambda)/(c+\lambda)$ on the others
  (Example~\ref{ex:real-R}), and the correction $\Phi_A$ is an Abel
  transform smoothed by an exponentially decaying kernel
  (Example~\ref{ex:Phi-real}). The only bad place is $2$, where the
  normalizing factor is not explicit. The discrete correction is
  $-\frac14\sum_{b=\pm1,\pm2}\operatorname{tr}I(\chi_{\psi_b,0},f)$
  (Example~\ref{ex:itf-Q}).
\end{itemize}

\subsection{Conventions}\label{sec:conventions-intro}

The revision treats the double cover only; the 1995 notes treated the
$n$-fold cover, and remarks in Sections~\ref{sec:background}
and~\ref{sec:rank0} say what changes for $n>2$. Comparisons with the
1995 notes and with the 1997 book manuscript, including the corrections
made, are collected in Appendix~\ref{sec:appendix-1995} and are not
needed to read the paper. Notation and terms are listed in the Glossary.

% SPDX-License-Identifier: CC-BY-4.0
% =====================================================================
% Section 2 of D. Joyner, "A hitch-hiker's guide to the metaplectic trace
% formula for SL(2)", revised for the 2-fold cover (n = 2), November 2026.
%
% Stage 1 of ../revision_log.md. Source: ../original-1995/cctrace_new.tex,
% Section 2 (lines 98-717), with ../book-1997/book1.tex for the global NAK
% lemma. Referee items I.1-I.10 and the Section 2 minor items are fixed here;
% see the "Stage 1" section of ../revision_log.md for the item-by-item list.
%
% This file is \input by the master; it compiles on its own through
% preview.tex. Lemma numbers 2.1-2.9 are those of the 1995 notes;
% Lemma 2.10 is new. Every explicit formula for n = 2 is checked by
% checks/sec2_n2_check.py (exact rational arithmetic).
% Comments beginning "% [Claude]" are editorial notes by Claude.
% =====================================================================

\section{Background on the metaplectic double cover}\label{sec:background}

This section fixes the double cover of $G=\Res_{k/\Q}SL(2)$ and the
conventions used in the rest of the paper. Section~\ref{sec:local}
treats a single place: Kubota's cocycle, the splittings over unipotent
and compact subgroups, the center, and the preimages of tori. These are
then assembled into the cover of $G(\Q_p)$. Section~\ref{sec:global}
builds the adelic cover and its splitting over $G(\Q)$, and
Section~\ref{sec:gl2} fixes a double cover of $GL(2)$ and the resulting
$GL(2)$-action on the cover of $SL(2)$, which is what Section~7 uses.

From now on $n=2$. The 1995 notes treated the $n$-fold cover; for
$n=2$ several statements become simpler, and one, the
commutativity of preimages of abelian subgroups
(Lemma~\ref{lem:local}(g)), becomes a two-line argument. What changes
for general $n$ is listed in Remark~\ref{rem:general-n}.

\subsection*{Notation}

Let $k$ be a number field. (For $n=2$ the hypothesis
$\mu_n\subset k$ of the 1995 notes always holds, and $k$ may have real
places.) Let $\Sigma_k$ be the set of places of $k$, $\Sigma_k^\infty$
the archimedean places and $\Sigma_k^f$ the finite places. Let $\A$ be
the adele ring of $\Q$ and $\A_k$ that of $k$. For $v\in\Sigma_k$ let
$k_v$ be the completion. For $F=k_v$ non-archimedean let $\OO=\OO_F$ be
its ring of integers, $\varpi=\varpi_F$ a uniformizer, $q$ the order of
the residue field and $\ord=\ord_F\colon F^\times\to\Z$ the normalized
valuation, so $\ord(\varpi)=1$ and $|\varpi|=q^{-1}$. We call $v$
\emph{dyadic} if $v\mid 2$, and put $e_v=\ord_v(2)$. Absolute values
are normalized; at a complex place $|z|_v=z\bar z$.

Let $(a,b)=(a,b)_v$ be the quadratic Hilbert symbol of
$F=k_v$~\cite{W1,Se}. We use only the following properties:
\begin{enumerate}[label=(H\arabic*)]
\item\label{H1} $(\cdot,\cdot)$ is bimultiplicative and symmetric,
  with values $\pm1$; hence $(a^{-1},b)=(a,b)$.
\item\label{H2} $(a,b)=1$ if $a$ or $b$ lies in $F^{\times2}$;
  $(a,-a)=1$; and $(a,1-a)=1$ for $a\neq 0,1$ (the Steinberg relation).
\item\label{H3} If $v$ is finite and not dyadic and $a,b\in\OO^\times$,
  then $(a,b)=1$. If $v$ is complex, $(\cdot,\cdot)\equiv1$. If $v$
  is real, $(a,b)=-1$ exactly when $a<0$ and $b<0$.
\item\label{H4} (Product formula) For $a,b\in k^\times$,
  $(a,b)_v=1$ for almost all $v$ and $\prod_v(a,b)_v=1$.
\end{enumerate}
By~\ref{H3}, $(-1,-1)_v=1$ at non-dyadic finite places and at complex
places, and $(-1,-1)_v=-1$ at real places. At a dyadic place both values
occur; $(-1,-1)_{\Q_2}=-1$.

Let $G$ be the group over $\Q$ obtained from $SL(2)$ by restriction of
scalars from $k$ to $\Q$, so $G(\Q)=SL(2,k)$ and
$G(\Q_p)=\prod_{v\mid p}SL(2,k_v)$.

\subsection{Local aspects}\label{sec:local}

Here $F=k_v$ for one place $v$. We define the cover, compute it on the
subgroups that matter later (unipotent, compact, diagonal, central,
anisotropic), and then pass to $G(\Q_p)$.

For $g=\begin{pmatrix}a&b\\c&d\end{pmatrix}\in SL(2,F)$ put
$x(g)=c$ if $c\neq0$ and $x(g)=d$ if $c=0$, and define
$\beta=\beta_v\colon SL(2,F)\times SL(2,F)\to\mu_2$ by
\begin{equation}\label{eq:beta}
\beta(g_1,g_2)
=\Bigl(\frac{x(g_1g_2)}{x(g_1)},\frac{x(g_1g_2)}{x(g_2)}\Bigr)
=\bigl(x(g_1),x(g_2)\bigr)\bigl(-x(g_1)^{-1}x(g_2),\,x(g_1g_2)\bigr).
\end{equation}
At a complex place $\beta\equiv1$ by~\ref{H3}.

\begin{lemma}[Kubota~\cite{K1}, Moore~\cite{M}]\label{lem:kubota}
$\beta$ is a continuous Borel $2$-cocycle, $\mu_2$ carrying the discrete
topology. Its class in $H^2(SL(2,F),\mu_2)$ defines a topological
central extension
\[
1\to\mu_2\to\overline{SL(2,F)}\xrightarrow{\;p\;}SL(2,F)\to1,
\]
which is nontrivial unless $F=\C$.
\end{lemma}

The two expressions in~\eqref{eq:beta} agree, and the cocycle identity
holds; both are checked on random elements over $\R$, $\Q_2$, $\Q_3$,
$\Q_5$ and $\Q_7$ in \texttt{checks/sec2\_n2\_check.py}. For $F=\R$ the
group $\overline{SL(2,\R)}$ is the metaplectic group $Mp(2,\R)$.

For a subgroup $H\subset SL(2,F)$ put $\overline H=p^{-1}(H)$. Elements
of $\overline{SL(2,F)}$ are written $(g,\zeta)$ with $g\in SL(2,F)$,
$\zeta\in\mu_2$, and multiplied by
\begin{equation}\label{eq:law}
(g_1,\zeta_1)(g_2,\zeta_2)=(g_1g_2,\ \beta(g_1,g_2)\zeta_1\zeta_2).
\end{equation}
Since $\zeta^{-1}=\zeta$, inverses and conjugates are
\begin{equation}\label{eq:inverse}
(g,\zeta)^{-1}=(g^{-1},\beta(g,g^{-1})\zeta),\qquad
(g,\zeta)(h,\zeta')(g,\zeta)^{-1}
=\bigl(ghg^{-1},\ \beta(g,h)\beta(gh,g^{-1})\beta(g,g^{-1})\zeta'\bigr).
\end{equation}

Let $B\subset SL(2)$ be the upper triangular Borel subgroup, $N$ its
unipotent radical and $A$ the diagonal torus. Write
$a(t)=\diag(t,t^{-1})$ and $n(u)=\begin{pmatrix}1&u\\0&1\end{pmatrix}$.
For $\theta\in F^\times\smallsetminus F^{\times2}$ let
\[
T_\theta=\Bigl\{\begin{pmatrix}a&b\theta\\b&a\end{pmatrix}
\;\Big|\;a,b\in F,\ a^2-\theta b^2=1\Bigr\}\cong E^1,
\qquad E=F(\sqrt\theta),
\]
the norm-one group of $E$. Every anisotropic maximal torus of
$SL(2,F)$ is $GL(2,F)$-conjugate to some $T_\theta$.
% [Claude] The 1995 text has "a^2+b^2\theta=1", which is not the
% determinant of the displayed matrix.

If $v$ is finite let $K=K_0=SL(2,\OO)$ and, for $m\ge1$,
\[
K_m=\Bigl\{\begin{pmatrix}a&b\\c&d\end{pmatrix}\in SL(2,\OO)
\;\Big|\;a-1,\ d-1,\ b,\ c\in\varpi^m\OO\Bigr\}.
\]
If $v$ is non-dyadic, define $\kappa=\kappa_v\colon K\to\mu_2$ by
\begin{equation}\label{eq:kappa}
\kappa\begin{pmatrix}a&b\\c&d\end{pmatrix}
=\begin{cases}(c,d)_v,&0<|c|_v<1,\\ 1,&\text{otherwise.}\end{cases}
\end{equation}
If $v$ is archimedean let $K=SO(2)$ ($v$ real) or $K=SU(2)$
($v$ complex).

\begin{lemma}\label{lem:local}
Let $F=k_v$.
\begin{enumerate}[label=(\alph*)]
\item $\beta\equiv1$ on $N(F)\times N(F)$, so $n\mapsto(n,1)$ is a
  homomorphism and $\overline{N(F)}=N(F)\times\mu_2$.
\item Let $v$ be finite and non-dyadic. Then $\beta(k_1,k_2)=
  \kappa(k_1)\kappa(k_2)\kappa(k_1k_2)$ for $k_1,k_2\in K$, so
  $k\mapsto(k,\kappa(k))$ is a homomorphism and
  $\overline K=K^*\times\mu_2$ with $K^*=\{(k,\kappa(k))\}$. The
  restriction of $\beta$ to $K\times K$ is not identically $1$, so the
  set $\{(k,1)\mid k\in K\}$ is not a subgroup.
\item Let $v$ be finite and $m\ge1$ with $1+\varpi^m\OO\subset
  F^{\times2}$. Then $\beta\equiv1$ on $K_m\times K_m$, so
  $k\mapsto(k,1)$ splits $\overline{K_m}$. The hypothesis holds for every
  $m\ge1$ if $v$ is non-dyadic, and for $m\ge 2e_v+1$ if $v$ is dyadic.
\item Let $v$ be archimedean. If $v$ is complex,
  $\overline{SL(2,F)}=SL(2,F)\times\mu_2$. If $v$ is real, $\beta\equiv1$
  on $A(\R)^0\times A(\R)^0$, so $\overline{A(\R)^0}=A(\R)^0\times\mu_2$,
  while $\overline{SO(2)}$ is connected, the nontrivial double cover of
  the circle~\cite{Ge2}.
\item The center of $\overline{SL(2,F)}$ is $\overline{\{\pm1\}}$, a
  group of order $4$ in which $(-1,\zeta)^2=(1,(-1,-1)_v)$. It is cyclic
  when $(-1,-1)_v=-1$ (for example $F=\R$ or $F=\Q_2$) and is
  $\{\pm1\}\times\mu_2$ otherwise.
\item On $A(F)$, $\beta(a(t_1),a(t_2))=(t_1,t_2)_v$. Hence
  $\overline{A(F)}$ is abelian, and
  $(a(t),\zeta)^2=(a(t^2),(t,-1)_v)$.
\item If $g_1,g_2\in SL(2,F)$ commute, then $\beta(g_1,g_2)=
  \beta(g_2,g_1)$. Consequently $\overline H$ is abelian for every abelian
  subgroup $H\subset SL(2,F)$; in particular $\overline{A(F)}$ and every
  $\overline{T_\theta}$ are abelian.
\item Let $U=gNg^{-1}$ with $g\in SL(2,F)$. Then $u\mapsto(u,1)$ is the
  unique homomorphic section of $\overline U\to U$. Consequently, for every
  $(g,\zeta)\in\overline{SL(2,F)}$ and $u\in N(F)$,
  \[
  (g,\zeta)(u,1)(g,\zeta)^{-1}=(gug^{-1},1).
  \]
\end{enumerate}
\end{lemma}

\begin{proof}
(a) For $n_1,n_2\in N(F)$ all three values of $x$ equal $1$.

(b) This is Kubota's computation~\cite{K2}; see also~\cite{Ge2}. The
identity is also checked numerically for $F=\Q_3,\Q_5,\Q_7$, together
with the existence of pairs in $K\times K$ with $\beta=-1$.

(c) Let $k_i=\begin{pmatrix}a_i&b_i\\c_i&d_i\end{pmatrix}\in K_m$ for
$i=1,2$, and $k_1k_2=\begin{pmatrix}a_3&b_3\\c_3&d_3\end{pmatrix}$. All
$a_i$ and $d_i$ lie in $1+\varpi^m\OO\subset F^{\times2}$, and
$c_3=c_1a_2+d_1c_2$. If $c_1,c_2,c_3\neq0$, then
$c_1a_2/c_3+d_1c_2/c_3=1$, so by~\ref{H2} and~\ref{H1}
\[
\beta(k_1,k_2)=\Bigl(\frac{c_3}{c_1},\frac{c_3}{c_2}\Bigr)
=\Bigl(\frac{c_1a_2}{c_3},\frac{d_1c_2}{c_3}\Bigr)=1,
\]
because $a_2$ and $d_1$ are squares. If $c_3=0\neq c_1c_2$, then
$x(k_1k_2)=d_3$ is a square and $c_2=-c_1a_2/d_1$, so
$\beta=(c_1,c_2)=(c_1,-c_1)(c_1,a_2/d_1)=1$. If $c_1=0\ne c_2$, then
$x(k_1)=d_1$ is a square and $c_3=d_1c_2$, so
$\beta=(c_3/d_1,\,d_1)=1$; the case $c_2=0\neq c_1$ is the same with
$a_2$ in place of $d_1$. If $c_1=c_2=0$, all $x$-values are squares.
Finally, $1+\varpi\OO\subset F^{\times2}$ for non-dyadic $F$ by Hensel's
lemma, and $1+4\varpi\OO\subset F^{\times2}$ for dyadic $F$ by the
binomial series for $(1+4\varpi y)^{1/2}$; and $4\varpi\OO=
\varpi^{2e_v+1}\OO$.

(d) At a complex place $\beta\equiv1$ by~\ref{H3}. For $t_1,t_2>0$,
$(t_1,t_2)_\R=1$ by~\ref{H3} and~(f). The statement about $SO(2)$ is
standard~\cite{Ge2}.

(e) Commuting elements give commuting lifts by~(g), so
$\overline{\{\pm1\}}$ is central; conversely, the image of a central
element is central in $SL(2,F)$. Since $x(-1)=-1$ and $x(1)=1$,
$(-1,\zeta)^2=(1,\beta(-1,-1))=(1,(-1,-1)_v)$ by~\eqref{eq:beta}. A
group of order $4$ with an element of order $4$ is cyclic, and one of
exponent $2$ is $\{\pm1\}\times\mu_2$.

(f) Here $x(a(t))=t^{-1}$, so~\eqref{eq:beta} gives
$\beta(a(t_1),a(t_2))=(t_2^{-1},t_1^{-1})=(t_1,t_2)$ by~\ref{H1}. The
square is $(a(t^2),(t,t))$, and $(t,t)=(t,-1)$ by~\ref{H2}.

(g) If $g_1g_2=g_2g_1$, then $x(g_1g_2)=x(g_2g_1)$, and the first
expression in~\eqref{eq:beta} is symmetric in $g_1,g_2$ because the
quadratic Hilbert symbol is symmetric~\ref{H1}. By~\eqref{eq:law}, lifts
$(g_1,\zeta_1)$ and $(g_2,\zeta_2)$ commute if and only if
$\beta(g_1,g_2)=\beta(g_2,g_1)$.

(h) If $g\in B(F)$, then $U=N$ and the claim is~(a). Otherwise $g$ has
lower left entry $\gamma\ne0$, and $u_s=g\,n(s)\,g^{-1}$ has lower left
entry $-\gamma^2s$. So for $s_1,s_2\ne0$ with $s_1+s_2\ne0$,
\[
\beta(u_{s_1},u_{s_2})
=\Bigl(\frac{s_1+s_2}{s_1},\frac{s_1+s_2}{s_2}\Bigr)
=\Bigl(\frac{s_1}{s_1+s_2},\frac{s_2}{s_1+s_2}\Bigr)=1
\]
by~\ref{H1} and the Steinberg relation~\ref{H2}. If $s_1+s_2=0$, then
$x(u_{s_1}u_{s_2})=1$ and $\beta=(1/s_1,-1/s_1)=1$. So $u\mapsto(u,1)$
is a homomorphism on $U$. Two homomorphic sections differ by a
continuous homomorphism $U\cong F\to\mu_2$, which is trivial because
$F=2F$. Conjugation by $(g,\zeta)$ carries the section $n\mapsto(n,1)$
of $\overline N$ to a homomorphic section of $\overline{gNg^{-1}}$, which
is therefore $u\mapsto(u,1)$.
\end{proof}

\begin{example}\label{ex:real}
Take $F=\R$. By~(e) and~(f), the element $(a(-1),1)=(-1,1)$ has order
$4$, so $\overline{\{\pm1\}}\cong\Z/4$ and $\overline{A(\R)}$ is an
abelian group which is not $A(\R)\times\mu_2$. By~(d), $\overline{A(\R)}$
is generated by $\overline{A(\R)^0}=A(\R)^0\times\mu_2$ and $(-1,1)$.
At a non-dyadic place, on the other hand, $(-1,-1)_v=1$, so the center
is $\{\pm1\}\times\mu_2$; and $\overline{A(\OO^\times)}$ is split by
$a\mapsto(a,1)$, because $(t_1,t_2)_v=1$ for units by~\ref{H3}. This is
consistent with~(b), since $\kappa\equiv1$ on $A(\OO^\times)$.
\end{example}

\begin{remarkrev}\label{rem:dyadic}
For $F=\Q_2$, Lemma~\ref{lem:local}(c) gives the splitting $k\mapsto
(k,1)$ over $K_3$, and $\beta$ is not identically $1$ on $K_2\times K_2$.
The cover nevertheless splits over $K_2$, by $k\mapsto(k,(c,d)_2)$ for
$c\ne0$ and $k\mapsto(k,1)$ for $c=0$. This is the local form at $2$ of
the classical splitting over $\Gamma_1(4)$~\cite{K2,Ge2}. We have
verified it numerically only (\texttt{checks/sec2\_n2\_check.py}), and
the same section fails on $K_1$. Nothing below depends on it: the rest
of the paper uses only that some $K_m$ splits at each dyadic place.
\end{remarkrev}

\paragraph{The cover of \texorpdfstring{$G(\Q_p)$}{G(Qp)}.}
Let $\overline{G(\Q_p)}$ be the extension of $G(\Q_p)=\prod_{v\mid p}
SL(2,k_v)$ by $\mu_2$ defined by the cocycle
\[
\beta_p(g_1,g_2)=\prod_{v\mid p}\beta_v(g_{1,v},g_{2,v}).
\]
It is the \emph{quotient} of $\prod_{v\mid p}\overline{SL(2,k_v)}$ by
the kernel $\{(\zeta_v)\mid\prod_v\zeta_v=1\}$ of the product map
$\prod_{v\mid p}\mu_2\to\mu_2$. (The 1995 notes call it a subgroup.)
For $H\subset G(\Q_p)$ put $\overline H=p^{-1}(H)$. Let $B$, $N$, $A$ be
the $\Q$-subgroups of $G$ obtained from those of $SL(2)$, and let
$T=\prod_{v\mid p}T_v$ with each $T_v$ a maximal torus of $SL(2,k_v)$.
Let $K_p=\prod_{v\mid p}K_v$. For $p$ odd, $K_p^*=\prod_{v\mid p}K_v^*$,
with $K_v^*$ as in Lemma~\ref{lem:local}(b), maps to a subgroup of
$\overline{G(\Q_p)}$ which we also denote by $K_p^*$.

\begin{lemma}\label{lem:Qp}
\begin{enumerate}[label=(\alph*)]
\item $n\mapsto(n,1)$ splits $\overline{N(\Q_p)}$.
\item If $p$ is odd, $\overline{K_p}=K_p^*\times\mu_2$.
\item For every $p$, $k\mapsto(k,1)$ splits $\overline{K_{m,p}}$, where
  $K_{m,p}=\prod_{v\mid p}K_{m,v}$, as soon as $m\ge 2e_v+1$ for every
  $v\mid p$; if $p$ is odd, $m\ge1$ suffices.
\item If $p=\infty$, $\overline{A(\R)^0}=A(\R)^0\times\mu_2$, where
  $A(\R)^0$ is the identity component of $A(k\otimes\R)$.
\item The center of $\overline{G(\Q_p)}$ is
  $\overline{\{\pm1\}^{S_p}}$, where $S_p=\{v\mid p\}$. It is abelian,
  and the square of a lift of $\varepsilon=(\varepsilon_v)$ is
  $\prod_{v\mid p,\ \varepsilon_v=-1}(-1,-1)_v$.
\item $\beta_p(a_1,a_2)=\prod_{v\mid p}(t_{1,v},t_{2,v})_v$ for
  $a_i=a(t_i)\in A(\Q_p)$, and $\overline{A(\Q_p)}$ is abelian.
\item The preimage of any abelian subgroup of $G(\Q_p)$, in particular
  $\overline T$, is abelian, and Lemma~\ref{lem:local}(h) holds for
  $G(\Q_p)$.
\end{enumerate}
\end{lemma}

\begin{proof}
Each statement holds componentwise by Lemma~\ref{lem:local}, and the
quotient map is compatible with all of them.
\end{proof}

\subsubsection{\texorpdfstring{$NAK$}{NAK} decompositions}

\begin{lemma}\label{lem:NAK-local}
$\overline{G(\Q_p)}=N(\Q_p)\cdot\overline{A(\Q_p)}\cdot\overline{K_p}$,
with $N(\Q_p)$ embedded by $n\mapsto(n,1)$. If
$(n,1)\,\tilde a\,\tilde k=(n',1)\,\tilde a'\,\tilde k'$ are two such
decompositions, then $p(\tilde a')\in p(\tilde a)\,(A(\Q_p)\cap K_p)$.
In particular $|t|_p:=\prod_{v\mid p}|t_v|_v$ is well defined for
$p(\tilde a)=a(t)$, and depends only on $(g,\zeta)$. For $p$ odd one may
take $\tilde k\in K_p^*$.
\end{lemma}

\begin{proof}
The Iwasawa decomposition $G(\Q_p)=N(\Q_p)A(\Q_p)K_p$ holds, with
$a$ determined modulo $A(\Q_p)\cap K_p$, which is
$\prod_{v\mid p}A(\OO_v^\times)$ for $p<\infty$ and a product of copies
of $\{\pm1\}$ and of $\{\diag(z,\bar z):|z|=1\}$ for $p=\infty$. Since
$\mu_2$ is central and lies in $\overline{A(\Q_p)}$, this lifts to
$\overline{G(\Q_p)}$. For $p$ odd, $\overline{K_p}=K_p^*\mu_2$ and the
$\mu_2$ factor can be moved into $\tilde a$. The absolute value $|t|_p$
is trivial on $A(\Q_p)\cap K_p$.
\end{proof}

% [Claude] The 1995 lemma says that a is determined modulo the center and
% identifies K_p with {(k,1)}; both are corrected here (items I.1, I.7).

\subsubsection{Heights}

For $g=(g_v)\in G(\Q_p)$ put $\|g\|_p=\prod_{v\mid p}\|g_v\|_v$, where
$\|g_v\|_v=\max(|a|_v,|b|_v,|c|_v,|d|_v)$ for $v$ finite and
$\|g_v\|_v=(|a|^2+|b|^2+|c|^2+|d|^2)^{1/2}$, raised to the power
$[k_v:\R]$, for $v$ archimedean (here $|\cdot|$ is the usual absolute
value). Extend to the cover by $\|(g,\zeta)\|_p=\|g\|_p$.

Define $H=H_B\colon\overline{G(\Q_p)}\to\R$ by
\begin{equation}\label{eq:H}
H(g,\zeta)=\log|t|_p=\sum_{v\mid p}\log|t_v|_v,
\end{equation}
where $(g,\zeta)=(n,1)\,\tilde a\,\tilde k$ as in
Lemma~\ref{lem:NAK-local} and $p(\tilde a)=a(t)$. This is the product
of the normalized absolute values over $v\mid p$, that is
$|N_{k\otimes\Q_p/\Q_p}(t)|_p$, and not the maximum used in the 1995
notes; see Section~\ref{sec:chambers}.

\begin{lemma}\label{lem:H}
\begin{enumerate}[label=(\alph*)]
\item $H(\tilde a\tilde a')=H(\tilde a)+H(\tilde a')$ for
  $\tilde a,\tilde a'\in\overline{A(\Q_p)}$.
\item $H(\tilde n\,x\,\tilde k)=H(x)$ for $\tilde n\in\overline{N(\Q_p)}$,
  $x\in\overline{G(\Q_p)}$ and $\tilde k\in\overline{K_p}$.
\item $|H(x)|\ll1+\log\|x\|_p$ for $x\in\overline{G(\Q_p)}$~\cite{A2}.
\end{enumerate}
\end{lemma}

The global analog, including invariance under $B(\Q)$ on the left, is
Lemma~\ref{lem:NAK-global}.
% [Claude] The 1995 Lemma 2.5 is stated over A(A) and G(A) before these
% are defined, and (b) asserts left K-invariance, which fails; H is left
% N-invariant and right K-invariant.

\subsubsection{Characteristic functions of Weyl chambers}\label{sec:chambers}

Let $P_0=B$; the standard parabolic subgroups are $B$ and $G$. For a
standard $P=MN_P$ let $A_P$ be the split component of the center of $M$
over $\Q$. Thus $A_G=1$, because the center $\{\pm1\}$ of $SL(2)$ is
finite, and $A_B\cong\mathbb G_m$ is the image of
$t\mapsto a(t)$, $t\in\Q^\times\subset k^\times$: the $\Q$-split part of
the torus $A\cong\Res_{k/\Q}\mathbb G_m$, which has $\Q$-rank one.

The $\Q$-rational characters of $A$ form $X(A)_\Q\cong\Z$, generated by
$\chi_1=N_{k/\Q}\circ t$. Put $\mathfrak a=\mathfrak a_B=
\Hom(X(A)_\Q,\R)$, identified with $\R$ by $X\mapsto X(\chi_1)$. Then
$H$ of~\eqref{eq:H} and its global analog are
$\langle H(x),\chi_1\rangle=\log|\chi_1(\tilde a)|$, which is why $H$
uses the norm. Let $d=[k:\Q]$. The simple root $\alpha$ of $(B,A_B)$ is
$a(s)\mapsto s^2$ ($s\in\Q^\times$), while $\chi_1(a(s))=s^d$; so
$\alpha=(2/d)\chi_1$ in $\mathfrak a^*$, and $\alpha(X)=2X/d$ under
$\mathfrak a\cong\R$. Let $\alpha^\vee=X_\alpha\in\mathfrak a$ be the
coroot $s\mapsto a(s)$, so $X_\alpha=\langle\chi_1,\alpha^\vee\rangle=d$
and $\alpha(\alpha^\vee)=2$, and let $\varpi=\tfrac12\alpha^\vee$ be the
fundamental coweight, so $\alpha(\varpi)=1$. The modulus character is
$\delta_B=e^{2H}$. (The 1995 notes call $\varpi$ the dual root.)
% [Claude] Stage 3 correction: the stage-1 text had alpha = 2 chi_1 and
% X_alpha = 1, which holds only for k = Q. See Remark 2026.8.

Let $\mathfrak a^+=\{X\in\mathfrak a\mid\alpha(X)>0\}$ and
$(\mathfrak a^*)^+=\{\Lambda\in\mathfrak a^*\mid\Lambda(\varpi)>0\}$,
with characteristic functions $\tau_B$ and $\hat\tau_B$, and set
$\tau_G=\hat\tau_G=1$. For standard $P_1\subset P_2$ let
$\tau_1^2$, $\hat\tau_1^2$ and $\sigma_1^2$ be as in~\cite{A2}. With
$\mathfrak a\cong\R$:
\begin{gather}
\tau_B=\tau_B^G=\hat\tau_B=\hat\tau_B^G
=\text{characteristic function of }\{u\in\R\mid u>0\},\qquad
\tau_1^2=\begin{cases}\tau_1,&P_2=G,\\1,&P_2\ne G,\end{cases}
\label{eq:tau}\\
\sigma_B^G=\tau_B,\qquad\sigma_B^B=0,\qquad\sigma_G^G=1.
\label{eq:sigma}
\end{gather}

\subsubsection{Algebras of functions}\label{sec:functions}

We follow Li~\cite{LiI,LiIV}. A representation $\pi$ of a cover is
\emph{genuine} if $\pi(1,\zeta)=\zeta$ for $\zeta\in\mu_2$, and a
function $f$ on the cover is \emph{anti-genuine} if
$f(g,\zeta)=\zeta^{-1}f(g,1)$. Test functions are anti-genuine. Then
$\pi(f)=\int f(\tilde g)\pi(\tilde g)\,d\tilde g$ picks up the factor
$\int_{\mu_2}\zeta^{-1}\zeta\,d\zeta=1$, so it is in general nonzero on
genuine representations, and it vanishes on representations that are
trivial on $\mu_2$.

For $n=2$, $\zeta^{-1}=\zeta$, so ``anti-genuine'' and ``genuine''
define the same functions, and the conventions of the 1995 notes (genuine
functions in~\S2.1.4, $\pi(g,\zeta)=\zeta^{-1}\pi(g,1)$ in~\S3.1) agree
with Li's. We keep Li's terminology so that the formulas remain correct
for general $n$; see Remark~\ref{rem:general-n}.

Let $C_c^\infty(\overline{G(\Q_p)})_-$ be the space of anti-genuine
$f$ such that $g\mapsto f(g,1)$ lies in $C_c^\infty(G(\Q_p))$. Let $p$
be odd. The \emph{spherical Hecke algebra}
$\mathcal H(\overline{G(\Q_p)},K_p^*)$ consists of the
$f\in C_c^\infty(\overline{G(\Q_p)})_-$ with
$f(\tilde k\,x\,\tilde k')=f(x)$ for $\tilde k,\tilde k'\in K_p^*$. Its
unit element is
\begin{equation}\label{eq:unit}
e_p(\tilde k\zeta)=\zeta^{-1}\ (\tilde k\in K_p^*,\ \zeta\in\mu_2),\qquad
e_p=0\ \text{off}\ \overline{K_p}.
\end{equation}
The function $(k,\zeta)\mapsto\zeta^{-1}$ on $\overline{K_p}$, which is
what the identification $k\mapsto(k,1)$ of the 1995 notes suggests, is
$\prod_v\kappa_v\cdot e_p$; it is not bi-$K_p^*$-invariant.

\subsubsection{Measures}\label{sec:measures-local}

Let $dk$ be the Haar measure on $K_p$ of total volume $1$. Fix a
nontrivial additive character $\psi_v$ of $k_v$ for each $v$, and let
$dn$ be the measure on $N(\Q_p)\cong\bigoplus_{v\mid p}k_v$ that is the
product of the self-dual measures for the $\psi_v$. Let $da$ be a Haar
measure on $A(\Q_p)$. If $T$ is a maximal torus of $G$ over $\Q_p$, and
$h\in G(\overline{\Q}_p)$ conjugates $T$ to $A$, let $dt$ be the
pull-back of $da$ as in~\cite{R}; stably conjugate tori get compatible
measures as in~\cite{R}. Put
\[
\delta_B(a(t))=|\alpha(a(t))|_p=|t|_p^2.
\]
Normalize the Haar measure $dg$ on $G(\Q_p)$ by
\begin{equation}\label{eq:iwasawa}
\int_{G(\Q_p)}\phi(g)\,dg
=\int_{N}\int_{A}\int_{K}\phi(nak)\,\delta_B(a)^{-1}\,dk\,da\,dn
=\int_{K}\int_{A}\int_{N}\phi(kan)\,\delta_B(a)\,dn\,da\,dk.
\end{equation}
The two expressions agree because $dg$ is invariant under
$g\mapsto g^{-1}$ (as $G$ is unimodular) and $(nak)^{-1}=k^{-1}a^{-1}
n^{-1}$; see~\cite[Ch.~VIII]{Kn} for the real case. (The 1995 notes
have $\delta_B(a)^{1/2}$ for $g=kan$; the exponent must be $1$ for
$g=kan$ and $-1$ for $g=nak$.)

On a cover we use $d\tilde g=dg\,d\zeta$, where $d\zeta$ gives
$\mu_2$ total volume~$1$. Then $\overline{K_p}$ has volume $1$, and
$e_p$ of~\eqref{eq:unit} is idempotent: for $x=\tilde k\in K_p^*$,
\[
(e_p*e_p)(x)=\int_{\overline{K_p}}e_p(y)\,e_p(y^{-1}x)\,dy
=\int_{K_p^*}\int_{\mu_2}\zeta^{-1}\zeta\,d\zeta\,dy=1.
\]

\subsection{Global aspects}\label{sec:global}

The adelic cover is not the cover defined by $\prod_v\beta_v$: that
product does not converge, because $\beta_v$ is not trivial on
$K_v\times K_v$ (Lemma~\ref{lem:local}(b)). The standard
remedy~\cite{K2,Ge2,KP}, used here, is to correct $\beta_v$ by the
coboundary of $\kappa_v$. This choice also fixes the splitting over
$G(\Q)$.

For each place $v$ of $k$ extend $\kappa_v$ to $SL(2,k_v)$ by
\[
\kappa_v(g)=\begin{cases}(c,d)_v,&v\ \text{finite non-dyadic},\ g\in K_v,\
0<|c|_v<1,\\ 1,&\text{otherwise,}\end{cases}
\]
so $\kappa_v\equiv1$ at dyadic and archimedean places, and put
\begin{equation}\label{eq:sigma-v}
\sigma_v(g_1,g_2)=\beta_v(g_1,g_2)\,\kappa_v(g_1)\kappa_v(g_2)\kappa_v(g_1g_2).
\end{equation}
This cocycle is cohomologous to $\beta_v$; it defines the same cover
$\overline{SL(2,k_v)}$, identified with the $\beta_v$-model by
$(g,\zeta)_{\sigma}\mapsto(g,\kappa_v(g)\zeta)_\beta$. By
Lemma~\ref{lem:local}(b), $\sigma_v\equiv1$ on $K_v\times K_v$ at every
finite non-dyadic $v$. Since $g\in G(\A)$ has $g_v\in K_v$ for almost all
$v$, the product
\begin{equation}\label{eq:sigma-A}
\sigma_\A(g_1,g_2)=\prod_v\sigma_v(g_{1,v},g_{2,v})
\end{equation}
has only finitely many factors different from $1$. It is a continuous
$2$-cocycle on $G(\A)=SL(2,\A_k)$ and defines
\[
1\to\mu_2\to\overline{G(\A)}\xrightarrow{\;p\;}G(\A)\to1.
\]
Equivalently, $\overline{G(\A)}$ is the restricted product of the
$\overline{SL(2,k_v)}$ with respect to the subgroups $K_v^*$, divided by
$\{(\zeta_v)\in\bigoplus_v\mu_2\mid\prod_v\zeta_v=1\}$. It is a
\emph{quotient} of the restricted product, not a subgroup. Factorization
of genuine representations of $\overline{G(\A)}$ (and of
$\overline{A(\A)}$ in Section~3) refers to this quotient. Group law,
inverses and conjugation are given by~\eqref{eq:law}
and~\eqref{eq:inverse} with $\sigma_\A$ in place of $\beta$. For
$H\subset G(\A)$ put $\overline H=p^{-1}(H)$.

For $\gamma\in SL(2,k)$ we have $\gamma\in K_v$ and
$\ord_v(c)\le0$ or $c=0$ for almost all $v$, so
\begin{equation}\label{eq:kappa-A}
\kappa_\A(\gamma)=\prod_v\kappa_v(\gamma)
=\prod_{\substack{v\ \text{finite non-dyadic}\\ \gamma\in K_v,\ 0<|c|_v<1}}(c,d)_v
\end{equation}
is a finite product.

\begin{lemma}\label{lem:global}
\begin{enumerate}[label=(\alph*)]
\item The map $\gamma\mapsto(\gamma,\kappa_\A(\gamma))$ is a homomorphism
  $G(\Q)\to\overline{G(\A)}$, so $\overline{G(\Q)}=G(\Q)\times\mu_2$.
\item $n\mapsto(n,1)$ splits $\overline{N(\A)}$, and on $B(\Q)$ the
  splitting of~(a) is $b\mapsto(b,1)$.
\item Let $S$ be a finite set of places containing the archimedean and
  dyadic places, and let $K_f=\prod_{v\in S\smallsetminus\Sigma_k^\infty}
  K_{m_v,v}\times\prod_{v\notin S}K_v$, with $m_v\ge1$, and $m_v\ge
  2e_v+1$ if $v$ is dyadic. Then $k\mapsto(k,1)$ is a homomorphism
  $K_f\to\overline{G(\A)}$; we write $K_f^*$ for its image. (At $v\notin S$,
  translated to the $\beta_v$-model, its component is $K_v^*$.)
\item The center of $\overline{G(\A)}$ is $\overline{Z(G(\A))}$,
  where $Z(G(\A))=\prod_v\{\pm1\}$. It is abelian, and a lift of
  $\varepsilon=(\varepsilon_v)$ has square
  $\prod_{v:\,\varepsilon_v=-1}(-1,-1)_v$, a finite product. It contains the
  subgroup $Z(G(\Q))\times\mu_2$ of order $4$, a direct product, embedded
  by~(a). It is not the direct product $Z(G(\A))\times\mu_2$ when $k$ has
  a real place.
\item $\sigma_\A(a(t_1),a(t_2))=\prod_v(t_{1,v},t_{2,v})_v$ for
  $t_1,t_2\in\A_k^\times$, and the preimage of every abelian subgroup of
  $G(\A)$, in particular $\overline{A(\A)}$, is abelian.
\end{enumerate}
\end{lemma}

\begin{proof}
(a) For $\gamma_1,\gamma_2\in SL(2,k)$, $\prod_v\beta_v(\gamma_1,\gamma_2)
=1$ by the product formula~\ref{H4}, because $\beta_v(\gamma_1,\gamma_2)$
is the local Hilbert symbol of two elements of $k^\times$. Hence
$\sigma_\A(\gamma_1,\gamma_2)=\kappa_\A(\gamma_1)\kappa_\A(\gamma_2)
\kappa_\A(\gamma_1\gamma_2)$, which is the condition for
$\gamma\mapsto(\gamma,\kappa_\A(\gamma))$ to be a homomorphism.

(b) On $N$ every $\beta_v$ and $\kappa_v$ is $1$. On $B(\Q)$ every
$\kappa_v$ is $1$ because $c=0$.

(c) It suffices that $\sigma_v\equiv1$ on the $v$-component of
$K_f\times K_f$ for each finite $v$. At $v\notin S$ this is
Lemma~\ref{lem:local}(b). At finite $v\in S$, $\kappa_v\equiv1$ on
$K_{m_v,v}$ (non-dyadic: $c\in\varpi\OO$ and $d$ is a square; dyadic: by
definition), so $\sigma_v=\beta_v\equiv1$ there by
Lemma~\ref{lem:local}(c).

(d) Lemma~\ref{lem:local}(g) applies to $\sigma_v$, since
$\kappa_v(g_1)\kappa_v(g_2)\kappa_v(g_1g_2)$ is symmetric in commuting
$g_1,g_2$; so preimages of abelian subgroups are abelian, and
$\overline{Z(G(\A))}$ is central. Its image under $p$ is central, so
equality holds. Since $\kappa_v(\pm1)=1$, the square of a lift of
$\varepsilon$ is $\prod_v\sigma_v(\varepsilon_v,\varepsilon_v)=
\prod_{\varepsilon_v=-1}(-1,-1)_v$, which is finite by~\ref{H3}. For
$\varepsilon=-1\in Z(G(\Q))$ it is $\prod_v(-1,-1)_v=1$ by~\ref{H4}. If
$v_0$ is real, the lift of the element that is $-1$ at $v_0$ and $1$
elsewhere has order $4$.

(e) As in Lemma~\ref{lem:local}(f), since $\kappa_v\equiv1$ on $A$.
\end{proof}

From now on we identify $G(\Q)$ with its image under
Lemma~\ref{lem:global}(a), and $N(\A)$ with $N(\A)\times1$. By~(b), the
two identifications agree on $N(\Q)$. The splitting over $B(\Q)$ by
$b\mapsto(b,1)$ is what the constant terms, Eisenstein series and the
truncation operator of Sections~4 and~5 use.

\begin{example}\label{ex:jacobi}
Let $k=\Q$ and $\gamma=\begin{pmatrix}a&b\\c&d\end{pmatrix}\in SL(2,\Z)$
with $c\ne0$. Then $\gamma\in K_v$ for every $v$, and for odd
$p\mid c$, $(c,d)_p=\bigl(\tfrac dp\bigr)^{\ord_p(c)}$ by~\ref{H3}. So
\[
\kappa_\A(\gamma)=\Bigl(\frac{d}{c'}\Bigr),\qquad
c'=\text{the odd part of }|c|,
\]
a Jacobi symbol: the factor behind the theta multiplier on
$\Gamma_0(4)$~\cite{K2,Ge2}. The script checks this on random elements
of $SL(2,\Z)$.
\end{example}

\begin{remarkrev}\label{rem:old-2.9}
The 1995 notes define the splitting as $\gamma\mapsto(\gamma,
s_\A(\gamma)^{-1})$, with $s_v(\gamma)^{-1}=(c,d)_v$ when $v$ is finite,
$cd\ne0$ and $n\nmid\ord_v(c)$, and with the cocycle $\prod_v\beta_v$.
That cocycle does not converge, and the formula belongs to a different
normalization. In the normalization~\eqref{eq:sigma-A},
$\kappa_\A$ of~\eqref{eq:kappa-A} replaces $s_\A^{-1}$. For $n=2$ the two
agree at every non-dyadic $v$ with $\gamma\in K_v$: if $\ord_v(c)>0$ is
even, or if $\ord_v(c)=0$, both $(c,d)_v$ and $\kappa_v(\gamma)$ are $1$.
They differ only at the finitely many places where $\gamma\notin K_v$,
and at dyadic places.
\end{remarkrev}

Let $K_0=\prod_vK_v$ be the standard maximal compact subgroup of
$G(\A)$, with $K_v$ as in Section~\ref{sec:local}; $K_f$ and $K_\infty$
denote its finite and archimedean parts.

\begin{lemma}[Strong approximation; compare~{\cite[Prop.~5]{K2}},
\cite{Ge2}]\label{lem:strong}
Let $K_f$ and $K_f^*$ be as in Lemma~\ref{lem:global}(c), and put
$\Gamma=G(\Q)\cap K_f$, a congruence subgroup of $SL(2,\OO_k)$, viewed in
$G(\R)=SL(2,k\otimes\R)$. For $\gamma\in\Gamma$ let $\tilde k_\gamma\in K_f^*$ be the lift of
$\gamma_f$. Then $\tilde\gamma=\gamma\,\tilde k_\gamma^{-1}$ lies in
$\overline{G(\R)}$, the map $\gamma\mapsto\tilde\gamma$ is an injective
homomorphism, its image $\widetilde\Gamma$ is discrete of finite
covolume, and
\[
G(\Q)\backslash\overline{G(\A)}/K_f^*\;\cong\;
\widetilde\Gamma\backslash\overline{G(\R)}.
\]
\end{lemma}

\begin{proof}
Since $SL(2)$ is simply connected and $G(\R)$ is not compact, strong
approximation gives $G(\A)=G(\Q)\,G(\R)\,K_f$. Lifting to the cover,
with $\mu_2\subset\overline{G(\R)}$, gives $\overline{G(\A)}=
G(\Q)\,\overline{G(\R)}\,K_f^*$, so the double coset space has a single
component. The stabilizer computation is the usual one: $\gamma\in
G(\Q)$ lies in $\overline{G(\R)}K_f^*$ exactly when $\gamma_f\in K_f$,
and then $\gamma=\tilde\gamma\,\tilde k_\gamma$. The map
$\gamma\mapsto\tilde\gamma$ is a homomorphism because lifts of elements
of $G(\A)$ with disjoint supports commute: $\sigma_v(1,h)=\sigma_v(h,1)=1$.
\end{proof}

% [Claude] The 1995 lemma uses K = prod_{p<=infty} K_p, which is not open,
% and a disjoint union over i; for simply connected G there is one term.

\subsubsection{\texorpdfstring{$NAK$}{NAK} decompositions}

% [Claude] Part (a) follows the corrected global lemma of book1.tex.
\begin{lemma}\label{lem:NAK-global}
\begin{enumerate}[label=(\alph*)]
\item $\overline{G(\A)}=N(\A)\cdot\overline{A(\A)}\cdot\overline{K_0}$. In
  a decomposition $(n,1)\,\tilde a\,\tilde k$, $p(\tilde a)$ is determined
  modulo $A(\A)\cap K_0$, and
  \[
  H(g,\zeta)=\log|t|_{\A_k}=\sum_v\log|t_v|_v,\qquad p(\tilde a)=a(t),
  \]
  is well defined.
\item $H$ satisfies Lemma~\ref{lem:H}(a)--(c) globally, and
  $H(\gamma x)=H(x)$ for $\gamma\in B(\Q)$ and $x\in\overline{G(\A)}$.
\end{enumerate}
\end{lemma}

\begin{proof}
(a) As for Lemma~\ref{lem:NAK-local}, place by place. (b) For
$\gamma=n'a(s)\in B(\Q)$, identified with $(\gamma,1)$ by
Lemma~\ref{lem:global}(b), and $x=(n,1)\tilde a\tilde k$,
$\gamma x=(n'',1)\,(a(s),1)\tilde a\,\tilde k$ with $n''\in N(\A)$, so
$H(\gamma x)=\log|s|_{\A_k}+H(x)=H(x)$ by the product formula.
\end{proof}

\subsubsection{Heights}

For $g=(g_v)\in G(\A)$ put $\|g\|=\prod_v\|g_v\|_v$ and
$\|(g,\zeta)\|=\|g\|$.

\begin{lemma}\label{lem:heights}
For $x,y\in\overline{G(\A)}$ and $\tilde k,\tilde k'\in\overline{K_0}$:
\begin{enumerate*}[label=(\alph*)]
\item $\|x\|\ge1$;
\item $\|\tilde k\,x\,\tilde k'\|=\|x\|$;
\item $\|xy\|\le\|x\|\,\|y\|$;
\item $\|x^{-1}\|=\|x\|$.
\end{enumerate*}
\end{lemma}

\begin{proof}
Each holds for $G(\A)$ place by place, since $\det g_v=1$ forces
$\|g_v\|_v\ge1$. For~(d), $g^{-1}$ has the same entries as $g$ up to
order and sign. The cover does not change $\|\cdot\|$.
\end{proof}

\subsubsection{Algebras of functions}

A function on $\overline{G(\A)}$ is \emph{anti-genuine} if
$f(g,\zeta)=\zeta^{-1}f(g,1)$. Let $\mathcal H(\overline{G(\A)})$ be the
span of the functions $f=\bigotimes_vf_v$, $f_v\in
C_c^\infty(\overline{SL(2,k_v)})_-$, with $f_v$ equal to the unit element
$e_v$ of~\eqref{eq:unit} (defined with $K_v^*$) for almost all $v$. Such
a product is a function on the restricted product of the local covers
which is invariant under $\{(\zeta_v)\mid\prod\zeta_v=1\}$, because each
$f_v$ is anti-genuine; so it descends to $\overline{G(\A)}$ and is
anti-genuine there.

\subsubsection{Measures}

Let $dk=\prod_vdk_v$ on $K_0$. Let $dg=\prod_vdg_v$ be the Tamagawa
measure on $G(\A)$~\cite{W2}, with $dg_v$ built from the self-dual
measures of a fixed nontrivial character $\psi=\prod\psi_v$ of
$\A_k/k$; then $\mathrm{vol}(G(\Q)\backslash G(\A))=1$. For a maximal
torus $T$ of $G$ over $\Q$ let $dt$ be the Tamagawa measure on $T(\A)$,
written $da$ for $T=A$, and let $dn$ be the product measure on $N(\A)$,
for which $\mathrm{vol}(N(\Q)\backslash N(\A))=1$. On the cover put
$d\tilde g=dg\,d\zeta$ with $\mathrm{vol}(\mu_2)=1$, so
$\mathrm{vol}(G(\Q)\backslash\overline{G(\A)})=1$.

\subsection{The double cover of \texorpdfstring{$GL(2)$}{GL(2)}}\label{sec:gl2}

Section~7 compares $SL(2)$-conjugacy with $GL(2)$-conjugacy on the cover,
so we need a double cover of $GL(2)$ which restricts to the one above.
There are two natural choices for $n=2$ (Remark~\ref{rem:gl2-twist}); we
fix one, and compute the conjugation action of $GL(2,F)$ on
$\overline{SL(2,F)}$ for later use.

For $g_1,g_2\in GL(2,F)$, with $x(g)$ defined as before, put
\begin{equation}\label{eq:sigma-gl2}
\sigma(g_1,g_2)=\Bigl(\frac{x(g_1g_2)}{x(g_1)},\
\frac{x(g_1g_2)}{x(g_2)\det g_1}\Bigr)_v,
\end{equation}
the $n=2$ case of Kubota's cocycle on $GL(2)$~\cite{K2,Ge2}. Let
$\overline{GL(2,F)}$ be the corresponding extension by $\mu_2$.

\begin{lemma}\label{lem:gl2}
\begin{enumerate}[label=(\alph*)]
\item $\sigma$ is a $2$-cocycle on $GL(2,F)$ whose restriction to
  $SL(2,F)$ is $\beta$. Hence $p^{-1}(SL(2,F))\subset\overline{GL(2,F)}$
  is the cover $\overline{SL(2,F)}$ of Section~\ref{sec:local}.
\item For $v$ finite non-dyadic, $k\mapsto(k,\kappa'(k))$ is a
  homomorphism on $GL(2,\OO)$, where $\kappa'(k)=(c,\,d/\det k)_v$ if
  $0<|c|_v<1$ and $\kappa'(k)=1$ otherwise. On $SL(2,\OO)$,
  $\kappa'=\kappa$.
\item Conjugation by any lift of $h\in GL(2,F)$ defines an automorphism
  $\mathrm{Ad}(h)$ of $\overline{SL(2,F)}$ lifting conjugation by $h$, and
  $h\mapsto\mathrm{Ad}(h)$ is an action of $GL(2,F)$. If $h$ commutes with
  $\gamma\in SL(2,F)$, then
  \begin{equation}\label{eq:gl2-commutator}
  \mathrm{Ad}(h)(\gamma,\zeta)
  =\Bigl(\gamma,\ \Bigl(\frac{x(h\gamma)}{x(h)},\ \det h\Bigr)_v\,\zeta\Bigr).
  \end{equation}
  In particular $\mathrm{Ad}(\diag(y,1))(a(t),\zeta)=(a(t),(y,t)_v\zeta)$,
  so the preimage of the diagonal torus of $GL(2,F)$ is not abelian
  unless $F=\C$. Likewise $E^\times\supset T_\theta$ (embedded as
  $\{\begin{pmatrix}a&b\theta\\b&a\end{pmatrix}\}$) acts on
  $\overline{T_\theta}$ through~\eqref{eq:gl2-commutator}.
\end{enumerate}
\end{lemma}

\begin{proof}
(a) On $SL(2)$, $\det g_1=1$ and~\eqref{eq:sigma-gl2}
becomes~\eqref{eq:beta}. The cocycle identity is due to
Kubota~\cite{K2}; we have also checked it numerically.

(b) We have checked this numerically for $F=\Q_3,\Q_5,\Q_7$; see
Remark~\ref{rem:verified}.

(c) Since $\mu_2$ is central, conjugation by $(h,\zeta)$ does not depend
on $\zeta$, and it preserves $p^{-1}(SL(2,F))$. If $h\gamma=\gamma h$,
then $\mathrm{Ad}(h)(\gamma,\zeta)=(\gamma,\zeta\,\sigma(h,\gamma)
\sigma(\gamma,h))$. With $x_3=x(h\gamma)$, bimultiplicativity gives
$\sigma(h,\gamma)=(x_3/x(h),x_3/x(\gamma))\,(x_3/x(h),\det h)$, and
$\sigma(\gamma,h)=(x_3/x(\gamma),x_3/x(h))$ since $\det\gamma=1$. By
symmetry~\ref{H1} the first factors cancel. For $h=\diag(y,1)$ and
$\gamma=a(t)$: $x(h\gamma)=t^{-1}$ and $x(h)=1$, giving $(t^{-1},y)=
(y,t)$. This is nontrivial for suitable $y,t$ unless $F=\C$.
\end{proof}

The consequences for Section~7 (the functions $\delta\mapsto
I_{\mathcal O}(f^\delta)$ are defined on $D(T/F)$ only up to the
character~\eqref{eq:gl2-commutator}, and the choice of $D(A/F)$) are
taken up there.

\begin{remarkrev}\label{rem:gl2-twist}
The cocycle $\sigma'(g_1,g_2)=\sigma(g_1,g_2)\,(\det g_1,\det g_2)_v$ is
another double cover of $GL(2,F)$ with the same restriction to $SL(2,F)$;
these are the two Kazhdan--Patterson covers for $n=2$~\cite{KP}. We fix
$\sigma$. Whether Section~7 needs $\sigma'$ instead has not been checked
at this stage. The commutator formula~\eqref{eq:gl2-commutator} is the
same for $\sigma'$: the extra factor $(\det h,\det\gamma)(\det\gamma,\det
h)$ is $1$ by~\ref{H1}. The
adelic cover of $GL(2,\A_k)$ is built as in Section~\ref{sec:global},
with $\kappa'_v$ in place of $\kappa_v$; Hilbert reciprocity for
$\prod_v\sigma_v$ on $GL(2,k)$ is checked numerically for $k=\Q$.
\end{remarkrev}

\begin{remarkrev}\label{rem:general-n}
For general $n$ (with $\mu_n\subset k$) the following change. The
conjugation action of $\overline{A(F)}$ on itself is through
$(t_1,t_2)_n^2$, so $\overline{A(F)}$ is abelian only for $n\le2$; the
argument of Lemma~\ref{lem:local}(g) uses the symmetry of the quadratic
symbol and does not extend. Lemma~\ref{lem:local}(c) holds with
$1+\varpi^m\OO\subset F^{\times n}$ (the same proof). Genuine
($\zeta$) and anti-genuine ($\zeta^{-1}$) differ, and with the 1995
convention $R(f)$ vanishes on $L^2_{\mathrm{gen}}$ unless $n\mid2$.
\end{remarkrev}

\begin{remarkrev}\label{rem:verified}
The script \texttt{checks/sec2\_n2\_check.py} verifies, in exact rational
arithmetic at $v=\infty,2,3,5,7$: the cocycle identity and the two forms
of~\eqref{eq:beta}; Lemma~\ref{lem:local}(a)--(c), (e)--(h);
Remark~\ref{rem:dyadic}; the product formula and the splitting of
Lemma~\ref{lem:global}(a),~(b),~(d); Example~\ref{ex:jacobi}; and
Lemma~\ref{lem:gl2}(a)--(c) with Remark~\ref{rem:gl2-twist}. The checks
are on random elements of $SL(2,\Q)$, $GL(2,\Q)$, $SL(2,\Z_{(p)})$ and
the tori $T_\theta(\Q)$; they are evidence, not proofs. Of the
statements in this section, Lemma~\ref{lem:local}(b) and
Lemma~\ref{lem:gl2}(a),~(b) are supported here only by citation and these
checks.
\end{remarkrev}

% SPDX-License-Identifier: CC-BY-4.0
% =====================================================================
% Section 3 of D. Joyner, "A hitch-hiker's guide to the metaplectic trace
% formula for SL(2)", revised for the 2-fold cover (n = 2), November 2026.
%
% Stage 2 of ../revision_log.md. Source: ../original-1995/cctrace2.tex
% (= ../book-1997/book7b.tex). For n = 2 the cover of a torus is abelian, so
% the Clifford-Mackey part of the 1995 section (Lemmas 3.3, 3.5, Remarks
% 3.4, 3.6, 3.7, 3.10) is not needed; the general-n situation is summarized
% in Remark 2026.6. The map from 1995 numbers to the numbers here is in the
% "Stage 2 notes" of ../revision_log.md. Explicit formulas are checked by
% checks/sec3_n2_check.py.
% Revised after the second referee report (ChatGPT, November 2026): test
% functions are defined on the topological cover, and the proof of
% Lemma 3.1(c) no longer claims that a continuous character has open kernel.
% Comments beginning "% [Claude]" are editorial notes by Claude.
% =====================================================================

\section{The rank 0 case}\label{sec:rank0}

Before the trace formula for $\overline{G(\A)}$ we treat the covers of
the tori of $G$, where it is elementary. For $n=2$ these covers are
abelian (Lemma~\ref{lem:global}(e)), but in general not split, so a
genuine automorphic character cannot be written as an ordinary one
times a sign (Example~\ref{ex:order4}). The trace formula is then
Poisson summation twisted by one fixed genuine automorphic character
(Theorem~\ref{thm:tf-rank0}, Lemma~\ref{lem:transfer}), and it gives a
bijection between genuine automorphic characters and Hecke characters
(Theorem~\ref{thm:corr}). For the split torus, Weil's character
$\chi_\psi$ (Lemma~\ref{lem:weil}) makes the bijection canonical once
$\psi$ is fixed. We also define the distributions $I^A_{\mathcal O}$ and
$I^A_{\bar\chi}$ used in Section~6. We work with an arbitrary maximal
torus because Section~7 applies the same formula to anisotropic tori.

\subsection{The cover of a torus}\label{sec:torus-cover}

Let $T$ be a maximal torus of $G$ defined over $\Q$. Then
$T=\Res_{k/\Q}T_k$ for a maximal torus $T_k$ of $SL(2)$ over $k$, which
is either split, and then $GL(2,k)$-conjugate to $A$, or of the form
$T_\theta$ of Section~\ref{sec:local} with $E=k(\sqrt\theta)$ a quadratic
extension of $k$. Put
\[
T(\A)^1=\{t\in T(\A)\mid |\chi(t)|=1\ \text{for all}\ \chi\in X(T)_\Q\}.
\]
For $T=A$ this is $A(\A)^1=\{a(t)\mid |t|_{\A_k}=1\}$, the kernel of
$H$ of Lemma~\ref{lem:NAK-global}; for $T_k=T_\theta$ the torus $T$ is
anisotropic over $\Q$ and $T(\A)^1=T(\A)$. In both cases
$T(\Q)\subset T(\A)^1$ and $T(\Q)\backslash T(\A)^1$ is
compact~\cite{PR}.

Let $\overline{T(\A)}$ and $\overline{T(\A)^1}$ be the preimages in
$\overline{G(\A)}$, and $\overline{T(k_v)}$ the preimage of $T_k(k_v)$
in $\overline{SL(2,k_v)}$. By Lemmas~\ref{lem:local}(g)
and~\ref{lem:global}(e) these groups are abelian. On $A$ the cocycle is
\begin{equation}\label{eq:cocycle-A}
\sigma_\A(a(t_1),a(t_2))=\prod_v(t_{1,v},t_{2,v})_v,\qquad
t_1,t_2\in\A_k^\times,
\end{equation}
since $\kappa_v\equiv1$ on $A$. Exactly as for $G$ in
Section~\ref{sec:global}, $\overline{T(\A)}$ is the \emph{quotient} of the
restricted product of the $\overline{T(k_v)}$, taken with respect to the
subgroups $T(\OO_v)^*=\{(t,\kappa_v(t))\mid t\in T_k(k_v)\cap K_v\}$, by
$\{(\zeta_v)\mid\prod_v\zeta_v=1\}$; it is not a subgroup of that
restricted product. Each $\overline{T(k_v)}$ maps injectively into
$\overline{T(\A)}$. These are closed subgroups of the covers of
Section~\ref{sec:background}, and they carry the induced topologies; for
$\overline{T(\A)}$ this is again the quotient topology of the restricted
product. We regard $T(\Q)$ as a subgroup of $\overline{T(\A)^1}$
by Lemma~\ref{lem:global}(a), $\gamma\mapsto(\gamma,\kappa_\A(\gamma))$;
for $T=A$ this is $a(t)\mapsto(a(t),1)$, $t\in k^\times$.

\paragraph{Conventions.}
As in Section~\ref{sec:functions}, a representation $\pi$ of one of these
covers is \emph{genuine} if $\pi(1,\zeta)=\zeta$, and a function $f$ is
\emph{anti-genuine} if $f(x\zeta)=\zeta^{-1}f(x)$ for $\zeta\in\mu_2$.
For $n=2$ the two conditions on functions coincide. Let
$C_c^\infty(\overline{T(k_v)})_-$ and $C_c^\infty(\overline{T(\A)^1})_-$
be the spaces of anti-genuine functions that are smooth and compactly
supported on these locally compact groups: locally constant in the
non-archimedean variables and $C^\infty$ in the archimedean ones. As in
Section~\ref{sec:functions}, this is a condition on the topological
group, not on the function $t\mapsto f(t,1)$. For the split torus the two
agree, because $a\mapsto(a,1)$ is continuous on $A(k_v)$ and the
cocycle~\eqref{eq:cocycle-A} is locally constant
(Lemma~\ref{lem:local}(i)); for an anisotropic torus they do not, for
example for $SO(2)\subset SL(2,\R)$, whose preimage is connected
(Lemma~\ref{lem:local}(d)). At places where $T(\OO_v)^*$ is defined
let $e_{T,v}$ be the anti-genuine function equal to $1$ on
$T(\OO_v)^*$ and supported on $\overline{T(\OO_v)}$. Let
$\mathcal H(\overline{T(\A)})$ be the span of the products
$\bigotimes_vf_v$, $f_v\in C_c^\infty(\overline{T(k_v)})_-$, with
$f_v=e_{T,v}$ for almost all $v$; these descend to $\overline{T(\A)}$ as
in Section~\ref{sec:global}. Measures on $T(\A)$ are as in
Section~\ref{sec:global}, and on covers $\mu_2$ has total volume $1$.

For a locally compact group $H$ let $\Pi(H)$ be the set of
(equivalence classes of) irreducible admissible genuine representations
of $H$, and $\Pi_u(H)$ the unitary ones.

\begin{lemma}\label{lem:genuine-chars}
Let $F=k_v$.
\begin{enumerate}[label=(\alph*)]
\item Every element of $\Pi(\overline{T(F)})$, $\Pi(\overline{T(\A)})$ or
  $\Pi(\overline{T(\A)^1})$ is a genuine character.
\item Genuine characters of $\overline{T(F)}$ exist, and any two differ
  by a character of $T(F)$: they form a torsor under $T(F)^\wedge$.
\item Every genuine character $\bar\chi$ of $\overline{T(\A)}$ factors
  uniquely as $\bar\chi=\prod_v\bar\chi_v$, where $\bar\chi_v$ is a
  genuine character of $\overline{T(k_v)}$, trivial on $T(\OO_v)^*$ for
  almost all $v$. Conversely, every such family defines a genuine
  character of $\overline{T(\A)}$. The same holds for unitary characters.
\end{enumerate}
\end{lemma}

\begin{proof}
(a) The groups are abelian, so Schur's lemma applies; for unitary
representations of locally compact abelian groups this is standard, and
for admissible representations Schur's lemma holds because they have
countable dimension. Since $\mu_2$ is central, it acts by a character,
and that character is the identity because the representation is
genuine.

(b) $\mu_2$ is a closed subgroup of the locally compact abelian group
$\overline{T(F)}$, so the character $\zeta\mapsto\zeta$ of $\mu_2$
extends~\cite{Ro}. The quotient of two genuine characters is trivial on
$\mu_2$, hence is the pull-back of a character of $T(F)$.

(c) Put $\bar\chi_v=\bar\chi\circ\iota_v$, where $\iota_v$ is the
inclusion of $\overline{T(k_v)}$. Since $\bar\chi$ is continuous, there
is a neighborhood $W$ of the identity in $\overline{T(\A)}$ with
$|\bar\chi(x)-1|<1$ for $x\in W$. By the definition of the topology, $W$
contains the image of $\prod_{v\notin S}T(\OO_v)^*$ for some finite set
$S$ of places. That image is a subgroup, so its image under $\bar\chi$ is
a subgroup of $\C^\times$ contained in the disc $\{|z-1|<1\}$, and the
only such subgroup is $\{1\}$. Hence $\bar\chi_v$ is trivial on
$T(\OO_v)^*$ for $v\notin S$. (The kernel of $\bar\chi$ need not be open:
at an archimedean place $\bar\chi_v$ can be nontrivial on every
neighborhood of the identity.) Every $x\in\overline{T(\A)}$ is the image of
some $(x_v)$ with $x_v\in T(\OO_v)^*$ for $v\notin S'$, and then
$\bar\chi(x)=\prod_{v\in S'}\bar\chi_v(x_v)$. Conversely, $\prod_v
\bar\chi_v$ is a character of the restricted product, and it is trivial
on $\{(\zeta_v)\mid\prod\zeta_v=1\}$ because each $\bar\chi_v$ is
genuine. So it descends to $\overline{T(\A)}$.
\end{proof}

Let $X_T=T(\Q)\backslash\overline{T(\A)^1}$. It is a compact abelian
group, an extension of the compact group $T(\Q)\backslash T(\A)^1$ by
$\mu_2$; $\mu_2$ injects into $X_T$ because $T(\Q)$ meets $\mu_2$
trivially in $\overline{G(\A)}$. Let $\Pi_{\mathrm{aut}}(X_T)$ be the set
of genuine characters of $\overline{T(\A)^1}$ that are trivial on
$T(\Q)$, i.e.\ the genuine characters of $X_T$.

\begin{theorem}\label{thm:corr}
\begin{enumerate}[label=(\alph*)]
\item $L^2(X_T)_{\mathrm{gen}}=\widehat\bigoplus_{\bar\chi\in
  \Pi_{\mathrm{aut}}(X_T)}\C\,\bar\chi$, each character occurring with
  multiplicity one.
\item $\Pi_{\mathrm{aut}}(X_T)$ is not empty. For any
  $\bar\chi_0\in\Pi_{\mathrm{aut}}(X_T)$, the map
  $\chi\mapsto\bar\chi_0\cdot(\chi\circ p)$ is a bijection from the
  characters of $T(\Q)\backslash T(\A)^1$ onto $\Pi_{\mathrm{aut}}(X_T)$.
  Replacing $\bar\chi_0$ by another element changes this bijection by a
  translation of $(T(\Q)\backslash T(\A)^1)^\wedge$.
\item For $T=A$, the character $\chi_\psi$ of Lemma~\ref{lem:weil} lies
  in $\Pi_{\mathrm{aut}}(X_A)$; it depends only on the fixed additive
  character $\psi$ of $\A_k/k$.
\end{enumerate}
\end{theorem}

\begin{proof}
(a) Peter--Weyl for the compact abelian group $X_T$: $L^2(X_T)$ is the
Hilbert sum of its characters, each once, and the genuine subspace is
spanned by the genuine ones. (b) Apply Lemma~\ref{lem:genuine-chars}(b)
to the closed subgroup $\mu_2$ of $X_T$; the second sentence follows
as in that lemma. (c) is Lemma~\ref{lem:weil}(c).
\end{proof}

The bijection in~(b) depends on the choice of $\bar\chi_0$, and for a
general torus there is no canonical one. The choice cannot be avoided:

\begin{example}\label{ex:order4}
Let $k=\Q$, and let $\varepsilon\in A(\A)^1$ be $a(-1)$ at the real place
and $1$ at every finite place. By Lemma~\ref{lem:local}(e) any lift
$\tilde\varepsilon$ satisfies $\tilde\varepsilon^2=(1,(-1,-1)_\infty)
=(1,-1)$. So every genuine character $\bar\chi$ of $\overline{A(\A)^1}$
has $\bar\chi(\tilde\varepsilon)=\pm i$, and no genuine character is of
the form $(\chi\circ p)\cdot\mathrm{sgn}$ for a splitting $A(\A)^1\to
\overline{A(\A)^1}$; indeed no such splitting exists.
\end{example}

For the split torus there is a canonical choice, given by Weil
indices~\cite{We64}. Let $\psi=\prod_v\psi_v$ be the nontrivial
character of $\A_k/k$ fixed in Section~\ref{sec:global}. For $a\in
k_v^\times$ let $\gamma_v(a)\in\mu_8$ be the Weil index of the character
$x\mapsto\psi_v(ax^2)$ of $k_v$, normalized by
\begin{equation}\label{eq:weil-index}
\int_{k_v}\psi_v(ax^2)\,dx=\gamma_v(a)\,|2a|_v^{-1/2}
\end{equation}
(self-dual measure; for $v$ finite the integral is the limit over the
lattices $\varpi^{-M}\OO_v$, $M\to\infty$).

\begin{lemma}[Weil~\cite{We64}]\label{lem:weil}
\begin{enumerate}[label=(\alph*)]
\item $\gamma_v(a)$ depends only on the class of $a$ in
  $k_v^\times/k_v^{\times2}$, and
  $\gamma_v(ab)\gamma_v(1)=\gamma_v(a)\gamma_v(b)(a,b)_v$.
\item If $v$ is finite and non-dyadic and $\psi_v$ is unramified, then
  $\gamma_v(u)=\gamma_v(1)=1$ for $u\in\OO_v^\times$. For $a\in k^\times$,
  $\prod_v\gamma_v(a)=1$.
\item Hence $\chi_{\psi,v}(a(t),\zeta)=\zeta\,\gamma_v(t)/\gamma_v(1)$ is
  a genuine character of $\overline{A(k_v)}$, trivial on $A(\OO_v)^*=
  A(\OO_v)\times\{1\}$ for almost all $v$, and
  $\chi_\psi=\prod_v\chi_{\psi,v}$ is a genuine character of
  $\overline{A(\A)}$ trivial on $A(\Q)$. Its restriction to
  $\overline{A(\A)^1}$ lies in $\Pi_{\mathrm{aut}}(X_A)$.
\item For $b\in k^\times$ and $\psi_b(x)=\psi(bx)$,
  $\chi_{\psi_b}=\chi_\psi\cdot\bigl(a(t)\mapsto\prod_v(b,t_v)_v\bigr)$.
\end{enumerate}
\end{lemma}

\begin{proof}
(a) and (b) are due to Weil~\cite[\S\S14, 28]{We64}. (c) By~(a),
$\chi_{\psi,v}(a(t_1),1)\chi_{\psi,v}(a(t_2),1)=(t_1,t_2)_v\,
\chi_{\psi,v}(a(t_1t_2),1)=\chi_{\psi,v}((a(t_1),1)(a(t_2),1))$, using the
group law~\eqref{eq:cocycle-A}. Unramifiedness is~(b), and triviality on
$(a(t),1)$, $t\in k^\times$, is the product formula in~(b) applied to
$\gamma(t)$ and to $\gamma(1)$. (d) $\psi_{b,v}(ax^2)=\psi_v(abx^2)$, so
$\gamma_{\psi_b,v}(t)/\gamma_{\psi_b,v}(1)=\gamma_v(bt)/\gamma_v(b)=
(\gamma_v(t)/\gamma_v(1))\,(b,t)_v$ by~(a).
\end{proof}

\begin{example}\label{ex:weil-Q}
Let $k=\Q$, $\psi_p(x)=e^{2\pi i\{x\}_p}$ and $\psi_\infty(x)=
e^{-2\pi ix}$, so $\psi$ is trivial on $\Q$. Then
$\gamma_\infty(a)=e^{-\pi i\,\mathrm{sgn}(a)/4}$, and for odd $p$,
$\gamma_p(u)=1$ and $\gamma_p(pu)=\bigl(\tfrac up\bigr)\epsilon_p$ for
$u\in\Z_p^\times$, where $\epsilon_p=1$ if $p\equiv1\pmod4$ and
$\epsilon_p=i$ if $p\equiv3\pmod4$. So
$\chi_{\psi,\infty}(a(-1),1)=i$, in accordance with
Example~\ref{ex:order4}, and $\chi_{\psi,p}(a(p),1)=\epsilon_p$ for odd
$p$. The dyadic values are tabulated by the check script, which computes
$\gamma_v$ from~\eqref{eq:weil-index} as a finite Gauss sum and verifies
Lemma~\ref{lem:weil}(a)--(c) at $v=\infty,2,3,5,7$ and the product
formula on random $a\in\Q^\times$.
\end{example}

\subsection{The trace formula}\label{sec:tf-rank0}

Let $R$ be the right regular representation of $\overline{T(\A)^1}$ on
$L^2(X_T)_{\mathrm{gen}}$, and for $f\in C_c^\infty(\overline{T(\A)^1})_-$
put $R(f)=\int_{\overline{T(\A)^1}}f(y)R(y)\,dy$ and
\[
\bar\chi(f)=\int_{\overline{T(\A)^1}}f(y)\bar\chi(y)\,dy,\qquad
\bar\chi\in\Pi(\overline{T(\A)^1}).
\]
The integrand $f\bar\chi$ is invariant under $\mu_2$, so $\bar\chi(f)$ is
an integral over $T(\A)^1$; if $\bar\chi$ were trivial on $\mu_2$ instead,
the integral would vanish.

\begin{theorem}\label{thm:tf-rank0}
Let $f\in C_c^\infty(\overline{T(\A)^1})_-$.
\begin{enumerate}[label=(\alph*)]
\item $R(f)$ is the integral operator on $L^2(X_T)_{\mathrm{gen}}$ with
  kernel
  \[
  K_f(x,y)=\sum_{\delta\in T(\Q)}f(x^{-1}\delta y),
  \]
  and $R(f)\bar\chi=\bar\chi(f)\,\bar\chi$ for
  $\bar\chi\in\Pi_{\mathrm{aut}}(X_T)$.
\item $R(f)$ is of trace class, and
  \begin{equation}\label{eq:tf-rank0}
  \sum_{\bar\chi\in\Pi_{\mathrm{aut}}(X_T)}\bar\chi(f)
  =\operatorname{tr}R(f)=\int_{X_T}K_f(x,x)\,dx
  =\mathrm{vol}\bigl(T(\Q)\backslash T(\A)^1\bigr)\sum_{\delta\in T(\Q)}f(\delta),
  \end{equation}
  the sum on the left converging absolutely.
\end{enumerate}
\end{theorem}

\begin{proof}
(a) For $\varphi\in L^2(X_T)_{\mathrm{gen}}$,
\[
(R(f)\varphi)(x)=\int_{\overline{T(\A)^1}}f(y)\varphi(xy)\,dy
=\int_{\overline{T(\A)^1}}f(x^{-1}y)\varphi(y)\,dy
=\int_{X_T}\sum_{\delta\in T(\Q)}f(x^{-1}\delta y)\varphi(y)\,dy,
\]
because $\varphi(\delta y)=\varphi(y)$. For $\varphi=\bar\chi$ the first
expression is $\bar\chi(x)\bar\chi(f)$.

(b) The group is abelian, so $f(x^{-1}\delta x)=f(\delta)$, and
$\int_{X_T}K_f(x,x)\,dx=\mathrm{vol}(X_T)\sum_\delta f(\delta)$ with
$\mathrm{vol}(X_T)=\mathrm{vol}(T(\Q)\backslash T(\A)^1)$ since $\mu_2$
has volume $1$. The function $F(x)=\sum_{\delta}f(\delta x)$ is smooth on
$X_T$ (the sum is locally finite, $T(\Q)$ being discrete), so its Fourier
series converges absolutely; its Fourier coefficient at $\bar\chi$ is
$\int_{X_T}F\bar\chi=\bar\chi(f)$, which vanishes unless $\bar\chi$ is
genuine. Evaluating the Fourier series at $x=1$ gives the outer equality
of~\eqref{eq:tf-rank0}, i.e.\ Poisson summation. Since $R(f)$ is diagonal
in the orthonormal basis $\mathrm{vol}(X_T)^{-1/2}\bar\chi$ with summable
eigenvalues $\bar\chi(f)$, it is of trace class with that trace.
\end{proof}

Every conjugacy class in $\overline{T(\A)^1}$ is a point and every
centralizer is the whole group, so the geometric side has no orbital
integrals, only the values $f(\delta)$. Here $f(\delta)$ means $f$ at
the image $(\delta,\kappa_\A(\delta))$ of $\delta$; for $T=A$ it is
$f(a(\delta),1)$.

The trace formula for $\overline{T(\A)^1}$ is literally Poisson summation
on $T(\Q)\backslash T(\A)^1$, after multiplying by a genuine automorphic
character:

\begin{lemma}\label{lem:transfer}
Fix $\bar\chi_0\in\Pi_{\mathrm{aut}}(X_T)$ with local components
$\bar\chi_{0,v}$, as in Lemma~\ref{lem:genuine-chars}(c). For an
anti-genuine $\bar f$ on $\overline{T(\A)^1}$ (resp.\
$\overline{T(k_v)}$) define $f$ on $T(\A)^1$ (resp.\ $T_k(k_v)$) by
$f(t)=\bar f(\tilde t)\bar\chi_0(\tilde t)$ (resp.\ $\bar\chi_{0,v}$),
$\tilde t$ any lift of $t$.
\begin{enumerate}[label=(\alph*)]
\item $\bar f\mapsto f$ is well defined, bijective from
  $C_c^\infty(\overline{T(\cdot)})_-$ onto $C_c^\infty(T(\cdot))$, and
  carries convolution to convolution.
\item At every $v$ where $\bar\chi_{0,v}$ is trivial on $T(\OO_v)^*$
  (almost all $v$), it carries $e_{T,v}$ to the characteristic function of
  $T_k(k_v)\cap K_v$, which is the unit of the Hecke algebra of
  $T_k(k_v)$ when $T_k(k_v)\cap K_v$ has volume $1$.
\item $f(\delta)=\bar f(\delta)$ for $\delta\in T(\Q)$, and
  $(\bar\chi_0\cdot(\chi\circ p))(\bar f)=\hat f(\chi):=\int_{T(\A)^1}
  f(t)\chi(t)\,dt$ for every character $\chi$ of $T(\A)^1$.
\end{enumerate}
Consequently~\eqref{eq:tf-rank0} for $\bar f$ is the Poisson summation
formula for $f$,
\[
\sum_\chi\hat f(\chi)=\mathrm{vol}\bigl(T(\Q)\backslash T(\A)^1\bigr)
\sum_{\delta\in T(\Q)}f(\delta),
\]
with $\chi$ running over the characters of $T(\Q)\backslash T(\A)^1$.
\end{lemma}

\begin{proof}
(a) $\bar f\bar\chi_0$ is $\mu_2$-invariant because $\bar f$ is
anti-genuine and $\bar\chi_0$ genuine, and it is smooth because
$\bar\chi_0$ is a continuous character; so it is the pull-back of a
smooth function on $T$. The inverse map is
$f\mapsto(f\circ p)\bar\chi_0^{-1}$. For convolution,
$(\bar f_1*\bar f_2)(x)\bar\chi_0(x)=\int\bar f_1(y)\bar\chi_0(y)\,
\bar f_2(y^{-1}x)\bar\chi_0(y^{-1}x)\,dy$, and the integrand is
$\mu_2$-invariant in $y$, so the integral over the cover equals the
integral over $T$. (b) is immediate. (c) $\bar\chi_0$ is trivial on
$T(\Q)$; the second identity is the definition, the integrand being
$\mu_2$-invariant.
\end{proof}

Together with Theorem~\ref{thm:tf-rank0}, the lemma shows that the
correspondence $\chi\leftrightarrow\bar\chi_0\cdot(\chi\circ p)$ of
Theorem~\ref{thm:corr} matches the two trace formulas term by term.

\subsection{Distributions for Section 6}\label{sec:IA}

Let $S$ be a finite set of places of $\Q$ and let $\overline{A(\A_S)}$ be
the cover of $A(\A_S)=\prod_{p\in S}A(\Q_p)$ defined by
$\prod_{p\in S}\beta_p$. Let $\mathcal H(\overline{A(\A_S)})$ be the
algebra of functions $\bigotimes_{p\in S}f_p$, $f_p\in
C_c^\infty(\overline{A(\Q_p)})_-$, and the analogous algebra
$\mathcal H(\overline{A(\A)})$ of Section~\ref{sec:torus-cover}. For a
class $\mathcal O=\{\delta\}$ of $A(\Q)$ (a single element, since the
cover is abelian) and $\bar\chi\in\Pi_u(\overline{A(\A_S)})$, put
\begin{equation}\label{eq:IA}
I^A_{\mathcal O}(f)=f(a(\delta),1),\qquad
I^A_{\bar\chi}(f)=\bar\chi(f)=\int_{\overline{A(\A_S)}}f(y)\bar\chi(y)\,dy,
\end{equation}
for $f\in\mathcal H(\overline{A(\A_S)})$, and likewise over $\A$. Here
$(a(\delta),1)$ is the image of $\delta$ under the rational splitting
(Lemma~\ref{lem:global}(b)). Both are invariant distributions: the cover
is abelian, so conjugation acts trivially.

\begin{remarkrev}\label{rem:general-n-rank0}
For general $n$ the cover of $A$ is not abelian: by the computation in
Lemma~\ref{lem:local}(f), $\overline{A(F)}$ is a Heisenberg-type group
whose commutator pairing is $(t_1,t_2)_n^{-2}$. Its genuine irreducible
representations are then parametrized by genuine characters of its
center, the preimage of $A^N$ with $N=n/\gcd(n,2)$. Matching them with
characters of $A^N$ requires a splitting over $A^N$, and the global
argument requires a Hasse principle for $N$-th powers, which holds when
$\mu_N\subset k$. These are referee items II.1--II.4; none of them arises
for $n=2$, and this section does not treat $n>2$.
\end{remarkrev}

\subsection{Status of the results}\label{sec:status3}

\begin{remarkrev}\label{rem:sec3-status}
\emph{What rests on citation.} Taken from the literature:
Lemma~\ref{lem:weil}(a)--(b)~\cite{We64}; the compactness of
$T(\Q)\backslash T(\A)^1$~\cite{PR}; and the extension of characters from
closed subgroups of locally compact abelian groups~\cite{Ro}. Everything
else in this section is proved in the text.
\end{remarkrev}

\begin{remarkrev}\label{rem:verified-rank0}
The script \texttt{checks/sec3\_n2\_check.py} verifies: that the 1995
cocycle on $A$ equals~\eqref{eq:cocycle-A}; that $\gamma_v$, computed
from~\eqref{eq:weil-index}, is an eighth root of unity, depends only on
the square class, is $1$ on units at $p=3,5,7$, and satisfies
Lemma~\ref{lem:weil}(a); that $\chi_{\psi,v}$ respects the group law
(Lemma~\ref{lem:weil}(c)) at $v=\infty,2,3,5,7$; the product formula of
Lemma~\ref{lem:weil}(b) for random $a\in\Q^\times$; the odd-$p$ values of
Example~\ref{ex:weil-Q}; and Example~\ref{ex:order4}. The rest of the
section is proved in the text.
\end{remarkrev}

% SPDX-License-Identifier: CC-BY-4.0
% =====================================================================
% Section 4 of D. Joyner, "A hitch-hiker's guide to the metaplectic trace
% formula for SL(2)", revised for the 2-fold cover (n = 2), November 2026.
%
% Stage 3 of ../revision_log.md. Sources: ../original-1995/cctrace3.tex
% (= ../book-1997/book7c.tex, Sections 4.1-4.2) and the missing cctrac3b.tex
% (= ../book-1997/book7d.tex, Sections 4.3-4.5). Referee items II.5-II.8,
% II.22 (local part), A.1-A.7, the Siegel-domain part of I.10 and the
% Section 4 minor items are fixed here; see the "Stage 3 notes" in
% ../revision_log.md, which also map the 1995 numbers to the numbers here.
%
% This file is \input by the master; it compiles on its own through
% preview.tex. Explicit formulas are checked by checks/sec4_n2_check.py.
% Comments beginning "% [Claude]" are editorial notes by Claude.
% Revised after the second referee report (ChatGPT, November 2026): roadmap
% added, comparisons with the 1995 notes moved to Appendix A.
% =====================================================================

\providecommand{\GA}{\overline{G(\A)}}

\section{The geometric side}\label{sec:geometric}

This section computes the geometric side of the trace formula for
$\GA$, class by class. The general theory is due to
Arthur~\cite{A2,A3,A12} for connected reductive groups and to
Li~\cite{LiI} for covering groups; we specialize it to the double cover of
$G=\Res_{k/\Q}SL(2)$. For this group one inequality between heights
(Lemma~\ref{lem:height-ineq}) makes the reduction theory explicit, with
truncation parameter $T_0=0$. With it, every step of Arthur's coarse
expansion can be checked directly, and each class contributes a
polynomial in $T$ of degree at most one that we write down
(Theorem~\ref{thm:coarse}, Propositions~\ref{prop:hyperbolic},
\ref{prop:elliptic} and~\ref{prop:central}). The new parts are the
corrected weight $2T+v_A$ for hyperbolic classes, the unipotent
contribution of the central classes, given as the finite part of a zeta
integral, and the local distribution $J_A(z,f)$ at a central element
(Lemma~\ref{lem:JA-central}). Section~\ref{sec:fine} states the fine
expansion~\cite{A12,LiI} in the form that Sections~5--7 use.

\paragraph{Roadmap.} The argument has four steps.
\begin{enumerate}[label=\arabic*.]
\item \emph{Classes} (Section~\ref{sec:classes}). $G(\Q)$ is sorted into
  classes $\mathcal O$ by semisimple part: two central classes, the
  hyperbolic classes and the elliptic classes.
\item \emph{Truncation} (Section~\ref{sec:kernels}). The kernel
  $K_G(x,x)=\sum_\gamma f(x^{-1}\gamma x)$ of $R(f)$ is not integrable
  over $G(\Q)\backslash\GA$. The divergence comes from the cusp, the
  region where $H(\delta x)>T$ for some $\delta\in G(\Q)$. There one
  subtracts the kernel $K_B$ of the Borel subgroup; the result is the
  truncated kernel $k^T(x,f)$. The height inequality
  (Lemma~\ref{lem:height-ineq}) says that for $T\ge0$ at most one coset
  $B(\Q)\delta$ is involved at each $x$, so no further reduction theory is
  needed.
\item \emph{Integration, class by class} (Section~\ref{sec:coarse}). For
  an elliptic class nothing is truncated and $J_{\mathcal O}$ is an
  orbital integral. For a hyperbolic class the divergent integral over
  $A(\Q)\backslash A(\A)$ is cut down to an interval of length
  $2T+v_A(y)$, and $J^T_{\mathcal O}$ is a weighted orbital integral. For
  a central class the unipotent elements give a zeta integral, of which
  $J^T_{\mathcal O}$ keeps the finite part at $s=1$ and a term linear in
  $T$.
\item \emph{Localization} (Sections~\ref{sec:local-dist}
  and~\ref{sec:fine}). For $f=f_S\otimes e^S$ the terms are rewritten as
  local distributions on $\overline{G(\Q_S)}$ with global coefficients.
\end{enumerate}
The objects that occur are the following.
\begin{center}
\begin{tabular}{p{0.24\textwidth}p{0.68\textwidth}}
$K_P$, $K_{P,\mathcal O}$ & kernel of $R_P(f)$ and its part on the class
  $\mathcal O$, \eqref{eq:KP}\\
$J_{P,\mathcal O}$ & variant of $K_{P,\mathcal O}$ with the same
  integral, \eqref{eq:JP}\\
$\hat\tau_B(H(\delta x)-T)$ & cutoff: $1$ in the cusp above height $T$,
  $0$ elsewhere\\
$k^T$, $k^T_{\mathcal O}$, $j^T_{\mathcal O}$ & truncated kernels,
  \eqref{eq:kT}\\
$J^T_{\mathcal O}(f)$, $J_{\mathcal O}(f)$ & their integrals, polynomials
  in $T$, and the value at $T=0$\\
$v_A$, $v_{A,v}$ & the weight $-H(x)-H(wx)$ and its local form,
  \eqref{eq:vA}\\
$Z(f,s)$, $\mathrm{FP}_{s=1}$ & zeta integral of the unipotent terms and
  its finite part, \eqref{eq:Zf}\\
$J_G(\gamma,f_S)$, $J_A(\gamma,f_S)$ & local orbital and weighted orbital
  integrals, \eqref{eq:JG}--\eqref{eq:JA}\\
$J_A(z,f_S)$ & the weighted orbital integral at a central element
  (Lemma~\ref{lem:JA-central})\\
$a^M(S,\gamma)$ & global coefficients (Theorem~\ref{thm:fine}); the
  $a^G(S,u_\xi)$ are defined but not evaluated
\end{tabular}
\end{center}

Throughout, $G(\Q)$ is identified with its image in $\GA$ under
Lemma~\ref{lem:global}(a), and $N(\A)$ with $N(\A)\times1$; on $B(\Q)$
the two identifications agree, and $b\mapsto(b,1)$
(Lemma~\ref{lem:global}(b)). Test functions are anti-genuine
(Section~\ref{sec:functions}). We write $w=\begin{pmatrix}0&-1\\1&0
\end{pmatrix}\in G(\Q)$, and $d=[k:\Q]$.

\subsection{Equivalence classes}\label{sec:classes}

Following~\cite[\S2]{A2}, call $\gamma,\gamma'\in G(\Q)$
\emph{equivalent} if their semisimple parts are $G(\Q)$-conjugate, and let
$\mathcal C$ be the set of equivalence classes. A class is
\emph{unramified}~\cite[p.~949]{A2} if it consists of semisimple elements.

\begin{lemma}\label{lem:classes}
The classes $\mathcal O\in\mathcal C$ are of three kinds.
\begin{enumerate}[label=(\alph*)]
\item \emph{Central}: for $z\in\{\pm1\}$, $\mathcal O_z=\{zu\mid
  u\in G(\Q)\ \text{unipotent}\}$. It is not unramified: it contains $z$
  and $z\,n(\nu)$, $\nu\in k^\times$.
\item \emph{Hyperbolic}: the $G(\Q)$-conjugacy class of $\mu=a(t)$,
  $t\in k^\times\smallsetminus\{\pm1\}$. It contains $a(t^{-1})=w\mu w^{-1}$
  and meets $A(\Q)$ in $\{a(t),a(t^{-1})\}$; it is determined by
  $\{t,t^{-1}\}$. The centralizer of $\mu$ is $A$.
\item \emph{Elliptic}: the $G(\Q)$-conjugacy class of a regular
  $\mu\in G(\Q)$ whose characteristic polynomial is irreducible over $k$.
  Its centralizer is a maximal torus $T_\mu=\Res_{k/\Q}T_k$, anisotropic
  over $\Q$, and $\mathcal O\cap A(\Q)=\emptyset$.
\end{enumerate}
Classes of kinds (b) and~(c) are unramified. In all three cases the
centralizer of (the image of) $\mu$ in $\GA$ is the full preimage
$\overline{G_\mu(\A)}$, which is abelian in cases~(b) and~(c).
\end{lemma}

\begin{proof}
An element of $SL(2,k)$ with semisimple part $\pm1$ is $\pm$ a unipotent
element. A regular semisimple element with eigenvalues in $k$ is
$GL(2,k)$-conjugate to $a(t)$, and then $SL(2,k)$-conjugate to it,
because $\diag(y,1)$ centralizes $a(t)$. The centralizers are standard.
For the cover: if $h\in G_\mu(\A)$, lifts of $h$ and $\mu$ commute by
Lemma~\ref{lem:global}(d)--(e) (the argument of
Lemma~\ref{lem:local}(g), applied to the abelian group generated by $h$
and $\mu$); conversely the centralizer in $\GA$ maps into $G_\mu(\A)$.
\end{proof}

For general $n$ the centralizer of a lift of a split $\gamma$ can be
smaller than the preimage of $A$, and then the genuine orbital integral
vanishes. For $n=2$ this does not happen, by the last sentence of the
lemma.

\subsection{Kernels, heights and truncation}\label{sec:kernels}

Let $f\in\mathcal H(\GA)$. We use Arthur's convention: $R$ is the right
regular representation on $L^2(G(\Q)\backslash\GA)_{\mathrm{gen}}$, the
space of $\varphi$ with $\varphi(\gamma x\zeta)=\zeta\varphi(x)$, and
$(R(f)\varphi)(x)=\int_{\GA}f(y)\varphi(xy)\,dy$. The integrand is
invariant under $y\mapsto y\zeta$ because $f$ is anti-genuine. For a
standard parabolic $P=MN_P\in\{B,G\}$ and $\mathcal O\in\mathcal C$ put
\begin{gather}
K_P(x,y)=\sum_{\gamma\in M(\Q)}\int_{N_P(\A)}f(x^{-1}\gamma ny)\,dn,
\qquad
K_{P,\mathcal O}(x,y)=\sum_{\gamma\in M(\Q)\cap\mathcal O}
\int_{N_P(\A)}f(x^{-1}\gamma ny)\,dn,\label{eq:KP}\\
J_{P,\mathcal O}(x,y)=\sum_{\mu\in M(\Q)\cap\mathcal O}\
\sum_{\nu\in N_{P,\mu_s}(\Q)\backslash N_P(\Q)}\ \int_{N_{P,\mu_s}(\A)}
f(x^{-1}\nu^{-1}\mu n\nu y)\,dn,\label{eq:JP}
\end{gather}
where $\mu=\mu_s\mu_u$ is the Jordan decomposition and $N_{P,\mu_s}$ the
centralizer of $\mu_s$ in $N_P$. For $P=G$ these are
$K_G=\sum_{\gamma\in G(\Q)}f(x^{-1}\gamma y)$ and
$K_{G,\mathcal O}=J_{G,\mathcal O}=\sum_{\gamma\in\mathcal O}
f(x^{-1}\gamma y)$.

\begin{lemma}\label{lem:kernels}
The sums~\eqref{eq:KP} and~\eqref{eq:JP} are finite, uniformly for $x,y$
in compact sets. $K_P(x,y)=\sum_{\mathcal O}K_{P,\mathcal O}(x,y)$, and
$K_P$ is the kernel of $R_P(f)$, where $R_P$ is the right regular
representation on $L^2(N_P(\A)M(\Q)\backslash\GA)_{\mathrm{gen}}$.
\end{lemma}

\begin{proof}
$G(\Q)$ is discrete in $G(\A)$ and $f$ has compact support modulo
$\mu_2$, so only finitely many $\gamma$ contribute for $x,y$ in a compact
set; the integrals over $N_P(\A)$ are over compact sets for the same
reason. For upper triangular $b_1,b_2\in B(\A)$ with lower right entries
$d_1,d_2$, \eqref{eq:beta} and $\kappa_v|_B\equiv1$ give
$\sigma_\A(b_1,b_2)=\prod_v(d_2,d_1)_v$, which is $1$ when
$d_1,d_2\in k^\times$. So $b\mapsto(b,1)$ is a homomorphism on
$N(\A)A(\Q)$, and $N_P(\A)M(\Q)$ is a subgroup of $\GA$. For
$\varphi$ in that space,
\[
(R_P(f)\varphi)(x)=\int_{\GA}f(x^{-1}y)\varphi(y)\,dy
=\int_{N_P(\A)M(\Q)\backslash\GA}\ \sum_{\gamma\in M(\Q)}
\int_{N_P(\A)}f(x^{-1}\gamma ny)\,dn\ \varphi(y)\,dy,
\]
using $\varphi(\gamma ny)=\varphi(y)$. The decomposition over
$\mathcal O$ is a regrouping of the finite sum.
\end{proof}

\paragraph{Heights.} For a row vector $u=(c,d)\in k_v^2$ put
$\|u\|_v=\max(|c|_v,|d|_v)$ if $v$ is finite and
$\|u\|_v=(|c|^2+|d|^2)^{[k_v:\R]/2}$ if $v$ is archimedean, and for
$u\in\A_k^2$ put $\|u\|=\prod_v\|u_v\|_v$. Then $\|\lambda u\|_v=
|\lambda|_v\|u\|_v$, and $\|e_2k\|_v=1$ for $k\in K_v$, where
$e_2=(0,1)$. Recall $H=H_B$ from Lemma~\ref{lem:NAK-global}.

\begin{lemma}\label{lem:height-ineq}
\begin{enumerate}[label=(\alph*)]
\item For $x\in\GA$ with image $g\in G(\A)$, $H(x)=-\log\|e_2g\|$.
  The same holds at a single place: $H_v(g)=-\log\|e_2g\|_v$ for
  $g\in SL(2,k_v)$.
\item If $\gamma\in G(\Q)\smallsetminus B(\Q)$, then
  \[
  H(x)+H(\gamma x)\le0\qquad(x\in\GA).
  \]
  Consequently, if $T\ge0$, then for each $x$ there is at most one
  $\delta\in B(\Q)\backslash G(\Q)$ with $H(\delta x)>T$.
\item $H(ax)=H(a)+H(x)$ and $H(wax)=-H(a)+H(wx)$ for
  $a\in\overline{A(\A)}$. Hence
  \begin{equation}\label{eq:vA}
  v_A(x)=-H(x)-H(wx)
  \end{equation}
  is left $\overline{A(\A)}$-invariant, right $\overline{K_0}$-invariant,
  and $v_A(x)\ge0$.
\end{enumerate}
\end{lemma}

\begin{proof}
(a) If $g=n\,a(t)\,k$, then $e_2g=t^{-1}e_2k$, so $\|e_2g\|=|t|^{-1}
=e^{-H(x)}$. The local statement is the same computation.

(b) Let $e_2\gamma=(c,d)$, with $c\ne0$ since $\gamma\notin B(\Q)$. The
matrix with rows $e_2g$ and $e_2\gamma g$ has determinant
$\det\begin{pmatrix}0&1\\c&d\end{pmatrix}\det g=-c$. At each place,
$|\det(u_1;u_2)|_v\le\|u_1\|_v\|u_2\|_v$: this is the ultrametric
inequality at finite places and Hadamard's inequality (squared at complex
places) at archimedean ones. Taking the product over $v$,
$1=|c|_{\A_k}\le\|e_2g\|\,\|e_2\gamma g\|=e^{-H(x)-H(\gamma x)}$ by~(a)
and the product formula. If $\delta_1\ne\delta_2$ in $B(\Q)\backslash
G(\Q)$ both had $H(\delta_ix)>T\ge0$, then (b), applied to $\delta_1x$ and
$\gamma=\delta_2\delta_1^{-1}\notin B(\Q)$, would give $0<2T<
H(\delta_1x)+H(\delta_2x)\le0$.

(c) The first identity is Lemma~\ref{lem:H}(a). For the second,
$wa\,w^{-1}=a^{-1}$, so $H(wax)=H(a^{-1}\cdot wx)=-H(a)+H(wx)$. Right
$K_0$-invariance is Lemma~\ref{lem:H}(b), and $v_A\ge0$ is~(b) with
$\gamma=w$.
\end{proof}

For Lemma~\ref{lem:height-ineq}(b) compare~\cite[Lemma~5.1]{A2}.
Reduction theory~\cite{Go} gives a Siegel set
$\mathfrak S=\{(n,1)\,\tilde a\,\tilde k\mid n\in\omega,\ H(\tilde a)>-c,\
p(\tilde a)\in\Omega\,A_B(\R)^0,\ \tilde k\in\overline{K_0}\}$, with
$\omega\subset N(\A)$ and $\Omega\subset A(\A)^1$ compact and $c\ge0$, such
that $\GA=G(\Q)\mathfrak S$. In the coordinates $x=nak$ the measure
is $\delta_B(a)^{-1}dn\,da\,dk$ by~\eqref{eq:iwasawa}, and $\mathfrak S$ has
finite volume.

\paragraph{Truncation.} Let $\hat\tau_B=\tau_B$ be the characteristic
function of $(0,\infty)$ (Section~\ref{sec:chambers}), $\hat\tau_G=1$, and
$\dim A_G=0$, $\dim A_B=1$. For $T\in\R\cong\mathfrak a$ put
\begin{equation}\label{eq:kT}
k^T(x,f)=K_G(x,x)-\sum_{\delta\in B(\Q)\backslash G(\Q)}
K_B(\delta x,\delta x)\,\hat\tau_B(H(\delta x)-T),
\end{equation}
and define $k^T_{\mathcal O}(x,f)$ and $j^T_{\mathcal O}(x,f)$ in the same
way with $K_{P,\mathcal O}$, respectively $J_{P,\mathcal O}$, in place of
$K_P$. By Lemma~\ref{lem:height-ineq}(b), for $T\ge0$ the sum over
$\delta$ has at most one nonzero term. Let $F^G(x,T)$ be the
characteristic function of $\{x\mid H(\delta x)\le T\ \text{for all}\
\delta\in G(\Q)\}$; its support is compact modulo $G(\Q)$ by reduction
theory~\cite{Go}.

\begin{lemma}\label{lem:partition}
Let $T\ge0$. Then $F^G(x,T)+\sum_{\delta\in B(\Q)\backslash G(\Q)}
\tau_B(H(\delta x)-T)=1$, and
\begin{equation}\label{eq:kT-split}
k^T(x,f)=F^G(x,T)\,K_G(x,x)+\sum_{\delta\in B(\Q)\backslash G(\Q)}
\tau_B(H(\delta x)-T)\bigl(K_G(\delta x,\delta x)-K_B(\delta x,\delta x)\bigr).
\end{equation}
\end{lemma}

\begin{proof}
The first identity says that either $H(\delta x)\le T$ for all $\delta$,
or $H(\delta x)>T$ for exactly one coset, which is
Lemma~\ref{lem:height-ineq}(b). Multiply it by $K_G(x,x)=
K_G(\delta x,\delta x)$ and subtract the $\delta$-sum of~\eqref{eq:kT}.
\end{proof}

In general~\cite[\S6]{A2}, \cite[\S8]{Ar05}, the
partition holds for $T$ sufficiently regular; here $T\ge0$ suffices.

\subsection{The coarse geometric expansion}\label{sec:coarse}

Let $dx$ be a Haar measure on $G(\A)$ with $dx=\delta_B(a)^{-1}dn\,da\,dk$
for $x=nak$ as in~\eqref{eq:iwasawa}, where $dk$ gives $K_0$ volume $1$,
$dn$ gives $N(\Q)\backslash N(\A)$ volume $1$, and $da$ is a Haar measure
on $A(\A)\cong\A_k^\times$, written $d^\times t$ for $a=a(t)$. On covers
$\mu_2$ has volume~$1$. Write $A(\A)=A(\A)^1\times A_B(\R)^0$, where
$A_B(\R)^0=\{a(u)\mid u\in\R_{>0}\}$ ($u$ embedded diagonally in
$k\otimes\R$), let $dX$ be Lebesgue measure on $\R\cong\mathfrak a$
transported to $A_B(\R)^0$ by $H$, and let $da^1$ be the Haar measure on
$A(\A)^1$ with $da=da^1\,dX$. Put
\[
c_k=\mathrm{vol}\bigl(A(\Q)\backslash A(\A)^1\bigr)
=\mathrm{vol}\bigl(k^\times\backslash\A_k^1\bigr).
\]

\begin{theorem}\label{thm:coarse}
Let $f\in\mathcal H(\GA)$ and $T\ge0$.
\begin{enumerate}[label=(\alph*)]
\item Only finitely many $\mathcal O\in\mathcal C$ have
  $k^T_{\mathcal O}(\cdot,f)\not\equiv0$, and $k^T(x,f)=
  \sum_{\mathcal O}k^T_{\mathcal O}(x,f)$.
\item Each $k^T_{\mathcal O}(\cdot,f)$ and each
  $j^T_{\mathcal O}(\cdot,f)$ is absolutely integrable over
  $G(\Q)\backslash\GA$, and
  \[
  J^T_{\mathcal O}(f):=\int_{G(\Q)\backslash\GA}k^T_{\mathcal O}(x,f)\,dx
  =\int_{G(\Q)\backslash\GA}j^T_{\mathcal O}(x,f)\,dx.
  \]
\item Hence $k^T(\cdot,f)$ is absolutely integrable, and
  $J^T(f):=\int k^T(x,f)\,dx=\sum_{\mathcal O}J^T_{\mathcal O}(f)$.
\item Each $J^T_{\mathcal O}(f)$ is a polynomial in $T$ of degree at
  most $1$, given by Propositions~\ref{prop:hyperbolic},
  \ref{prop:elliptic} and~\ref{prop:central}.
\end{enumerate}
\end{theorem}

We put $J_{\mathcal O}(f)=J^0_{\mathcal O}(f)$ and $J(f)=J^0(f)=
\sum_{\mathcal O}J_{\mathcal O}(f)$. In general one defines $J_{\mathcal
O}$ as the value at $T=0$ of the polynomial $J^T_{\mathcal O}$, which is
known only for $T$ sufficiently regular~\cite{A3,Ar05}. Here the integral
itself converges at $T=0$, so the two definitions agree.
Theorem~\ref{thm:coarse} is the case of $\GA$ of Arthur's coarse
expansion~\cite[\S\S7--8]{A2}, \cite{A3}, extended to covers by
Li~\cite{LiI}; for this group we prove it directly.

\begin{proposition}[Hyperbolic classes]\label{prop:hyperbolic}
Let $\mathcal O$ be the class of $\mu=a(t)$, $t\in k^\times
\smallsetminus\{\pm1\}$, and $T\ge0$. Then
\begin{equation}\label{eq:jT-hyp}
j^T_{\mathcal O}(x,f)=\sum_{\eta\in A(\Q)\backslash G(\Q)}
f(x^{-1}\eta^{-1}\mu\eta x)\bigl(1-\hat\tau_B(H(\eta x)-T)
-\hat\tau_B(H(w\eta x)-T)\bigr),
\end{equation}
and
\begin{equation}\label{eq:JT-hyp}
J^T_{\mathcal O}(f)=c_k\int_{\overline{A(\A)}\backslash\GA}
f(y^{-1}\mu y)\,\bigl(2T+v_A(y)\bigr)\,dy,
\end{equation}
with $v_A(y)=-H(y)-H(wy)\ge0$ as in~\eqref{eq:vA}. The integral converges
absolutely.
\end{proposition}

\begin{proof}
Here $\mathcal O\cap A(\Q)=\{\mu,\mu^{-1}\}$ and $N_{\mu_s}=1$, so
\[
J_{B,\mathcal O}(y,y)=\sum_{\mu'\in\{\mu,\mu^{-1}\}}\ \sum_{\nu\in N(\Q)}
f(y^{-1}\nu^{-1}\mu'\nu y).
\] As $\delta$ runs over $B(\Q)\backslash
G(\Q)$ and $\nu$ over $N(\Q)$, $\eta=\nu\delta$ runs over
$A(\Q)\backslash G(\Q)$, and $H(\nu\delta x)=H(\delta x)$. Writing
$\mu^{-1}=w^{-1}\mu w$ and replacing $\eta$ by $w^{-1}\eta$ in the
$\mu^{-1}$-term (note that $H(w^{-1}y)=H(-wy)=H(wy)$)
gives~\eqref{eq:jT-hyp}, since $K_{G,\mathcal O}(x,x)=\sum_{\eta}f(x^{-1}\eta^{-1}\mu\eta x)$.

By Lemma~\ref{lem:height-ineq}(b), for $T\ge0$ the two $\hat\tau_B$ in
\eqref{eq:jT-hyp} are never both $1$, so the bracket is $0$ or $1$.
Unfolding, and writing $x=\tilde a y$ with $\tilde a\in\overline{A(\A)}$,
\[
\int|j^T_{\mathcal O}|\le\int_{\overline{A(\A)}\backslash\GA}
|f(y^{-1}\mu y)|\int_{A(\Q)\backslash\overline{A(\A)}}
\bigl(1-\hat\tau_B(H(\tilde ay)-T)-\hat\tau_B(H(w\tilde ay)-T)\bigr)
\,d\tilde a\,dy,
\]
with equality if $f\ge0$, and we used $f(y^{-1}\tilde a^{-1}\mu\tilde
ay)=f(y^{-1}\mu y)$, because $\overline{A(\A)}$ is abelian
(Lemma~\ref{lem:global}(e)). Put $X=H(\tilde a)$, $h_1=H(y)$, $h_2=
H(wy)$; by Lemma~\ref{lem:height-ineq}(c) the bracket is $1-\mathbf
1[X+h_1>T]-\mathbf 1[-X+h_2>T]$, whose integral over $X\in\R$ is
$2T-h_1-h_2=2T+v_A(y)$. The integral over $A(\Q)\backslash A(\A)^1$
gives $c_k$. This proves~\eqref{eq:JT-hyp} for $\int j^T_{\mathcal O}$
and its absolute convergence: in the coordinates $y=n(z)k$ on
$\overline{A(\A)}\backslash\GA$ (measure $dz\,dk$), $y^{-1}\mu y=
k^{-1}\mu\,n((1-t^{-2})z)k$, so $z$ runs over a compact set and
$v_A(n(z)k)$ is bounded there (Lemma~\ref{lem:weight-local}).

It remains to show $\int k^T_{\mathcal O}=\int j^T_{\mathcal O}$. By
definition
\[
k^T_{\mathcal O}(x)-j^T_{\mathcal O}(x)=\sum_{\delta}
\hat\tau_B(H(\delta x)-T)\,\bigl(J_{B,\mathcal O}-K_{B,\mathcal O}\bigr)
(\delta x,\delta x).
\]
Let $x=(n(u_0),1)\,\tilde a\,\tilde k$ with $p(\tilde a)=a(s)$. For
$\mu'=a(t')$ one has $n(-\nu)\mu'n(\nu)=\mu'n((1-t'^{-2})\nu)$, a
bijection of $N(\Q)$, and $\tilde a^{-1}\mu'(n(z),1)\tilde a=
\mu'(n(s^{-2}z),1)$ (Lemma~\ref{lem:local}(h)). So with
$g_{\mu'}(z)=f(\tilde k^{-1}\mu'(n(z),1)\tilde k)$ and $c=(1-t'^{-2})u_0$,
\[
(J_{B,\mathcal O}-K_{B,\mathcal O})(x,x)=\sum_{\mu'}\Bigl(\sum_{\nu\in k}
g_{\mu'}(s^{-2}(\nu+c))-\int_{\A_k}g_{\mu'}(s^{-2}(z+c))\,dz\Bigr).
\]
By Poisson summation over $k\subset\A_k$~\cite{Ta}, this is
$|s|^2\sum_{\mu'}\sum_{\nu\in k^\times}\hat g_{\mu'}(s^2\nu)\psi(\nu c)$,
which is $O(|s|^{2-M})$ for every $M$, uniformly for $\tilde k\in
\overline{K_0}$ and $s$ in $\Omega\times\R_{\ge1}$, $\Omega\subset
\A_k^1$ compact. Since $\delta_B(a)^{-1}=|s|^{-2}$, the function
$\hat\tau_B(H(x)-T)(J_{B,\mathcal O}-K_{B,\mathcal O})(x,x)$ is
absolutely integrable over $B(\Q)\backslash\GA$, and unfolding is
justified. Integrating first over $u_0\in k\backslash\A_k$ gives $0$,
because $\int_{k\backslash\A_k}\sum_{\nu\in k}g(s^{-2}(\nu+c))\,du_0=
|1-t'^{-2}|_{\A_k}^{-1}\int_{\A_k}g(s^{-2}z)\,dz$ and
$|1-t'^{-2}|_{\A_k}=1$ by the product formula.
\end{proof}

In the form~\eqref{eq:JT-hyp}, $J_{\mathcal O}(f)$ is $c_k$ times Arthur's
weighted orbital integral, $v_A$ being the volume of the convex hull of
$\{-H_B(y),-H_{\bar B}(y)\}$ with $H_{\bar B}(y)=-H(wy)$. For the measure
on $\mathfrak a$ see Remark~\ref{rem:alpha}.

\begin{proposition}[Elliptic classes]\label{prop:elliptic}
Let $\mathcal O$ be elliptic, with $\mu\in\mathcal O$ and $T_\mu=G_\mu$.
Then
$k^T_{\mathcal O}=j^T_{\mathcal O}=K_{G,\mathcal O}(x,x)$ for every $T$,
and
\[
J_{\mathcal O}(f)=J^T_{\mathcal O}(f)=\mathrm{vol}\bigl(T_\mu(\Q)\backslash
T_\mu(\A)\bigr)\int_{\overline{T_\mu(\A)}\backslash\GA}f(y^{-1}\mu y)\,dy,
\]
the integral converging absolutely.
\end{proposition}

\begin{proof}
$\mathcal O\cap A(\Q)=\emptyset$, so the $B$-terms vanish. Unfold
$\sum_{\eta\in T_\mu(\Q)\backslash G(\Q)}$ and use that
$\overline{T_\mu(\A)}$ is abelian, so $f(y^{-1}\tilde t^{-1}\mu\tilde ty)=
f(y^{-1}\mu y)$. The quotient $T_\mu(\Q)\backslash T_\mu(\A)$ is compact
($T_\mu$ is anisotropic, so $T_\mu(\A)^1=T_\mu(\A)$), and the orbit of $\mu$ is closed, so the orbital
integral of the compactly supported $|f|$ converges.
\end{proof}

\paragraph{Central classes.} For $z\in\{\pm1\}$ let $f_z(x)=f(zx)$, with
$z$ identified with its image $(z,1)$ in $\GA$ (Lemma~\ref{lem:global}(a),
$\kappa_\A(\pm1)=1$). It is central, so $f_z\in\mathcal H(\GA)$, and at
almost all $v$ its local component is again $e_v$, because $(-1,1)\in
K_v^*$. For $f\in\mathcal H(\GA)$ define a Schwartz--Bruhat function on
$\A_k$ by
\begin{equation}\label{eq:Phi}
\Phi_f(y)=\int_{\overline{K_0}}f\bigl(\tilde k^{-1}(n(y),1)\tilde k\bigr)
\,d\tilde k,
\end{equation}
let $\hat\Phi_f$ be its Fourier transform with respect to $\psi$, and for
$\mathrm{Re}\,s>1$ put
\begin{equation}\label{eq:Zf}
Z(f,s)=\int_{k^\times\backslash\A_k^\times}|t|^{-2s}\sum_{\nu\in k^\times}
\Phi_f(t^{-2}\nu)\,d^\times t
=\frac12\sum_{\xi\in k^\times/k^{\times2}}\ \int_{\A_k^\times}
\Phi_f(\xi t^2)\,|t|^{2s}\,d^\times t .
\end{equation}
For $f=\bigotimes_vf_v$ the last integral is the product of the local
integrals $Z_v(\Phi_{f,v},\xi,s)=\int_{k_v^\times}\Phi_{f,v}(\xi t^2)
|t|_v^{2s}d^\times t$, which are, at $s=1$, twice the local unipotent orbital
integrals of $f_v$ at $n(\xi)$ (Lemma~\ref{lem:JA-central}(c)).

\begin{proposition}[Central classes]\label{prop:central}
Let $z\in\{\pm1\}$. Then $k^T_{\mathcal O_z}=j^T_{\mathcal O_z}$, and
$J^T_{\mathcal O_z}(f)=J^T_{\mathrm{unip}}(f_z)$, where for every $T\in\R$
\begin{equation}\label{eq:Junip}
J^T_{\mathrm{unip}}(f)=\mathrm{vol}\bigl(G(\Q)\backslash\GA\bigr)f(1)
+c_k\,\hat\Phi_f(0)\,T+\mathop{\mathrm{FP}}_{s=1}Z(f,s).
\end{equation}
Here $Z(f,s)$, absolutely convergent for $\mathrm{Re}\,s>1$, extends to a
meromorphic function on $\C$ whose only possible poles are simple ones at
$s=1$ and $s=0$; the residue at $s=1$ is $c_k\hat\Phi_f(0)/2$, and
$\mathrm{FP}_{s=1}$ denotes the constant term of the Laurent expansion.
In particular
\[
J_{\mathcal O_z}(f)=J_{\mathrm{unip}}(f_z),\qquad
J_{\mathrm{unip}}(f)=\mathrm{vol}\bigl(G(\Q)\backslash\GA\bigr)f(1)
+\mathop{\mathrm{FP}}_{s=1}Z(f,s).
\]
\end{proposition}

\begin{proof}
Since $z$ is central, $K_{P,\mathcal O_z}(x,y)$ for $f$ equals
$K_{P,\mathcal O_1}(x,y)$ for $f_z$, and likewise for $J_P$; and
$J_{P,\mathcal O_1}=K_{P,\mathcal O_1}$ because $\mu_s=1$ and
$N_{P,1}=N_P$. So it suffices to treat $z=1$. The nontrivial unipotent
elements of $G(\Q)$ are the $\delta^{-1}n(\nu)\delta$,
$\delta\in B(\Q)\backslash G(\Q)$, $\nu\in k^\times$, each written once,
so
\[
k^T_{\mathcal O_1}(x)=f(1)+\sum_{\delta}\Bigl(F(\delta x)
-\hat\tau_B(H(\delta x)-T)\,C(\delta x)\Bigr),
\]
where $F(x)=\sum_{\nu\in k^\times}f(x^{-1}n(\nu)x)$ and
$C(x)=\int_{N(\A)}f(x^{-1}nx)\,dn$.
For $x=(n_0,1)\tilde a\tilde k$ with $p(\tilde a)=a(t)$ we have
$x^{-1}(n(\nu),1)x=\tilde k^{-1}(n(t^{-2}\nu),1)\tilde k$, so $F$ and $C$
depend only on $t$ and $\tilde k$. With $g_{\tilde k}(y)=
f(\tilde k^{-1}(n(y),1)\tilde k)$, Poisson summation gives
\[
\delta_B(a)^{-1}\bigl(F-C\bigr)(x)=\sum_{\nu\in k^\times}
\hat g_{\tilde k}(t^2\nu)-|t|^{-2}g_{\tilde k}(0),
\]
which is integrable over $|t|>e^T$, uniformly in $\tilde k$; and
$F(x)=0$ once $|t|$ is small, because $g_{\tilde k}$ has compact support
uniformly in $\tilde k$ and $k$ is discrete in $\A_k$. Hence the
three pieces $f(1)$, $\sum_\delta(1-\hat\tau_B)F$ and
$\sum_\delta\hat\tau_B(F-C)$ of $k^T_{\mathcal O_1}$ are absolutely
integrable, and unfolding each over $B(\Q)\backslash\GA=N(\Q)\backslash
N(\A)\times A(\Q)\backslash A(\A)\times\overline{K_0}$ gives, with
$G(t)=|t|^{-2}\sum_{\nu\in k^\times}\Phi_f(t^{-2}\nu)$ and
$\delta_B(a)^{-1}C=\hat\Phi_f(0)$ after the $\tilde k$-integration,
\[
J^T_{\mathrm{unip}}(f)=\mathrm{vol}\bigl(G(\Q)\backslash\GA\bigr)f(1)
+\int_{|t|\le e^T}G(t)\,d^\times t+\int_{|t|>e^T}\bigl(G(t)-\hat\Phi_f(0)
\bigr)\,d^\times t.
\]
Now $Z(f,s)=\int G(t)|t|^{2-2s}d^\times t$. Split it at $|t|=e^T$ and
write $G=(G-\hat\Phi_f(0))+\hat\Phi_f(0)$ on $|t|>e^T$. Since
$\int_{|t|>e^T}|t|^{2-2s}d^\times t=c_k\,e^{(2-2s)T}/(2s-2)$,
\[
Z(f,s)=\int_{|t|\le e^T}G|t|^{2-2s}+\int_{|t|>e^T}\bigl(G-\hat\Phi_f(0)
\bigr)|t|^{2-2s}+c_k\hat\Phi_f(0)\frac{e^{(2-2s)T}}{2s-2}.
\]
The first integral is entire in $s$ (compact range of $|t|$). By the
Poisson formula above, $G-\hat\Phi_f(0)=\sum_{\nu\ne0}\hat\Phi_f(t^2\nu)
-|t|^{-2}\Phi_f(0)$, so the second integral is an entire function minus
$c_k\Phi_f(0)e^{-2sT}/(2s)$. This gives the continuation and the poles.
Expanding $e^{(2-2s)T}/(2s-2)=1/(2(s-1))-T+O(s-1)$ gives~\eqref{eq:Junip}.
The second form of~\eqref{eq:Zf} follows by writing $k^\times$ as the
union of the cosets $\xi k^{\times2}$ and unfolding
$\sum_{c\in k^\times/\{\pm1\}}$ against $k^\times\backslash\A_k^\times$,
using $|ct|=|t|$.
\end{proof}

The linear term of~\eqref{eq:Junip} agrees with the fine expansion: the
term $\tfrac12a^A(S,1)J^T_A(1,f)$ there has weight $v_A+2T$, and
$a^A(S,1)=c_k$, $J_G$ of the class induced from $1$ tends to
$\hat\Phi_f(0)$ (Lemma~\ref{lem:JA-central}).

\begin{example}\label{ex:unip-Q}
Let $k=\Q$ and $\Phi=\Phi_\infty\otimes\mathbf 1_{\hat\Z}$ with
$\Phi_\infty(y)=e^{-\pi y^2}$, so $\hat\Phi(0)=1$, and take $d^\times t$
with $\mathrm{vol}(\hat\Z^\times)=1$, so $c_\Q=1$. Representing
$\Q^\times\backslash\A^\times$ by $\R_{>0}\times\hat\Z^\times$,
$Z(s)=\int_0^\infty u^{-2s}\sum_{\nu\ne0}e^{-\pi\nu^2/u^4}\,du/u=
\zeta(s)M(s)$ with $M(s)=\int_0^\infty e^{-\pi w^2}w^{s-1}dw=
\tfrac12\pi^{-s/2}\Gamma(s/2)$. The residue at $s=1$ is $M(1)=\tfrac12=
c_\Q\hat\Phi(0)/2$, and
\[
\mathop{\mathrm{FP}}_{s=1}Z(s)=\gamma_EM(1)+M'(1)
=\frac{\gamma_E}4-\frac{\log4\pi}4\approx-0.48845 .
\]
In the second form of~\eqref{eq:Zf}, $\xi$ runs over $\pm m$ with $m$
squarefree, each contributing $m^{-s}M(s)\zeta(2s)$, and
$\sum_m m^{-s}=\zeta(s)/\zeta(2s)$. The check script computes the
truncated integrals of the proof of Proposition~\ref{prop:central} at
$T=-0.3,0,0.5,1.2$ and finds that
$J^T_{\mathrm{unip}}-\mathrm{vol}\cdot f(1)-T$ equals this constant to
$14$ digits.
\end{example}

\begin{proof}[Proof of Theorem~\ref{thm:coarse}]
(a) $k^T(x)=\sum_{\mathcal O}k^T_{\mathcal O}(x)$ by
Lemma~\ref{lem:kernels}. If $k^T_{\mathcal O}\not\equiv0$, then some
element of $\mathcal O$, or some $\mu n$ with $\mu\in\mathcal O\cap A(\Q)$,
is $G(\A)$-conjugate into the support of $f$; its trace lies in $k$ and
in a compact subset of $\A_k$, so in a finite set. This leaves the two
central classes, finitely many hyperbolic classes, and the elliptic
classes with finitely many characteristic polynomials. The
$SL(2,k)$-classes with a given regular elliptic characteristic polynomial,
splitting field $E$, form a torsor under $k^\times/N_{E/k}E^\times$,
which injects into $\bigoplus_vk_v^\times/N E_v^\times$ by the Hasse
norm theorem~\cite{CF}. Let $v$ be a finite place outside a finite set,
where $f_v=e_v$, $E_v/k_v$ is unramified and $\mu$ is a unit with
$|D(\mu)|_v=1$. If $v$ splits in $E$, the local class is unique. If $v$
is inert, an element $\gamma\in K_v$ with the characteristic polynomial
of $\mu$ has irreducible reduction, so it fixes no edge of the
Bruhat--Tits tree and exactly one vertex, that of $K_v$. If $\gamma'=
m\gamma m^{-1}\in K_v$ with $m\in GL(2,k_v)$, then $m$ fixes that vertex,
so $m\in k_v^\times GL(2,\OO_v)$ and $\det m\in k_v^{\times2}\OO_v^\times
=N E_v^\times$; correcting $m$ by an element of the centralizer $E_v^\times$
makes it lie in $SL(2,k_v)$. So one local class meets $K_v$. Hence only
finitely many local classes, and so finitely many global classes, meet the
support of $f$.

(b) and (d) are Propositions~\ref{prop:hyperbolic}, \ref{prop:elliptic}
and~\ref{prop:central}. (c) follows from (a) and (b).
\end{proof}

\subsection{Local distributions}\label{sec:local-dist}

\paragraph{Test functions.} Let $S$ be a finite set of places of $\Q$
containing $\infty$ and $2$, and let $S_k$ be the set of places of $k$
above $S$. Let $\overline{G(\Q_S)}$ be the cover of $G(\Q_S)=
\prod_{v\in S_k}SL(2,k_v)$ defined by $\prod_{v\in S_k}\sigma_v$, the
quotient of $\prod_{v\in S_k}\overline{SL(2,k_v)}$ as in
Section~\ref{sec:local}. Let $\mathcal H(\overline{G(\Q_S)})$ be the span
of the products $\bigotimes_{v\in S_k}f_v$, $f_v\in
C_c^\infty(\overline{SL(2,k_v)})_-$, with $f_v$ bi-$\overline{K_v}$-finite
at archimedean $v$; here $C_c^\infty(\cdot)_-$ is the space of
Section~\ref{sec:functions}, defined on the topological cover. Let
$\mathcal H_N(\overline{G(\Q_S)})_\Gamma$, for $N>0$ and a finite set $\Gamma$ of irreducible
representations of $\overline{K_S}$, be the subspace of functions
supported on $\{\|g\|\le e^N\}$ whose left and right
$\overline{K_S}$-translates span representations with constituents in
$\Gamma$. With the seminorms $\sup|Df|$, $D$ a differential operator on
$\overline{G(\R)}$, it is a Fr\'echet space, and
$\mathcal H(\overline{G(\Q_S)})$ is the union over $N$ and $\Gamma$. For
$f_S\in\mathcal H(\overline{G(\Q_S)})$ put $f=f_S\otimes e^S$, $e^S=
\bigotimes_{v\notin S_k}e_v$. For $\gamma\in G(\Q)$ let
$\gamma_S=(\gamma,\kappa_\A(\gamma))\in\overline{G(\Q_S)}$; since
$e_v((k,\zeta)_\sigma)=\zeta$ for $k\in K_v$, $f(\gamma)=f_S(\gamma_S)$
when $\gamma\in K_v$ for all $v\notin S_k$.

Arthur's spaces of almost compactly supported and moderate
functions~\cite{A9} coincide with $\mathcal H$ for $G$, because $A_G=1$.
For $A$ they do not: $\mathfrak a_A=\mathfrak a\ne0$, and a function that
is compactly supported on each fiber of $H$ need not be compactly
supported. This matters in Section~\ref{sec:IAgamma}.

\paragraph{Orbital integrals.} At each $v\in S_k$ we use the local
measures of Section~\ref{sec:measures-local}, and on $G(\Q_S)$ their
product. For $\gamma=a(t)\in A(\Q_S)$ with $t_v^2\ne1$ for all $v$, and
$\gamma$ also denoting a lift in $\overline{A(\Q_S)}$, let
$D(\gamma)=\prod_v(t_v-t_v^{-1})^2$, and
\begin{align}
J_G(\gamma,f_S)&=|D(\gamma)|^{1/2}\int_{A(\Q_S)\backslash G(\Q_S)}
f_S(x^{-1}\gamma x)\,dx,\label{eq:JG}\\
J_A(\gamma,f_S)&=|D(\gamma)|^{1/2}\int_{A(\Q_S)\backslash G(\Q_S)}
f_S(x^{-1}\gamma x)\,v_A(x)\,dx,\qquad
v_A(x)=\sum_{v\in S_k}v_{A,v}(x_v),\label{eq:JA}
\end{align}
where $v_{A,v}(x)=-H_v(x)-H_v(wx)$ is the local weight. For
$\gamma=\mu_S$ regular elliptic, $J_G(\gamma,f_S)$ is defined
by~\eqref{eq:JG} with the centralizer $T_\mu(\Q_S)$ in place of
$A(\Q_S)$.
For a unipotent $u\in G(\Q_S)$ and $z\in\{\pm1\}$, $J_G(zu,f_S)=
\int_{G_u(\Q_S)\backslash G(\Q_S)}f_S(x^{-1}zux)\,dx$. These integrals
converge absolutely, and the quotient measures are those of~\eqref{eq:iwasawa}:
on $A(\Q_S)\backslash G(\Q_S)$, $dx=dn\,dk$ for $x=nk$.

\begin{lemma}\label{lem:weight-local}
Let $v$ be a place of $k$. Then $v_{A,v}\ge0$, it is left
$A(k_v)$-invariant and right $K_v$-invariant, and
\[
v_{A,v}(n(y)k)=\log\|(1,y)\|_v=\begin{cases}
\log\max(1,|y|_v),&v\ \text{finite},\\
\tfrac12\log(1+y^2),&v\ \text{real},\\
\log(1+|y|^2),&v\ \text{complex}.\end{cases}
\]
For $x\in G(\A)$, $v_A(x)=\sum_vv_{A,v}(x_v)$, a finite sum.
\end{lemma}

\begin{proof}
As for Lemma~\ref{lem:height-ineq}, at one place: $|\det(e_2x;e_2wx)|_v
=|-1|_v=1$ gives $v_{A,v}\ge0$, and $H_v(n(y))=0$, $e_2wn(y)=(1,y)$. The
global identity is $H=\sum_vH_v$.
\end{proof}

\begin{lemma}\label{lem:JA-central}
Let $f_S=\bigotimes_{v\in S_k}f_v$, $z\in\{\pm1\}$, and let
$\Phi_v(y)=\int_{\overline{K_v}}f_v(\tilde k^{-1}z(n(y),1)\tilde k)\,
d\tilde k$, a Schwartz--Bruhat function on $k_v$, with $\hat\Phi_v(0)=
\int_{k_v}\Phi_v$. For $t\in k_S^\times$ near $1$ let $a_t=(a(t),1)z$
and
\[
r(t)=\sum_{v\in S_k}\log\bigl|t_v^2-t_v^{-2}\bigr|_v .
\]
\begin{enumerate}[label=(\alph*)]
\item $J_G(a_t,f_S)\to\prod_v\hat\Phi_v(0)$ as $t\to1$.
\item The limit
  \[
  J_A(z,f_S):=\lim_{t\to1}\bigl(J_A(a_t,f_S)+r(t)\,J_G(a_t,f_S)\bigr)
  =\sum_{v\in S_k}\Bigl(\int_{k_v}\Phi_v(y)\log|2y|_v\,dy\Bigr)
  \prod_{w\in S_k\smallsetminus\{v\}}\hat\Phi_w(0)
  \]
  exists. Without the term $r(t)J_G$, $J_A(a_t,f_S)=-\log|1-t^{-2}|_S
  \prod_v\hat\Phi_v(0)+O(1)$, which diverges unless
  $\prod_v\hat\Phi_v(0)=0$.
\item For $\xi\in k_S^\times$, $J_G(z\,n(\xi),f_S)=\prod_{v\in S_k}
  \tfrac12\int_{k_v^\times}\Phi_v(\xi t^2)|t|_v^2\,d^\times t$.
\end{enumerate}
\end{lemma}

\begin{proof}
Everything factors over $v\in S_k$, $J_A$ by the Leibniz rule
$J_A=\sum_vJ_{A,v}\prod_{w\ne v}J_{G,w}$, and $r(t)=\sum_vr_v(t_v)$; so
it suffices to treat one place. For $x=n(y)k$,
$x^{-1}a_tx=\tilde k^{-1}a_t(n(\varepsilon y),1)\tilde k$ with
$\varepsilon=1-t^{-2}$, and $|D(a(t))|^{1/2}=|t|\,|\varepsilon|$. Put
$y'=\varepsilon y$ and $h_t(y')=\int f_v(\tilde k^{-1}a_t(n(y'),1)
\tilde k)\,d\tilde k$. By Lemma~\ref{lem:weight-local} and homogeneity of
$\|\cdot\|_v$,
\[
J_{G,v}=|t|\int h_t(y')\,dy',\qquad
J_{A,v}=|t|\int h_t(y')\bigl(\log\|(\varepsilon,y')\|_v-\log|\varepsilon|_v
\bigr)\,dy'.
\]
This gives~(a), since $h_t\to\Phi_v$ with uniformly compact support, and
the asymptotics in~(b). Since $t^2-t^{-2}=t^2\varepsilon(1+t^{-2})$,
\[
J_{A,v}+r_vJ_{G,v}=|t|\int h_t(y')\bigl(\log\|(\varepsilon,y')\|_v+
\log|t^2(1+t^{-2})|_v\bigr)\,dy'\to\int\Phi_v(y)\bigl(\log|y|_v+
\log|2|_v\bigr)\,dy
\]
by dominated convergence ($\log|y|_v$ is locally integrable). For~(c),
$G_{n(\xi)}=\{\pm1\}N$, so the integral is over $\{\pm1\}\backslash
A(k_v)\times K_v$ with measure $\tfrac12\delta_B(a)^{-1}da\,dk$, and
$a(t)^{-1}n(\xi)a(t)=n(t^{-2}\xi)$; substitute $t\mapsto t^{-1}$.
\end{proof}

Lemma~\ref{lem:JA-central}(b) is Arthur's definition~\cite[(8.1)]{A7} of
$J_M(\gamma,f)$ at a non-regular point, as the limit of
$\sum_{L\in\mathcal L(M)}r_M^L(\gamma,a)J_L(a\gamma,f)$, specialized to
$M=A$, $\gamma=z$: $r^A_A=1$, and $r^G_A(1,a)$ is the function $r(t)$
above (Remark~\ref{rem:alpha}). The term $L=G$ is present because the
class induced from $1\in A$ to $G$ is the regular unipotent class; without
it the limit diverges, as~(b) shows.

\subsection{The fine expansion}\label{sec:fine}

For $S$ as above let $\mathcal U_S$ be the set of $(G,S)$-equivalence
classes of unipotent elements of $G(\Q)$, that is, the classes
$\{1\}$ and $u_\xi=n(\xi)$, $\xi\in k_S^\times/k_S^{\times2}$ (each
meets $G(\Q)$, by weak approximation). For $v\notin S_k$ we use Haar
measures giving $K_v$, $A(\OO_v)$ and $T_\mu(\OO_v)$ volume $1$, and global
measures that are products of these with the measures at $S_k$; the
volumes below are taken with respect to them.

\begin{theorem}[Arthur~{\cite[Thm.~8.1]{A12}}; Li~\cite{LiI} for
covers]\label{thm:fine}
Let $f_S\in\mathcal H(\overline{G(\Q_S)})$ and $f=f_S\otimes e^S$. If
$S$ is large enough (depending on the support of $f_S$), then:
\begin{enumerate}[label=(\alph*)]
\item for $z\in\{\pm1\}$,
  \[
  J_{\mathcal O_z}(f)=a^G(S,1)f_S(z)+\sum_{\xi\in k_S^\times/k_S^{\times2}}
  a^G(S,u_\xi)\,J_G(z\,u_\xi,f_S)+\tfrac12\,a^A(S,1)\,J_A(z,f_S),
  \]
  with $a^G(S,1)=\mathrm{vol}(G(\Q)\backslash\GA)$, $a^A(S,1)=c_k$, and
  $J_A(z,f_S)$ as in Lemma~\ref{lem:JA-central}(b); the coefficients
  $a^G(S,u_\xi)$ are those of~\cite{A12,LiI} and are not evaluated in this
  paper;
\item for $\mathcal O$ elliptic with $\mu\in\mathcal O$,
  $J_{\mathcal O}(f)=a^G(S,\mu)J_G(\mu_S,f_S)$ with
  $a^G(S,\mu)=\mathrm{vol}(T_\mu(\Q)\backslash T_\mu(\A))$;
\item for $\mathcal O$ hyperbolic with $\mu=a(t)\in\mathcal O$,
  $J_{\mathcal O}(f)=a^A(S,\mu)J_A(\mu_S,f_S)$ with $a^A(S,\mu)=c_k$.
\end{enumerate}
\end{theorem}

In Arthur's notation the right-hand sides are
$\sum_{M\in\{A,G\}}|W^M_0||W^G_0|^{-1}\sum_{\gamma}a^M(S,\gamma)
J_M(\gamma,f_S)$, with $\gamma$ running over the $(M,S)$-classes in
$M(\Q)\cap\mathcal O$; in~(c) the factor $\tfrac12$ cancels against the
two elements $\mu^{\pm1}$, since $J_A(\mu^{-1},f_S)=J_A(\mu,f_S)$. (Substitute
$x=wy$, and use $v_A(wy)=-H(wy)-H(-y)=v_A(y)$.) Parts~(b) and~(c) follow from
Propositions~\ref{prop:elliptic} and~\ref{prop:hyperbolic}: for $S$
large, $|D(\mu)|_S=|D(\mu)|_{\A_k}=1$, and at $v\notin S_k$ the orbital
integrals of $e_v$ equal $1$ while $v_{A,v}=0$ on their support. For the
same reason $f(\gamma)=f_S(\gamma_S)$, which is why the lift $\mu_S$
carries the global factor $\kappa_\A(\mu)$. In~(a) the coefficient of
$f_S(z)$ and the linear term agree with Proposition~\ref{prop:central}.
The coefficients $a^G(S,u_\xi)$ exist by the theorem and are determined
by~(a), but we do not compute them (Remark~\ref{rem:sec4-status}).
Wherever they occur below, in Theorems~\ref{thm:tf} and~\ref{thm:itf}, they
are defined, not evaluated.

The coefficients are defined with centralizers in the cover, which for
$n=2$ are the full preimages (Lemma~\ref{lem:classes}), and the sums run
over $(G,S)$-classes, not over the rational class $\mathcal O$.

\begin{remarkrev}\label{rem:alpha}
\emph{Normalizations on $\mathfrak a$.} Identify $\mathfrak a$ with $\R$ by
$X\mapsto X(\chi_1)$, $\chi_1=N_{k/\Q}\circ t$, as in
Section~\ref{sec:chambers}. The coroot is $\alpha^\vee\colon s\mapsto
a(s)$, $s\in\Q^\times$, so $\alpha^\vee=\langle\chi_1,\alpha^\vee\rangle=
d$, and the simple root $\alpha$ of $(B,A_B)$, $a(s)\mapsto s^2$, is
$\alpha=(2/d)\chi_1$, i.e.\ $\alpha(X)=2X/d$. The modulus is $\delta_B=e^{2H}$ for every $k$. We use
Lebesgue measure $dX$ on $\mathfrak a$, for which $\mathfrak
a/\Z\alpha^\vee$ has volume $d$; Arthur normalizes it to volume $1$. With
Arthur's normalization, $\theta_B(\lambda)=\lambda(\alpha^\vee)$, the
weight $v_A$ and the function $r_A^G$ are divided by $d$, and
$\mathrm{vol}(A(\Q)\backslash A(\A)^1)$ is multiplied by $d$; the products
in~\eqref{eq:JT-hyp} and Theorem~\ref{thm:fine} do not change. With
$\theta_P(\lambda)$ in our normalization, Arthur's
$r^G_A(1,a)=\lim_{\lambda\to0}\sum_{P\in\mathcal P(A)}r_P(\lambda,1,a)
\theta_P(\lambda)^{-1}$~\cite[\S3]{A7} with $r_P(\lambda,1,a)=
|a^\beta-a^{-\beta}|^{\lambda(\beta^\vee)/2}$ evaluates to $r(t)$ for
$a=a(t)$, $t\in\Q_S^\times$. We have checked this only for the formula
of~\cite{A7} as quoted in \texttt{book7d.tex}, not against the general
definition of the exponents $\rho(\beta,u)$ there. The independent check is
that the limit in Lemma~\ref{lem:JA-central}(b) exists.
\end{remarkrev}

\subsection{Status of the results}\label{sec:status4}

\begin{remarkrev}\label{rem:sec4-status}
\emph{What rests on citation.} Theorem~\ref{thm:coarse} and
Propositions~\ref{prop:hyperbolic}--\ref{prop:central} are proved here.
Theorem~\ref{thm:fine}(a) is taken from~\cite{A12,LiI}: the coefficients
$a^G(S,u_\xi)$ are defined by it, and we have not derived~(a) from
Proposition~\ref{prop:central}. Doing so amounts to computing
$a^G(S,u_\xi)$ from the Laurent expansion at $s=1$ of the second form
of~\eqref{eq:Zf}, grouped by the image of $\xi$ in
$k_S^\times/k_S^{\times2}$. Each fiber is infinite, and its contribution
is a Dirichlet series in $s$, so a regularization is needed. This is
left to the revision of Section~7, where the same divergence occurs. The unramified computations used for (b)--(c) of
Theorem~\ref{thm:fine} are standard and not written out. For an
elliptic $\mu$ they use the uniqueness of the local class meeting $K_v$
from the proof of Theorem~\ref{thm:coarse}(a).
\end{remarkrev}

\begin{remarkrev}\label{rem:verified-sec4}
The script \texttt{checks/sec4\_n2\_check.py} verifies, for $k=\Q$: that
the Iwasawa height computed from an explicit decomposition equals
$-\log\|e_2g\|_v$ at $v=\infty,2,3,5,7$; the inequality
$H(x)+H(\gamma x)\le0$ on random adelic $x$ and $\gamma\in
SL(2,\Q)\smallsetminus B(\Q)$, and $H(\gamma x)=H(x)$ for $\gamma\in
B(\Q)$; Lemma~\ref{lem:height-ineq}(c) and Lemma~\ref{lem:weight-local};
the weight integral in the proof of Proposition~\ref{prop:hyperbolic};
the conjugation identities used there (symbolically);
Lemma~\ref{lem:JA-central}(b) at $v=\R$ (quadrature) and $v=\Q_3$
(exactly), including the divergence without $r(t)$; the symmetry
$J_A(\mu^{-1})=J_A(\mu)$ at $\R$; and Example~\ref{ex:unip-Q}.
\end{remarkrev}

% SPDX-License-Identifier: CC-BY-4.0
% =====================================================================
% Section 5 of D. Joyner, "A hitch-hiker's guide to the metaplectic trace
% formula for SL(2)", revised for the 2-fold cover (n = 2), November 2026.
%
% Stage 4 of ../revision_log.md. Sources: ../original-1995/cctrace4.tex
% (= ../book-1997/book7e.tex, Sections 5.1-5.2) and cctrace5.tex
% (= ../book-1997/book7f.tex, Sections 5.3-5.4). Referee items II.9-II.13,
% the Section 5 part of II.6 and the Section 5 minor items are fixed here;
% see the "Stage 4 notes" in ../revision_log.md, which also map the 1995
% numbers to the numbers here.
%
% This file is \input by the master; it compiles on its own through
% preview.tex. Explicit formulas are checked by checks/sec5_n2_check.py.
% Comments beginning "% [Claude]" are editorial notes by Claude.
% Revised after the second referee report (ChatGPT, November 2026): diagram
% of the objects, Lemma 5.7 (the Weil representation on the cover of A), which
% settles b' = -b in Theorem 5.8(c); comparisons with the 1995 notes moved to
% Appendix A.
% =====================================================================

\providecommand{\GA}{\overline{G(\A)}}

\section{The spectral side}\label{sec:spectral}

This section computes the spectral side of the trace formula for $\GA$,
class by class, and assembles the trace formula. The general theory is
due to Langlands~\cite{L2} and Arthur~\cite{A3,A14} for connected
reductive groups, and to M\oe glin--Waldspurger~\cite{MW2} and
Li~\cite{LiIII} for covering groups; as in Section~\ref{sec:geometric},
we state it for the double cover of $G=\Res_{k/\Q}SL(2)$ and make the
terms explicit. Three things are new compared with the 1995 notes. The
residual spectrum is not empty: it consists of one even Weil
representation for each class in $k^\times/k^{\times2}$
(Theorem~\ref{thm:residual}). These classes are exactly the cuspidal
data $(A,\mu)$ with $\mu^w=\mu$, which the 1995 notes omit
(Lemma~\ref{lem:w-action}, Theorem~\ref{thm:Jchi}(c)). And the operator
$M_B^T$ is $2T-M^{-1}M'$ (Theorem~\ref{thm:Jchi}), whose $T$-term we check
against the geometric side (Proposition~\ref{prop:Tcoef}). The
normalizing factor of the intertwining operator, $L^S(2\lambda,\chi^2)/
L^S(2\lambda+1,\chi^2)$, is computed in Lemma~\ref{lem:intertwining}.

\paragraph{Roadmap.} The objects attached to a genuine automorphic
character $\mu$ of the cover of $A$ are related as follows.
\[
\begin{array}{ccc}
\varphi\in\mathcal H_B(\mu)
 & \xrightarrow{\ \text{sum over }B(\Q)\backslash G(\Q)\ }
 & E(x,\varphi,\lambda)\\[6pt]
 & & \Big\downarrow{\scriptstyle\ \text{constant term}}\\[6pt]
\theta_\varphi=\mathrm{Res}_{\lambda=1/2}E(x,\varphi,\lambda)\in V_b
 & \xleftarrow{\ \text{pole at }\lambda=\frac12\ (\mu^w=\mu)\ }
 & \varphi_\lambda+(M(w,\lambda)\varphi)_{-\lambda}
\end{array}
\]
For $\lambda=it$ the Eisenstein series are not square integrable; their
wave packets span the continuous part $L^2_{\chi,\mathrm{cont}}$ of the
spectrum.
Section~\ref{sec:decomp} sorts the spectrum by cuspidal data $\chi$.
Section~\ref{sec:eisenstein} constructs the Eisenstein series $E$, whose
constant term defines the intertwining operator $M(w,\lambda)$. The operator factors as a quotient of partial
$L$-functions times finitely many local operators
(Lemma~\ref{lem:intertwining}); its only pole in $\mathrm{Re}\,\lambda>0$
is at $\lambda=\frac12$ and occurs exactly for the ramified classes, and
the residues of $E$ there span the residual spectrum, which consists of
even Weil representations (Theorem~\ref{thm:residual}).
Section~\ref{sec:Jchi} truncates the kernels as on the geometric side and
computes the integral $J^T_\chi(f)$ of each truncated kernel from the
inner products of truncated Eisenstein series (the Maass--Selberg
relations): the diagonal terms give $2T-M^{-1}M'$, and for a ramified
class the cross terms give $\frac14\operatorname{tr}(M(w,0)I(\mu_0,f))$.
Section~\ref{sec:tf} equates the two sides.

\paragraph{Conventions.} $L^2_{\mathrm{gen}}=L^2(G(\Q)\backslash\GA)_
{\mathrm{gen}}$ is the space of $\varphi$ with $\varphi(\gamma x\zeta)=
\zeta\varphi(x)$, with the measure of Section~\ref{sec:global} (total
volume $1$), and $R$ is the right regular representation; test functions
$f\in\mathcal H(\GA)$ are anti-genuine and $R(f)=\int f(y)R(y)\,dy$, as in
Section~\ref{sec:kernels}. As in Section~\ref{sec:chambers} and
Remark~\ref{rem:alpha}, $\mathfrak a\cong\R$ by $X\mapsto X(\chi_1)$, with
Lebesgue measure $dX$, and we identify $\mathfrak a^*_\C$ with $\C$ so that
$\lambda(X)=\lambda X$. Then $e^{\lambda(H(x))}=e^{\lambda H(x)}$,
\[
\rho=\rho_B=1,\qquad\alpha=2/d,\qquad\alpha^\vee=d,\qquad d=[k:\Q],
\]
since $\delta_B=e^{2H}$, and $\lambda\in i\R$ is written $\lambda=it$ with
the measure $dt/2\pi$, dual to $dX$. For $k=\Q$ and $x=n(u)a(t)k$ the
point $x\cdot i$ of the upper half plane has imaginary part $y=t^2$, so
$e^{(\lambda+1)H(x)}=y^{s}$ with $s=(\lambda+1)/2$. We write $c_k=
\mathrm{vol}(A(\Q)\backslash A(\A)^1)$ as in Section~\ref{sec:coarse}, and
$X_A=A(\Q)\backslash\overline{A(\A)^1}$ and $\Pi_{\mathrm{aut}}(X_A)$ as
in Section~\ref{sec:torus-cover}; $\Pi_u(H)$ is the set of irreducible
unitary genuine representations of $H$, as there.

\subsection{The spectral decomposition and cuspidal data}\label{sec:decomp}

\begin{lemma}\label{lem:spec-decomp}
$L^2_{\mathrm{gen}}=L^2_{\mathrm{disc}}\oplus L^2_{\mathrm{cont}}$, where
$L^2_{\mathrm{disc}}$ is the closed span of the irreducible closed
invariant subspaces. Each $\pi\in\Pi_u(\GA)$ occurs in
$L^2_{\mathrm{disc}}$ with finite multiplicity $m(\pi)$, and
$L^2_{\mathrm{disc}}=L^2_{\mathrm{cusp}}\oplus L^2_{\mathrm{res}}$, where
$L^2_{\mathrm{cusp}}$ is the space of cusp forms,
\[
\int_{N(\Q)\backslash N(\A)}\varphi(nx)\,dn=0\quad\text{for almost all }x,
\]
and $L^2_{\mathrm{res}}$ is its orthogonal complement in
$L^2_{\mathrm{disc}}$.
\end{lemma}

\begin{proof}
By Lemma~\ref{lem:strong}, $G(\Q)\backslash\GA/K_f^*\cong
\widetilde\Gamma\backslash\overline{G(\R)}$ with $\widetilde\Gamma$ a
lattice, so this is Langlands' theory~\cite{L2} for
$\widetilde\Gamma\backslash\overline{G(\R)}$; in adelic form,
see~\cite{MW2}, whose setting includes finite central covers. The finiteness of multiplicities in
$L^2_{\mathrm{disc}}$ is Harish-Chandra's finiteness theorem for
automorphic forms.
\end{proof}

\begin{lemma}\label{lem:factor}
Every $\pi\in\Pi_u(\GA)$ is a restricted tensor product
$\pi\cong\bigotimes'_v\pi_v$, with $\pi_v\in\Pi_u(\overline{SL(2,k_v)})$
genuine and $\pi_v$ having a nonzero $K_v^*$-fixed vector for almost all
$v$; the restricted product is taken with respect to these vectors.
Conversely every such family defines an element of $\Pi_u(\GA)$.
\end{lemma}

\begin{proof}
A representation of $\GA$ is a representation of the restricted product
of the $\overline{SL(2,k_v)}$ (with respect to the $K_v^*$) that is
trivial on $Z_0=\{(\zeta_v)\mid\prod_v\zeta_v=1\}$
(Section~\ref{sec:global}). Flath's theorem~\cite{F} applies to the
restricted product and gives $\pi\cong\bigotimes'\pi_v$. Each $\mu_2$ at
$v$ acts by a sign $\varepsilon_v$; triviality on $Z_0$ forces all
$\varepsilon_v$ to be equal, and genuineness forces them to be the
identity. Conversely, if every $\pi_v$ is genuine, $\bigotimes'\pi_v$ is
trivial on $Z_0$ and descends.
\end{proof}

\paragraph{Cuspidal data.} The \emph{cuspidal data} are the pairs
$(G,\pi)$, with $\pi\in\Pi_u(\GA)$ occurring in $L^2_{\mathrm{cusp}}$, and
the pairs $(A,\mu)$, with $\mu\in\Pi_{\mathrm{aut}}(X_A)$ a genuine
character of $X_A$ (Theorem~\ref{thm:corr}). Among the standard
parabolic subgroups, $B$ is associate only to itself, with
$W(\mathfrak a)=\Omega(\mathfrak a,\mathfrak a)=\{1,w\}$. (The opposite
$\bar B=wBw^{-1}$ also has Levi $A$ and is associate to $B$, but it is not
standard.) For $\mu$ as above put
$\mu^w(\tilde a)=\mu(w^{-1}\tilde aw)$, where $w$ stands for its image
$(w,\kappa_\A(w))=(w,1)$ in $\GA$. Two data are \emph{equivalent} if they
are equal or of the form $(A,\mu)$, $(A,\mu^w)$. Let $\mathcal X$ be the
set of classes, $\mathcal X_G$ the classes of the $(G,\pi)$ and
$\mathcal X_A$ those of the $(A,\mu)$; for $\chi\in\mathcal X_A$ we write
$\mu\in\chi$ for the members of the set $\{\mu,\mu^w\}$, which has one or
two elements. A class $\chi\in\mathcal X_A$ is \emph{unramified} if
$\mu^w\ne\mu$, and \emph{ramified} if $\mu^w=\mu$~\cite[\S4]{A3}.

\begin{lemma}\label{lem:w-action}
\begin{enumerate}[label=(\alph*)]
\item For $t\in\A_k^\times$ and $\zeta\in\mu_2$,
  $(w,1)\,(a(t),\zeta)\,(w,1)^{-1}=(a(t^{-1}),\zeta)$ in $\GA$, and the same
  holds in each $\overline{SL(2,k_v)}$.
\item $(a(t),1)^{-1}=(a(t^{-1}),\prod_v(t_v,-1)_v)$; so $w$-conjugation is
  not inversion on $\overline{A(\A)}$, and
  $\mu^w(a(t),1)=\mu(a(t),1)^{-1}\prod_v(t_v,-1)_v$.
\item $\chi_\psi^w=\chi_\psi$ for Weil's character $\chi_\psi$
  (Lemma~\ref{lem:weil}). Hence, writing $\mu=\chi_\psi\cdot(\chi\circ p)$
  with $\chi$ a character of $A(\Q)\backslash A(\A)^1$
  (Theorem~\ref{thm:corr}), $\mu^w=\chi_\psi\cdot(\chi^{-1}\circ p)$.
\item $\mu^w=\mu$ if and only if $\chi^2=1$, and then
  $\mu=\chi_{\psi_b}$ on $\overline{A(\A)^1}$ for a unique $b\in
  k^\times/k^{\times2}$. So the ramified classes in $\mathcal X_A$ are the
  classes $\chi_b$ of $(A,\chi_{\psi_b})$, $b\in k^\times/k^{\times2}$,
  each with one member, and every other class in $\mathcal X_A$ has two
  members.
\end{enumerate}
\end{lemma}

\begin{proof}
(a) At one place, with $x(\cdot)$ as in~\eqref{eq:beta}: $x(w)=1$,
$x(a(t))=t^{-1}$, $wa(t)=\begin{pmatrix}0&-t^{-1}\\t&0\end{pmatrix}$ has
$x=t$, $wa(t)w^{-1}=a(t^{-1})$ has $x=t$, and $x(w^{-1})=-1$. So by the
conjugation formula~\eqref{eq:inverse},
$\beta(w,a(t))\beta(wa(t),w^{-1})\beta(w,w^{-1})=(t,t^2)(1,-t)(1,-1)=1$.
At non-dyadic finite $v$ the correction $\kappa_v$ is $1$ on $w$,
$w^{-1}$, $a(t)$, $a(t^{-1})$ and $wa(t)$ (the lower left entry is $0$ or
of absolute value $\ge1$), so $\sigma_v=\beta_v$ on these pairs. Taking
the product over $v$ gives the global statement.

(b) $(a(t),1)(a(t^{-1}),\varepsilon)=(1,\sigma_\A(a(t),a(t^{-1}))
\varepsilon)$ and $(t_v,t_v^{-1})_v=(t_v,t_v)_v=(t_v,-1)_v$ by~\ref{H1}
and~\ref{H2}. The formula for $\mu^w$ follows from~(a).

(c) $t^{-1}=t\cdot t^{-2}$ is in the square class of $t$, so
$\gamma_v(t^{-1})=\gamma_v(t)$ by Lemma~\ref{lem:weil}(a), and
$\chi_\psi(a(t^{-1}),\zeta)=\chi_\psi(a(t),\zeta)$. By~(a),
$\mu^w(a(t),\zeta)=\mu(a(t^{-1}),\zeta)=\chi_\psi(a(t),\zeta)\,
\chi(a(t))^{-1}$.

(d) By~(c), $\mu^w=\mu$ if and only if $\chi^2=1$. A quadratic character
of $k^\times\backslash\A_k^1$ extends uniquely to a quadratic character of
$k^\times\backslash\A_k^\times$ trivial on $\R_{>0}$ (embedded diagonally
at the archimedean places). By class field theory the quadratic
characters of $k^\times\backslash\A_k^\times$ are the
$t\mapsto\prod_v(b,t_v)_v$, $b\in k^\times$, and two values of $b$ give
the same character exactly when they differ by an element of
$k^\times\cap\prod_vk_v^{\times2}=k^{\times2}$; each such character is
trivial on $\R_{>0}$, since $(b,u)_v=1$ for $u>0$ at a real $v$. By
Lemma~\ref{lem:weil}(d), $\chi_\psi\cdot(b,\cdot)=\chi_{\psi_b}$.
\end{proof}

For $SL(2)$ itself the ramified classes are $(A,\mu)$ with $\mu^2=1$,
because there $w$ acts on $A$ by inversion~\cite{LL}. For the cover,
inversion is twisted by $(t,-1)$, and Weil's character absorbs the twist.

\begin{example}\label{ex:ramified-Q}
Let $k=\Q$. Then $k^\times/k^{\times2}$ is represented by the squarefree
integers $b$, positive and negative, and the ramified classes are the
$\chi_{\psi_b}$. If $f=f_S\otimes e^S$ as in Section~\ref{sec:local-dist},
only the classes with $\chi_{\psi_b,p}$ trivial on $A(\Z_p)$ for every
$p\notin S$ can contribute; since $(b,u)_p=1$ for all units $u$ exactly
when $\ord_p(b)$ is even, these are the $b$ whose prime factors lie in
$S$. For $S=\{\infty,2\}$ they are $b\in\{\pm1,\pm2\}$.
\end{example}

\begin{proposition}\label{prop:Lchi}
$L^2_{\mathrm{gen}}=\widehat\bigoplus_{\chi\in\mathcal X}L^2_\chi$, an
orthogonal decomposition into invariant subspaces, where:
\begin{enumerate}[label=(\alph*)]
\item for $\chi$ the class of $(G,\pi)$, $L^2_\chi$ is the
  $\pi$-isotypic subspace of $L^2_{\mathrm{cusp}}$, of multiplicity
  $m(\pi)$;
\item for $\chi\in\mathcal X_A$ with $\mu\in\chi$,
  $L^2_\chi=L^2_{\chi,\mathrm{res}}\oplus L^2_{\chi,\mathrm{cont}}$, where
  $L^2_{\chi,\mathrm{cont}}$ is the closed span of the wave packets
  $\int_\R E(x,\Phi(t),it)\,dt$ of the Eisenstein series of
  Section~\ref{sec:eisenstein}, $\Phi\in C_c^\infty(\R,\mathcal H_B(\mu))$
  with values in a finite sum of $\overline{K_0}$-types, and
  $L^2_{\chi,\mathrm{res}}=L^2_{\mathrm{res}}\cap L^2_\chi$ is spanned by
  residues of these Eisenstein series at poles with $\mathrm{Re}\,\lambda
  >0$. In particular $L^2_{\mathrm{cusp}}=\widehat\bigoplus_{\chi\in
  \mathcal X_G}L^2_\chi$.
\end{enumerate}
\end{proposition}

This is Langlands' decomposition~\cite[Ch.~7]{L2} along cuspidal
data~\cite[Ch.~II, V]{MW2}. By Lemma~\ref{lem:w-action}(d) and
Lemma~\ref{lem:intertwining}(c) below, $L^2_{\chi,\mathrm{res}}=0$ unless
$\chi$ is ramified.

\subsection{Eisenstein series and intertwining operators}\label{sec:eisenstein}

For $\mu\in\Pi_{\mathrm{aut}}(X_A)$ let $\mathcal H_B(\mu)$ be the space of
genuine functions $\varphi$ on $N(\A)A(\Q)A_B(\R)^0\backslash\GA$ that are
right $\overline{K_0}$-finite and satisfy
\begin{equation}\label{eq:HBmu}
\varphi(\tilde a x)=\mu(\tilde a)\varphi(x)\quad(\tilde a\in
\overline{A(\A)^1}),\qquad
\langle\varphi,\varphi'\rangle=\int_{X_A}\int_{\overline{K_0}}
\varphi(\tilde a\tilde k)\overline{\varphi'(\tilde a\tilde k)}\,d\tilde k\,
d\tilde a=c_k\int_{\overline{K_0}}\varphi\overline{\varphi'}.
\end{equation}
Here $A(\Q)$ and $N(\A)$ are embedded by $b\mapsto(b,1)$
(Lemma~\ref{lem:global}(b)), $\mu$ is extended to $\overline{A(\A)}$
trivially on $A_B(\R)^0$, and restriction to $\overline{K_0}$ is injective
by Lemma~\ref{lem:NAK-global}. For $\lambda\in\C$ put
$\varphi_\lambda(x)=\varphi(x)e^{(\lambda+1)H(x)}$, and let
$I(\mu_\lambda)$ be the representation on $\mathcal H_B(\mu)$ given by
$(I(\mu_\lambda,y)\varphi)_\lambda(x)=\varphi_\lambda(xy)$, and
$I(\mu_\lambda,f)=\int_{\GA}f(y)I(\mu_\lambda,y)\,dy$. It is the induced
representation $\mathrm{Ind}_{\overline{B(\A)}}^{\GA}(\mu\,e^{\lambda H})$
(normalized induction), unitary for $\lambda\in i\R$. Note that
$\varphi$ is $A_B(\R)^0$-invariant and that $\lambda$ enters only through
$\varphi_\lambda$.

For $\varphi\in\mathcal H_B(\mu)$ and $\mathrm{Re}\,\lambda>1$ put
\begin{gather}
E(x,\varphi,\lambda)=\sum_{\delta\in B(\Q)\backslash G(\Q)}
\varphi_\lambda(\delta x),\label{eq:Eis}\\
\bigl(M(w,\lambda)\varphi\bigr)_{-\lambda}(x)=\int_{N(\A)}
\varphi_\lambda(w^{-1}nx)\,dn .\label{eq:Mw}
\end{gather}
The sum is well defined because $\varphi_\lambda$ is left
$B(\Q)$-invariant (Lemma~\ref{lem:NAK-global}(b), and $\mu$ is trivial
on $A(\Q)$). By Lemma~\ref{lem:w-action}(a), $M(w,\lambda)$ maps
$\mathcal H_B(\mu)$ to $\mathcal H_B(\mu^w)$, and it intertwines
$I(\mu_\lambda)$ with $I(\mu^w_{-\lambda})$. The constant term of a
function $\phi$ on $G(\Q)\backslash\GA$ is $\phi_B(x)=\int_{N(\Q)
\backslash N(\A)}\phi(nx)\,dn$.

\begin{theorem}[Langlands~\cite{L2}; M\oe glin--Waldspurger~\cite{MW2}
for covers]\label{thm:eisenstein}
Let $\mu\in\Pi_{\mathrm{aut}}(X_A)$ and $\varphi\in\mathcal H_B(\mu)$.
\begin{enumerate}[label=(\alph*)]
\item The series~\eqref{eq:Eis} and the integral~\eqref{eq:Mw} converge
  absolutely for $\mathrm{Re}\,\lambda>1$, and $E(x,\varphi,\lambda)$
  and $M(w,\lambda)\varphi$ continue to meromorphic functions of
  $\lambda\in\C$.
\item $E_B(x,\varphi,\lambda)=\varphi_\lambda(x)+(M(w,\lambda)
  \varphi)_{-\lambda}(x)$.
\item $E(x,M(w,\lambda)\varphi,-\lambda)=E(x,\varphi,\lambda)$ and
  $M(w,-\lambda)M(w,\lambda)=1$.
\item On $i\R$, $E(x,\varphi,\lambda)$ and $M(w,\lambda)$ are holomorphic
  and $M(w,\lambda)$ is unitary. In $\mathrm{Re}\,\lambda>0$ they have at
  most finitely many poles, all simple and in $(0,1]$, and the poles of
  $E$ are those of $M(w,\lambda)$.
\item $E(x,I(\mu_\lambda,f)\varphi,\lambda)=\int_{\GA}f(y)
  E(xy,\varphi,\lambda)\,dy$, and $M(w,\lambda)I(\mu_\lambda,f)=
  I(\mu^w_{-\lambda},f)M(w,\lambda)$.
\end{enumerate}
\end{theorem}

The measure on $i\mathfrak a^*$ is $dt/2\pi$, and the number $n(A_B)$ of
chambers in $\mathfrak a$ is $2$; it produces the factor $\frac12$ in
Proposition~\ref{prop:kernels-spec}. The convergence of~\eqref{eq:Eis} for $\mathrm{Re}\,\lambda>\rho=1$ is
Godement's criterion; the integral~\eqref{eq:Mw} in fact converges for
$\mathrm{Re}\,\lambda>\frac12$, by Lemma~\ref{lem:intertwining}.

\begin{lemma}\label{lem:intertwining}
Let $\mu=\chi_\psi\cdot(\chi\circ p)\in\Pi_{\mathrm{aut}}(X_A)$. Let $S$
be a finite set of places of $k$ containing the archimedean and dyadic
places and those where $\psi_v$ has conductor $\ne\OO_v$ or $\chi_v$ is
ramified. For $v\notin S$ let $\varphi_v^0\in I(\mu_{v,\lambda})$ be the
$K_v^*$-fixed vector with $\varphi_v^0=1$ on $K_v^*$, and let $M_v(\lambda)$
be the local intertwining operator, $(M_v(\lambda)\varphi_v)(x)=
\int_{k_v}\varphi_{v,\lambda}(w^{-1}n(y)x)\,dy$ with the self-dual
measure.
\begin{enumerate}[label=(\alph*)]
\item For $\varphi=\varphi_S\otimes\bigotimes_{v\notin S}\varphi_v^0$,
  $M(w,\lambda)\varphi=(M_S(\lambda)\varphi_S)\otimes\bigotimes_{v\notin S}
  M_v(\lambda)\varphi_v^0$, with $M_S=\bigotimes_{v\in S}M_v$.
\item For $v\notin S$, with $q=q_v$ and $\chi_v(\varpi)$ the value at a
  uniformizer,
  \[
  M_v(\lambda)\varphi_v^0=\frac{1-\chi_v(\varpi)^2q^{-1-2\lambda}}
  {1-\chi_v(\varpi)^2q^{-2\lambda}}\,\varphi_v^{0,w}
  =\frac{L_v(2\lambda,\chi_v^2)}{L_v(2\lambda+1,\chi_v^2)}\,
  \varphi_v^{0,w},
  \]
  where $\varphi_v^{0,w}$ is the $K_v^*$-fixed vector of
  $I(\mu^w_{v,-\lambda})$. Hence $M(w,\lambda)=m^S(\mu,\lambda)\,M_S(\lambda)
  \otimes1$ on such $\varphi$, with
  \begin{equation}\label{eq:mS}
  m^S(\mu,\lambda)=\frac{L^S(2\lambda,\chi^2)}{L^S(2\lambda+1,\chi^2)} .
  \end{equation}
\item In $\mathrm{Re}\,\lambda>0$, $M(w,\lambda)$ is holomorphic unless
  $\chi^2=1$ ($\chi$ extended to $k^\times\backslash\A_k^\times$ trivially
  on $\R_{>0}$). If $\chi^2=1$, its only pole there is a simple pole at
  $\lambda=\frac12$, with residue $\mathrm{Res}_{s=1}\zeta_k^S(s)\,
  \bigl(2\zeta_k^S(2)\bigr)^{-1}M_S(\tfrac12)\otimes1$, which is not $0$.
\end{enumerate}
\end{lemma}

\begin{proof}
(a) The lift $(w,1)$ of $w\in G(\Q)$ is the product of the local lifts
$(w,1)_v$ ($\kappa_v(w)=1$ for all $v$), $N(\A)$ is the restricted
product of the $N(k_v)$ with $dn=\prod dn_v$, and a factorizable
$\varphi$ is a product of local functions (Lemma~\ref{lem:factor}).

(b) Since $\chi_{\psi,v}$ is trivial on $A(\OO_v)^*$
(Lemma~\ref{lem:weil}(b)), $\varphi_v^0((n,1)(a(t),\zeta)\tilde k)=\zeta
\mu_v(a(t),1)|t|^{\lambda+1}$ for $\tilde k\in K_v^*$ is well defined. If
$|y|\le1$, then $w^{-1}n(y)\in K_v$, its lift lies in $K_v^*$ because
$\sigma_v\equiv1$ on $K_v\times K_v$, and the integrand is $1$. If
$|y|=q^m$ with $m\ge1$, write $y=\varpi^{-m}u$, $u\in\OO^\times$; then
\[
w^{-1}n(y)=n(-y^{-1})\,a(-y^{-1})\,\begin{pmatrix}1&0\\y^{-1}&1
\end{pmatrix},
\]
the last factor in $K_v$. Computing the cocycle of this decomposition, the
integrand is $\chi_v(\varpi)^m|y|^{-\lambda-1}$ times $1$ if $m$ is even,
and times $c\,(u,\varpi)_v$ with a constant $c$ of absolute value $1$ if
$m$ is odd. (The odd case uses $\gamma_v(\varpi u)=(u,\varpi)_v
\gamma_v(\varpi)$, Lemma~\ref{lem:weil}(a)--(b); both cases are checked in
\texttt{checks/sec5\_n2\_check.py} at $p=3,5,7$.) The shell
$\{|y|=q^m\}$ has volume $q^m(1-q^{-1})$, and $u\mapsto(u,\varpi)_v$ is a
nontrivial character of $\OO^\times$, so the odd shells contribute $0$
and
\[
M_v(\lambda)\varphi_v^0(1)=1+(1-q^{-1})\sum_{j\ge1}\bigl(\chi_v(\varpi)^2
q^{-2\lambda}\bigr)^j=\frac{1-\chi_v(\varpi)^2q^{-1-2\lambda}}
{1-\chi_v(\varpi)^2q^{-2\lambda}} .
\]
$M_v(\lambda)\varphi_v^0$ is $K_v^*$-fixed, so it is this multiple of
$\varphi_v^{0,w}$.

(c) The local integrals at $v\in S$ converge absolutely for
$\mathrm{Re}\,\lambda>0$, since the integrand is $O(|y|^{-\mathrm{Re}\,
\lambda-1})$, so $M_S(\lambda)$ is holomorphic there. If $\chi^2\ne1$,
$L^S(s,\chi^2)$ is entire (the removed Euler factors have no poles in
$\mathrm{Re}\,s>0$) and $L^S(s,\chi^2)\ne0$ for $\mathrm{Re}\,s\ge1$.
If $\chi^2=1$, $m^S=\zeta_k^S(2\lambda)/\zeta_k^S(2\lambda+1)$, whose only
pole in $\mathrm{Re}\,\lambda>0$ is at $2\lambda=1$. The operator
$M_S(\frac12)$ is not $0$: if $\varphi_S$ is supported in
$\overline{B(\Q_S)}\,w^{-1}N(U)$ for a small compact neighborhood $U$ of
$0$, then $(M_S(\frac12)\varphi_S)_{-1/2}(1)=\int_U\varphi_{S,1/2}
(w^{-1}n(y))\,dy$, and on the big cell $\varphi_S$ may be any smooth
compactly supported function of $y$.
\end{proof}

For $SL(2)$ itself the same computation gives
$L_v(\lambda,\chi_v)/L_v(\lambda+1,\chi_v)$, with the pole at
$\lambda=1=\rho$ for $\chi=1$ (the script checks this too). On the double
cover the variable is doubled and the character squared; this is the
Gindikin--Karpelevich formula for the double cover~\cite{KP}; for
general covering groups see~\cite{McN,Gao}. The factor~\eqref{eq:mS} is
what produces the logarithmic derivative $m'/m$ in Section~6.

\paragraph{The Weil representation.} The residual spectrum consists of
Weil representations, and to say which ones we fix a normalization. Let
$v$ be a place of $k$ and $\psi'$ a nontrivial character of $k_v$. The
\emph{Weil representation} $\omega_{\psi'}$ attached to $\psi'$ and the
quadratic form $x^2$~\cite{We64,Ge2} is a genuine representation of
$\overline{SL(2,k_v)}$ on the Schwartz--Bruhat space $\mathcal S(k_v)$ in
which
\begin{equation}\label{eq:weil-rep}
\begin{gathered}
\omega_{\psi'}(n(u),1)\phi(x)=\psi'(ux^2)\phi(x),\qquad
\omega_{\psi'}(w,1)\phi(x)=c\,(\mathcal F\phi)(-x),\\
\mathcal F\phi(x)=\int_{k_v}\phi(y)\psi'(2xy)\,dy,
\end{gathered}
\end{equation}
with a constant $c$ of absolute value $1$ and the measure that is
self-dual for $\psi'(2xy)$. (In Kubota's model $c=\gamma'(1)^{-1}$, with
$\gamma'$ as in Lemma~\ref{lem:weil-rep}; see Lemma~\ref{lem:weil-kernel}.)
It is the sum of its restrictions
$r^+_{\psi'}$ and $r^-_{\psi'}$ to the even and the odd functions, which
are irreducible. In some references $n(u)$ acts by $\psi'(\frac12ux^2)$;
that representation is $\omega_{\psi'_{1/2}}$ in the present notation.

\begin{lemma}\label{lem:weil-rep}
Let $\omega=\omega_{\psi'}$ be as in~\eqref{eq:weil-rep}, and let
$\gamma'(a)$ be the Weil index of $x\mapsto\psi'(ax^2)$, as
in~\eqref{eq:weil-index}. In the model of Section~\ref{sec:local}, for
$t\in k_v^\times$ and $\zeta\in\mu_2$,
\[
\omega(a(t),\zeta)\phi(x)=\zeta\,\frac{\gamma'(1)}{\gamma'(t)}\,
|t|_v^{1/2}\phi(tx)=\zeta\,\chi_{\psi'_{-1}}(a(t),1)\,|t|_v^{1/2}\phi(tx),
\]
where $\psi'_{-1}(x)=\psi'(-x)$. Consequently $\phi\mapsto\bigl(x\mapsto
(\omega(x)\phi)(0)\bigr)$ maps $r^+_{\psi'}$ into the induced
representation $I(\chi_{\psi'_{-1},v},-\frac12)$: the function
$F(x)=(\omega(x)\phi)(0)$ satisfies $F((n,1)\tilde ax)=\chi_{\psi'_{-1}}
(\tilde a)e^{H_v(\tilde a)/2}F(x)$.
\end{lemma}

\begin{proof}
Let $\bar n(u)=w\,n(-u)\,w^{-1}$, the lower triangular unipotent matrix
with entry $u$. By Lemma~\ref{lem:local}(h),
$(\bar n(u),1)=(w,1)(n(-u),1)(w,1)^{-1}$, so by~\eqref{eq:weil-rep}
$\omega(\bar n(u),1)=\mathcal F m(-u)\mathcal F^{-1}$, where $m(u)$ is
multiplication by $\psi'(ux^2)$; the constant $c$ and the reflection
$x\mapsto-x$ of~\eqref{eq:weil-rep} cancel, because the reflection
commutes with $m(u)$. Completing the square and using~\eqref{eq:weil-index} for
$\psi'$ (with $\gamma'(1/t)=\gamma'(t)$; the measure self-dual for
$\psi'(2xy)$ is $|2|_v^{1/2}$ times the one in~\eqref{eq:weil-index}),
\[
\bigl(\mathcal Fm(1/t)\mathcal F^{-1}\phi\bigr)(x)=\gamma'(t)\,|t|_v^{1/2}
\int_{k_v}\phi(z)\,\psi'\bigl(-t(x-z)^2\bigr)\,dz .
\]
Let $w(t)=n(t)\bar n(-t^{-1})n(t)=\begin{pmatrix}0&t\\-t^{-1}&0
\end{pmatrix}$. Multiplying by $m(t)$ on both sides, the terms
$\psi'(tx^2)$ and $\psi'(tz^2)$ cancel against the square, and
\[
\omega(n(t),1)\,\omega(\bar n(-t^{-1}),1)\,\omega(n(t),1)\phi(x)
=\gamma'(t)\,|t|_v^{1/2}\,(\mathcal F\phi)(tx).
\]
Now compute in the cover with~\eqref{eq:beta}. The values of $x(\cdot)$
on $n(t)$, $\bar n(-t^{-1})$, $n(t)\bar n(-t^{-1})$ and $w(t)$ are $1$,
$-t^{-1}$, $-t^{-1}$, $-t^{-1}$, so both cocycles in the triple product are
$1$ and $(n(t),1)(\bar n(-t^{-1}),1)(n(t),1)=(w(t),1)$. Next
$(w(1),1)^{-1}=(w(1)^{-1},1)$, and $\beta(w(t),w(1)^{-1})=
\bigl(\frac{t^{-1}}{-t^{-1}},t^{-1}\bigr)_v=(-1,t)_v$, because
$w(t)w(1)^{-1}=a(t)$ has $x=t^{-1}$. Hence $(w(t),1)(w(1),1)^{-1}=
(a(t),(-1,t)_v)$ and
\[
\omega\bigl(a(t),(-1,t)_v\bigr)\phi(x)=\frac{\gamma'(t)}{\gamma'(1)}\,
|t|_v^{1/2}\phi(tx).
\]
By Lemma~\ref{lem:weil}(a) with $a=b=t$, $(\gamma'(t)/\gamma'(1))^2=
(t,t)_v=(t,-1)_v$. So $(-1,t)_v\gamma'(t)/\gamma'(1)=(\gamma'(t)/
\gamma'(1))^3=\gamma'(1)/\gamma'(t)$, which is the first formula. The
second is Lemma~\ref{lem:weil}(d) with $b=-1$:
$\chi_{\psi'_{-1}}=\chi_{\psi'}\cdot(-1,\cdot)_v=\chi_{\psi'}^{-1}$. For the
last statement, $(\omega((n(u),1)x)\phi)(0)=\psi'(0)(\omega(x)\phi)(0)$,
and $e^{(-\frac12+1)H_v(\tilde a)}=|t|_v^{1/2}$.
\end{proof}

For example, let $v$ be real and $\psi'(x)=e^{2\pi iax}$. The operator
$d\omega\begin{pmatrix}0&1\\-1&0\end{pmatrix}=i\bigl(2\pi ax^2-
\frac1{8\pi a}\frac{d^2}{dx^2}\bigr)$ generates the action of the
subgroup $\tilde k_\theta$ of Lemma~\ref{lem:explicit-r}(b); for $a>0$ its
eigenvalues on even functions are $i(\frac12+2j)$, $j\ge0$, with ground
state $e^{-2\pi ax^2}$, and for $a<0$ they are $-i(\frac12+2j)$. So
$r^+_{\psi'}$ has lowest weight $\frac12$ when $a>0$ and highest weight
$-\frac12$ when $a<0$. This agrees with the lemma:
$\chi_{\psi'_{-1}}(a(-1),1)=i$ for $a>0$, which is the condition
$m_0=\frac12$ of Lemma~\ref{lem:explicit-r}(b).

Globally, for a nontrivial character $\psi'$ of $\A_k/k$ let
$\omega_{\psi'}=\bigotimes_v\omega_{\psi'_v}$ on $\mathcal S(\A_k)$, a
genuine representation of $\GA$, and for $\phi\in\mathcal S(\A_k)$ let
\begin{equation}\label{eq:theta}
\theta^{\psi'}_\phi(x)=\sum_{\xi\in k}\bigl(\omega_{\psi'}(x)\phi\bigr)(\xi).
\end{equation}
By Weil~\cite{We64}, $\theta^{\psi'}_\phi$ is left invariant under
$G(\Q)$. By~\eqref{eq:weil-rep} its constant term is
$x\mapsto(\omega_{\psi'}(x)\phi)(0)$, because $\int_{k\backslash\A_k}
\psi'(u\xi^2)\,du=0$ for $\xi\ne0$.

\begin{theorem}\label{thm:residual}
\begin{enumerate}[label=(\alph*)]
\item $L^2_{\chi,\mathrm{res}}=0$ for every $\chi\in\mathcal X_A$ that is
  unramified.
\item For $b\in k^\times$, $V_b:=L^2_{\chi_b,\mathrm{res}}$ is the closed
  span of the residues $\theta_\varphi(x)=\mathrm{Res}_{\lambda=1/2}
  E(x,\varphi,\lambda)$, $\varphi\in\mathcal H_B(\chi_{\psi_b})$; it is
  not $0$, and $\theta_{\varphi,B}(x)=(\mathrm{Res}_{\lambda=1/2}
  M(w,\lambda)\varphi)_{-1/2}(x)$, of order $e^{H(x)/2}$.
\item $L^2_{\mathrm{res}}=\widehat\bigoplus_{b\in k^\times/k^{\times2}}V_b$,
  and each $V_b$ is irreducible. With the
  normalization~\eqref{eq:weil-rep},
  \[
  V_b\cong r^+_{\psi_{-b}}=\bigotimes\nolimits'_v r^+_{\psi_{-b,v}},
  \]
  the even Weil representation attached to $\psi_{-b}(x)=\psi(-bx)$, and
  the elements of $V_b$ are the theta functions
  $\theta^{\psi_{-b}}_\phi$ of~\eqref{eq:theta}, $\phi$ an even
  Schwartz--Bruhat function on $\A_k$.
\end{enumerate}
\end{theorem}

\begin{proof}
(a) and the first part of (b): by Proposition~\ref{prop:Lchi}(b) and
Theorem~\ref{thm:eisenstein}(d), $L^2_{\chi,\mathrm{res}}$ is spanned by
residues of $E$ at poles of $M(w,\lambda)$ with $\mathrm{Re}\,\lambda>0$,
and by Lemma~\ref{lem:intertwining}(c) there are none unless
$\chi^2=1$, i.e.\ $\mu=\chi_{\psi_b}$ (Lemma~\ref{lem:w-action}(d)), and
then only $\lambda=\frac12$. The constant term of the residue is the
residue of Theorem~\ref{thm:eisenstein}(b), in which $\varphi_\lambda$ is
holomorphic; it is nonzero by Lemma~\ref{lem:intertwining}(c), and it is
$O(e^{(-\frac12+1)H(x)})$. Since $\int_T^\infty e^{X}e^{-2X}\,dX<\infty$
for the measure $\delta_B^{-1}=e^{-2X}$, $\theta_\varphi$ is square
integrable (its difference with the constant term is rapidly decreasing
on a Siegel set).

(c) By Gelbart~\cite{Ge2}, Gelbart and
Piatetski-Shapiro~\cite{GPS80} and Waldspurger~\cite{Wa80},
$L^2_{\mathrm{res}}$ is the orthogonal sum, over the classes $b'\in
k^\times/k^{\times2}$, of the spaces $\Theta_{b'}$ of theta functions
$\theta^{\psi_{b'}}_\phi$, $\phi$ even, and $\Theta_{b'}$ is irreducible
and isomorphic to $r^+_{\psi_{b'}}$; the identification of the residues
with theta functions is the Siegel--Weil formula in this case. It remains
to match $b'$ with $b$. The constant term of an element of $V_b$
transforms on the left under $\overline{A(\A)}$ by $\chi_{\psi_b}e^{H/2}$,
by~(b) and $\chi_{\psi_b}^w=\chi_{\psi_b}$. The constant term of
$\theta^{\psi_{b'}}_\phi$ is $x\mapsto(\omega_{\psi_{b'}}(x)\phi)(0)$, which
by Lemma~\ref{lem:weil-rep}, applied at each place (on $A$ and $N$ the
cocycles $\sigma_v$ and $\beta_v$ agree), transforms by
$\chi_{\psi_{-b'}}e^{H/2}$. So $\Theta_{b'}\subset L^2_{\chi_{-b'}}$, and
therefore $\Theta_{b'}=V_{-b'}$, by the uniqueness in
Lemma~\ref{lem:w-action}(d).
\end{proof}

\begin{example}\label{ex:theta}
Let $k=\Q$, $b=1$, and $\psi$ as in Example~\ref{ex:weil-Q}, so
$\psi_{-1,\infty}(x)=e^{2\pi ix}$. The residual representation
$V_1\cong r^+_{\psi_{-1}}$ contains the theta function of
$\phi=e^{-2\pi\xi^2}\otimes\mathbf 1_{\hat\Z}$: at $x=n(u)a(y^{1/2})$ in
the real component it is $y^{1/4}\theta(u+iy)$ with the classical
$\theta(z)=\sum_{n\in\Z}e^{2\pi in^2z}$, of weight
$\frac12$ on $\Gamma_0(4)$, in agreement with the lowest weight
$\frac12$ found above. The space $V_{-1}$ contains its complex conjugate.
The multiplier of $\theta$ involves the Jacobi symbol of
Example~\ref{ex:jacobi}. In the classical variable $s=(\lambda+1)/2$ the
pole is at $s=\frac34$, and the constant term of $\theta$ is a multiple
of $y^{1/4}$. For $SL(2)$ the only residual representation is the trivial
one, the residue of $E$ at $\lambda=1$ ($s=1$), whose constant term is a
constant. With $S=\{\infty,2\}$ and $\mu=\chi_\psi$,
$m^S(\mu,\lambda)=(1-2^{-2\lambda})\zeta(2\lambda)/\bigl((1-2^{-1-2\lambda})
\zeta(2\lambda+1)\bigr)$, with residue $2/\pi^2$ at $\lambda=\frac12$.
\end{example}

The residual term is the analog for the cover of the
term~\cite[(5.5)]{LL} for $SL(2)$. The odd Weil representations, the
products $\bigotimes'_vr^{\varepsilon_v}_{\psi_{b,v}}$ with some
$\varepsilon_v=-$, are not residual: those that are automorphic occur in
$L^2_{\mathrm{cusp}}$~\cite{GPS80,Wa80}, and they are counted among the
$(G,\pi)$.

\subsection{Kernels, truncation and the terms \texorpdfstring{$J_\chi$}{J\_chi}}\label{sec:Jchi}

For $\chi\in\mathcal X$ let $K_{G,\chi}(x,y)$ be the kernel of $R(f)$ on
$L^2_\chi$, and $K_{B,\chi}(x,y)$ the kernel of $R_B(f)$
(Lemma~\ref{lem:kernels}) on the $\chi$-part of
$L^2(N(\A)A(\Q)\backslash\GA)_{\mathrm{gen}}$, which is $0$ for
$\chi\in\mathcal X_G$ and the closed span of the $\mathcal H_B(\mu)\,
e^{(\lambda+1)H}$, $\mu\in\chi$, $\lambda\in i\R$, for $\chi\in\mathcal X_A$.
For $\mu\in\Pi_{\mathrm{aut}}(X_A)$ let $\mathcal B(\mu)$ be an orthonormal
basis of $\mathcal H_B(\mu)$, for~\eqref{eq:HBmu}, consisting of
$\overline{K_0}$-finite vectors.

\begin{proposition}\label{prop:kernels-spec}
\begin{enumerate}[label=(\alph*)]
\item If $\chi$ is the class of $(G,\pi)$, then $K_{B,\chi}=0$ and
  $K_{G,\chi}(x,y)=\sum_\phi(R(f)\phi)(x)\overline{\phi(y)}$, $\phi$
  running over an orthonormal basis of $L^2_\chi$ of
  $\overline{K_0}$-finite vectors.
\item If $\chi\in\mathcal X_A$, then
  \begin{align}
  K_{B,\chi}(x,y)&=\sum_{\mu\in\chi}\frac1{2\pi}\int_\R\sum_{\varphi\in
  \mathcal B(\mu)}\bigl(I(\mu_{it},f)\varphi\bigr)_{it}(x)\,
  \overline{\varphi_{it}(y)}\,dt,\label{eq:KBchi}\\
  K_{G,\chi}(x,y)&=K_{\chi,\mathrm{res}}(x,y)+\frac12\sum_{\mu\in\chi}
  \frac1{2\pi}\int_\R\sum_{\varphi\in\mathcal B(\mu)}
  E\bigl(x,I(\mu_{it},f)\varphi,it\bigr)\,\overline{E(y,\varphi,it)}\,dt,
  \label{eq:KGchi}
  \end{align}
  where $K_{\chi,\mathrm{res}}(x,y)=\sum_\theta(R(f)\theta)(x)
  \overline{\theta(y)}$ over an orthonormal basis of
  $L^2_{\chi,\mathrm{res}}$; it is $0$ unless $\chi=\chi_b$. If
  $\mu^w\ne\mu$, the two terms of the sum over $\mu\in\chi$
  in~\eqref{eq:KGchi} are equal.
\item $\sum_{\chi\in\mathcal X}K_{P,\chi}(x,y)=K_P(x,y)=
  \sum_{\mathcal O\in\mathcal C}K_{P,\mathcal O}(x,y)$ for $P\in\{B,G\}$, the
  sum over $\chi$ converging absolutely and uniformly on compact sets.
\end{enumerate}
\end{proposition}

\begin{proof}
(a) is the spectral theorem on the discrete space $L^2_\chi$, and
$L^2_\chi$ has no constant terms. For~\eqref{eq:KBchi}, write $x=n\,
\tilde a^1e^Xk$ with $\tilde a^1\in\overline{A(\A)^1}$; for a function
$F$ on $N(\A)A(\Q)\backslash\GA$, $e^{-X}F$ is a function on
$X_A\times\R\times\overline{K_0}$ with $\int|F|^2dx=\int|e^{-X}F|^2\,
d\tilde a^1\,dX\,dk$, because $dx=\delta_B^{-1}dn\,da\,dk$. Fourier
analysis on $X_A$ (Theorem~\ref{thm:corr}(a); only genuine characters
occur) and on $\R$ (with $dX$ and $dt/2\pi$) writes $F$ as
$\sum_\mu\frac1{2\pi}\int\varphi_{\mu,t,it}\,dt$ with $\varphi_{\mu,t}\in
\mathcal H_B(\mu)$ and $\|F\|^2=\sum_\mu\frac1{2\pi}\int
\|\varphi_{\mu,t}\|^2dt$, the norm being~\eqref{eq:HBmu}. So
$R_B(f)$ is the direct integral of the $I(\mu_{it},f)$, which
gives~\eqref{eq:KBchi}. Formula~\eqref{eq:KGchi} is Langlands' spectral
decomposition of $L^2_{\chi,\mathrm{cont}}$ by wave packets of
Eisenstein series~\cite[Ch.~7]{L2}, \cite[Ch.~V]{MW2}, written as
in~\cite[\S4]{A3}; the factor $\frac12=n(A_B)^{-1}$ is the reciprocal of
the number of chambers $\pm\mathfrak a^+$ in $\mathfrak a$~\cite[p.~927]{A2}.
(Proposition~\ref{prop:Tcoef} checks this factor against the geometric
side.) The equality of the two terms is Theorem~\ref{thm:eisenstein}(c),(e)
and the unitarity of $M(w,it)$. (c) Both sides are kernels of the same
operator; the absolute convergence is~\cite[\S3]{A3}, \cite{MW2},
\cite{LiIII}.
\end{proof}

\paragraph{Truncation.} For a continuous function $\phi$ on
$G(\Q)\backslash\GA$ and $T\ge0$ put
\begin{equation}\label{eq:trunc}
\Lambda^T\phi(x)=\phi(x)-\sum_{\delta\in B(\Q)\backslash G(\Q)}
\hat\tau_B\bigl(H(\delta x)-T\bigr)\,\phi_B(\delta x).
\end{equation}
By Lemma~\ref{lem:height-ineq}(b) the sum has at most one nonzero term.

\begin{lemma}\label{lem:truncation}
Let $T\ge0$.
\begin{enumerate}[label=(\alph*)]
\item $(\Lambda^T\phi)_B(x)=0$ if $H(x)>T$.
\item $\Lambda^T\Lambda^T=\Lambda^T$, and $\langle\Lambda^T\phi,\phi'\rangle
  =\langle\phi,\Lambda^T\phi'\rangle$ whenever both sides converge
  absolutely.
\item If $\phi$ is an automorphic form, for example $E(\cdot,\varphi,
  \lambda)$ at a regular point $\lambda$, then $\Lambda^T\phi$ is rapidly
  decreasing on $G(\Q)\backslash\GA$. It is not compactly supported in
  general.
\end{enumerate}
\end{lemma}

\begin{proof}
(a) Integrating~\eqref{eq:trunc} over $n\in N(\Q)\backslash N(\A)$, the
term $\delta=1$ gives $\hat\tau_B(H(x)-T)\phi_B(x)=\phi_B(x)$. For
$\delta\ne1$ and every $n$, $H(\delta nx)\le-H(nx)=-H(x)<-T\le T$ by
Lemma~\ref{lem:height-ineq}(b), so these terms vanish. (b) Idempotence
follows from~(a), since $\Lambda^T\phi$ then has no constant term where
$\hat\tau_B(H(x)-T)\ne0$. For self-adjointness, unfold:
$\langle\phi-\Lambda^T\phi,\phi'\rangle=\int_{B(\Q)\backslash\GA}
\hat\tau_B(H(x)-T)\phi_B(x)\overline{\phi'(x)}\,dx=\int_{N(\A)B(\Q)
\backslash\GA}\hat\tau_B(H(x)-T)\phi_B(x)\overline{\phi'_B(x)}\,dx$, because
$\hat\tau_B(H(x)-T)\phi_B(x)$ is left $N(\A)$-invariant; this is symmetric.
(c) is in~\cite[\S1]{A3}; the cover splits over $N(\A)$
(Lemma~\ref{lem:global}(b)), so the proof applies verbatim.
\end{proof}

Define, for $\chi\in\mathcal X$ and $T\ge0$,
\begin{equation}\label{eq:kchi}
k^T_\chi(x,f)=K_{G,\chi}(x,x)-\sum_{\delta\in B(\Q)\backslash G(\Q)}
K_{B,\chi}(\delta x,\delta x)\,\hat\tau_B\bigl(H(\delta x)-T\bigr).
\end{equation}

\begin{lemma}\label{lem:kchi}
For $T\ge0$, $k^T_\chi(x,f)=F^G(x,T)K_{G,\chi}(x,x)+\sum_{\delta}
\tau_B(H(\delta x)-T)\bigl(K_{G,\chi}-K_{B,\chi}\bigr)(\delta x,\delta x)$,
and $\sum_\chi k^T_\chi(x,f)=k^T(x,f)$.
\end{lemma}

\begin{proof}
As for Lemma~\ref{lem:partition}, whose first identity is used; the sum
over $\chi$ is Proposition~\ref{prop:kernels-spec}(c) and~\eqref{eq:kT}.
\end{proof}

\begin{theorem}[Arthur~{\cite[\S\S2--4]{A3}}; Li~\cite{LiIII} for
covers]\label{thm:Jchi-poly}
Let $f\in\mathcal H(\GA)$.
\begin{enumerate}[label=(\alph*)]
\item For $T$ sufficiently large, $\sum_{\chi\in\mathcal X}\int_{G(\Q)
  \backslash\GA}|k^T_\chi(x,f)|\,dx<\infty$.
\item $J^T_\chi(f):=\int_{G(\Q)\backslash\GA}k^T_\chi(x,f)\,dx$ is, for
  such $T$, a polynomial in $T$, and $\sum_\chi J^T_\chi(f)=J^T(f)$, the
  geometric side of Theorem~\ref{thm:coarse}.
\item $J^T_\chi(f)-\int_{G(\Q)\backslash\GA}\Lambda^T_2K_{G,\chi}(x,x)\,dx
  \to0$ as $T\to\infty$, where $\Lambda^T_2$ is $\Lambda^T$ applied to the
  second variable.
\end{enumerate}
We put $J_\chi(f)=J^0_\chi(f)$, the value of the polynomial at $T=0$.
\end{theorem}

By Theorem~\ref{thm:Jchi} below each $J^T_\chi$ has
degree at most $1$, so $J^T(f)=\sum_\chi J^T_\chi(f)$, with
Theorem~\ref{thm:coarse}, is an identity of polynomials of degree at most
$1$ and holds for every $T$, in particular at $T=0$.

\begin{lemma}[Maass--Selberg relations]\label{lem:MS}
Let $\mu,\mu'\in\Pi_{\mathrm{aut}}(X_A)$, $\varphi\in\mathcal H_B(\mu)$,
$\varphi'\in\mathcal H_B(\mu')$, $T\ge0$, and let $\lambda,\lambda'$ be
regular points with $\lambda+\bar\lambda'\ne0$ and
$\lambda-\bar\lambda'\ne0$. Then, with $M=M(w,\lambda)$ and
$\check M=M(w,\lambda')$,
\begin{multline}\label{eq:MS}
\bigl\langle\Lambda^TE(\varphi,\lambda),\Lambda^TE(\varphi',\lambda')
\bigr\rangle
=\frac{\langle\varphi,\varphi'\rangle e^{(\lambda+\bar\lambda')T}
-\langle M\varphi,\check M\varphi'\rangle e^{-(\lambda+\bar\lambda')T}}
{\lambda+\bar\lambda'}\\
+\frac{\langle M\varphi,\varphi'\rangle e^{(-\lambda+\bar\lambda')T}}
{-\lambda+\bar\lambda'}
+\frac{\langle\varphi,\check M\varphi'\rangle e^{(\lambda-\bar\lambda')T}}
{\lambda-\bar\lambda'} .
\end{multline}
The inner products are those of~\eqref{eq:HBmu} and vanish unless the
characters agree; in particular the last two terms vanish unless
$\mu'=\mu^w$.
\end{lemma}

\begin{proof}[Sketch]
For $\mathrm{Re}\,\lambda'$ large, Lemma~\ref{lem:truncation}(b) gives
$\langle\Lambda^TE,\Lambda^TE'\rangle=\langle\Lambda^TE,E'\rangle$. Unfold
$E'$ and integrate over $N(\Q)\backslash N(\A)$ first. By the proof of
Lemma~\ref{lem:truncation}, $(\Lambda^TE)_B=E_B$ on $\{H\le T\}$ minus the
terms $\delta\ne1$; folding these back into a sum over $G(\Q)$ and
unfolding the other way, they contribute
$-\int_{H>T}E_B\,\overline{(\check M\varphi')_{-\lambda'}}$, because the
constant term of $E'-\varphi'_{\lambda'}$ is $(\check M\varphi')_{-\lambda'}$.
So
\[
\langle\Lambda^TE,E'\rangle=\int_{H(x)\le T}E_B\,\overline{\varphi'_{\lambda'}}
-\int_{H(x)>T}E_B\,\overline{(\check M\varphi')_{-\lambda'}},
\]
integrals over $N(\A)A(\Q)\backslash\GA$. Insert
Theorem~\ref{thm:eisenstein}(b); in the coordinates $x=\tilde a^1e^Xk$ each
of the four products is $e^{\nu X}$ times a product of values on
$\overline{K_0}$, with $\nu\in\{\pm\lambda\pm\bar\lambda'\}$, and
$\int_{-\infty}^Te^{\nu X}dX=e^{\nu T}/\nu$, $\int_T^\infty e^{\nu X}dX=
-e^{\nu T}/\nu$ in the appropriate half planes. The integral over
$X_A\times\overline{K_0}$ gives the inner products~\eqref{eq:HBmu}. The
general case follows by analytic continuation. In general the relations
are due to Langlands; see~\cite[\S4]{A3}, \cite{MW2}, \cite{LiIII}.
\end{proof}

\begin{theorem}\label{thm:Jchi}
Let $f\in\mathcal H(\GA)$, and for $\mu\in\Pi_{\mathrm{aut}}(X_A)$ and
$\lambda\in i\R$ put
\begin{equation}\label{eq:MT}
M_B^T(\mu,\lambda)=2T-M(w,\lambda)^{-1}\frac{d}{d\lambda}M(w,\lambda),
\end{equation}
an operator on $\mathcal H_B(\mu)$, self-adjoint for $\lambda\in i\R$.
\begin{enumerate}[label=(\alph*)]
\item If $\chi$ is the class of $(G,\pi)$, then $J^T_\chi(f)=
  m(\pi)\operatorname{tr}\pi(f)$.
\item If $\chi\in\mathcal X_A$ is unramified and $\mu\in\chi$, then
  \[
  J^T_\chi(f)=\frac1{2\pi}\int_\R\operatorname{tr}\bigl(M_B^T(\mu,it)\,
  I(\mu_{it},f)\bigr)\,dt .
  \]
\item If $\chi=\chi_b$ is ramified and $\mu=\chi_{\psi_b}$, then
  \[
  J^T_\chi(f)=\operatorname{tr}\bigl(R(f)|_{V_b}\bigr)
  +\frac12\cdot\frac1{2\pi}\int_\R\operatorname{tr}\bigl(M_B^T(\mu,it)\,
  I(\mu_{it},f)\bigr)\,dt
  +\frac14\operatorname{tr}\bigl(M(w,0)\,I(\mu_0,f)\bigr).
  \]
\end{enumerate}
\end{theorem}

\begin{proof}
By Theorem~\ref{thm:Jchi-poly}(c) it suffices to compute
$\int\Lambda^T_2K_{G,\chi}(x,x)\,dx$ up to $o(1)$ as $T\to\infty$. By
Lemma~\ref{lem:truncation}(b) and Proposition~\ref{prop:kernels-spec}
this is $\sum_\phi\langle R(f)\phi,\Lambda^T\phi\rangle$ in case~(a) and
\[
\sum_\theta\langle R(f)\theta,\Lambda^T\theta\rangle+\frac12\sum_{\mu\in\chi}
\frac1{2\pi}\int_\R\sum_{\varphi\in\mathcal B(\mu)}\bigl\langle\Lambda^T
E(I(\mu_{it},f)\varphi,it),\Lambda^TE(\varphi,it)\bigr\rangle\,dt
\]
in cases~(b)--(c), the inner product on $i\R$ being the value at
$\lambda'=\lambda$ of~\eqref{eq:MS}, which is regular there; the
interchange of the integrals is justified as in~\cite[\S4]{A3}. (a) Cusp
forms satisfy $\Lambda^T\phi=\phi$. For the residual part,
$\Lambda^T\theta\to\theta$ in $L^2$, so the first sum tends to
$\operatorname{tr}(R(f)|_{V_b})$.

Now let $\lambda=it$, $\lambda'=\lambda+\varepsilon$ with
$\varepsilon\in i\R$, so $\lambda+\bar\lambda'=-\varepsilon$. Expanding
$M(w,\lambda')=M+\varepsilon\,\frac{dM}{d\lambda}+O(\varepsilon^2)$ and
using $M^*=M^{-1}$ on $i\R$, the first fraction of~\eqref{eq:MS} tends to
$\langle\varphi,M_B^T(\mu,\lambda)\varphi'\rangle$ as $\varepsilon\to0$.
With $\varphi=I(\mu_\lambda,f)\varphi_j$, $\varphi'=\varphi_j$ and the sum
over $j$ this gives $\operatorname{tr}(M_B^T(\mu,\lambda)
I(\mu_\lambda,f))$, using that $M^{-1}\frac{dM}{d\lambda}$ is self-adjoint on
$i\R$ (differentiate $M(it)^*M(it)=1$). For $\mu^w\ne\mu$ the two $\mu\in
\chi$ give equal contributions, and the factor $\frac12$ cancels; this
proves~(b), the expression being already linear in $T$.

For $\mu^w=\mu$ the last two terms of~\eqref{eq:MS} also contribute, with
$-\lambda+\bar\lambda'\to-2it$ and $\lambda-\bar\lambda'\to2it$. Put
$F(t)=\operatorname{tr}(M(w,it)I(\mu_{it},f))$; by
Theorem~\ref{thm:eisenstein}(c),(e),
$\operatorname{tr}(M(w,-it)I(\mu_{it},f))=F(-t)$. The contribution is
\[
\frac12\cdot\frac1{2\pi}\,\mathrm{PV}\!\int_\R\frac{F(-t)e^{2itT}-F(t)
e^{-2itT}}{2it}\,dt
=\frac1{4\pi}\Bigl(\int_\R F(t)\frac{\sin2tT}{t}\,dt
+i\,\mathrm{PV}\!\int_\R F(t)\frac{\cos2tT}{t}\,dt\Bigr),
\]
which tends to $\frac1{4\pi}\cdot\pi F(0)=\frac14F(0)$ as $T\to\infty$, by
Dirichlet's integral and the Riemann--Lebesgue lemma applied to
$(F(t)-F(-t))/t$. Since $J^T_\chi$ is a polynomial, it equals the
polynomial part of these limits, which is~(c).
\end{proof}

Two remarks on the constants. First, $M_B^T$ is the case of
$SL(2)$ of Arthur's formula
\[
M_B^T=\mathrm{vol}(\mathfrak a/\Z\alpha^\vee)\lim_{\zeta\to0}
\sum_{s\in\{1,w\}}\frac{e^{(s\zeta)(T)}M(s,\lambda)^{-1}
M(s,\lambda+\zeta)}{\theta(s\zeta)},
\]
in which $\theta(\zeta)=\zeta(\alpha^\vee)=d\zeta$ and
$\mathrm{vol}(\mathfrak a/\Z\alpha^\vee)=d$ for our $dX$
(Remark~\ref{rem:alpha}); the factor $d$ cancels and the limit
is~\eqref{eq:MT} for every $k$. Proposition~\ref{prop:Tcoef} confirms the
coefficient of $T$ against the geometric side. Second, in Arthur's fine
expansion~\cite[Thm.~8.2]{A14}, \cite{LiIII} the last term of~(c) is
$|W^A_0||W^G_0|^{-1}|\det(w-1)_{\mathfrak a}|^{-1}
\operatorname{tr}(M(w,0)I(\mu_0,f))$, with $|W^A_0||W^G_0|^{-1}=\frac12$
and $|\det(w-1)_{\mathfrak a}|=2$. In the trace formula for $GL(2)$
of~\cite[\S6]{GJ} the same term is written $-\frac14\operatorname{tr}
(M(0)\pi_0(f))$, next to $+\frac1{4\pi}\int\operatorname{tr}(M(-iy)M'(iy)
\pi_{iy}(f))\,dy$: that formula has the trace of $R(f)$ on the discrete
spectrum alone on one side, and moving both terms to the spectral side
gives the signs of Theorem~\ref{thm:tf}(e). The sign of the cross terms
in~\eqref{eq:MS}, from which it comes, is also checked numerically on
the Eisenstein series of $SL(2,\Z)$ (Remark~\ref{rem:verified-sec6}),
and $M(w,0)$ is computed in Proposition~\ref{prop:Mw0}.

\begin{corollary}\label{cor:mprime}
Let $S$ and $\mu=\chi_\psi\cdot(\chi\circ p)$ be as in
Lemma~\ref{lem:intertwining}, $f=f_S\otimes e^S$, and $\lambda\in i\R$.
Then $I(\mu_\lambda,f)=0$ unless $\mu_v$ is trivial on $A(\OO_v)^*$ for
all $v\notin S$, and in that case
\begin{multline*}
\operatorname{tr}\Bigl(M(w,\lambda)^{-1}\tfrac{dM}{d\lambda}(w,\lambda)
I(\mu_\lambda,f)\Bigr)\\
=\frac{(m^S)'}{m^S}(\mu,\lambda)\operatorname{tr}
I_S(\mu_{S,\lambda},f_S)+\operatorname{tr}\Bigl(M_S(\lambda)^{-1}
\tfrac{dM_S}{d\lambda}(\lambda)I_S(\mu_{S,\lambda},f_S)\Bigr),
\end{multline*}
where $I_S=\bigotimes_{v\in S}I_v$ and
$\frac{(m^S)'}{m^S}(\mu,\lambda)=2\frac{(L^S)'}{L^S}(2\lambda,\chi^2)
-2\frac{(L^S)'}{L^S}(2\lambda+1,\chi^2)$.
\end{corollary}

\begin{proof}
$I(\mu_\lambda,e_v)$ is the projection onto the $K_v^*$-fixed vectors,
which are $0$ unless $\mu_v$ is trivial on $A(\OO_v)^*$. Then apply
Lemma~\ref{lem:intertwining}(a)--(b) and the Leibniz rule.
\end{proof}

\begin{example}\label{ex:mprime-Q}
For $k=\Q$, $S=\{\infty,2\}$ and $\mu=\chi_\psi$,
$\frac{(m^S)'}{m^S}(\lambda)=2\frac{\zeta'}{\zeta}(2\lambda)
-2\frac{\zeta'}{\zeta}(2\lambda+1)+\frac{2\log2\cdot2^{-2\lambda}}
{1-2^{-2\lambda}}-\frac{2\log2\cdot2^{-1-2\lambda}}{1-2^{-1-2\lambda}}$.
For $SL(2)$ the corresponding factor is
$\frac{\zeta'}{\zeta}(\lambda)-\frac{\zeta'}{\zeta}(\lambda+1)$ plus the
$2$-adic terms.
\end{example}

\subsection{The trace formula}\label{sec:tf}

\begin{theorem}\label{thm:tf}
Let $S$ be a finite set of places of $\Q$ as in
Section~\ref{sec:local-dist}, large enough for Theorem~\ref{thm:fine} and
containing the places below which $\psi$ is ramified; let $f_S\in
\mathcal H(\overline{G(\Q_S)})$ and $f=f_S\otimes e^S$. Then
$\sum_{\mathcal O}J_{\mathcal O}(f)=\sum_{\chi}J_\chi(f)$, where both
sides converge absolutely and:
\begin{enumerate}[label=(\alph*)]
\item for the central classes $\mathcal O_z$, $z\in\{\pm1\}$,
  \[
  J_{\mathcal O_z}(f)=a^G(S,1)f_S(z)+\sum_{\xi\in k_S^\times/k_S^{\times2}}
  a^G(S,u_\xi)\,J_G(zu_\xi,f_S)+\tfrac12\,c_k\,J_A(z,f_S),
  \]
  with the coefficients $a^G(S,u_\xi)$ of Theorem~\ref{thm:fine}, which
  are computed in Proposition~\ref{prop:unipotent};
\item for $\mathcal O$ elliptic, $J_{\mathcal O}(f)=
  \mathrm{vol}(T_\mu(\Q)\backslash T_\mu(\A))\,J_G(\mu_S,f_S)$;
\item for $\mathcal O$ hyperbolic, the class of $a(t)$,
  $J_{\mathcal O}(f)=c_k\,J_A(a(t)_S,f_S)$;
\item for $\chi$ the class of $(G,\pi)$, $J_\chi(f)=m(\pi)
  \operatorname{tr}\pi(f)$;
\item summed over $\chi\in\mathcal X_A$,
  \begin{multline*}
  \sum_{\chi\in\mathcal X_A}J_\chi(f)=\sum_{b}\operatorname{tr}
  \bigl(R(f)|_{V_b}\bigr)
  -\frac12\sum_{\mu}\frac1{2\pi}\int_\R\operatorname{tr}\Bigl(M(w,it)^{-1}
  \tfrac{dM}{d\lambda}(w,it)\,I(\mu_{it},f)\Bigr)\,dt\\
  +\frac14\sum_b\operatorname{tr}\bigl(M(w,0)\,I(\chi_{\psi_b,0},f)
  \bigr),
  \end{multline*}
  where $b$ runs over $k^\times/k^{\times2}$ and $\mu$ over
  $\Pi_{\mathrm{aut}}(X_A)$; only finitely many $b$, and only $\mu$
  unramified outside $S$, contribute (Example~\ref{ex:ramified-Q},
  Corollary~\ref{cor:mprime}).
\end{enumerate}
\end{theorem}

\begin{proof}
At $T=0$, Theorem~\ref{thm:coarse}, Theorem~\ref{thm:fine} and
Theorem~\ref{thm:Jchi-poly}(b) give $J(f)=\sum_{\mathcal O}J_{\mathcal O}(f)
=\sum_\chi J_\chi(f)$. (a)--(c) are Theorem~\ref{thm:fine}. (d)--(e) are
Theorem~\ref{thm:Jchi} at $T=0$: in~(b) of that theorem, write the
integral as half the sum over the two members of $\chi$; in~(c) the class
has one member. The absolute convergence of the spectral side
is~\cite{A3}, \cite{LiIII}.
\end{proof}

\begin{proposition}\label{prop:Tcoef}
The coefficients of $T$ in $\sum_\chi J^T_\chi(f)$ and in
$\sum_{\mathcal O}J^T_{\mathcal O}(f)$ agree. Both equal
$c_k\sum_{\gamma\in A(\Q)}F_f(\gamma)$, where $F_f(x)=\int_{\overline{K_0}}
\int_{N(\A)}f(\tilde k^{-1}xn\tilde k)\,dn\,d\tilde k$.
\end{proposition}

\begin{proof}
Spectral side. By Theorem~\ref{thm:Jchi}, only the term $2T$ of $M_B^T$
contributes, and summing over $\chi$ the coefficient is
$\sum_\mu\frac1{2\pi}\int_\R\operatorname{tr}I(\mu_{it},f)\,dt$ (each
unramified class contributes twice $\frac12$ through its two members).
Computing $I(\mu_\lambda,f)$ on $\overline{K_0}$ in the coordinates
$y=n\tilde ak'$ of~\eqref{eq:iwasawa} gives the kernel
$\int_{\overline{A(\A)}}\int_{N(\A)}f(\tilde k^{-1}\tilde an\tilde k')\mu(\tilde
a)e^{(\lambda+1)H(\tilde a)}\,dn\,d\tilde a$, so $\operatorname{tr}
I(\mu_\lambda,f)=\int_{\overline{A(\A)}}e^{H(\tilde a)}F_f(\tilde a)
\mu(\tilde a)e^{\lambda H(\tilde a)}\,d\tilde a$. Fourier inversion on $\R$
($dX$, $dt/2\pi$) and then Theorem~\ref{thm:tf-rank0} on $X_A$, applied
to the anti-genuine function $F_f$ on $\overline{A(\A)^1}$, give
$\sum_\mu\frac1{2\pi}\int\operatorname{tr}I(\mu_{it},f)\,dt=c_k\sum_{\gamma\in
A(\Q)}F_f(\gamma)$.

Geometric side. By Propositions~\ref{prop:hyperbolic}--\ref{prop:central},
the coefficient is $c_k\bigl(\hat\Phi_f(0)+\hat\Phi_{f_{-1}}(0)\bigr)
+2c_k\sum_{\mathcal O\ \mathrm{hyp}}\int_{\overline{A(\A)}\backslash\GA}
f(y^{-1}\mu y)\,dy$. Here $\hat\Phi_{f_z}(0)=F_f(z)$ for $z=\pm1$. For
$\mu=a(t)$, $t\in k^\times\smallsetminus\{\pm1\}$, the coordinates
$y=n(z)\tilde k$ and $y^{-1}\mu y=\tilde k^{-1}\mu\,n((1-t^{-2})z)\tilde k$
give $\int f(y^{-1}\mu y)\,dy=|1-t^{-2}|_{\A_k}^{-1}F_f(\mu)=F_f(\mu)$ by the
product formula. The integral is the same for $a(t^{-1})=w\,a(t)\,w^{-1}$
(Lemma~\ref{lem:w-action}(a)), so $2F_f(a(t))=F_f(a(t))+F_f(a(t^{-1}))$,
and the hyperbolic class of $a(t)$ accounts for the two elements
$a(t^{\pm1})$ of $A(\Q)$.
\end{proof}

Proposition~\ref{prop:Tcoef} is a consistency check between
Sections~\ref{sec:geometric} and~\ref{sec:spectral}: with $4T$ in place
of $2T$ in $M_B^T$, or with $n(A_B)=1$ in place of $2$, the spectral
coefficient would be twice the geometric one.

\begin{remarkrev}\label{rem:sec5-norm}
\emph{Normalizations.} Everything here uses $\lambda(X)=\lambda X$ on
$\mathfrak a\cong\R$, Lebesgue $dX$ and $dt/2\pi$ on $i\R$, and the inner
product~\eqref{eq:HBmu}, which includes $c_k$. Arthur normalizes the
measure on $\mathfrak a$ so that $\mathfrak a/\Z\alpha^\vee$ has volume
$1$, i.e.\ $dX/d$, and takes the dual measure on $i\mathfrak a^*$; the
operator $M_B^T$ then carries the factor $\mathrm{vol}(\mathfrak a/
\Z\alpha^\vee)\theta(\zeta)^{-1}$, and the products in
Theorem~\ref{thm:Jchi} do not change (Remark~\ref{rem:alpha}). In the
classical variable $s=(\lambda+1)/2$ of the Conventions above, for
$k=\Q$, $\frac{d}{d\lambda}=\frac12\frac{d}{ds}$, and $2T=\log y_T$, where
$y_T=e^{2T}$ is the height at which $\Lambda^T$ truncates.
\end{remarkrev}

\subsection{Status of the results}\label{sec:status5}

\begin{remarkrev}\label{rem:sec5-status}
\emph{What rests on citation.} Proved here: Lemma~\ref{lem:factor} (given Flath's
theorem); Lemmas~\ref{lem:w-action}, \ref{lem:truncation}(a)--(b)
and~\ref{lem:kchi}; Lemma~\ref{lem:intertwining}, with~(b) also checked
numerically;
Theorem~\ref{thm:residual}(a)--(b), the computations in the proofs of
Lemma~\ref{lem:MS} and Theorem~\ref{thm:Jchi} from~\eqref{eq:MS}, and
Proposition~\ref{prop:Tcoef}. Taken from the literature:
Lemma~\ref{lem:spec-decomp}, Proposition~\ref{prop:Lchi},
Theorem~\ref{thm:eisenstein}~\cite{L2,MW2}; the wave-packet formula~\eqref{eq:KGchi} and the
absolute convergence in Proposition~\ref{prop:kernels-spec};
Lemma~\ref{lem:truncation}(c); Theorem~\ref{thm:Jchi-poly}~\cite{A3,
LiIII}, including the interchange of integrals used for
Theorem~\ref{thm:Jchi}; the analytic continuation in Lemma~\ref{lem:MS};
and Theorem~\ref{thm:residual}(c)~\cite{Ge2,GPS80,Wa80}, except for
the matching $V_b\cong r^+_{\psi_{-b}}$ of the parameters, which is proved
here from Lemma~\ref{lem:weil-rep}. The matching depends on the
normalization~\eqref{eq:weil-rep}: if $n(u)$ acts by $\psi'(\frac12ux^2)$,
it reads $V_b\cong r^+_{\psi_{-2b}}$. We use from the cited papers that
the residual spectrum is the sum of the spaces of even theta
functions~\eqref{eq:theta} and that these are irreducible. The range ``$T$ sufficiently large'' in
Theorem~\ref{thm:Jchi-poly} is Arthur's; for this group we expect $T\ge0$
to suffice, as on the geometric side, but have not checked it.
\end{remarkrev}

\begin{remarkrev}\label{rem:verified-sec5}
The script \texttt{checks/sec5\_n2\_check.py} verifies, for $k=\Q$:
Lemma~\ref{lem:w-action}(a)--(c) at $v=\infty,2,3,5,7$ (exactly, with
$\sigma_v$ and the Weil index); the unramified intertwining integral of
Lemma~\ref{lem:intertwining}(b) at $p=3,5,7$, by an explicit
decomposition of the lift of $w^{-1}n(y)$, shell by shell, including the
shape of the integrand on odd and even shells, and the same integral for
$SL(2)$; Arthur's limit for $M_B^T$ for $d=1,2,3$ and the 1995
normalization (symbolically); the diagonal limit of~\eqref{eq:MS} and the
self-adjointness of $M^{-1}M'$ for a model unitary family; the limit
$\frac14F(0)$ of the cross terms; the residue $2/\pi^2$ of
Example~\ref{ex:theta}; at $v=\R$, the two local identities behind
Proposition~\ref{prop:Tcoef}; and, for Lemma~\ref{lem:weil-rep}, the
product $(w(t),1)(w(1),1)^{-1}=(a(t),(-1,t)_v)$ of unipotent lifts and the
identity $(\gamma_v(t)/\gamma_v(1))^2=(t,-1)_v$ at $v=\infty,2,3,5,7$, and
the weights $\pm\frac12$ of the ground state at $v=\R$.
\end{remarkrev}

% SPDX-License-Identifier: CC-BY-4.0
% =====================================================================
% Section 6 of D. Joyner, "A hitch-hiker's guide to the metaplectic trace
% formula for SL(2)", revised for the 2-fold cover (n = 2), November 2026.
%
% Stage 5 of ../revision_log.md. Sources: ../original-1995/cctrace6.tex
% (= ../book-1997/book7g.tex; ../book-1997/book7gnw.tex is a later version
% of the same chapter). Referee items II.14-II.17 and the Section 6 minor
% items are fixed here; see the "Stage 5 notes" in ../revision_log.md,
% which also map the 1995 numbers to the numbers here.
%
% This file is \input by the master; it compiles on its own through
% preview.tex. Explicit formulas are checked by checks/sec6_n2_check.py.
% Comments beginning "% [Claude]" are editorial notes by Claude.
% Revised after the second referee report (ChatGPT, November 2026): an
% explanation of the correction term before the formal construction, real
% normalizing factors for all unitary characters, the three senses of
% "normalizing factor", a fuller proof of M(w,0) = -1, and an accurate
% statement of what is explicit; comparisons with the 1995 notes moved to
% Appendix A.
% =====================================================================

\providecommand{\GA}{\overline{G(\A)}}
\providecommand{\GS}{\overline{G(\Q_S)}}
\providecommand{\AS}{\overline{A(\Q_S)}}
\providecommand{\KS}{\overline{K_S}}

\section{The invariant form}\label{sec:invariant}

The trace formula of Theorem~\ref{thm:tf} is an identity $\sum_{\mathcal O}
J_{\mathcal O}(f)=\sum_\chi J_\chi(f)$ between distributions, but not all
of the terms are invariant: $J(f^y)\ne J(f)$ in general, where
$f^y(x)=f(yxy^{-1})$. The elliptic terms $J_{\mathcal O}$ and the cuspidal
terms $J_\chi$ are invariant; the hyperbolic and central terms, and the
terms $J_\chi$ with $\chi\in\mathcal X_A$, are not. This section computes
the failure of invariance, defines the invariant distributions
$I_A(\gamma,f)$ that replace the weighted orbital integrals
$J_A(\gamma,f)$, following Arthur~\cite{A6,A9,A10}, and puts the trace
formula in invariant form (Theorem~\ref{thm:itf} and
Corollary~\ref{cor:itf-final}).

For covering groups the general theory is due to Li. The normalization
of the local intertwining operators, which the construction needs, is
in~\cite{LiII}. The invariant trace formula is in~\cite{LiIV}, under the
hypothesis that the invariant Paley--Wiener theorem holds at the real
places for the Levi subgroups; Li verifies the hypothesis for the
two-fold covers of $Sp(2n)$, hence for the group here. For
$G=\Res_{k/\Q}SL(2)$ much less is needed. The only proper Levi subgroup is
$A$, and its cover is abelian (Lemma~\ref{lem:global}(e)), so the one
Paley--Wiener theorem that enters is Fourier inversion on the abelian
group $\AS$. Given the normalized intertwining operators, we can
therefore prove the variance formulas, the invariance of $I_A(\gamma,f)$
and the invariant trace formula directly; we do so, and cite
Arthur and Li for the general statements.

\paragraph{Why a correction term, and why this one.} The weighted
orbital integral $J_A(\gamma,f)$ integrates $f(x^{-1}\gamma x)$ against the
weight $v_A(x)=-H(x)-H(wx)$. The height $H$ is defined through the
Iwasawa decomposition, so it depends on the choice of the maximal compact
subgroup: $H(xy)\ne H(x)+H(y)$ in general, but $H(xy)=H(x)+H(\tilde ky)$
when $x=n\tilde a\tilde k$. Replacing $f$ by $f^y$ replaces $v_A(x)$ by
$v_A(xy)$, and the difference $-H(\tilde ky)-H(\tilde k'y)$ does not
vanish. This is the geometric source of the non-invariance
(Lemma~\ref{lem:variance-local}). Since the difference depends only on the
$K$-components, the change in $J_A(\gamma,f)$ is not a new kind of
object: it is the value at $\gamma$ and at $w\gamma w^{-1}$ of one function
$f_{B,y}$ on the cover of $A$.

On the spectral side the same function appears. The logarithmic
derivative $\mathcal M^{-1}\mathcal M'$ of an intertwining operator
between $I(\pi_\lambda)$ and $I(\pi^w_{-\lambda})$ is not an intertwining
operator, because the two representations depend on $\lambda$ through
the same height: differentiating $e^{(\lambda+1)H(\tilde ky)}$ brings down
$H(\tilde ky)$. Conjugating $\mathcal M^{-1}\mathcal M'$ by $I(\pi_\lambda,
y)$ therefore changes it by multiplication by $H(\tilde ky)$, once
directly and once through $\mathcal M$ (Lemma~\ref{lem:conj-identity}),
and the trace against $I(\pi_\lambda,f)$ changes by the Fourier transform
of $f_{B,y}$ at $\pi_\lambda$ and at $\pi^w_{-\lambda}$.

So the function $\varphi_A(f,\pi)=-\operatorname{tr}(R^{-1}R'\,I(\pi,f))$
on the unitary dual of the cover of $A$, built from the normalized
operators $R$, has a variation which is the Fourier transform of the
variation of $\gamma\mapsto J_A(\gamma,f)$. Its inverse Fourier transform
$\Phi_A(\gamma,f)$ then varies exactly as $J_A(\gamma,f)$ does, and
\[
I_A(\gamma,f)=J_A(\gamma,f)-\Phi_A(\gamma,f)
\]
is invariant (Theorem~\ref{thm:IA}). The two functions live on opposite
sides of Fourier inversion: $\varphi_A(f,\cdot)$ on characters $\pi$ of
the cover of $A(\Q_S)$, $\Phi_A(\cdot,f)$ on its elements. Globally the
corrections on the two sides are matched by Poisson summation on the
cover of $A$, which is the rank-$0$ trace formula of
Section~\ref{sec:rank0}, so the identity is preserved. The normalized
operators $R_v$ are used, rather than $M$, because they are holomorphic
and unitary on $i\R$ and trivial on the unramified vectors, so that
$\varphi_A$ is given by finitely many local terms. The scalar that
remains, $m$ in $M=m\cdot R$, contributes $\frac{m'}m\operatorname{tr}
I(\pi,f)$, which is invariant; this is the term that survives on the
spectral side.

\paragraph{Results.} The main points of the section are these.
\begin{itemize}
\item The weight in the variance formulas is $u'_B(k,y)=-H(ky)$
  (Lemma~\ref{lem:uprime}).
\item The normalization of the intertwining operators is a theorem of
  Li~\cite{LiII}, and no trace Paley--Wiener theorem is needed for this
  group.
\item The invariant spectral terms attached to $\chi\in\mathcal X_A$ do
  not vanish. The logarithmic derivative $m'/m$ of the global normalizing
  factor survives, and $m$ is identified up to finitely many local
  factors: $m'/m=2\frac{L'}{L}(2\lambda,\chi^2)-2\frac{L'}{L}(2\lambda+1,
  \chi^2)$ plus the terms of the places in $S_{\mathrm{bad}}(\mu)$
  (Corollary~\ref{cor:global-m}).
\item The residual terms $\operatorname{tr}(R(f)|_{V_b})$ and the ramified
  terms of Theorem~\ref{thm:Jchi}(c) are already invariant. The ramified
  term is evaluated: $M(w,0)=-1$ on $\mathcal H_B(\chi_{\psi_b})$, so it
  equals $-\frac14\operatorname{tr}I(\chi_{\psi_b,0},f)$
  (Proposition~\ref{prop:Mw0}).
\item The function $\varphi_A(f)$ that carries the correction is not the
  Fourier transform of a compactly supported function on $\AS$
  (Example~\ref{ex:Phi-real}).
\end{itemize}

\paragraph{Conventions.} $S$ is a finite set of places of $\Q$ containing
$\infty$, $2$ and the places below which $\psi$ is ramified, $S_k$ is the
set of places of $k$ above $S$, and $\GS$, $\mathcal H(\GS)$,
$f=f_S\otimes e^S$ and $\gamma_S$ are as in Section~\ref{sec:local-dist}.
Let $K_S=\prod_{v\in S_k}K_v$, with preimage $\KS$ in $\GS$, and let
$\AS$ be the preimage of $A(\Q_S)=\prod_{v\in S_k}A(k_v)$, an abelian
group (it is the group $\overline{A(\A_S)}$ of Section~\ref{sec:IA}). Put
$H_S(x)=\sum_{v\in S_k}H_v(x_v)$ for $x\in\GS$; on the image of $\GS$ in
$\GA$ it is the height $H$ of Lemma~\ref{lem:NAK-global}. Let
$\Pi_u(\AS)$ be the group of genuine unitary characters of $\AS$. It is
a principal homogeneous space under the unitary dual of $A(\Q_S)$
(Lemma~\ref{lem:genuine-chars}), and $d\pi$ is the measure on it that
corresponds to the Plancherel measure dual to the Haar measure $da$ on
$A(\Q_S)$. For an anti-genuine integrable function $h$ on $\AS$ put
\[
\hat h(\pi)=\int_{\AS}h(\tilde a)\pi(\tilde a)\,d\tilde a,\qquad\text{so
that}\qquad h(\tilde a)=\int_{\Pi_u(\AS)}\hat h(\pi)\,\pi(\tilde a)^{-1}\,d\pi
\]
when $\hat h$ is integrable and $h$ continuous. For $\pi\in\Pi_u(\AS)$ and
$\lambda\in\C$ put $\pi_\lambda=\pi\,e^{\lambda H_S}$, and
$\pi^w(\tilde a)=\pi(w^{-1}\tilde aw)$, where $w$ stands for
$\prod_{v\in S_k}(w,1)_v$; by Lemma~\ref{lem:w-action}(a),
$(\pi_\lambda)^w=(\pi^w)_{-\lambda}$. Let $V(\pi)$ be the space of smooth
functions $\varphi$ on $\KS$ with $\varphi(\tilde m\tilde k)=\pi(\tilde
m)\varphi(\tilde k)$ for $\tilde m\in\AS\cap\KS$, and let $I(\pi_\lambda)$
be the induced representation of $\GS$ on $V(\pi)$:
$(I(\pi_\lambda,y)\varphi)(\tilde k)=\varphi_\lambda(\tilde ky)$, where
$\varphi_\lambda(n\tilde a\tilde k)=\pi(\tilde a)e^{(\lambda+1)H_S(\tilde
a)}\varphi(\tilde k)$. These are the local analogs of
$\mathcal H_B(\mu)$ and $I(\mu_\lambda)$ in
Section~\ref{sec:eisenstein}: for $\mu\in\Pi_{\mathrm{aut}}(X_A)$ with
$\mu_v$ trivial on $A(\OO_v)^*$ for all $v\notin S_k$, the restriction
$\mu_S$ of $\mu$ to $\AS$ lies in $\Pi_u(\AS)$, and
$\operatorname{tr}I(\mu_\lambda,f)=\operatorname{tr}I(\mu_{S,\lambda},f_S)$.

Invariance refers to conjugation by all $y\in\GS$. At an archimedean
place $f^y$ need not be $K$-finite when $f$ is, so in
Sections~\ref{sec:noninv} and~\ref{sec:IAgamma} test functions are taken
in $C_c^\infty(\GS)_-$, the space of smooth, compactly supported,
anti-genuine functions, which is stable under $f\mapsto f^y$ and to which
the distributions $J_G(\gamma,\cdot)$, $J_A(\gamma,\cdot)$ and
$\operatorname{tr}I(\pi,\cdot)$ extend by the same formulas.

\subsection{Non-invariance}\label{sec:noninv}

This subsection computes $J(f^y)-J(f)$ for each of the distributions in
the trace formula. In every case the difference is a rank-$0$
distribution on the cover of $A$, evaluated at the function $f_{B,y}$
defined below. The geometric and the spectral differences are the two
sides of the trace formula for $\overline{A(\A)^1}$
(Theorem~\ref{thm:tf-rank0}), which is why the correction terms of
Section~\ref{sec:itf} cancel in the sum.

As in~\cite[\S\S2--3]{A6}, let
\begin{equation}\label{eq:Gammaprime}
\Gamma'_B(X,Y)=\tau_B(X)-\hat\tau_B(X-Y),\qquad X,Y\in\mathfrak a\cong\R,
\end{equation}
the case $Q=B$ of $\Gamma'_Q(X,Y)=\sum_{R\supset Q}(-1)^{\dim(A_R/A_G)}
\tau_Q^R(X)\hat\tau_R(X-Y)$, in the notation of~\eqref{eq:tau}; for fixed
$Y$ it is compactly supported in $X$. For $\tilde k\in\KS$, $y\in\GS$ and
$f\in C_c^\infty(\GS)_-$ put
\begin{gather}
u'_B(\tilde k,y)=\int_{\mathfrak a}\Gamma'_B\bigl(X,-H_S(\tilde ky)\bigr)
\,dX,\label{eq:uprime}\\
f_{B,y}(\tilde a)=\delta_B(\tilde a)^{1/2}\int_{\KS}\int_{N(\Q_S)}
f(\tilde k^{-1}\tilde an\tilde k)\,u'_B(\tilde k,y)\,dn\,d\tilde k,\qquad
\tilde a\in\AS,\label{eq:fBy}
\end{gather}
and likewise for $y\in\GA$, $f\in\mathcal H(\GA)$ and $\tilde a\in
\overline{A(\A)}$, with $\overline{K_0}$, $N(\A)$ and $H$. Here $dX$ is
Lebesgue measure, as in Remark~\ref{rem:alpha}.

\begin{lemma}\label{lem:uprime}
\begin{enumerate}[label=(\alph*)]
\item $u'_B(\tilde k,y)=-H_S(\tilde ky)$.
\item $f_{B,y}$ is smooth, compactly supported and anti-genuine on $\AS$,
  and $f_{B,1}=0$.
\item For $\pi\in\Pi_u(\AS)$ and $\lambda\in\C$ let $U_y$ be the operator
  of multiplication by $\tilde k\mapsto u'_B(\tilde k,y)$ on $V(\pi)$.
  Then $\operatorname{tr}\bigl(U_y\,I(\pi_\lambda,f)\bigr)
  =\widehat{f_{B,y}}(\pi_\lambda)$.
\end{enumerate}
\end{lemma}

\begin{proof}
(a) $\int_\R\bigl(\tau_B(X)-\tau_B(X-Y)\bigr)\,dX=Y$ for every real $Y$:
the integrand is the characteristic function of $(0,Y]$ if $Y>0$ and
minus that of $(Y,0]$ if $Y<0$. (b) $u'_B(\cdot,y)$ is smooth and left
invariant under $\AS\cap\KS$, and $u'_B(\tilde k,1)=-H_S(\tilde k)=0$.
(c) As in the proof of Proposition~\ref{prop:Tcoef}, $I(\pi_\lambda,f)$
has the kernel $\int_{\AS}\int_{N(\Q_S)}f(\tilde k^{-1}\tilde
an\tilde k')\pi(\tilde a)e^{(\lambda+1)H_S(\tilde a)}\,dn\,d\tilde a$ on
$\KS$; multiply by $u'_B(\tilde k,y)$ and integrate over the diagonal,
using $\delta_B^{1/2}=e^{H_S}$.
\end{proof}

The sign matters: $u'_B(\cdot,y)$ is not of constant sign
(Example~\ref{ex:uprime}), and with $|H|$ in place of $-H$ the formulas of
Lemmas~\ref{lem:variance-global} and~\ref{lem:variance-local} fail.

\begin{example}\label{ex:uprime}
Let $v$ be a real place, $k_\theta=\begin{pmatrix}\cos\theta&\sin\theta\\
-\sin\theta&\cos\theta\end{pmatrix}$ and $y=a(s)$ with $s>0$. Since
$H_v(g)=-\log\|e_2g\|$ (Lemma~\ref{lem:height-ineq}),
\[
u'_B(k_\theta,a(s))=\tfrac12\log\bigl(s^2\sin^2\theta+s^{-2}\cos^2\theta
\bigr),
\]
which is $-\log s$ at $\theta=0$ and $+\log s$ at $\theta=\frac\pi2$.
\end{example}

For a class $\mathcal O\in\mathcal C$, a class $\chi\in\mathcal X$ and an
anti-genuine function $h$ on $\overline{A(\A)}$ whose restriction to
$\overline{A(\A)^1}$ is smooth and compactly supported, put
\begin{equation}\label{eq:JAo}
J^A_{\mathcal O}(h)=c_k\sum_{\gamma\in\mathcal O\cap A(\Q)}h(\gamma),\qquad
J^A_\chi(h)=\sum_{\mu\in\chi}\int_{\overline{A(\A)^1}}h(\tilde a)\mu(\tilde
a)\,d\tilde a,
\end{equation}
with $J^A_\chi=0$ for $\chi\in\mathcal X_G$. These are the
distributions~\eqref{eq:IA}, summed over $\mathcal O\cap A(\Q)$ and over
$\chi$, and Theorem~\ref{thm:tf-rank0} says that
$\sum_{\mathcal O}J^A_{\mathcal O}(h)=\sum_\chi J^A_\chi(h)$. By Fourier
inversion on $\mathfrak a\cong\R$, $\int_{\overline{A(\A)^1}}h\mu=\frac1{2\pi}
\int_\R\hat h(\mu_{it})\,dt$ when $\hat h(\mu_{it})=\int_{\overline{A(\A)}}
h\,\mu\,e^{itH}$ is integrable in $t$.

\begin{lemma}\label{lem:variance-global}
Let $f\in\mathcal H(\GA)$ and $y\in\GA$ with $f^y\in\mathcal H(\GA)$. Then,
for all $\mathcal O\in\mathcal C$ and $\chi\in\mathcal X$,
\[
J_{\mathcal O}(f^y)=J_{\mathcal O}(f)+J^A_{\mathcal O}(f_{B,y}),\qquad
J_\chi(f^y)=J_\chi(f)+J^A_\chi(f_{B,y}).
\]
In particular $J_{\mathcal O}$ is invariant when $\mathcal O$ is elliptic,
and $J_\chi$ is invariant when $\chi\in\mathcal X_G$.
\end{lemma}

\begin{proof}
This is the formula of~\cite[\S3]{A6} for this group; the proof is
short. Since
$f^y(x^{-1}\gamma nx)=f\bigl((xy^{-1})^{-1}\gamma n(xy^{-1})\bigr)$, the
kernels satisfy $K_{P,\mathcal O}^{f^y}(x,x)=K_{P,\mathcal O}^f(xy^{-1},
xy^{-1})$. Substituting $x\mapsto xy$ in $J^T_{\mathcal O}(f^y)=\int
k^T_{\mathcal O}(x,f^y)\,dx$ (Theorem~\ref{thm:coarse}) replaces
$\hat\tau_B(H(\delta x)-T)$ in~\eqref{eq:kT} by $\hat\tau_B(H(\delta
xy)-T)$. Write $\delta x=n\tilde a\tilde k$; then $H(\delta xy)=H(\delta
x)+H(\tilde ky)$, and by~\eqref{eq:Gammaprime},
\[
\hat\tau_B\bigl(H(\delta x)+H(\tilde ky)-T\bigr)=\hat\tau_B\bigl(H(\delta
x)-T\bigr)-\Gamma'_B\bigl(H(\delta x)-T,-H(\tilde ky)\bigr).
\]
So $J^T_{\mathcal O}(f^y)-J^T_{\mathcal O}(f)=\int_{B(\Q)\backslash\GA}
K_{B,\mathcal O}(x,x)\,\Gamma'_B(H(x)-T,-H(\tilde ky))\,dx$. In the
coordinates $x=n\tilde a\tilde k$, with $dx=\delta_B(\tilde a)^{-1}dn\,
d\tilde a\,d\tilde k$,
\[
K_{B,\mathcal O}(x,x)=\delta_B(\tilde a)\sum_{\gamma\in\mathcal O\cap A(\Q)}
\int_{N(\A)}f(\tilde k^{-1}\gamma n\tilde k)\,dn,
\]
and $N(\Q)\backslash N(\A)$ has volume $1$. The integral over
$A(\Q)\backslash\overline{A(\A)}$ of $\Gamma'_B(H(\tilde a)-T,-H(\tilde
ky))$ is $c_k\,u'_B(\tilde k,y)$, for every $T$. Since $\delta_B(\gamma)=1$
for $\gamma\in A(\Q)$, the difference is $c_k\sum_\gamma f_{B,y}(\gamma)$.
All integrals converge absolutely: the sum over $\gamma$ is finite and
$\Gamma'_B(\cdot,Y)$ is compactly supported.

For $\chi$ the argument is the same, with~\eqref{eq:kchi} and
$K_{B,\chi}$. By~\eqref{eq:KBchi}, and because an orthonormal basis
for~\eqref{eq:HBmu} satisfies $c_k\int_{\overline{K_0}}|\varphi|^2=1$, the
difference is
\[
\sum_{\mu\in\chi}\frac1{2\pi}\int_\R\operatorname{tr}\bigl(U_y\,
I(\mu_{it},f)\bigr)\,dt=\sum_{\mu\in\chi}\frac1{2\pi}\int_\R
\widehat{f_{B,y}}(\mu_{it})\,dt=J^A_\chi(f_{B,y})
\]
by Lemma~\ref{lem:uprime}(c); the interchange of integrals is justified
as in~\cite[\S3]{A6}, by Theorem~\ref{thm:Jchi-poly}.
\end{proof}

Written as a sum over all parabolic subgroups $Q\supset A$, as
in~\cite{A9,A10}, the formula is $J(f^y)=\sum_Q|W_0^{M_Q}||W_0|^{-1}
J^{M_Q}(f_{Q,y})$. Here $Q$ runs over $G$, $B$ and $\bar B$; the last
two contribute $\frac12J^A(f_{B,y})$ each, because
$J^A(f_{\bar B,y})=J^A(f_{B,y})$.

The local distributions vary in the same way. Recall $J_G$, $J_A$ and the
weight $v_A$ from~\eqref{eq:JG}--\eqref{eq:JA}, and $J_A(z,f)$ from
Lemma~\ref{lem:JA-central}.

\begin{lemma}\label{lem:variance-local}
Let $f\in C_c^\infty(\GS)_-$ and $x,y\in\GS$.
\begin{enumerate}[label=(\alph*)]
\item If $x=n\tilde a\tilde k$ and $wx=n'\tilde a'\tilde k'$ are Iwasawa
  decompositions, then
  $v_A(xy)=v_A(x)-H_S(\tilde ky)-H_S(\tilde k'y)$.
\item If $\gamma\in\AS$ is regular (that is, $\gamma$ lies over $a(t)$
  with $t_v^2\ne1$ for all $v\in S_k$), then $J_G(\gamma,f^y)=J_G(\gamma,f)$
  and
  \[
  J_A(\gamma,f^y)=J_A(\gamma,f)+f_{B,y}(\gamma)+f_{B,y}(w\gamma w^{-1}).
  \]
\item For $z\in\{\pm1\}$, $J_A(z,f^y)=J_A(z,f)+2f_{B,y}(z)$.
\end{enumerate}
\end{lemma}

\begin{proof}
(a) $H_S(xy)=H_S(x)+H_S(\tilde ky)$ and $H_S(wxy)=H_S(wx)+H_S(\tilde
k'y)$; subtract from $v_A(x)=-H_S(x)-H_S(wx)$. (b) Substituting $x\mapsto
xy$ in~\eqref{eq:JA} gives $J_A(\gamma,f^y)=|D(\gamma)|^{1/2}\int
f(x^{-1}\gamma x)v_A(xy)\,dx$, and likewise for $J_G$ without the weight.
By~(a) the difference from $J_A(\gamma,f)$ is the sum of two terms. In the
first, with $x=n\tilde k$ and $n^{-1}\gamma n=\gamma n'$, where
$dn=|1-t^{-2}|^{-1}dn'$ and $|D(\gamma)|^{1/2}=\delta_B(\gamma)^{1/2}
|1-t^{-2}|$, we get $f_{B,y}(\gamma)$. In the second, substitute
$x\mapsto w^{-1}x$; this replaces $\gamma$ by $w\gamma w^{-1}$ and
$\tilde k'$ by the $\KS$-component of $x$, and
$D(w\gamma w^{-1})=D(\gamma)$. (c) Apply~(b) to $a_t$ in
Lemma~\ref{lem:JA-central}(b): $J_G(a_t,\cdot)$ and $r(t)$ do not change,
and $f_{B,y}(a_t)+f_{B,y}(wa_tw^{-1})\to2f_{B,y}(z)$ as $t\to1$, because
$f_{B,y}$ is continuous and $z$ is central.
\end{proof}

Parts~(b) and~(c) are the case $M=A$ of Arthur's formula
$J_M(\gamma,f^y)=\sum_{Q\in\mathcal F(M)}J_M^{M_Q}(\gamma,f_{Q,y})$
\cite{A6,A9}, with $\mathcal F(A)=\{G,B,\bar B\}$ and
$f_{\bar B,y}(\gamma)=f_{B,y}(w\gamma w^{-1})$.

On the spectral side everything follows from one identity. It holds for
the global operator $M(w,\lambda)$, for the local operators and for their
normalized versions alike, so we state it for any intertwining family.

\begin{lemma}\label{lem:conj-identity}
Let $\pi\in\Pi_u(\AS)$, and let $\mathcal M(\lambda)\colon V(\pi)\to
V(\pi^w)$ be a meromorphic family of invertible operators with
$\mathcal M(\lambda)I(\pi_\lambda,y)=I(\pi^w_{-\lambda},y)\mathcal
M(\lambda)$ for all $y\in\GS$. For $y\in\GS$ let $H_y$ be multiplication
by $\tilde k\mapsto H_S(\tilde ky)$, on $V(\pi)$ and on $V(\pi^w)$. Then,
with $\mathcal M'=d\mathcal M/d\lambda$,
\begin{equation}\label{eq:conj-identity}
I(\pi_\lambda,y)\,\mathcal M(\lambda)^{-1}\mathcal M'(\lambda)\,
I(\pi_\lambda,y)^{-1}=\mathcal M(\lambda)^{-1}\mathcal M'(\lambda)+H_y
+\mathcal M(\lambda)^{-1}H_y\,\mathcal M(\lambda),
\end{equation}
and for $f\in C_c^\infty(\GS)_-$,
\begin{equation}\label{eq:trace-variance}
-\operatorname{tr}\bigl(\mathcal M^{-1}\mathcal M'\,I(\pi_\lambda,f^y)\bigr)
=-\operatorname{tr}\bigl(\mathcal M^{-1}\mathcal M'\,I(\pi_\lambda,f)\bigr)
+\widehat{f_{B,y}}(\pi_\lambda)+\widehat{f_{B,y}}(\pi^w_{-\lambda}).
\end{equation}
The same holds over $\A$, for $\mathcal H_B(\mu)$, $I(\mu_\lambda)$ and
$\GA$.
\end{lemma}

\begin{proof}
In the compact picture, $(I(\pi_\lambda,y)\varphi)(\tilde k)=\pi(\tilde
a)e^{(\lambda+1)H_S(\tilde ky)}\varphi(\tilde k_1)$ for $\tilde
ky=n\tilde a\tilde k_1$, so $\frac{d}{d\lambda}I(\pi_\lambda,y)=H_y
I(\pi_\lambda,y)$ and $\frac{d}{d\lambda}I(\pi^w_{-\lambda},y)=-H_y
I(\pi^w_{-\lambda},y)$. Differentiating the intertwining relation,
\[
\mathcal M'I(\pi_\lambda,y)+\mathcal MH_yI(\pi_\lambda,y)
=-H_yI(\pi^w_{-\lambda},y)\mathcal M+I(\pi^w_{-\lambda},y)\mathcal M'.
\]
Multiply by $\mathcal M^{-1}$ on the left and by $I(\pi_\lambda,y)^{-1}$
on the right, and use $\mathcal M^{-1}I(\pi^w_{-\lambda},y)=I(\pi_\lambda,
y)\mathcal M^{-1}$; this is~\eqref{eq:conj-identity}. Since
$I(\pi_\lambda,f^y)=I(\pi_\lambda,y)^{-1}I(\pi_\lambda,f)I(\pi_\lambda,y)$,
the left side of~\eqref{eq:trace-variance} is $-\operatorname{tr}\bigl(
I(\pi_\lambda,y)\mathcal M^{-1}\mathcal M'I(\pi_\lambda,y)^{-1}
I(\pi_\lambda,f)\bigr)$. Insert~\eqref{eq:conj-identity}; the last term
contributes $-\operatorname{tr}(H_y\,\mathcal MI(\pi_\lambda,f)\mathcal
M^{-1})=-\operatorname{tr}(H_yI(\pi^w_{-\lambda},f))$, and $-H_y=U_y$, so
Lemma~\ref{lem:uprime}(c) applies to both terms.
\end{proof}

With $\mathcal M=M(w,\lambda)$, formula~\eqref{eq:trace-variance},
Theorem~\ref{thm:Jchi} and the invariance of $\operatorname{tr}(R(f)|_{V_b})$
and of $\operatorname{tr}(M(w,0)I(\mu_0,f))$ give a second proof of the
formula for $J_\chi(f^y)$ in Lemma~\ref{lem:variance-global}. (The
operator $M(w,0)$ commutes with $I(\mu_0,y)$ when $\mu^w=\mu$, so
$f\mapsto\operatorname{tr}(M(w,0)I(\mu_0,f))$ is invariant.)

\subsection{Normalized intertwining operators}\label{sec:normalized}

The correction terms are built from the logarithmic derivative of the
intertwining operator. The global operator $M(w,\lambda)$ is not a
product of local operators that are holomorphic on $i\R$: the local
operators have poles there, and the Euler product of their unramified
eigenvalues does not converge. One divides each local operator by a
scalar normalizing factor; the product of these factors is the global
function $m(\mu,\lambda)$, and what remains is a finite product of
normalized local operators, which are unitary and holomorphic on $i\R$.

Let $v$ be a place of $k$, $\pi_v$ a genuine unitary character of
$\overline{A(k_v)}$, and
\[
\bigl(M_v(\pi_v,\lambda)\varphi\bigr)_{-\lambda}(x)=\int_{k_v}
\varphi_\lambda\bigl((w,1)^{-1}(n(u),1)x\bigr)\,du,\qquad
M_v(\pi_v,\lambda)\colon V(\pi_v)\to V(\pi_v^w),
\]
as in Lemma~\ref{lem:intertwining}; the integral converges for
$\mathrm{Re}\,\lambda>0$ and continues meromorphically.

\begin{theorem}[Arthur~{\cite[Thm.~2.1]{A11}}; Li~{\cite[\S3]{LiII}} for
covers]\label{thm:normalization}
There are meromorphic functions $r_v(\pi_v,\lambda)$ of $\lambda\in\C$,
with $r_v(\pi_{v,\nu},\lambda)=r_v(\pi_v,\lambda+\nu)$ for $\nu\in i\R$,
such that the operators $R_v(\pi_v,\lambda)=r_v(\pi_v,\lambda)^{-1}
M_v(\pi_v,\lambda)$ have the following properties.
\begin{enumerate}[label=(N\arabic*)]
\item\label{N1} $R_v(\pi_v,\lambda)I(\pi_{v,\lambda},x)=I(\pi^w_{v,-\lambda},x)
  R_v(\pi_v,\lambda)$ for $x\in\overline{SL(2,k_v)}$.
\item\label{N2} $R_v(\pi_v^w,-\lambda)R_v(\pi_v,\lambda)=1$, and for
  $\lambda\in i\R$ the operator $R_v(\pi_v,\lambda)$ is holomorphic and
  unitary.
\item\label{N3} On each $\overline{K_v}$-isotypic subspace,
  $R_v(\pi_v,\lambda)$ is a rational function of $q_v^{-\lambda}$ if $v$
  is finite, and a rational function of $\lambda$ if $v$ is archimedean.
  In the archimedean case the norms of $R_v^{-1}\frac{d}{d\lambda}R_v$ on
  the $\overline{K_v}$-types, for $\lambda\in i\R$, are bounded by a
  polynomial in the highest weight of the type and in $|\lambda|$ (for
  real $v$ this is explicit in~\eqref{eq:rho-m}).
\item\label{N4} If $v$ is finite and non-dyadic, $\psi_v$ is unramified
  and $\pi_v$ is trivial on $A(\OO_v)^*$, then
  $R_v(\pi_v,\lambda)\varphi^0_v=\varphi^{0,w}_v$ for all $\lambda$.
\end{enumerate}
\end{theorem}

These are Arthur's conditions $(R_1)$--$(R_8)$ of~\cite[\S2]{A11} for the
pair $(B,\bar B)$, written with $w$: $M_v(\pi_v,\lambda)$ is the
composition of the operator $J_{\bar B|B}(\pi_{v,\lambda})$ of~\cite{A11}
with left translation by $w$, which does not depend on $\lambda$.
At the real places everything is explicit for the double cover
(Lemma~\ref{lem:explicit-r}(b),(c)); for the universal cover of
$SL(2,\R)$ see Sally~\cite{Sa}. The factors $r_v$ are not
unique. Three statements about them should be kept apart.
\begin{enumerate}[label=(\roman*)]
\item \emph{Existence.} At every place there are factors $r_v$ with
  \ref{N1}--\ref{N4}. This is Theorem~\ref{thm:normalization}, and it is
  all that the construction of the invariant distributions uses.
\item \emph{An explicit formula at a given place.} We have one at the
  finite places covered by~\ref{N4}, at the real places and at the
  complex places (Lemma~\ref{lem:explicit-r} and the remark after
  Example~\ref{ex:real-R}). We do not have one at the places of
  $S_{\mathrm{bad}}(\mu)$ defined below (Remark~\ref{rem:bad-places}).
\item \emph{Compatibility with the global operator.} For any choice of
  local factors, the product $m(\mu,\lambda)$ in~\eqref{eq:global-m} is
  defined, $M(w,\lambda)=m\cdot R_S$, and $m$ has the properties of
  Corollary~\ref{cor:global-m}(b). So no compatibility condition has to be
  imposed on the local choices; but $m$, $m'/m$ and the distributions
  $I_A(\gamma,f)$ depend on them, and $m$ is a known function only up to
  the unknown factors at $S_{\mathrm{bad}}(\mu)$.
\end{enumerate}
At most places there is a natural choice:

\begin{lemma}\label{lem:explicit-r}
Write $\pi_v=\chi_{\psi,v}\cdot(\chi_v\circ p)$.
\begin{enumerate}[label=(\alph*)]
\item Let $v$, $\psi_v$ and $\pi_v$ be as in~\ref{N4}. Then
  \[
  r_v(\pi_v,\lambda)=\frac{L_v(2\lambda,\chi_v^2)}{L_v(2\lambda+1,\chi_v^2)}
  \]
  satisfies~\ref{N1}--\ref{N4}.
\item Let $v$ be real and $\chi_v=\mathrm{sgn}^\epsilon$,
  $\epsilon\in\{0,1\}$. Let $\theta\mapsto\tilde k_\theta$ be the
  one-parameter subgroup of $\overline{K_v}$ above $k_\theta$
  (Example~\ref{ex:uprime}). The $\overline{K_v}$-types of $I(\pi_v)$ are
  the characters $\tilde k_\theta\mapsto e^{im\theta}$, $m\in m_0+2\Z$,
  each with multiplicity one, where $m_0\in\{\frac12,-\frac12\}$ is
  determined by $e^{i\pi m_0}=\pi_v(a(-1),1)$. On the line of weight $m$,
  with $du$ Lebesgue measure,
  \begin{equation}\label{eq:real-M}
  M_v(\pi_v,\lambda)=e^{i\pi m/2}\,\frac{\pi\,2^{1-\lambda}\,\Gamma(\lambda)}
  {\Gamma\bigl(\frac{\lambda+1+m}2\bigr)\Gamma\bigl(\frac{\lambda+1-m}2\bigr)} .
  \end{equation}
  With $r_v(\pi_v,\lambda)=\sqrt2\,e^{i\pi m_0/2}\,L_v(2\lambda)/L_v(2\lambda+1)$,
  where $L_v(s)=\pi^{-s/2}\Gamma(s/2)$, the operator $R_v(\pi_v,\lambda)$
  is the identity on the weight $m_0$, and on the weight $m_0+2j$ it is
  \begin{equation}\label{eq:real-R}
  \prod_{i=0}^{j-1}\frac{a_i-\lambda}{a_i+\lambda}\quad(j>0,\ a_i=1+m_0+2i),
  \qquad
  \prod_{i=0}^{|j|-1}\frac{b_i-\lambda}{b_i+\lambda}\quad(j<0,\ b_i=1-m_0+2i).
  \end{equation}
  These operators satisfy~\ref{N1}--\ref{N3}.
\item Let $v$ be real. Every genuine unitary character of
  $\overline{A(k_v)}$ is of the form $\pi_{v,i\tau}=\pi_v\,e^{i\tau H_v}$
  with $\pi_v$ as in~(b) and a unique $\tau\in\R$; that is,
  $\chi_v=\mathrm{sgn}^\epsilon|\cdot|^{i\tau}$. Then
  $V(\pi_{v,i\tau})=V(\pi_v)$, $I(\pi_{v,i\tau},\lambda)=I(\pi_v,\lambda+
  i\tau)$ and $M_v(\pi_{v,i\tau},\lambda)=M_v(\pi_v,\lambda+i\tau)$. With
  $r_v(\pi_{v,i\tau},\lambda)=r_v(\pi_v,\lambda+i\tau)$, the normalized
  operator $R_v(\pi_{v,i\tau},\lambda)$ is given by~\eqref{eq:real-R} with
  $\lambda+i\tau$ in place of $\lambda$, and satisfies~\ref{N1}--\ref{N3}.
  In terms of $\chi_v$,
  \[
  r_v(\pi_{v,i\tau},\lambda)=\sqrt2\,e^{i\pi m_0/2}\,
  \frac{L_v(2\lambda,\chi_v^2)}{L_v(2\lambda+1,\chi_v^2)},\qquad
  L_v(s,|\cdot|^{2i\tau})=L_v(s+2i\tau).
  \]
\end{enumerate}
\end{lemma}

\begin{proof}
(a) \ref{N1} holds for $M_v$, and~\ref{N4} is
Lemma~\ref{lem:intertwining}(b). The induced representation
$I(\pi_{v,\lambda})$ is irreducible for $\lambda$ outside a discrete
set, so $M_v(\pi_v^w,-\lambda)M_v(\pi_v,\lambda)$, which commutes with
it, is a scalar; on $\varphi_v^0$ it is $r_v(\pi_v^w,-\lambda)
r_v(\pi_v,\lambda)$ by Lemma~\ref{lem:intertwining}(b). So
$R_v(\pi^w_v,-\lambda)R_v(\pi_v,\lambda)=1$. For $\lambda\in i\R$ the
adjoint of $M_v(\pi_v,\lambda)$ is $M_v(\pi_v^w,-\lambda)$~\cite[\S1]{A11}. (This
uses that $(w,1)_v^2$ acts trivially on $I(\pi_{v,\lambda})$. Here it
does: $(-1,-1)_v=1$ at a non-dyadic finite place, so $(w,1)_v^2=(-1,1)$,
a central element lying in $A(\OO_v)^*$, and it acts by the scalar
$\pi_v(-1,1)=1$. At a real place $(-1,1)$ has order $4$ and acts by $\pm
i$ on a genuine representation, by Lemma~\ref{lem:local}(e); part~(b)
does not use the adjoint formula.) Moreover
$\overline{r_v(\pi_v,\lambda)}=r_v(\pi_v^w,-\lambda)$ for $\lambda\in i\R$;
hence $R_v$ is unitary on $i\R$ wherever it is defined, so it is bounded
there and has no poles. Rationality is~\cite{LiII}.

(b) In the model of Section~\ref{sec:local}, $\tilde k_\theta=(k_\theta,1)$
for $0<\theta<\pi$: for $\theta_1,\theta_2,\theta_1+\theta_2\in(0,\pi)$
the three values of $x(\cdot)$ in~\eqref{eq:beta} are $-\sin$ of the
angles, all negative, so $\beta(k_{\theta_1},k_{\theta_2})=1$. Then
$\tilde k_\pi=(k_{\pi/2},1)^2=(-1,1)=(a(-1),1)$, which gives the
$\overline{K_v}$-types, since $\overline{A(k_v)}\cap\overline{K_v}$ is
generated by $\tilde k_\pi$. Also $(w,1)^{-1}=\tilde k_{\pi/2}$, and
$\tilde k_{\pi/2}(n(u),1)=(n',1)(a(s),1)\tilde k_{\theta(u)}$ with
$s=(1+u^2)^{-1/2}>0$ and $\theta(u)=\frac\pi2+\arctan u$, by
continuity in $u$. So on a vector of weight $m$ the
integral is $e^{i\pi m/2}\int_\R(1+u^2)^{-(\lambda+1)/2}e^{im\arctan u}\,du$,
and $\int_{-\pi/2}^{\pi/2}\cos^{\lambda-1}\vartheta\,e^{im\vartheta}\,
d\vartheta$ is the classical integral in~\eqref{eq:real-M}. By the
duplication formula, $\Gamma(\frac{\lambda}2+\frac34)\Gamma(\frac\lambda2+
\frac14)=2^{\frac12-\lambda}\sqrt\pi\,\Gamma(\lambda+\frac12)$,
so~\eqref{eq:real-M} at $m=\pm\frac12$ is $\sqrt2\,e^{i\pi m/2}\sqrt\pi\,
\Gamma(\lambda)/\Gamma(\lambda+\frac12)$. The quotient of~\eqref{eq:real-M}
at $m+2$ and at $m$ is $-(\lambda-1-m)/(\lambda+1+m)$,
which gives~\eqref{eq:real-R}. Each factor in~\eqref{eq:real-R} has
absolute value $1$ on $i\R$ and no pole there, because $a_i$ and $b_i$
are nonzero real numbers; \ref{N2} and~\ref{N3} follow, and~\ref{N1}
holds for $M_v$.

(c) The unitary characters of $\R^\times$ are the
$\mathrm{sgn}^\epsilon|\cdot|^{i\tau}$, and $H_v(a(t),\zeta)=\log|t|$. The
space $V(\pi)$ depends only on the restriction of $\pi$ to
$\overline{A(k_v)}\cap\overline{K_v}$, on which $e^{i\tau H_v}=1$, and
$\varphi_\lambda$ formed with $\pi_{v,i\tau}$ is $\varphi_{\lambda+i\tau}$
formed with $\pi_v$; this gives the first three statements. The factors
in~\eqref{eq:real-R} with $\lambda+i\tau$ in place of $\lambda$ still have
absolute value $1$ and no poles for $\lambda\in i\R$. The last formula is
$L_v(s,|\cdot|^{2i\tau})=\pi^{-(s+2i\tau)/2}\Gamma((s+2i\tau)/2)$.
\end{proof}

\begin{example}\label{ex:real-R}
In Lemma~\ref{lem:explicit-r}(b) take $m_0=\frac12$, which is the case
$\pi_v=\chi_{\psi,v}$ for $\psi_v(x)=e^{-2\pi ix}$
(Example~\ref{ex:weil-Q}). The weights are $\frac12,\frac52,\frac92,\dots$
and $-\frac32,-\frac72,\dots$, and
\[
R_v(\lambda)\big|_{5/2}=\frac{\frac32-\lambda}{\frac32+\lambda},\qquad
R_v(\lambda)\big|_{-3/2}=\frac{\frac12-\lambda}{\frac12+\lambda},\qquad
R_v^{-1}R_v'(\lambda)\big|_{5/2}=-\frac{3}{\frac94-\lambda^2}.
\]
$R_v(\lambda)$ has a zero at $\lambda=\frac12$ on every weight below
$\frac12$ and is $1$ on the weight $\frac12$: this is the reducibility of
$I(\chi_{\psi,v,1/2})$, whose quotient with lowest weight $\frac12$ is
the even Weil representation of Theorem~\ref{thm:residual}. In
general, by~\eqref{eq:real-R}, on the weight $m_0+2j$ and for $\lambda=it$,
\begin{equation}\label{eq:rho-m}
-R_v^{-1}R_v'(it)=\rho_j(t):=\sum_i\frac{2c_i}{c_i^2+t^2}\;\ge0,
\end{equation}
the sum over $c_i=a_i$, $0\le i<j$, if $j>0$, and over $c_i=b_i$,
$0\le i<|j|$, if $j<0$; and $\rho_0=0$. For $\pi_{v,i\tau}$ as in
Lemma~\ref{lem:explicit-r}(c) the value is $\rho_j(t+\tau)$. For $SL(2,\R)$ the weights are
integers and the same computation gives the familiar formula with
$L_v(\lambda)/L_v(\lambda+1)$.
\end{example}

At a complex place the cover splits and $R_v$ is the normalized operator
of $SL(2,\C)$~\cite{A11}; the factor
$L_v(\lambda,\chi_v)/L_v(\lambda+1,\chi_v)$ there equals $2L_v(2\lambda,
\chi_v^2)/L_v(2\lambda+1,\chi_v^2)$. So the factor of
Lemma~\ref{lem:explicit-r} is available at every place outside the finite
set
\[
S_{\mathrm{bad}}(\mu)=\{v\text{ dyadic}\}\cup\{v\mid\psi_v\text{ ramified}\}
\cup\{v\text{ finite}\mid\mu_v\text{ nontrivial on }A(\OO_v)^*\}.
\]

\paragraph{The global normalizing factor.} Let
$\mu=\chi_\psi\cdot(\chi\circ p)\in\Pi_{\mathrm{aut}}(X_A)$, and fix
factors $r_v(\mu_v,\lambda)$ as in Theorem~\ref{thm:normalization} for
$v\in S_{\mathrm{bad}}(\mu)$ and as in Lemma~\ref{lem:explicit-r} (and
the remark on complex places) for all other $v$. At a real place the
component $\mu_v$ is $\chi_{\psi,v}\cdot(\chi_v\circ p)$ with
$\chi_v=\mathrm{sgn}^\epsilon|\cdot|^{i\tau_v}$, and
Lemma~\ref{lem:explicit-r}(c) applies; for $k=\Q$ one has $\tau_\infty=0$,
because $\mu$ is trivial on $A_B(\R)^0$, but for other fields the
$\tau_v$ need not vanish. For a finite set
$S_1\supset S_{\mathrm{bad}}(\mu)$ of places of $k$ containing the
archimedean ones, put
\begin{equation}\label{eq:global-m}
m(\mu,\lambda)=m^{S_1}(\mu,\lambda)\prod_{v\in S_1}r_v(\mu_v,\lambda),
\qquad m^{S_1}(\mu,\lambda)=\frac{L^{S_1}(2\lambda,\chi^2)}
{L^{S_1}(2\lambda+1,\chi^2)}
\end{equation}
as in~\eqref{eq:mS}. By Lemma~\ref{lem:explicit-r}(a) this does not
depend on $S_1$.

\begin{corollary}\label{cor:global-m}
Let $f=f_S\otimes e^S$ and $\mu\in\Pi_{\mathrm{aut}}(X_A)$ with $\mu_v$
trivial on $A(\OO_v)^*$ for $v\notin S_k$, and let $R_S(\mu_S,\lambda)=
\bigotimes_{v\in S_k}R_v(\mu_v,\lambda)$.
\begin{enumerate}[label=(\alph*)]
\item On $\bigotimes_{v\in S_k}V(\mu_v)\otimes\bigotimes_{v\notin S_k}
  \varphi_v^0$, $M(w,\lambda)=m(\mu,\lambda)\,R_S(\mu_S,\lambda)\otimes1$,
  and for $\lambda\in i\R$
  \[
  \operatorname{tr}\bigl(M(w,\lambda)^{-1}M'(w,\lambda)I(\mu_\lambda,f)\bigr)
  =\frac{m'}{m}(\mu,\lambda)\operatorname{tr}I(\mu_{S,\lambda},f_S)+
  \operatorname{tr}\bigl(R_S^{-1}R_S'(\mu_S,\lambda)I(\mu_{S,\lambda},f_S)
  \bigr).
  \]
\item $m(\mu^w,-\lambda)m(\mu,\lambda)=1$; for $\lambda\in i\R$,
  $m(\mu,\lambda)$ is holomorphic of absolute value $1$, $\frac{m'}m(\mu,
  \lambda)$ is real, and $\frac{m'}m(\mu^w,-\lambda)=\frac{m'}m(\mu,\lambda)$.
\item With $\Lambda^{\mathrm{bad}}(s,\chi^2)=\prod_{v\notin
  S_{\mathrm{bad}}(\mu)}L_v(s,\chi_v^2)$, the product over all places
  outside $S_{\mathrm{bad}}(\mu)$, archimedean ones included,
  \[
  \frac{m'}m(\mu,\lambda)=2\frac{(\Lambda^{\mathrm{bad}})'}
  {\Lambda^{\mathrm{bad}}}(2\lambda,\chi^2)-2\frac{(\Lambda^{\mathrm{bad}})'}
  {\Lambda^{\mathrm{bad}}}(2\lambda+1,\chi^2)+\sum_{v\in S_{\mathrm{bad}}
  (\mu)}\frac{r_v'}{r_v}(\mu_v,\lambda).
  \]
\end{enumerate}
\end{corollary}

\begin{proof}
(a) Lemma~\ref{lem:intertwining}(a)--(b), with $S_1=S_k$, and
$M_v=r_vR_v$; then the Leibniz rule, as for Corollary~\ref{cor:mprime}.
(b) $M(w,-\lambda)M(w,\lambda)=1$ (Theorem~\ref{thm:eisenstein}(c))
and~\ref{N2} give the first statement, and the unitarity of $M(w,\lambda)$
and of $R_S$ on $i\R$ the second; differentiate. (c) is the logarithmic
derivative of~\eqref{eq:global-m}.
\end{proof}

\begin{remarkrev}\label{rem:bad-places}
\emph{The factors at the bad places.} For $v\in S_{\mathrm{bad}}(\mu)$
(dyadic places, and places where $\psi_v$ or $\mu_v$ is ramified) we use
the factors $r_v$ of~\cite{LiII} without making them explicit. The
expected choice is again $L_v(2\lambda,\chi_v^2)/L_v(2\lambda+1,\chi_v^2)$
times an exponential factor, which would give
$m(\mu,\lambda)=c\,q^{-\lambda}\Lambda(2\lambda,\chi^2)/\Lambda(2\lambda+1,
\chi^2)$ with the completed Hecke $L$-function and constants $c$, $q>0$,
as the Shimura correspondence suggests; we have not verified~\ref{N2}
for this choice at those places. The distributions $I_A(\gamma,f)$ and the function $m'/m$
below depend on the choice of the $r_v$; the invariant trace formula
holds for each choice. For local coefficients and Plancherel measures of
the genuine principal series of covers of $SL(2,k_v)$ see
Szpruch~\cite{Sz09} and Goldberg--Szpruch~\cite{GS16}; we have not derived
the factors $r_v$ from those results.
\end{remarkrev}

\subsection{The map \texorpdfstring{$\varphi_A$}{phi\_A} and the
distributions \texorpdfstring{$I_A(\gamma,f)$}{I\_A(gamma,f)}}
\label{sec:IAgamma}

We now define the invariant distributions. Fix normalizing factors as in
Section~\ref{sec:normalized}, and for $\pi\in\Pi_u(\AS)$ let
$R_S(\pi,\lambda)=\bigotimes_{v\in S_k}R_v(\pi_v,\lambda)$. For
$f\in C_c^\infty(\GS)_-$ put
\begin{equation}\label{eq:phiA}
\varphi_A(f,\pi)=-\operatorname{tr}\Bigl(R_S(\pi,0)^{-1}\,\frac{dR_S}
{d\lambda}(\pi,0)\,I(\pi,f)\Bigr),\qquad\pi\in\Pi_u(\AS).
\end{equation}
By Theorem~\ref{thm:normalization}, $\varphi_A(f,\pi_\lambda)=
-\operatorname{tr}(R_S^{-1}R_S'(\pi,\lambda)I(\pi_\lambda,f))$ for
$\lambda\in i\R$, and $R_S^{-1}R_S'=\sum_{v\in S_k}R_v^{-1}R_v'$. This is
Arthur's map $\phi_M(f,\pi)=\operatorname{tr}(\mathcal R_M(\pi,P)I_P(\pi,
f))$~\cite[\S7]{A6}, \cite{A9} for $M=A$, $P=B$: the $(G,A)$-family is
$\mathcal R_B(\Lambda)=1$, $\mathcal R_{\bar B}(\Lambda)=R_S(\pi,0)^{-1}
R_S(\pi,\Lambda)$, and, as for~\eqref{eq:MT},
\[
\mathcal R_A(\pi,B)=\lim_{\Lambda\to0}\Bigl(\frac{1}{\theta_B(\Lambda)}+
\frac{R_S(\pi,0)^{-1}R_S(\pi,\Lambda)}{\theta_{\bar B}(\Lambda)}\Bigr)
=-R_S(\pi,0)^{-1}R_S'(\pi,0),
\]
with $\theta_B(\Lambda)=\Lambda=-\theta_{\bar B}(\Lambda)$ in the
normalization of Remark~\ref{rem:alpha}.

\begin{proposition}\label{prop:phiA}
Let $f\in C_c^\infty(\GS)_-$ and $y\in\GS$.
\begin{enumerate}[label=(\alph*)]
\item $\varphi_A(f,\cdot)$ is continuous and integrable on $\Pi_u(\AS)$.
  More precisely, it vanishes unless the conductors of the $\pi_v$, $v$
  finite, are bounded in terms of $f$, and it decreases rapidly in the
  continuous parameters of $\pi$.
\item $\varphi_A(f,\pi^w)=\varphi_A(f,\pi)$.
\item $\varphi_A(f^y,\pi)=\varphi_A(f,\pi)+\widehat{f_{B,y}}(\pi)+
  \widehat{f_{B,y}}(\pi^w)$.
\item If $v\notin S_k$ is as in~\ref{N4} and $S'=S\cup\{v\}$, then
  $\varphi_A(f\otimes e_v,\pi\otimes\pi_v)=\varphi_A(f,\pi)$ if $\pi_v$ is
  trivial on $A(\OO_v)^*$, and $=0$ otherwise.
\end{enumerate}
\end{proposition}

\begin{proof}
(a) $f$ is bi-invariant under an open subgroup at the finite places, which
bounds the conductors. $\pi\mapsto I(\pi,f)$ is the operator-valued
Fourier transform of a compactly supported smooth function, so its trace
norm on each $\KS$-type decreases rapidly in $\pi$ and in the
$\KS$-type; by~\ref{N2} and~\ref{N3}, $R_S^{-1}R_S'$ is continuous in
$\pi$ and of polynomial growth. (b) By~\ref{N2}, $R_S(\pi^w,\nu)=R_S(\pi,-\nu)^{-1}$, so
$R_S(\pi^w,\nu)^{-1}\frac{d}{d\nu}R_S(\pi^w,\nu)$ at $\nu=0$ is
$R_S'R_S^{-1}$ at $(\pi,0)$, and
$\operatorname{tr}(R_S'R_S^{-1}I(\pi^w,f))=\operatorname{tr}(R_S^{-1}R_S'
R_S^{-1}I(\pi^w,f)R_S)=\operatorname{tr}(R_S^{-1}R_S'I(\pi,f))$
by~\ref{N1}. (c) is~\eqref{eq:trace-variance} with $\mathcal
M=R_S(\pi,\lambda)$ at $\lambda=0$. (d) $I(\pi_v,e_v)$ is the projection
onto $\C\varphi^0_v$ (or $0$), and $R_v^{-1}R_v'\varphi_v^0=0$
by~\ref{N4}.
\end{proof}

For $\tilde a\in\AS$ put
\begin{equation}\label{eq:PhiA}
\Phi_A(\tilde a,f)=\int_{\Pi_u(\AS)}\varphi_A(f,\pi)\,\pi(\tilde a)^{-1}\,
d\pi,
\end{equation}
the function on $\AS$ whose Fourier transform is $\varphi_A(f,\cdot)$,
and for $\gamma\in\AS$ regular and $z\in\{\pm1\}$ define
\begin{equation}\label{eq:IAgamma}
I_A(\gamma,f)=J_A(\gamma,f)-\Phi_A(\gamma,f),\qquad
I_A(z,f)=J_A(z,f)-\Phi_A(z,f).
\end{equation}

\begin{theorem}\label{thm:IA}
Let $f\in C_c^\infty(\GS)_-$.
\begin{enumerate}[label=(\alph*)]
\item $\Phi_A(\cdot,f)$ is a continuous, bounded, anti-genuine function
  on $\AS$, smooth in the archimedean coordinates and locally constant in
  the others, with
  $\Phi_A(w\tilde aw^{-1},f)=\Phi_A(\tilde a,f)$. If $f=\bigotimes_vf_v$,
  its restriction to each set $\{\tilde a\in\AS\mid H_S(\tilde a)=X\}$ is
  compactly supported. It is not compactly supported on $\AS$ in general.
\item The distributions $I_A(\gamma,\cdot)$ and $I_A(z,\cdot)$ are
  invariant: $I_A(\gamma,f^y)=I_A(\gamma,f)$ for all $y\in\GS$.
\item $I_A(w\gamma w^{-1},f)=I_A(\gamma,f)$.
\item If $v\notin S_k$ is as in~\ref{N4}, $S'=S\cup\{v\}$, and
  $\gamma_v\in A(\OO_v)^*$ lies over $a(t)$ with $|t^2-1|_v=1$, then
  $I_A(\gamma\gamma_v,f\otimes e_v)=I_A(\gamma,f)$.
\end{enumerate}
\end{theorem}

\begin{proof}
(a) Continuity and boundedness follow from
Proposition~\ref{prop:phiA}(a), smoothness from the rapid decrease, local
constancy from the bound on the conductors, and
the symmetry from Proposition~\ref{prop:phiA}(b) and the invariance of
$d\pi$ under $\pi\mapsto\pi^w$. For $f=\bigotimes f_v$, by the Leibniz
rule $\varphi_A(f,\pi)=\sum_{v\in S_k}\varphi_{A,v}(f_v,\pi_v)\prod_{u\ne
v}\operatorname{tr}I(\pi_u,f_u)$, and $\operatorname{tr}I(\pi_u,f_u)=
\widehat{F_u}(\pi_u)$ with $F_u(\tilde a)=\delta_B(\tilde a)^{1/2}\int
\int f_u(\tilde k^{-1}\tilde an\tilde k)\,dn\,d\tilde k$ compactly
supported on $\overline{A(k_u)}$. So $\Phi_A=\sum_v\Phi_{A,v}\otimes
\bigotimes_{u\ne v}F_u$ with $\Phi_{A,v}$ continuous. On
$\{H_S=X\}$ the $v$-th summand is supported where $\tilde a_u\in
\mathrm{supp}\,F_u$ for $u\ne v$, and then $H_v(\tilde a_v)=X-\sum_{u\ne
v}H_u(\tilde a_u)$ is bounded; this set is compact. The last statement is
Example~\ref{ex:Phi-real}.

(b) By Proposition~\ref{prop:phiA}(c) and Fourier inversion for
$f_{B,y}\in C_c^\infty(\AS)_-$ (Lemma~\ref{lem:uprime}(b)),
\[
\Phi_A(\tilde a,f^y)=\Phi_A(\tilde a,f)+f_{B,y}(\tilde a)+f_{B,y}(w\tilde
aw^{-1}),
\]
since the inverse transform of $\pi\mapsto\hat h(\pi^w)$ is $\tilde
a\mapsto h(w\tilde aw^{-1})$. Subtract this from
Lemma~\ref{lem:variance-local}(b), and at $\tilde a=z$ from
Lemma~\ref{lem:variance-local}(c).

(c) $J_A(w\gamma w^{-1},f)=J_A(\gamma,f)$ by the substitution $x\mapsto
wx$ in~\eqref{eq:JA}, as after Theorem~\ref{thm:fine}; and~(a). (d) For
such $\gamma_v$ the orbital integral of $e_v$ is $1$ and its weighted
orbital integral is $0$, because $v_{A,v}=0$ on the support (the proof of
Theorem~\ref{thm:fine}); and Proposition~\ref{prop:phiA}(d) gives
$\Phi_A(\gamma\gamma_v,f\otimes e_v)=\Phi_A(\gamma,f)$.
\end{proof}

For connected reductive groups, $I_M(\gamma,f)$ is defined in~\cite[\S2]{A9}
by $I_M(\gamma,f)=J_M(\gamma,f)-\sum_{L\in\mathcal L(M),\,L\ne G}
\hat I^L_M(\gamma,\phi_L(f))$, and for covers in~\cite[\S4.5]{LiIV}. Here
$\mathcal L(A)=\{A,G\}$, the only term is $L=A$, and
$\hat I_A^A(\gamma,\phi_A(f))$ is the value at $\gamma$ of the function
with Fourier transform $\phi_A(f)$, which is $\Phi_A(\gamma,f)$. For a
general group the corresponding step needs the trace Paley--Wiener
theorem for $L$; for $L=A$ it is Fourier inversion on an abelian group.
Two further remarks.

\emph{The space containing $\varphi_A(f)$.} Arthur's
$\phi_M$ maps $\mathcal H(G(F_S))$ to a space $\mathcal I_{\mathrm{ac}}
(M(F_S))$ of functions $\phi$ on $\Pi_{\mathrm{temp}}(M(F_S))$ for which
\[
\phi(\pi,X)=\int_{i\mathfrak a_M^*}\phi(\pi_\Lambda)\,e^{-\Lambda(X)}\,
d\Lambda,\qquad X\in\mathfrak a_M,
\]
is, for each $X$, a Paley--Wiener function of $\pi$~\cite[\S11]{A11}.
In our
case $\phi(\pi,X)$ is the Fourier transform of the restriction of
$\Phi_A(\cdot,f)$ to $\{H_S=X\}$, and Theorem~\ref{thm:IA}(a) is the
statement that $\varphi_A(f)\in\mathcal I_{\mathrm{ac}}$. This space is
strictly larger than the space $\mathcal I(\overline{A(\A_S)})$ of Fourier
transforms of compactly supported functions, for the reason given in
Section~\ref{sec:local-dist} for the spaces of test functions:
$\mathfrak a_A\ne0$ (Example~\ref{ex:Phi-real}).

\emph{Weighted characters.} The weighted character of~\cite{A9} is
$J_A(\pi,f)=\operatorname{tr}(\mathcal R_A(\pi,B)I(\pi,f))=\varphi_A(f,
\pi)$, so the invariant distribution $I_A(\pi,X,f)=J_A(\pi,X,f)-
\hat I^A_A(\pi,X,\varphi_A(f))$ vanishes for every $\pi\in\Pi_u(\AS)$.
This does not
make the spectral side of the invariant formula vanish: the terms that
remain are characters, with the coefficients of Theorem~\ref{thm:itf}.

\begin{example}\label{ex:Phi-real}
Let $v$ be real, $\pi_v$ and $m_0$ as in Lemma~\ref{lem:explicit-r}(b),
and let $P_j$ be the projection of $V(\pi_v)$ onto the weight $m_0+2j$.
For $f_v\in C_c^\infty(\overline{SL(2,\R)})_-$ let $F_j$ be the function
on $\overline{A(\R)}$ with $\widehat{F_j}(\pi)=\operatorname{tr}(P_jI(\pi,
f_v))$; then $\sum_jF_j=F_{f_v}$, the Abel transform in the proof of
Theorem~\ref{thm:IA}(a). By~\eqref{eq:rho-m}, $\varphi_{A,v}(f_v,\pi_{v,it})
=\sum_j\rho_j(t)\widehat{F_j}(\pi_{v,it})$, and since
$\frac1{2\pi}\int_\R\frac{2c}{c^2+t^2}e^{-itX}\,dt=e^{-c|X|}$ for $c>0$,
\[
\Phi_{A,v}(\tilde a,f_v)=\sum_j\int_\R\kappa_j(X)\,F_j\bigl(\tilde a\,a(e^X)
\bigr)\,dX,\qquad\kappa_j(X)=\sum_ie^{-c_i|X|},
\]
with $c_i$ as in~\eqref{eq:rho-m}. So $\Phi_{A,v}(\cdot,f_v)$ is the Abel
transform smoothed, on each $\overline{K_v}$-type, by a kernel
$\kappa_j\ge0$ that decays exponentially but is not compactly supported;
it vanishes on the weight $m_0$. For example, if $f_v$ transforms on
both sides under $\overline{K_v}$ by the character of weight $m_0+2j$
with $j\ne0$, and $F_j=F_{f_v}\ge0$ is not zero, then $\Phi_{A,v}(\tilde
a,f_v)>0$ at every $\tilde a$ in the cosets of $\{a(s)\mid s>0\}$ that
meet the support of $F_{f_v}$, although $F_{f_v}$ is compactly supported.
\end{example}

\subsection{The invariant trace formula}\label{sec:itf}

We return to $f=f_S\otimes e^S$ with $f_S\in\mathcal H(\GS)$, and to the
situation of Theorem~\ref{thm:tf}. For $\gamma\in A(\Q)$ put
$\Phi_A(\gamma,f)=\Phi_A(\gamma_S,f_S)$ if $\gamma\in K_v$ for all
$v\notin S_k$, and $\Phi_A(\gamma,f)=0$ otherwise; for
$\mu\in\Pi_{\mathrm{aut}}(X_A)$ put $\varphi_A(f,\mu_\lambda)=\varphi_A(f_S,
\mu_{S,\lambda})$ if $\mu_v$ is trivial on $A(\OO_v)^*$ for all
$v\notin S_k$, and $0$ otherwise.

\begin{lemma}\label{lem:poisson-phi}
$\Phi_A(\gamma,f)=0$ for all but finitely many $\gamma\in A(\Q)$, and
\[
\sum_{\mu\in\Pi_{\mathrm{aut}}(X_A)}\frac1{2\pi}\int_\R\varphi_A(f,\mu_{it})
\,dt=c_k\sum_{\gamma\in A(\Q)}\Phi_A(\gamma,f),
\]
the left side converging absolutely.
\end{lemma}

\begin{proof}
We may take $f_S=\bigotimes_vf_v$. Let $h=\Phi_A(\cdot,f_S)\otimes
e_A^S$ on $\overline{A(\A)}$, where $e_A^S=\bigotimes_{v\notin S_k}e_{A,v}$
is the product of the units of Lemma~\ref{lem:transfer}(b). By
Theorem~\ref{thm:IA}(a) the restriction of $h$ to $\overline{A(\A)^1}$ is
compactly supported; it is smooth at the archimedean places and locally
constant at the finite ones, because $\varphi_A(f_S,\cdot)$ decreases
rapidly and is supported on finitely many components. So
$\gamma\mapsto h(\gamma)=\Phi_A(\gamma,f)$ has finite support on the
discrete subgroup $A(\Q)$, and Theorem~\ref{thm:tf-rank0} applies to
$h|_{\overline{A(\A)^1}}$: $\sum_\mu\int_{\overline{A(\A)^1}}h\mu=c_k
\sum_\gamma h(\gamma)$. Finally $\int_{\overline{A(\A)^1}}h\mu=\frac1{2\pi}
\int_\R\hat h(\mu_{it})\,dt$ and $\hat h(\mu_{it})=\varphi_A(f,\mu_{it})$,
by Fourier inversion on $\R$. The left side converges absolutely by
Proposition~\ref{prop:phiA}(a): the $\mu$ that contribute form finitely
many cosets of a lattice in the unitary dual of $A(\Q_S)$ modulo
$i\R$.
\end{proof}

For $\mathcal O\in\mathcal C$ and $\chi\in\mathcal X$ define
\begin{equation}\label{eq:Iochi}
I_{\mathcal O}(f)=J_{\mathcal O}(f)-\frac{c_k}2\sum_{\gamma\in\mathcal O\cap
A(\Q)}\Phi_A(\gamma,f),\qquad
I_\chi(f)=J_\chi(f)-\frac12\sum_{\mu\in\chi}\frac1{2\pi}\int_\R
\varphi_A(f,\mu_{it})\,dt,
\end{equation}
the second sum being empty for $\chi\in\mathcal X_G$. That is, $I=J-J'$
with
$J'=|W_G|^{-1}(I^A)^\wedge(\varphi_A(f))$, where $(I^A_{\mathcal O})^\wedge$
and $(I^A_\chi)^\wedge$ are the distributions~\eqref{eq:JAo} transported
to the Fourier transform side.

\begin{theorem}[The invariant trace formula]\label{thm:itf}
Let $S$ and $f=f_S\otimes e^S$ be as in Theorem~\ref{thm:tf}. Then
\[
\sum_{\mathcal O\in\mathcal C}I_{\mathcal O}(f)=\sum_{\chi\in\mathcal X}
I_\chi(f),
\]
both sides converging absolutely, and:
\begin{enumerate}[label=(\alph*)]
\item for the central classes $\mathcal O_z$, $z\in\{\pm1\}$,
  \[
  I_{\mathcal O_z}(f)=a^G(S,1)f_S(z)+\sum_{\xi\in k_S^\times/k_S^{\times2}}
  a^G(S,u_\xi)\,J_G(zu_\xi,f_S)+\tfrac12\,c_k\,I_A(z,f_S);
  \]
\item for $\mathcal O$ elliptic, $I_{\mathcal O}(f)=J_{\mathcal O}(f)=
  \mathrm{vol}(T_\mu(\Q)\backslash T_\mu(\A))\,J_G(\mu_S,f_S)$;
\item for $\mathcal O$ hyperbolic, the class of $\gamma=a(t)$,
  $I_{\mathcal O}(f)=c_k\,I_A(\gamma_S,f_S)$ if $t\in\OO_v^\times$ for all
  $v\notin S_k$, and $I_{\mathcal O}(f)=0$ otherwise;
\item for $\chi$ the class of $(G,\pi)$, $I_\chi(f)=m(\pi)\operatorname{tr}
  \pi(f)$;
\item for $\chi\in\mathcal X_A$ unramified and $\mu\in\chi$,
  \[
  I_\chi(f)=-\frac1{2\pi}\int_\R\frac{m'}m(\mu,it)\,\operatorname{tr}
  I(\mu_{it},f)\,dt;
  \]
\item for $\chi=\chi_b$ ramified and $\mu=\chi_{\psi_b}$,
  \[
  I_\chi(f)=\operatorname{tr}\bigl(R(f)|_{V_b}\bigr)-\frac12\cdot\frac1{2\pi}
  \int_\R\frac{m'}m(\mu,it)\,\operatorname{tr}I(\mu_{it},f)\,dt
  +\frac14\operatorname{tr}\bigl(M(w,0)\,I(\mu_0,f)\bigr).
  \]
\end{enumerate}
Each of the distributions (a)--(f) is the restriction to $\mathcal H(\GS)$
of an invariant distribution in $f_S\in C_c^\infty(\GS)_-$. Here
$a^G(S,u_\xi)$ are the coefficients of Theorem~\ref{thm:fine}, and $m$ and
$I_A$ are formed with the normalizing factors fixed in
Section~\ref{sec:normalized}.
\end{theorem}

\begin{proof}
By Theorem~\ref{thm:tf} and Lemma~\ref{lem:poisson-phi}, $\sum_{\mathcal
O}I_{\mathcal O}(f)-\sum_\chi I_\chi(f)=-\frac12\bigl(c_k\sum_\gamma
\Phi_A(\gamma,f)-\sum_\mu\frac1{2\pi}\int\varphi_A(f,\mu_{it})\,dt\bigr)=0$.
(a) is Theorem~\ref{thm:tf}(a) and~\eqref{eq:IAgamma}, since $\mathcal O_z
\cap A(\Q)=\{z\}$. (b) $\mathcal O\cap A(\Q)$ is empty. (c) $\mathcal
O\cap A(\Q)=\{\gamma,\gamma^{-1}\}$ and $\Phi_A(\gamma^{-1},f)=\Phi_A(\gamma,
f)$ by Theorem~\ref{thm:IA}(a), because $\gamma^{-1}=w\gamma w^{-1}$; with
Theorem~\ref{thm:tf}(c) this gives $c_k(J_A-\Phi_A)(\gamma_S,f_S)$. If $t$
is not a unit at some $v\notin S_k$, both $J_{\mathcal O}(f)$ and
$\Phi_A(\gamma,f)$ vanish. (d) is Theorem~\ref{thm:tf}(d). (e), (f) By
Theorem~\ref{thm:Jchi} at $T=0$ and Corollary~\ref{cor:global-m}(a),
\[
-\operatorname{tr}\bigl(M^{-1}M'(w,it)I(\mu_{it},f)\bigr)=-\frac{m'}m(\mu,
it)\operatorname{tr}I(\mu_{it},f)+\varphi_A(f,\mu_{it})
\]
when $\mu$ is unramified outside $S$, and both sides vanish otherwise. In
the unramified case the two members of $\chi$ contribute equally
to~\eqref{eq:Iochi}, by Proposition~\ref{prop:phiA}(b) and $t\mapsto-t$,
so $\varphi_A$ cancels; in the ramified case $J_\chi$ has the factor
$\frac12$ and $\chi$ has one member. Invariance: (b) and (d) are an
orbital integral and a character; the terms of~(a) and~(c) are invariant
by Theorem~\ref{thm:IA}(b), and $f_S(z)$ and $J_G(zu_\xi,\cdot)$ are
orbital integrals. In (e) and (f), $\operatorname{tr}I(\mu_{it},\cdot)$
and $\operatorname{tr}(R(\cdot)|_{V_b})$ are characters, and
$\operatorname{tr}(M(w,0)I(\mu_0,\cdot))$ is invariant because $M(w,0)$
commutes with $I(\mu_0)$. The absolute convergence is that of
Theorem~\ref{thm:tf} and Lemma~\ref{lem:poisson-phi}.
\end{proof}

Theorem~\ref{thm:itf} is the invariant trace formula of~\cite{LiIV} for
this group. Its terms are explicit with two exceptions, which are
defined but not evaluated here. The unipotent coefficients
$a^G(S,u_\xi)$ in~(a) are those of Theorem~\ref{thm:fine}
(Remark~\ref{rem:sec4-status}). And the normalizing factors $r_v$ at the
places of $S_{\mathrm{bad}}(\mu)$, which always include the dyadic places,
are those of Theorem~\ref{thm:normalization}; the function $m'/m$
in~(e) and~(f) and the distributions $I_A$ in~(a) and~(c) depend on them
(Remark~\ref{rem:bad-places}). The elliptic terms~(b), the discrete
terms~(d), the residual term and the ramified term in~(f) do not involve
either.

It remains to evaluate $M(w,0)$ in~(f). For $a\in k^\times$ let $\eta_a$
be the quadratic character $t\mapsto\prod_v(-a,t_v)_v$ of
$k^\times\backslash\A_k^\times$.

\begin{lemma}\label{lem:whittaker}
Let $\mu=\chi_\psi\cdot(\chi\circ p)\in\Pi_{\mathrm{aut}}(X_A)$,
$\varphi=\bigotimes_v\varphi_v\in\mathcal H_B(\mu)$ and $a\in k^\times$.
For $\mathrm{Re}\,\lambda>1$ put
\[
W_a(x,\varphi,\lambda)=\int_{k\backslash\A_k}E\bigl(n(u)x,\varphi,\lambda
\bigr)\,\psi(-au)\,du .
\]
\begin{enumerate}[label=(\alph*)]
\item $E(x,\varphi,\lambda)=E_B(x,\varphi,\lambda)+\sum_{a\in k^\times}
  W_a(x,\varphi,\lambda)$, and
  $W_a(x,\varphi,\lambda)=\prod_vW_{a,v}(x_v,\varphi_v,\lambda)$ with
  $W_{a,v}(x_v,\varphi_v,\lambda)=\int_{k_v}\varphi_{v,\lambda}\bigl((w,1)^{-1}
  (n(u),1)x_v\bigr)\psi_v(-au)\,du$.
\item Each $W_{a,v}(x_v,\varphi_v,\lambda)$ is an entire function of
  $\lambda$.
\item If $v$ is as in~\ref{N4}, $|a|_v=1$, $\varphi_v=\varphi_v^0$ and
  $x_v\in K_v^*$, then
  \[
  W_{a,v}(x_v,\varphi^0_v,\lambda)=1+\eta_{a,v}(\varpi)\chi_v(\varpi)
  q^{-\lambda-\frac12}=\frac{L_v(\lambda+\frac12,\chi_v\eta_{a,v})}
  {L_v(2\lambda+1,\chi_v^2)} .
  \]
\item Hence, for a finite set $S_1$ of places containing those not covered
  by~(c),
  \[
  W_a(x,\varphi,\lambda)=\frac{L^{S_1}(\lambda+\frac12,\chi\eta_a)}
  {L^{S_1}(2\lambda+1,\chi^2)}\prod_{v\in S_1}W_{a,v}(x_v,\varphi_v,\lambda),
  \]
  which continues $W_a$ meromorphically to $\C$.
\end{enumerate}
\end{lemma}

\begin{proof}
(a) The first formula is the Fourier expansion on $k\backslash\A_k$. By
the Bruhat decomposition $G(\Q)=B(\Q)\sqcup B(\Q)w^{-1}N(\Q)$, for
$a\ne0$ the series~\eqref{eq:Eis} unfolds to $\int_{\A_k}\varphi_\lambda
(w^{-1}n(u)x)\psi(-au)\,du$, which is the product. (b) For $v$ finite and
$|u|$ large, $w^{-1}n(u)=n(-u^{-1})a(-u^{-1})\bar n(u^{-1})$ with $\bar
n(u^{-1})$ in a subgroup that fixes $I(x_v)\varphi_v$, so the integrand
is $|u|^{-\lambda-1}$ times a function of $u$ that is invariant under
multiplication by a subgroup $1+\varpi^c\OO$; its integral against
$\psi_v(-au)$ over a shell $|u|=q^m$ vanishes for $m$ large, and
$W_{a,v}$ is a polynomial in $q^{\pm\lambda}$. For $v$ archimedean the
integrand is smooth in $u$ with an asymptotic expansion in powers of
$u^{-1}$ times $|u|^{-\lambda-1}$, and repeated integration by parts
against $\psi_v(-au)$ continues the integral to all $\lambda$. (c)
Proceed as for Lemma~\ref{lem:intertwining}(b). On $|u|\le1$ the
integrand is $1$. On the shell $|u|=q^m$, $m\ge1$, write
$u=\varpi^{-m}u_0$; the decomposition $(w,1)^{-1}(n(u),1)=(n(-u^{-1}),1)
(a(-u^{-1}),1)(\bar n(u^{-1}),1)$ holds in the cover, because every
value of Kubota's cocycle that occurs is of the form $(1,\cdot)_v$ or
$(u,-u)_v$, hence $1$, and $\kappa_v=1$ on all factors. So the integrand is
$\mu_v(a(-u^{-1}),1)|u|^{-\lambda-1}=\chi_v(\varpi)^m|u|^{-\lambda-1}
\gamma_v(-\varpi^mu_0)$, and by Lemma~\ref{lem:weil}(a)--(b)
$\gamma_v(-\varpi^mu_0)$ is $1$ for $m$ even and $(-u_0,\varpi)_v
\gamma_v(\varpi)$ for $m$ odd. For $m$ even the integral of $\psi_v(-au)$
over the shell is $0$ if $m\ge2$. For $m$ odd the integral of
$(u_0,\varpi)_v\psi_v(-au)$ is $0$ if $m\ge3$, and for $m=1$ it is the
Gauss sum $\sum_{u_0\in(\OO/\varpi)^\times}(u_0,\varpi)_v\psi_v(-a\varpi^{-1}
u_0)=(-a,\varpi)_v\,q^{1/2}\gamma_v(\varpi)$, by~\eqref{eq:weil-index}. The
shell $m=1$ therefore contributes $(-1,\varpi)_v(-a,\varpi)_v\gamma_v
(\varpi)^2\chi_v(\varpi)q^{-\lambda-\frac12}$, and
$\gamma_v(\varpi)^2=(-1,\varpi)_v$. (d) follows.
\end{proof}

Part~(c) is the unramified Whittaker function of a genuine principal
series at the identity; compare~\cite{GPS80}. For $SL(2)$ the same
computation gives $1-\chi_v(\varpi)q^{-\lambda-1}=L_v(\lambda+1,\chi_v)
^{-1}$. The script checks both at $p=3,5,7,11$.

\begin{proposition}\label{prop:Mw0}
Let $b\in k^\times$ and $\mu=\chi_{\psi_b}$. Then $E(x,\varphi,0)=0$ for
all $\varphi\in\mathcal H_B(\mu)$, and
\[
M(w,0)=-1\quad\text{on }\mathcal H_B(\chi_{\psi_b}).
\]
Consequently the ramified term of Theorem~\ref{thm:Jchi}(c) is
$\frac14\operatorname{tr}(M(w,0)I(\mu_0,f))=-\frac14\operatorname{tr}
I(\chi_{\psi_b,0},f)$.
\end{proposition}

\begin{proof}
Write $\mu=\chi_\psi\cdot(\chi\circ p)$ with $\chi=\prod_v(b,\cdot)_v$
(Lemma~\ref{lem:weil}(d)), so $\chi^2=1$. It suffices to treat
factorizable $\varphi$. By
Theorem~\ref{thm:eisenstein}(d), $E(x,\varphi,\lambda)$ and
$M(w,\lambda)$ are holomorphic at $\lambda=0$.

\emph{Step 1: the Whittaker coefficients vanish at $\lambda=0$.} In
Lemma~\ref{lem:whittaker}(d) the denominator is
$L^{S_1}(2\lambda+1,\chi^2)=\zeta_k^{S_1}(2\lambda+1)$, because $\chi^2=1$.
The Dedekind zeta function has a pole at $s=1$, and removing finitely
many Euler factors does not remove it; so
$\zeta_k^{S_1}(2\lambda+1)^{-1}$ has a zero at $\lambda=0$. The numerator
$L^{S_1}(\lambda+\frac12,\chi\eta_a)$ is the $L$-function of a quadratic
Hecke character at the center of the critical strip, which is finite
(for $\chi\eta_a=1$ it is $\zeta_k^{S_1}(\frac12)$), and the local factors
at $S_1$ are entire. So $W_a(x,\varphi,0)=0$ for every $a\in k^\times$ and
every $x$. The zero is specific to the cover: the unramified Whittaker
function has the denominator $L_v(2\lambda+1,\chi_v^2)$, in which the
character is squared, so every ramified class sees the pole of
$\zeta_k$.

\emph{Step 2: the Eisenstein series is its constant term.} By
Lemma~\ref{lem:whittaker}(a) and Theorem~\ref{thm:eisenstein}(b),
\[
F(x):=E(x,\varphi,0)=E_B(x,\varphi,0)=\bigl(\varphi+M(w,0)\varphi\bigr)(x)
\,e^{H(x)}.
\]

\emph{Step 3: $|F|$ is constant.} $F$ is left invariant under $G(\Q)$.
Being equal to its constant term $\int_{N(\Q)\backslash N(\A)}F(nx)\,dn$,
it is left invariant under $N(\A)$. Since $w\in G(\Q)$, it is then left
invariant under $wN(\A)w^{-1}$ as well: $F(wnw^{-1}x)=F(nw^{-1}x)=
F(w^{-1}x)=F(x)$. Here $wN(\A)w^{-1}$ is the image of the lower triangular
unipotent group $\bar N(\A)$ under its canonical section
(Lemma~\ref{lem:local}(h)). Now $N(\A)$ and $\bar N(\A)$ generate
$G(\A)$: over a field or a local ring every element of $SL(2)$ is a
product of a bounded number of elementary matrices, the bound being
independent of the ring, so the factorization can be carried out at all
places simultaneously, with integral factors at almost all places. Hence
the subgroup of $\GA$ generated by the lifts of $N(\A)$ and $\bar N(\A)$
maps onto $G(\A)$, and $F$ is left invariant under it. Since $F$ is
genuine, $|F|$ is invariant under $\mu_2$ too. So $|F|$ is left invariant
under all of $\GA$, that is, constant.

\emph{Step 4: the constant is $0$.} For $\tilde a\in\overline{A(\A)}$ the
formula of Step~2 gives $|F(\tilde ax)|=e^{H(\tilde a)}|F(x)|$, because
$\varphi$ and $M(w,0)\varphi$ both lie in $\mathcal H_B(\mu)$ when
$\mu^w=\mu$. Choosing $H(\tilde a)\ne0$ we get $F=0$, and restricting to
$\overline{K_0}$ gives $\varphi+M(w,0)\varphi=0$.
\end{proof}

The statement is about representations, not only about $L$-functions:
$M(w,0)$ is a self-intertwining operator of the unitary representation
$I(\mu_0)$ with $M(w,0)^2=1$, and the proposition says that it is the
scalar $-1$ for every ramified class. The proof finds the sign from a
global fact: an automorphic function that is equal to its constant term
and homogeneous of degree one under $A$ must vanish.

For $SL(2)$ the same argument applies to the trivial character, where
$L^{S_1}(\lambda+1,1)^{-1}$ vanishes at $\lambda=0$: this is the classical
fact that the scattering matrix of $SL(2,\Z)$ is $-1$ at $s=\frac12$ and
$E(z,\frac12)=0$. For a quadratic character $\eta\ne1$ of $SL(2)$ it does
not apply, and $M(w,0)$ is then a product of nonscalar local
operators~\cite{LL}; that is the source of the multiplicity formula
in~\cite{LL}. On the double cover the denominator is always
$\zeta_k^{S_1}(2\lambda+1)$, and every ramified class behaves like the
trivial character of $SL(2)$. This agrees with the irreducibility of the
unitary principal series of $\overline{SL(2,k_v)}$~\cite{Mo2}. For
$k=\Q$ and $b\in\{\pm1,\pm2\}$ the script confirms $M(w,0)=-1$ by a
separate computation, from the residues at $\lambda=0$ of the real and
the $2$-adic local operators and the zero of order $|S|$ of
$m^S(\mu,\lambda)$, $S=\{\infty,2\}$.

\begin{corollary}\label{cor:itf-final}
In Theorem~\ref{thm:itf},
\begin{multline*}
\sum_{\chi\in\mathcal X}I_\chi(f)=\sum_{\pi}m_{\mathrm{disc}}(\pi)
\operatorname{tr}\pi(f)-\frac14\sum_{b\in k^\times/k^{\times2}}
\operatorname{tr}I(\chi_{\psi_b,0},f)\\
-\frac12\sum_{\mu\in\Pi_{\mathrm{aut}}(X_A)}\frac1{2\pi}\int_\R\frac{m'}m
(\mu,it)\operatorname{tr}I(\mu_{it},f)\,dt,
\end{multline*}
where $\pi$ runs over the irreducible constituents of
$L^2_{\mathrm{disc}}=L^2_{\mathrm{cusp}}\oplus\bigoplus_bV_b$ with their
multiplicities $m_{\mathrm{disc}}(\pi)$. Every term is a character of
$\GS$ evaluated at $f_S$.
\end{corollary}

\begin{proof}
Theorem~\ref{thm:itf}(d)--(f), Proposition~\ref{prop:Mw0}, and
Corollary~\ref{cor:global-m}(b) to write~(e) as half the sum over the two
members of $\chi$.
\end{proof}

In Arthur's notation~\cite[\S4]{A10} this is $\sum_t\int_{\Pi(G,t)}a^G(\pi)
I_G(\pi,f)\,d\pi$: the coefficient $a^G(\pi)$ is $m_{\mathrm{disc}}(\pi)$ on
the discrete spectrum, $-\frac14$ at $I(\chi_{\psi_b,0})$, and
$-\frac1{4\pi}\frac{m'}m(\mu,it)\,dt$ on the unitary principal series; the
first two make up the discrete part $I_{\mathrm{disc}}$. Since
$V_b\cong r^+_{\psi_{-b}}$ (Theorem~\ref{thm:residual}(c)),
$\operatorname{tr}(R(f)|_{V_b})=\prod_v\operatorname{tr}r^+_{\psi_{-b},v}
(f_v)$, the character of the even Weil representation.

\begin{example}\label{ex:itf-Q}
Let $k=\Q$ and $S=\{\infty,2\}$, so $f=f_\infty\otimes f_2\otimes e^S$.
By Example~\ref{ex:ramified-Q} the ramified classes that contribute are
$\chi_{\psi_b}$, $b\in\{\pm1,\pm2\}$, and the discrete correction is
$-\frac14\sum_{b=\pm1,\pm2}\operatorname{tr}I(\chi_{\psi_b,0},f)$. For
$\mu=\chi_{\psi_b}$ we have $S_{\mathrm{bad}}(\mu)=\{2\}$ and, with
$\xi(s)=\pi^{-s/2}\Gamma(\frac s2)\zeta(s)$,
\[
m(\mu,\lambda)=c\,\frac{\xi(2\lambda)}{\xi(2\lambda+1)}\cdot
\frac{1-2^{-2\lambda}}{1-2^{-1-2\lambda}}\cdot r_2(\mu_2,\lambda),\qquad
\frac{\xi(2\lambda)}{\xi(2\lambda+1)}\Big|_{\lambda=0}=-1,
\]
where $r_2$ is the $2$-adic factor of Theorem~\ref{thm:normalization}
(Remark~\ref{rem:bad-places}) and $c=\sqrt2\,e^{\pm i\pi/4}$ comes from
Lemma~\ref{lem:explicit-r}(b), the sign being that of $b$. The
hyperbolic classes that contribute are those of $a(t)$, $t=\pm2^j$,
$j\ge1$, and their terms are $c_\Q\,I_A(a(t)_S,f_S)$. The archimedean part of the correction $\Phi_A$ is
given by Example~\ref{ex:Phi-real}.
\end{example}

\subsection{Status of the results}\label{sec:status6}

\begin{remarkrev}\label{rem:sec6-status}
\emph{What rests on citation.} Proved here, given
Sections~\ref{sec:rank0}--\ref{sec:spectral} and
Theorem~\ref{thm:normalization}: Lemmas~\ref{lem:uprime}--%
\ref{lem:conj-identity}; Lemma~\ref{lem:explicit-r}(b), and~(a) except
for rationality; Corollary~\ref{cor:global-m}; Proposition~\ref{prop:phiA};
Theorem~\ref{thm:IA}; Lemma~\ref{lem:poisson-phi}; Theorem~\ref{thm:itf};
Lemma~\ref{lem:whittaker}; Proposition~\ref{prop:Mw0}; and
Corollary~\ref{cor:itf-final}. Taken from the literature:
Theorem~\ref{thm:normalization}~\cite{A11,LiII}, of which only the bad
places are not made explicit (Remark~\ref{rem:bad-places}); the generic
irreducibility of $I(\pi_{v,\lambda})$ and the adjoint formula used in
Lemma~\ref{lem:explicit-r}(a); the estimates behind
Proposition~\ref{prop:phiA}(a) for functions that are not $\KS$-finite;
and, in the proof of Lemma~\ref{lem:variance-global}, the interchange of
integrals on the spectral side. We have not compared the
definition~\eqref{eq:IAgamma} with the definitions of~\cite{LiIV} beyond
the remark after Theorem~\ref{thm:IA}. As in
Theorem~\ref{thm:tf}, the set $S$ must be large for the given $f_S$, and
the coefficients $a^G(S,u_\xi)$ are not computed
(Remark~\ref{rem:sec4-status}).
\end{remarkrev}

\begin{remarkrev}\label{rem:verified-sec6}
The script \texttt{checks/sec6\_n2\_check.py} verifies, for $k=\Q$: the
sign of the cross terms in the Maass--Selberg relation~\eqref{eq:MS},
hence of the ramified term in Theorem~\ref{thm:Jchi}(c), on the
Eisenstein series of $SL(2,\Z)$ (by numerical integration over the
modular surface); Lemma~\ref{lem:uprime}(a) and
Example~\ref{ex:uprime}; Lemma~\ref{lem:variance-local}(a)--(b) at
$v=\R$, with explicit Iwasawa decompositions and by quadrature for a
Gaussian $f$; the identity~\eqref{eq:conj-identity} for the principal
series of $SL(2,\R)$ in the compact picture;
formulas~\eqref{eq:real-M}--\eqref{eq:rho-m} and the Fourier transform in
Example~\ref{ex:Phi-real}; Lemma~\ref{lem:whittaker}(c) at $p=3,5,7,11$,
shell by shell in the cover, and its analog for $SL(2)$;
Proposition~\ref{prop:Mw0} for $b=\pm1,\pm2$, from local residues; and the
value $-1$ in Example~\ref{ex:itf-Q}.
\end{remarkrev}

% SPDX-License-Identifier: CC-BY-4.0
% =====================================================================
% Section 7 of D. Joyner and Claude, "A hitch-hiker's guide to the
% metaplectic trace formula for SL(2)", revised for the 2-fold cover
% (n = 2), November 2026.
%
% Stage 6 of ../revision_log.md. Sources: ../original-1995/cctrace7.tex
% (= ../book-1997/book7h.tex) and ../book-1997/book3b.tex. Referee items
% II.18-II.23, C.1, C.3, C.4 and the Section 7 part of I.1 and II.22 are
% treated here; see the "Stage 6 notes" in ../revision_log.md, which also
% map the 1995 numbers to the numbers here.
%
% This file is \input by the master; it compiles on its own through
% sl2bar-trace-2026.tex. Explicit formulas are checked by checks/sec7_n2_check.py.
% Comments beginning "% [Claude]" are editorial notes by Claude.
% =====================================================================

\providecommand{\GA}{\overline{G(\A)}}
\providecommand{\GS}{\overline{G(\Q_S)}}
\providecommand{\AS}{\overline{A(\Q_S)}}
\providecommand{\KS}{\overline{K_S}}
\providecommand{\Gb}{\overline{SL(2,F)}}

\section{The trace formula in stable(r) form}\label{sec:stable}

The invariant trace formula of Theorem~\ref{thm:itf} is an identity
between distributions that are invariant under conjugation by the cover
itself. For $SL(2)$, Labesse and Langlands~\cite{LL} went one step
further: they regrouped the terms according to conjugacy under $GL(2)$,
which for $SL(2)$ is stable conjugacy, and identified the part of the
formula that is not stably invariant with contributions of the elliptic
tori. This section carries out the first half of that program for the
double cover: it determines how $GL(2)$ acts on the cover, what stable
conjugacy is, and how each term of the trace formula behaves. We call the
result the \emph{stable(r) form}: the terms are organized by stable
classes, and the elliptic regular terms are written as stable orbital
integrals, but nothing is transferred to another group.

Three properties of a distribution on the cover occur below and should be
kept apart. \emph{Equivariance}: it is multiplied by a character of $\det
h$ when $f$ is replaced by $f^h=f\circ\mathrm{Ad}(h)$, $h\in GL(2)$; this
is about ordinary conjugation, and it is what
Section~\ref{sec:hyperbolic-gl2} proves for the weighted terms.
\emph{Stability}: it is invariant under the twisted action $f\mapsto
f^{[h]}$ of Section~\ref{sec:stable-conj}; this is proved for the regular
elliptic terms, the characters of the principal series and the theta
functions. \emph{Endoscopic transfer}: it is identified with a stable
distribution on $SO(3)$ applied to a transfer of $f$; this is Li's
theorem, and nothing of the kind is proved here.

For the double cover of $Sp(2n)$ the full theory exists and is due to
Li: stable conjugacy and the transfer of orbital
integrals~\cite{Li11}, following Adams~\cite{Ad}, Renard~\cite{Re} and,
for the group here over a $p$-adic field, Schultz~\cite{Sch} (see
also~\cite{How}); the stabilization of the elliptic
terms~\cite{Li15}; and the stabilization of the whole invariant trace
formula~\cite{Li21}. The elliptic endoscopic groups of the double cover of
$Sp(2n)$ are the groups $SO(2n'+1)\times SO(2n''+1)$, $n'+n''=n$. For
$n=1$ these are two copies of $SO(3)=PGL(2)$, attached to the pairs
$(1,0)$ and $(0,1)$; the tori of $SL(2)$, which are the endoscopic groups
of $SL(2)$ in~\cite{LL}, do not occur. We do not reprove any of this.
What can be done by hand for the group of this paper, and is done below,
is the following.

\paragraph{Results.}
\begin{itemize}
\item $GL(2,F)$ acts on $\Gb$ by automorphisms $\mathrm{Ad}(h)$, with an
  explicit formula (Lemma~\ref{lem:Ad}). An element $h$ that commutes
  with a semisimple $\gamma$ multiplies the lifts of $\gamma$ by the sign
  $(\det h,\det(1+\gamma))_v$ (Lemma~\ref{lem:eta}). So conjugation by
  $GL(2)$ can carry a lift of $\gamma$ to the other lift, and it is not
  stable conjugacy.
\item Stable conjugacy on the cover is $GL(2)$-conjugacy \emph{twisted}
  by that sign: $h\star x=(\det h,\det(1+x))_v\,\mathrm{Ad}(h)x$. Its
  classes are the level sets, over the stable classes of $SL(2,F)$, of an
  explicit genuine function $\Delta_\psi$ with values in $\mu_4$
  (Lemma~\ref{lem:Delta}). This function is $|\det(1+x)|^{1/2}$ times the
  character of the difference $r^+_\psi-r^-_\psi$ of the two halves of
  the Weil representation (Proposition~\ref{prop:weil-char}), so the
  definition agrees with that of Adams and Li.
\item For a regular elliptic $\gamma\in G(\Q)$ with splitting field $E$,
  the terms of the trace formula attached to the $GL(2,k)$-conjugates of
  $\gamma$ add up to $\frac12$ of a stable orbital integral plus $\frac12$
  of a $\kappa$-orbital integral, as for $SL(2)$. For the cover the
  $\kappa$-orbital integral at $\gamma$ is again a stable orbital
  integral: that of the translate $f_-$ of $f$ by the central element
  $-1$, at $-\gamma$ (Lemma~\ref{lem:kappa-orbital}). Hence the elliptic
  regular part of the trace formula is $S_{\mathrm{ell}}(f)+
  S_{\mathrm{ell}}(f_-)$ with one stable distribution $S_{\mathrm{ell}}$
  (Theorem~\ref{thm:elliptic-stable}). This is the two-term pattern, with
  coefficients $\frac12$ and $\frac12$, of the elliptic part of Li's
  stable trace formula for $n=1$ with its two endoscopic data; the terms
  are not matched with the transferred distributions on $SO(3)$.
\item The remaining geometric terms transform under $GL(2,k_S)$ by
  explicit characters (Proposition~\ref{prop:IA-gl2}), and the
  unipotent coefficients $a^G(S,u_\xi)$, left open in
  Section~\ref{sec:geometric}, are computed:
  $a^G(S,u_\xi)=\frac12\bigl(c_S+\sum_{\kappa\ne1}\kappa_S(\xi)L^S(1,\kappa)
  \bigr)$, the sum over the quadratic Hecke characters unramified outside
  $S$ (Proposition~\ref{prop:unipotent}).
\item On the spectral side, $GL(2,k)$ permutes the pieces of the
  spectrum, and $V_b\mapsto V_{ab}$. The theta functions attached to
  $\psi_{-b}$ span a subspace of the discrete spectrum whose trace is
  $\frac12\Theta(f)+\frac12\Theta(f_-)$ with $\Theta$ a stable character
  (Proposition~\ref{prop:theta-packet}): the simplest instance of the
  same pattern on the spectral side.
\end{itemize}
The hyperbolic terms, the unipotent terms and the cuspidal spectrum are
\emph{not} stabilized here; Section~\ref{sec:stable-form} says what is
missing, and Remark~\ref{rem:sec7-status} what rests on citation.

\paragraph{Conventions.} $F=k_v$ is a local field of characteristic $0$,
as in Section~\ref{sec:local}; we write $G_1=GL(2)$, $Z$ for its center,
and $\delta_a=\diag(a,1)$. For $h\in G_1(F)$, $\mathrm{Ad}(h)$ is the
automorphism of $\Gb$ of Lemma~\ref{lem:gl2}(c), and for a function $f$
on the cover
\[
f^h(x)=f(\mathrm{Ad}(h)x),
\]
which for $h\in SL(2,F)$ is the function $f^y$ of
Section~\ref{sec:invariant}, $y$ a lift of $h$. For a central element
$z$ of a cover, $f_z(x)=f(zx)$, as in Section~\ref{sec:coarse}. The
element $-1\in G(\Q)$ is identified with its image $s(-1)=(-1,1)$ in
$\GA$, and
\[
f_-(x)=f_{-1}(x)=f\bigl(s(-1)\,x\bigr),\qquad f\in\mathcal H(\GA).
\]
Since $s(-1)^2=1$, $(f_-)_-=f$. For a quadratic extension $E=F(\sqrt\theta)$
let $\kappa_E=(\cdot,\theta)_v$ be the quadratic character of $F^\times$
with kernel $N E^\times=N_{E/F}(E^\times)$; for $E=F\times F$ put
$N E^\times=F^\times$ and $\kappa_E=1$.

\subsection{Conjugation by \texorpdfstring{$GL(2)$}{GL(2)}}\label{sec:Ad}

Conjugation by $h\in G_1(F)$ is an automorphism of $SL(2,F)$, and it
lifts to the cover in exactly one way (for the metaplectic cover of
$Sp(V)$ and the similitude group in general, see~\cite[\S3]{Th13}). This subsection makes the lift
explicit and computes it on the subgroups that matter: $N$, $A$, the
center, the maximal compact subgroup, the tori, and globally on $G(\Q)$.

\begin{lemma}\label{lem:Ad}
Let $h\in G_1(F)$.
\begin{enumerate}[label=(\alph*)]
\item $\mathrm{Ad}(h)$ is the only automorphism of the group $\Gb$ that
  induces $g\mapsto hgh^{-1}$ on $SL(2,F)$. It is a homeomorphism for
  the topology of Lemma~\ref{lem:kubota}(c), and it fixes $\mu_2$
  pointwise.
\item $\mathrm{Ad}(h)$ depends only on the image of $h$ in $PGL(2,F)$.
  For $h\in SL(2,F)$ it is conjugation by either lift of $h$. So
  $\mathrm{Ad}(h)$ is inner when $\det h\in F^{\times2}$, and the group
  $F^\times/F^{\times2}$ acts through the $\mathrm{Ad}(\delta_a)$ modulo
  inner automorphisms.
\item For $a\in F^\times$ and $g=\begin{pmatrix}\alpha&b\\c&d
  \end{pmatrix}\in SL(2,F)$,
  \begin{equation}\label{eq:Ad-delta}
  \mathrm{Ad}(\delta_a)(g,\zeta)=\bigl(\delta_ag\delta_a^{-1},\ v_a(g)\,
  \zeta\bigr),\qquad v_a(g)=\begin{cases}1,&c\ne0,\\(a,d)_v,&c=0.
  \end{cases}
  \end{equation}
\item In particular
  \begin{align*}
  \mathrm{Ad}(\delta_a)(n(u),1)&=(n(au),1),&
  \mathrm{Ad}(\delta_a)(a(t),\zeta)&=(a(t),(a,t)_v\zeta),\\
  \mathrm{Ad}(\delta_a)(-1,\zeta)&=(-1,(a,-1)_v\zeta),&
  \mathrm{Ad}(\delta_a)(w,1)&=(a(a),1)(w,1).
  \end{align*}
\item Let $v$ be finite and non-dyadic. If $h\in Z(F)GL(2,\OO)$, then
  $\mathrm{Ad}(h)K^*=K^*$ and $e_v^h=e_v$ for the unit element $e_v$
  of~\eqref{eq:unit}.
\item (Global.) For $h=(h_v)\in G_1(\A_k)$ the automorphisms
  $\mathrm{Ad}(h_v)$ define an automorphism $\mathrm{Ad}(h)$ of $\GA$,
  and $f^h\in\mathcal H(\GA)$ for $f\in\mathcal H(\GA)$. If
  $h\in G_1(k)$, then $\mathrm{Ad}(h)$ preserves the image of $G(\Q)$:
  $\mathrm{Ad}(h)s(\gamma)=s(h\gamma h^{-1})$ for $\gamma\in G(\Q)$, where
  $s(\gamma)=(\gamma,\kappa_\A(\gamma))$ is the splitting of
  Lemma~\ref{lem:global}(a).
\end{enumerate}
\end{lemma}

\begin{proof}
(a) Two automorphisms that induce the same automorphism of $SL(2,F)$
differ by a homomorphism $SL(2,F)\to\mu_2$, which is trivial because
$SL(2,F)$ is perfect. $\mathrm{Ad}(h)$ fixes $\mu_2$ because $\mu_2$ is
central in the cover of $GL(2,F)$. For continuity it suffices to look
near the identity. Let $U$ be as in Lemma~\ref{lem:kubota}(b), and let
$V\subset U$ be a neighborhood of $1$ with $hVh^{-1}\subset U$. For
$g\in V$ write $\mathrm{Ad}(h)(g,1)=(hgh^{-1},\varepsilon(g))$. Since
$\beta\equiv1$ on $U\times U$, $\varepsilon(g_1g_2)=\varepsilon(g_1)
\varepsilon(g_2)$ whenever $g_1,g_2,g_1g_2\in V$; and every $g$ in a
smaller neighborhood is a square of an element of $V$, as in the proof of
Lemma~\ref{lem:kubota}(c). So $\varepsilon=1$ near $1$, and
$\mathrm{Ad}(h)$ maps the basic neighborhoods $\{(g,1)\mid g\in V'\}$ to
basic neighborhoods.

(b) A scalar matrix induces the identity on $SL(2,F)$, so by~(a) it acts
trivially on the cover; and conjugation by a lift of $h\in SL(2,F)$ is an
automorphism over $g\mapsto hgh^{-1}$. If $\det h=c^2$, then $c^{-1}h\in
SL(2,F)$.

(c) By definition $\mathrm{Ad}(\delta_a)(g,\zeta)=(\delta_a,1)(g,\zeta)
(\delta_a,1)^{-1}$ in $\overline{GL(2,F)}$, which is
\[
\bigl(\delta_ag\delta_a^{-1},\ \zeta\,\sigma(\delta_a,g)\,\sigma(\delta_ag,
\delta_a^{-1})\,\sigma(\delta_a,\delta_a^{-1})\bigr)
\]
with the cocycle~\eqref{eq:sigma-gl2}. Here $x(\delta_a)=x(\delta_a^{-1})=1$ and
$x(\delta_ag)=x(g)$, so $\sigma(\delta_a,g)=(x(g),a)_v$ and
$\sigma(\delta_a,\delta_a^{-1})=1$. If $c\ne0$, then $x(\delta_ag
\delta_a^{-1})=c/a$ and $\sigma(\delta_ag,\delta_a^{-1})=(a^{-1},c/a^2)_v
=(a,c)_v$, so the product is $1$. If $c=0$, then
$x(\delta_ag\delta_a^{-1})=d=x(g)$ and $\sigma(\delta_ag,\delta_a^{-1})=
(1,d/a)_v=1$, so the product is $(d,a)_v$. (The script also checks
directly that~\eqref{eq:Ad-delta} is a homomorphism for $\beta$, which
by~(a) proves~(c) without reference to the cover of $GL(2)$.)

(d) The first three are~\eqref{eq:Ad-delta} with $c=0$ and $d=1$,
$t^{-1}$, $-1$. For the last, $\delta_aw\delta_a^{-1}=a(a)w$ has lower
left entry $a^{-1}\ne0$, so $\mathrm{Ad}(\delta_a)(w,1)=(a(a)w,1)$, and
$\beta(a(a),w)=(1,a^{-1})_v=1$.

(e) By~(b) we may take $h\in GL(2,\OO)$, so $hKh^{-1}=K$. Then
$k\mapsto\mathrm{Ad}(h)^{-1}(hkh^{-1},\kappa(hkh^{-1}))$ is a
homomorphic section of $\overline K\to K$. It differs from
$k\mapsto(k,\kappa(k))$ by a homomorphism $K=SL(2,\OO)\to\mu_2$, and there
is none except the trivial one: $K$ is generated by elementary matrices,
which are squares.

(f) By~(e), $\mathrm{Ad}(h_v)$ preserves $K_v^*$ for almost all $v$, so
the product acts on the restricted product of the local covers; it fixes
every $\mu_2$, hence the subgroup $Z_0$ of Section~\ref{sec:global}, and
descends to $\GA$. For $f=\bigotimes f_v$ we get
$f^h=\bigotimes f_v^{h_v}$ with $f_v^{h_v}=e_v$ for almost all $v$. If
$h\in G_1(k)$, then $\gamma\mapsto\mathrm{Ad}(h)^{-1}s(h\gamma h^{-1})$ is a
homomorphic section of $\overline{G(\Q)}\to G(\Q)$; it equals $s$ because
$SL(2,k)$ is perfect.
\end{proof}

Part~(d) already shows that $\mathrm{Ad}(h)$ is not inner in general: it
acts on the abelian group $\overline{A(F)}$ by the character
$a(t)\mapsto(a,t)_v$, and it exchanges the two lifts of $-1$ when
$(\det h,-1)_v=-1$. The next lemma gives the general statement. For a
semisimple $\gamma\in SL(2,F)$, $\gamma\ne-1$, put
\begin{equation}\label{eq:eta}
D^+(\gamma)=\det(1+\gamma)=2+\operatorname{tr}\gamma\ \in F^\times,\qquad
\eta_\gamma(a)=\bigl(a,D^+(\gamma)\bigr)_v\quad(a\in F^\times).
\end{equation}
The definition of $D^+$ and $\eta_\gamma$ makes sense for every $\gamma$
with $\operatorname{tr}\gamma\ne-2$, and depends only on the trace.

\begin{lemma}\label{lem:eta}
Let $\gamma\in SL(2,F)$ be semisimple, $\gamma\ne-1$, and let $h\in
G_1(F)$ commute with $\gamma$. Then
\[
\mathrm{Ad}(h)(\gamma,\zeta)=\bigl(\gamma,\ \eta_\gamma(\det h)\,\zeta
\bigr).
\]
For $\gamma=-1$ and every $h\in G_1(F)$,
$\mathrm{Ad}(h)(-1,\zeta)=(-1,(\det h,-1)_v\zeta)$.
\end{lemma}

\begin{proof}
The second statement is Lemma~\ref{lem:Ad}(b),(d). For $\gamma=1$ there is
nothing to prove ($D^+=4$). Let $\gamma$ be regular. Its centralizer in
$M_2(F)$ is the commutative algebra $E=F[\gamma]$, and $h\in E^\times$.
Put $z=1+\gamma\in E$. Then $\det z=D^+(\gamma)\ne0$, and $z\gamma^{-1}=
1+\gamma^{-1}$ is the conjugate $\bar z$ of $z$ in $E$, so
$z\bar z=\det z$ and
\[
\gamma\cdot\bigl(D^+(\gamma)\cdot1\bigr)=z^2\qquad\text{in }E^\times
\subset G_1(F).
\]
For commuting $x,y\in G_1(F)$ let $[x,y]\in\mu_2$ be the commutator of
any two lifts to $\overline{GL(2,F)}$. On the commutative group
$E^\times$ it is bimultiplicative, so $[h,z^2]=[h,z]^2=1$ and therefore
$[h,\gamma]=[h,D^+(\gamma)\cdot1]$. For a scalar $\lambda$ and any $h$,
\eqref{eq:sigma-gl2} gives $\sigma(h,\lambda)=(\lambda,x(h)/\det h)_v$ and
$\sigma(\lambda,h)=(x(h),\lambda)_v$, so $[h,\lambda]=(\lambda,\det h)_v$.
Hence $\mathrm{Ad}(h)(\gamma,\zeta)=(\gamma,[h,\gamma]\zeta)$ with
$[h,\gamma]=(D^+(\gamma),\det h)_v$.
\end{proof}

By~\eqref{eq:gl2-commutator} the sign is also $(x(h\gamma)/x(h),\det
h)_v$; the lemma says that this equals $(\det h,2+\operatorname{tr}
\gamma)_v$, which the script confirms on random elements of the tori
$T_\theta(\Q)$. If $\gamma$ is regular with centralizer $E^\times$ in
$G_1(F)$, then $\det h$ runs through $N E^\times$, and $D^+(\gamma)=
N_{E/F}(1+\gamma)$ is itself a norm. So the lifts of $\gamma$ are
exchanged by the centralizer of $\gamma$ in $G_1(F)$ exactly when
$\eta_\gamma$ is not trivial on $N E^\times$.

\begin{example}\label{ex:eta}
(a) \emph{Split torus.} For $\gamma=a(t)$, $t\ne\pm1$:
$D^+=(1+t)^2/t$, so $\eta_\gamma=(\cdot,t)_v$ and $N E^\times=F^\times$.
The two lifts of $a(t)$ are conjugate under $G_1(F)$ if and only if
$t\notin F^{\times2}$.

(b) \emph{$F=\R$, $\gamma$ elliptic.} $N E^\times=\R_{>0}$ and
$D^+=2+2\cos\theta>0$, so $\eta_\gamma=1$. In $Mp(2,\R)$ the automorphism
$\mathrm{Ad}(\delta_{-1})$ sends the element $\tilde k_\theta$ of
Lemma~\ref{lem:explicit-r}(b) to $\tilde k_{-\theta}$; on the center
$\{\tilde k_{j\pi}\}\cong\Z/4$ it is inversion, in agreement with
$(-1,-1)_\R=-1$.

(c) \emph{$F$ non-dyadic, $E/F$ unramified, $\gamma\in K$.} The norms are
the elements of even valuation, and $\eta_\gamma$ is trivial on them if
$D^+(\gamma)$ is a unit. But $\eta_\gamma(\varpi)=(\varpi,D^+(\gamma))_v$
is $+1$ or $-1$ according as the unit $D^+(\gamma)$ is a square or not;
both occur.

(d) \emph{$F=\Q_3$, $E=\Q_3(\sqrt3)$.} Let $h=\begin{pmatrix}0&3\\1&0
\end{pmatrix}$, which is $\sqrt3\in E^\times$, of determinant $-3$. Then
$\mathrm{Ad}(h)(-1,\zeta)=(-1,(-3,-1)_3\zeta)=(-1,-\zeta)$. By continuity
the elements of $\overline{T_3}$ close to a lift of $-1$ are conjugate
under $G_1(\Q_3)$, indeed under the centralizer $E^\times$, to their
negatives; here $D^+(\gamma)\in-3\,\Q_3^{\times2}$ and
$(-3,-3)_3=-1$.

(e) \emph{A rational example at $2$ and $\infty$.} Let
\[
\gamma=\begin{pmatrix}3/5&-4/5\\4/5&3/5\end{pmatrix},\qquad
h=\begin{pmatrix}1&-1\\1&1\end{pmatrix}.
\]
Both lie in $\Q[w]$, so they commute; $\det h=2$ and $D^+(\gamma)=16/5$.
Over $\Q_2$, $(2,16/5)_2=(2,5)_2=-1$: ordinary conjugation by $h$ exchanges
the two lifts of $\gamma$, and $h\star x=x$ for both. Over $\R$ the sign is
$(2,16/5)_\infty=+1$, and over $\Q_5$ it is $(2,5)_5=-1$, in agreement with
the product formula. (This example is due to the referee.)
\end{example}

\subsection{Stable conjugacy}\label{sec:stable-conj}

Two elements of $SL(2,F)$ are \emph{stably conjugate} if they are
conjugate in $SL(2,\bar F)$; for semisimple elements this is the same as
conjugacy under $G_1(F)$~\cite{LL}. On the cover, conjugacy under
$G_1(F)$ is too coarse, by Example~\ref{ex:eta}: it can identify the two
lifts of an element, on which every genuine function takes opposite
values. The remedy is to twist the action by the sign of
Lemma~\ref{lem:eta}.

Let $\Gb^\flat$ be the set of $x\in\Gb$ whose image $p(x)$ has trace
different from $-2$; it is open, dense, invariant under every
$\mathrm{Ad}(h)$, and contains the regular semisimple elements. For
$x\in\Gb^\flat$ write $\eta_x=\eta_{p(x)}$, and for $h\in G_1(F)$ put
\begin{equation}\label{eq:star}
h\star x=\eta_x(\det h)\cdot\mathrm{Ad}(h)x,\qquad
\eta_x(\det h)=\bigl(\det h,\ \det(1+p(x))\bigr)_v\in\mu_2 .
\end{equation}
Since $\eta_x$ depends only on the trace of $p(x)$, which
$\mathrm{Ad}(h)$ preserves, $(h_1h_2)\star x=h_1\star(h_2\star x)$: this
is an action of $G_1(F)$ on $\Gb^\flat$. Scalars act trivially, and
$h\in SL(2,F)$ acts by conjugation. Two elements of $\Gb^\flat$ are
called \emph{stably conjugate} if they lie in the same orbit of this
action. For a function $f$ on $\Gb$ put
\[
f^{[h]}(x)=f(h\star x)=\eta_x(\det h)\,f^h(x)\qquad(x\in\Gb^\flat),
\]
the last equality for anti-genuine $f$. The complement of $\Gb^\flat$ has
measure $0$, so $f^{[h]}$ is defined almost everywhere; it is smooth on
$\Gb^\flat$ but in general not across $\operatorname{tr}=-2$. A locally
integrable class function $\Theta$ on $\Gb$, and the invariant
distribution $f\mapsto\int f\,\Theta$ it defines, are called
\emph{stable} if $\Theta(h\star x)=\Theta(x)$ for all $h\in G_1(F)$ and
almost all $x$. Since $x\mapsto h\star x$ preserves the Haar measure,
this means $\int f^{[h]}\Theta=\int f\,\Theta$.

By Lemma~\ref{lem:eta}, a regular semisimple $x$ over $\gamma$ is fixed
by $h\star$ for every $h$ in the centralizer $E^\times$ of $\gamma$. So
the stable class of $x$ maps bijectively onto the $G_1(F)$-conjugacy
class of $\gamma$; it never contains the other lift of $\gamma$; and it
is the union of the conjugacy classes of the $\delta_a\star x$, $a\in
F^\times/N E^\times$: two classes if $E$ is a field, and one if $\gamma$
is hyperbolic.

That this is the right notion is shown by the following function, which
separates the two lifts and is constant on stable classes. Let
$\gamma_v(a)$ be the Weil index of~\eqref{eq:weil-index}, for the fixed
character $\psi=\psi_v$.

\begin{lemma}\label{lem:Delta}
Let $\Gb_{\mathrm{reg}}$ be the set of regular semisimple elements of
$\Gb$.
\begin{enumerate}[label=(\alph*)]
\item There is exactly one function $\Delta_\psi\colon
  \Gb_{\mathrm{reg}}\to\mu_4$ that is invariant under conjugation by
  $\Gb$ and satisfies
  \begin{equation}\label{eq:Delta}
  \Delta_\psi(g,\zeta)=\zeta\,\bigl(c,D^+(g)\bigr)_v\,
  \frac{\gamma_v(D^+(g))}{\gamma_v(1)}\qquad\text{for }
  g=\begin{pmatrix}\alpha&b\\c&d\end{pmatrix},\ c\ne0.
  \end{equation}
  It is genuine: $\Delta_\psi(x\zeta)=\zeta\,\Delta_\psi(x)$.
\item $\Delta_\psi(\mathrm{Ad}(h)x)=\eta_x(\det h)\,\Delta_\psi(x)$ for
  $h\in G_1(F)$; that is, $\Delta_\psi(h\star x)=\Delta_\psi(x)$. So two
  elements of $\Gb_{\mathrm{reg}}$ are stably conjugate if and only if
  their images in $SL(2,F)$ are conjugate under $G_1(F)$ and
  $\Delta_\psi$ takes the same value on them.
\item On the split torus, $\Delta_\psi(a(t),\zeta)=\zeta\,\gamma_v(1)/
  \gamma_v(t)=\chi_{\psi_{-1},v}(a(t),\zeta)$, the complex conjugate of
  Weil's character $\chi_{\psi,v}$ (Lemma~\ref{lem:weil}).
\item For $a\in F^\times$, $\Delta_{\psi_a}=\eta_x(a)\,\Delta_\psi$, where
  $\psi_a(u)=\psi(au)$.
\end{enumerate}
\end{lemma}

\begin{proof}
Write $\Delta$ for the right side of~\eqref{eq:Delta}, a function on the
set $X$ of $(g,\zeta)$ with $c\ne0$ and $\operatorname{tr}g\ne-2$; and
$D=D^+(g)$, which is invariant under conjugation. We show that
$\Delta(yxy^{-1})=\Delta(x)$ when $x$ and $yxy^{-1}$ lie in $X$.

\emph{Conjugation by $(n(u),1)$.} The lower left entry is unchanged, and
by~\eqref{eq:inverse} the sign is $\beta(n(u),g)\beta(n(u)g,n(-u))
\beta(n(u),n(-u))=(c,1)_v(1,c)_v=1$.

\emph{Conjugation by $(a(s),1)$.} The lower left entry becomes
$c/s^2$, and the sign is
\[
\beta(a(s),g)\,\beta(a(s)g,a(s)^{-1})\,\beta(a(s),a(s)^{-1})
=(c,s)_v\cdot(s,cs)_v\cdot(s,-1)_v=1 .
\]
So $\Delta$ is unchanged.

\emph{Conjugation by $(w,1)$, when $b\ne0$.} Here $wgw^{-1}=
\begin{pmatrix}d&-c\\-b&\alpha\end{pmatrix}$. Suppose $\alpha\ne0$.
Then $x(wg)=\alpha$, and the sign is $\beta(w,g)\beta(wg,w^{-1})
\beta(w,w^{-1})=(\alpha,\alpha/c)_v(-b/\alpha,b)_v=
(\alpha,-c)_v(\alpha,b)_v=(\alpha,-b c)_v$. So invariance amounts
to $(-b,D)_v(\alpha,-b c)_v=(c,D)_v$, that is, to
$(-b c,\ \alpha D)_v=1$. Put $X_0=-b c=1-\alpha d$ and
$Y_0=\alpha+1$. Then $\alpha D=\alpha^2+\alpha d+2\alpha=Y_0^2-X_0$, and
$(X_0,Y_0^2-X_0)_v=1$ because $Y_0^2-X_0$ is the norm of $Y_0+\sqrt{X_0}$.
If $\alpha=0$, then $b c=-1$, $x(wg)=b$, the sign is
$(b,b/c)_v(-1,b)_v=(b,-b^2)_v(-1,b)_v=1$, and the
identity to be checked is $(-b,D)_v=(c,D)_v$, i.e.\ $(-b c,D)_v=1$,
which holds because $-b c=1$.

Now let $y\in\Gb$ be arbitrary with $x,yxy^{-1}\in X$. If $p(y)\in B$,
then $y$ is a product of elements of the first two kinds and of $\mu_2$.
Otherwise $p(y)=b_1wn$ with $b_1\in B(F)$ and $n\in N(F)$ (Bruhat). The
element $x'=(n,1)x(n,1)^{-1}$ is in $X$, and its upper right entry is not
$0$: otherwise $(w,1)x'(w,1)^{-1}$, hence $yxy^{-1}$, would be upper
triangular. So the three steps apply in turn. Finally, every regular
semisimple element is conjugate to an element of $X$ (if $c=0\ne b$,
conjugate by $w$; if $c=b=0$, first by $n(1)$), which gives existence
and uniqueness of $\Delta_\psi$. The values lie in $\mu_4$, because $(\gamma_v(D)/\gamma_v(1))^2=(D,-1)_v$ by Lemma~\ref{lem:weil}(a), and
genuineness is clear.

(b) For $h\in SL(2,F)$ this is~(a). For $h=\delta_a$ and $x=(g,\zeta)\in
X$, \eqref{eq:Ad-delta} gives $\mathrm{Ad}(\delta_a)x=(\delta_ag
\delta_a^{-1},\zeta)$ with lower left entry $c/a$, so
$\Delta(\mathrm{Ad}(\delta_a)x)=(a,D)_v\Delta(x)$. Since $G_1(F)$ is
generated by $Z(F)$, $SL(2,F)$ and the $\delta_a$, the formula follows.
For the last statement, if $p(x')=hp(x)h^{-1}$, then $x'=\pm\,h\star x$,
and $\Delta_\psi(x')=\pm\Delta_\psi(x)$ accordingly.

(c) $(a(t),1)$ is conjugate, by $(n(1),1)$ and then $(w,1)$, to
$(g_2,(t^2-1,t)_v)$ with $g_2=\begin{pmatrix}t^{-1}&0\\t-t^{-1}&t
\end{pmatrix}$: the first conjugation has sign $1$, and the second has
sign $\beta(w,g_1)\beta(wg_1,w^{-1})=(1-t^{-2},t)_v=(t^2-1,t)_v$, where
$g_1=n(1)a(t)n(-1)$. Here
$D=(t+1)^2/t$, so $(c,D)_v=((t^2-1)/t,t)_v=(t^2-1,t)_v(t,-1)_v$ and
$\Delta_\psi(a(t),1)=(t,-1)_v\gamma_v(t)/\gamma_v(1)$. By
Lemma~\ref{lem:weil}(a), $(\gamma_v(t)/\gamma_v(1))^2=(t,-1)_v$, so this
is $\gamma_v(1)/\gamma_v(t)$; and $\chi_{\psi_{-1}}=\chi_\psi\cdot
(-1,\cdot)_v$ by Lemma~\ref{lem:weil}(d).

(d) $\gamma_{\psi_a,v}(D)/\gamma_{\psi_a,v}(1)=\gamma_v(aD)/\gamma_v(a)=
(a,D)_v\,\gamma_v(D)/\gamma_v(1)$, by Lemma~\ref{lem:weil}(a).
\end{proof}

So stable conjugacy can be described without $GL(2)$: \emph{$x$ and $x'$
are stably conjugate when $p(x)$, $p(x')$ are stably conjugate in
$SL(2,F)$ and $\Delta_\psi(x)=\Delta_\psi(x')$}; by~(d) the condition does
not depend on $\psi$. This is the definition of Adams~\cite{Ad} and
Li~\cite[\S5]{Li11}, who use, in place of $\Delta_\psi$, the character of
the difference of the two halves of the Weil representation. The two
functions agree up to a positive factor:

\begin{lemma}\label{lem:weil-kernel}
Let $\omega=\omega_{\psi'}$ be the Weil representation
of~\eqref{eq:weil-rep}, and $\gamma'$ the Weil index of $\psi'$. Then
$\omega(w,1)\phi(x)=\gamma'(1)^{-1}\int\phi(y)\psi'(-2xy)\,dy$, and for
$g=\begin{pmatrix}\alpha&b\\c&d\end{pmatrix}$ with $c\ne0$,
\begin{equation}\label{eq:weil-kernel}
\omega(g,1)\phi(x)=\gamma'(c)^{-1}\,|c|_v^{-1/2}\int_{F}\psi'\Bigl(
\frac{\alpha x^2-2xy+dy^2}{c}\Bigr)\phi(y)\,dy,
\end{equation}
with the measure that is self-dual for $\psi'(2xy)$. Moreover
$\omega_{\psi'}\circ\mathrm{Ad}(\delta_a)=\omega_{\psi'_a}$ for $a\in
F^\times$, as representations on $\mathcal S(F)$; in particular
$r^\pm_{\psi'}\circ\mathrm{Ad}(\delta_a)=r^\pm_{\psi'_a}$.
\end{lemma}

\begin{proof}
In the proof of Lemma~\ref{lem:weil-rep} it is shown that
$\omega(w(t),1)\phi(x)=\gamma'(t)|t|^{1/2}(\mathcal F\phi)(tx)$ for
$w(t)=\begin{pmatrix}0&t\\-t^{-1}&0\end{pmatrix}$; for $t=-1$ this is the
first formula, since $w(-1)=w$ and $\gamma'(-1)=\gamma'(1)^{-1}$. For
$c\ne0$, $g=n(\alpha/c)\,a(c^{-1})\,w\,n(d/c)$, and the same holds for
the lifts $(\cdot,1)$, because the three values of Kubota's cocycle are
$(c,1)_v$, $(1,c)_v$ and $(1,c)_v$. Apply~\eqref{eq:weil-rep},
Lemma~\ref{lem:weil-rep} with $t=c^{-1}$, and $\gamma'(-1)\gamma'(1)=1$.
For the last statement it suffices to compare the two representations on
$(n(u),1)$ and $(w,1)$, which generate $\Gb$ together with $\mu_2$. By
Lemma~\ref{lem:Ad}(d), $\omega_{\psi'}(\mathrm{Ad}(\delta_a)(n(u),1))$ is
multiplication by $\psi'(aux^2)$, and
$\omega_{\psi'}(\mathrm{Ad}(\delta_a)(w,1))\phi(x)=\frac{\gamma'(1)}
{\gamma'(a)}|a|^{1/2}\cdot\gamma'(1)^{-1}\int\phi(y)\psi'(-2axy)\,dy$. This
is $\omega_{\psi'_a}(w,1)\phi(x)$: the Weil index of $\psi'_a$ at $1$ is
$\gamma'(a)$, and the measure self-dual for $\psi'_a(2xy)$ is $|a|^{1/2}$
times the one for $\psi'(2xy)$.
\end{proof}

As a check on the constant in~\eqref{eq:weil-kernel}: composing the
kernels of $(g_1,1)$ and $(g_2,1)$ by a Gaussian integral gives the
kernel of $g_1g_2$ times $\gamma'(c_3/c_1c_2)\gamma'(c_3)/\gamma'(c_1)
\gamma'(c_2)$, where $c_3$ is the lower left entry of $g_1g_2$; by
Lemma~\ref{lem:weil}(a) this is $(c_3/c_1,c_3/c_2)_v=\beta(g_1,g_2)$. So
\eqref{eq:weil-kernel} is a genuine representation of the cover defined
by Kubota's cocycle, with no further constant.

\begin{proposition}\label{prop:weil-char}
Let $\Theta^\pm_{\psi}$ be the characters of $r^\pm_{\psi}$. On
$\Gb_{\mathrm{reg}}$ the distribution $\Theta^+_\psi-\Theta^-_\psi$ is the
function
\begin{equation}\label{eq:weil-char}
\Theta^+_\psi(x)-\Theta^-_\psi(x)=\frac{\Delta_\psi(x)}{|D^+(x)|_v^{1/2}} ;
\end{equation}
that is, $\operatorname{tr}r^+_\psi(f)-\operatorname{tr}r^-_\psi(f)=\int
f(x)\Delta_\psi(x)|D^+(x)|_v^{-1/2}dx$ for $f\in C_c^\infty(\Gb)$ with
support in $\Gb_{\mathrm{reg}}$. Hence $\Theta^+_\psi-\Theta^-_\psi$ is
invariant under $\star$ on $\Gb_{\mathrm{reg}}$, it does not vanish there,
and stable conjugacy as defined above coincides with stable conjugacy in
the sense of~\cite{Ad,Li11}. The right side of~\eqref{eq:weil-char} is
locally integrable on $\Gb$, and by the regularity theorem for the
character of the Weil
representation~\cite{Mak,Th08} the formula holds for all $f\in
C_c^\infty(\Gb)$: $\Theta^+_\psi-\Theta^-_\psi$ is a stable character.
\end{proposition}

\begin{proof}
Let $P\phi(x)=\phi(-x)$, so that $\operatorname{tr}r^+_\psi(f)-
\operatorname{tr}r^-_\psi(f)=\operatorname{tr}(\omega_\psi(f)P)$. Both
sides of the identity to be proved are unchanged when $f$ is replaced by a
conjugate $f^y$. Let $X$ be the set of the proof of
Lemma~\ref{lem:Delta}; it is open, and every element of
$\Gb_{\mathrm{reg}}$ is conjugate to an element of it. So each point of the
support of $f$ has an open neighborhood of the form $yUy^{-1}$ with
$U\subset X$ open. Finitely many of these cover the support, and a
partition of unity subordinate to them writes $f$ as a finite sum of
functions $f_i\in C_c^\infty(\Gb)$ with $f_i^{y_i}$ supported in $X$. We
may therefore assume that $f$ has compact support in $X$. Then $c$ and
$D=D^+(g)$ are bounded away from $0$ and $\infty$ on the support of $f$.

For $x=(g,\zeta)\in X$ let $k_x(\xi,\eta)$ be $\zeta$ times the kernel
of~\eqref{eq:weil-kernel}, so that $|k_x|=|c|_v^{-1/2}$ and
\[
k_x(\xi,-\xi)=\zeta\,\gamma_v(c)^{-1}|c|_v^{-1/2}\,\psi\bigl(D\xi^2/c\bigr).
\]
By Fubini's theorem $\omega_\psi(f)$ has the kernel $K_f(\xi,\eta)=\int
f(x)k_x(\xi,\eta)\,dx$, which is continuous and bounded.

\emph{Let $v$ be finite.} Then $f$ is bi-invariant under the image $K'$ of
some $K_m$, and $\omega_\psi(f)=\omega_\psi(e')\omega_\psi(f)\omega_\psi(e')$
for the idempotent $e'$ of $K'$. The space $V'$ of the vectors of
$L^2(F)$ fixed by $K'$ is finite-dimensional and consists of
Schwartz--Bruhat functions, by a two-lattice argument. The elements $(n(u),1)$, $u\in\varpi^m\OO$, act by multiplication
by $\psi(u\xi^2)$, so a fixed vector is supported in the bounded set
$L_1=\{\xi\mid\psi(u\xi^2)=1\text{ for all }u\in\varpi^m\OO\}$. The
conjugates by $(w,1)$ of the $(n(u),1)$ with $u\in2\varpi^m\OO$ lie in $K'$
(they lie over $K_m$, and the sign is a character of $u\in\varpi^m\OO$,
hence trivial on $2\varpi^m\OO$), and they act in the same way on the
Fourier transform, by~\eqref{eq:weil-rep}; so the Fourier
transform of a fixed vector is supported in a bounded set as well, and the
vector is constant on the cosets of a lattice $L_2$. Functions supported in a
bounded set and constant modulo a lattice form a finite-dimensional space.
So $\omega_\psi(f)$ has finite rank: $\omega_\psi(f)\phi=\sum_i\langle
\phi,w_i\rangle u_i$ with finitely many $u_i,w_i\in V'$ (the functionals
are of this form because $\omega_\psi(e')$ is the orthogonal projection
onto $V'$). Its kernel is $\sum_iu_i(\xi)\overline{w_i(\eta)}$, a
Schwartz--Bruhat function, which equals the continuous function $K_f$,
and $\operatorname{tr}(\omega_\psi(f)P)=\sum_i\langle Pu_i,w_i\rangle=
\int_FK_f(\xi,-\xi)\,d\xi$. Hence
\[
\operatorname{tr}(\omega_\psi(f)P)=\lim_{M\to\infty}\int f(x)\Bigl(
\int_{\varpi^{-M}\OO}k_x(\xi,-\xi)\,d\xi\Bigr)dx .
\]
By Lemma~\shellref below, applied to $a=D/c$, the inner integral
does not depend on $M$ for $M\ge M_0$, with $M_0$ uniform on the support
of $f$, and by~\eqref{eq:weil-index} (the measure being $|2|_v^{1/2}$ times
the one there) its value is $\zeta$ times
\[
\gamma_v(c)^{-1}|c|_v^{-1/2}\int_F\psi\Bigl(\frac{D\,\xi^2}{c}\Bigr)d\xi
=\frac{\gamma_v(cD)}{\gamma_v(c)}\,|D|_v^{-1/2}
=(c,D)_v\,\frac{\gamma_v(D)}{\gamma_v(1)}\,|D|_v^{-1/2};
\]
the last step is Lemma~\ref{lem:weil}(a). This is~\eqref{eq:weil-char}.

\emph{Let $v$ be archimedean.} We show that $K_f$ is a Schwartz function
on $F\times F$. On $X$ use the coordinates $(\alpha,c,d)$. The section
$g\mapsto(g,1)$ is continuous on each connected component of $\{c\ne0\}$:
$g=n(\alpha/c)\,a(-c^{-1})\,w\,n(d/c)$, the sections of
Lemma~\ref{lem:local}(i) are continuous on the factors, and the cocycle
of the factors is a Hilbert symbol in $c$, which is locally constant. The Haar
measure is a smooth
multiple of $d\alpha\,dc\,dd$, and $\gamma_v(c)$ is locally constant. So
\[
K_f(\xi,\eta)=\int_{F^\times}\widehat\Phi_c\Bigl(\frac{\xi^2}c,\frac{\eta^2}c
\Bigr)\,\psi\Bigl(-\frac{2\xi\eta}c\Bigr)\,dc,
\]
where $\Phi(\alpha,c,d)$ is smooth with compact support in $\{c\ne0\}$ and
$\widehat\Phi_c$ is its Fourier transform in $(\alpha,d)$. The functions
$\widehat\Phi_c$ and their derivatives in $c$ are Schwartz functions,
uniformly for $c$ in the compact set in question, and $|\xi|^2+|\eta|^2$
is bounded by a constant times $|\xi^2/c|+|\eta^2/c|$ there. So $K_f$ decays
faster than any power of $1+|\xi|+|\eta|$, and so do its derivatives, which
are integrals of the same kind with polynomial factors. An integral
operator on $L^2(F)$ with a Schwartz kernel is of trace class, and its
trace is the integral of the kernel over the diagonal.
So $\operatorname{tr}(\omega_\psi(f)P)=\int K_f(\xi,-\xi)\,d\xi=\lim_{\epsilon
\to0}\int f(x)\bigl(\int
k_x(\xi,-\xi)e^{-\epsilon|\xi|^2}d\xi\bigr)dx$. The inner integrals are
Gaussian; they are bounded by a constant times $|D|_v^{-1/2}$, uniformly
in $\epsilon$, and converge to the value above. Dominated convergence
gives~\eqref{eq:weil-char}.

The remaining statements follow from Lemma~\ref{lem:Delta}(b) and
$\Delta_\psi\ne0$. Local integrability: $|D^+|_v^{-1/2}$ is the inverse
square root of the function $2+\operatorname{tr}$, which vanishes to first
order on the smooth part of its zero set and is a nondegenerate quadratic
form in three variables at $-1$.
\end{proof}

\begin{lemmaprime}
Let $F$ be non-archimedean, let $\psi$ be trivial on $\varpi^{n_0}\OO$ and
not on $\varpi^{n_0-1}\OO$, and let $e=\mathrm{ord}(2)$. For $a\in F^\times$
and $j\in\Z$ with
\[
2j\ \ge\ \mathrm{ord}(a)-n_0+2+2e,
\]
the integral of $\psi(a\xi^2)$ over the shell $\{\mathrm{ord}(\xi)=-j\}$
vanishes. Consequently $\int_{\varpi^{-M}\OO}\psi(a\xi^2)\,d\xi$ does not
depend on $M$ once $2(M+1)\ge\mathrm{ord}(a)-n_0+2+2e$, and the bound is
uniform for $a$ in a compact subset of $F^\times$.
\end{lemmaprime}

\begin{proof}
Put $m=\mathrm{ord}(a)-2j$, the order of $a\xi^2$ on the shell, and
$r=n_0-1-m-e$. The hypothesis is $m\le n_0-2-2e$; it gives $r\ge1+e\ge1$
and $m+2r\ge n_0$. For $u\in\varpi^r\OO$ the substitution
$\xi\mapsto\xi(1+u)$ preserves the shell and the measure, and
\[
\psi\bigl(a\xi^2(1+u)^2\bigr)=\psi(a\xi^2)\,\psi(2a\xi^2u)\,\psi(a\xi^2u^2)
=\psi(a\xi^2)\,\psi(2a\xi^2u),
\]
because $\mathrm{ord}(a\xi^2u^2)\ge m+2r\ge n_0$. Averaging over
$u\in\varpi^r\OO$ multiplies the integrand by the average of the character
$u\mapsto\psi(2a\xi^2u)$ of $\varpi^r\OO$, which is not trivial because
$2a\xi^2\varpi^r\OO=\varpi^{m+e+r}\OO=\varpi^{n_0-1}\OO$. So the integral
vanishes.
\end{proof}

For $F=\Q_2$ and $\psi(x)=e^{2\pi i\{x\}}$ the lemma says that the
integrals over $2^{-M}\Z_2$ are constant from $M=1$ on for $a=1$ and from
$M=2$ on for $a=2$; the values are $1+i$ and $\sqrt2(1+i)$.

The same computation without $P$ gives the character of $\omega_\psi$
itself on the set where $D^-=\operatorname{tr}g-2\ne0$:
\begin{equation}\label{eq:weil-char-full}
\Theta_{\omega_\psi}(g,\zeta)=\zeta\,(c,D^-)_v\,\frac{\gamma_v(D^-)}
{\gamma_v(1)}\,|D^-|_v^{-1/2},\qquad c\ne0 .
\end{equation}
Two comparisons. First, $\det(g-1)=-D^-$ and $\gamma_v(1)\gamma_v(-D^-)=
(-1,D^-)_v\gamma_v(D^-)/\gamma_v(1)$, so~\eqref{eq:weil-char-full} is
$(-c,D^-)_v\cdot\gamma_v(1)\gamma_v(\det(g-1))|\det(g-1)|_v^{-1/2}$. This is
the formula of Thomas~\cite[Theorem~1C]{Th08}. That theorem is stated up
to a sign, so the comparison confirms the absolute value and the Weil
index and says nothing about the sign $(-c,D^-)_v$, which rests on the
computation above. Second,
$\omega_\psi(-1,1)=\gamma_v(1)^2P$ (Lemma~\ref{lem:weil-rep}), so
$\Theta^+_\psi(x)-\Theta^-_\psi(x)=\gamma_v(1)^{-2}\,\Theta_{\omega_\psi}
\bigl((-1,1)x\bigr)$; and indeed~\eqref{eq:weil-char-full} at
$(-1,1)(g,1)=(-g,(c,-1)_v)$ gives back~\eqref{eq:weil-char}, by
Lemma~\ref{lem:weil}(a) and $\gamma_v(1)^4=(-1,-1)_v$. The character of the
difference in~\cite[Theorem~D]{Th13} stands in the same relation to the
character of the Weil representation at $-g$; that theorem is written in
other coordinates on the metaplectic group and with a Weil index of a
quadratic form, and we have not verified an equality between its factors
and ours, nor matched the normalizations of the Weil index
in~\cite{Th08,Th13} with ours. Lemma~\ref{lem:Delta}, on which
the rest of the section depends, is algebraic and uses none of this
(Remark~\ref{rem:sec7-status}).

\begin{example}\label{ex:Delta-real}
Let $F=\R$. For $0<\theta<\pi$ the element $\tilde k_\theta=(k_\theta,1)$
has $c=-\sin\theta<0$ and $D^+=4\cos^2\frac\theta2>0$, so
$\Delta_\psi(\tilde k_\theta)=1$ and~\eqref{eq:weil-char} gives
$\Theta^+-\Theta^-=1/(2\cos\frac\theta2)$. Directly: if $r^+_\psi$ has
lowest weight $\frac12$, its weights are $\frac12+2j$ and those of
$r^-_\psi$ are $\frac32+2j$, $j\ge0$, and
\[
\sum_{j\ge0}e^{i(\frac12+2j)\theta}-\sum_{j\ge0}e^{i(\frac32+2j)\theta}
=\frac{e^{i\theta/2}}{1+e^{i\theta}}=\frac1{2\cos\frac\theta2}
\]
(Abel summation). The function is even in $\theta$, in agreement with
Example~\ref{ex:eta}(b): $\tilde k_\theta$ and $\tilde k_{-\theta}=
\mathrm{Ad}(\delta_{-1})\tilde k_\theta$ are stably conjugate, and they
are not conjugate. On the hyperbolic elements over $a(t)$, $t<0$, the
stable class is the conjugacy class, and $\mathrm{Ad}(\delta_{-1})$
exchanges the two lifts.
\end{example}

\begin{example}\label{ex:Delta-finite}
Let $v$ be finite and non-dyadic, $\psi$ unramified, and $x=(g,\kappa(g))
\in K^*$ with $c$ and $D^+(g)$ units. Then $\Delta_\psi(x)=1$ and
$\Delta_{\psi_\varpi}(x)=(\varpi,D^+(g))_v$, the quadratic residue symbol
of $\det(1+g)$. The second value can be seen in the Weil representation
of $SL(2,\mathbb F_q)$ on the functions on $\mathbb F_q$, which is the
space of vectors of $\omega_{\psi_\varpi}$ fixed by $K_1$ (the functions on
$\varpi^{-1}\OO/\OO$): there the trace of $\omega(g)P$ is
$\bigl(\frac{\det(1+g)}{q}\bigr)$ for every $g$ with $\det(1+g)\ne0$. The
script verifies this for $q=3,5,7,11$ by computing the matrices. (The
character of $\omega_{\psi_\varpi}$ also involves the vectors of higher
level, which we have not computed separately.)
\end{example}

\paragraph{The dual group.} The unramified genuine principal series of
$\Gb$ are parametrized by $\chi_v(\varpi)\in\C^\times$ modulo inversion,
through $\pi_v=\chi_{\psi,v}\cdot(\chi_v\circ p)$, and the intertwining
operator has the normalizing factor $L_v(2\lambda,\chi_v^2)/L_v(2\lambda+1,
\chi_v^2)$ (Lemma~\ref{lem:intertwining}). This is the normalizing factor
of the principal series of $PGL(2,F)$ induced from $\diag(t_1,t_2)\mapsto
\chi_v(t_1/t_2)$, whose Satake parameter is $\diag(\chi_v(\varpi),
\chi_v(\varpi)^{-1})\in SL(2,\C)$. So the dual group of the double cover
of $SL(2)$ is $SL(2,\C)$, the dual group of $SO(3)=PGL(2)$, and not
$PGL(2,\C)$, the dual group of $SL(2)$~\cite{Sav,Li11,GGW,We18}. The
elliptic endoscopic data correspond to the two central elements $\pm1$
of $SL(2,\C)$. This is why the endoscopic groups are $SO(3)$ twice, why
the translation by $-1$ appears below, and why no tori occur.

\subsection{Stable and \texorpdfstring{$\kappa$}{kappa}-orbital
integrals}\label{sec:kappa}

We now sort the orbital integrals in a $G_1(F)$-conjugacy class by the
characters of $F^\times/N E^\times$. For $SL(2)$ this gives the stable
orbital integral and one $\kappa$-orbital integral~\cite{LL}. On the
cover there are again two combinations, and the second is a stable
orbital integral at another point.

Fix Haar measures on the centralizers so that they are compatible under
every $\mathrm{Ad}(h)$, $h\in G_1(F)$ (Section~\ref{sec:measures-local}).
For $x\in\Gb_{\mathrm{reg}}$ over $\gamma$, with centralizer $T_x=T_\gamma$
in $SL(2,F)$, and $f\in C_c^\infty(\Gb)_-$, let
\[
J_G(x,f)=|D(\gamma)|_v^{1/2}\int_{T_x\backslash SL(2,F)}f(y^{-1}xy)\,dy,
\qquad D(\gamma)=(\operatorname{tr}\gamma)^2-4,
\]
as in~\eqref{eq:JG}. Then $J_G(x,f^h)=J_G(\mathrm{Ad}(h)x,f)$ and
$J_G(x,f^{[h]})=J_G(h\star x,f)$ for $h\in G_1(F)$. Let $E=F[\gamma]$,
and define
\begin{equation}\label{eq:Jst}
J^{\mathrm{st}}_G(x,f)=\sum_{a\in F^\times/N E^\times}J_G(\delta_a\star x,f),
\qquad
J^{\kappa}_G(x,f)=\sum_{a\in F^\times/N E^\times}\kappa_E(a)\,
J_G(\delta_a\star x,f).
\end{equation}
The terms depend only on $a$ modulo $N E^\times$, because $\delta_{aN(e)}
\star x$ is conjugate to $\delta_a\star(e\star x)=\delta_a\star x$ for
$e\in E^\times$. If $\gamma$ is hyperbolic, both are $J_G(x,f)$.

\begin{lemma}\label{lem:kappa-orbital}
Let $x\in\Gb_{\mathrm{reg}}$ lie over $\gamma$, $E=F[\gamma]$, and
$f\in C_c^\infty(\Gb)_-$.
\begin{enumerate}[label=(\alph*)]
\item $J^{\mathrm{st}}_G(h\star x,f)=J^{\mathrm{st}}_G(x,f)$ and
  $J^\kappa_G(h\star x,f)=\kappa_E(\det h)J^\kappa_G(x,f)$ for $h\in
  G_1(F)$. Equivalently, $J^{\mathrm{st}}_G(x,f^{[h]})=J^{\mathrm{st}}_G(x,
  f)$: the distribution $J^{\mathrm{st}}_G(x,\cdot)$ is stable.
\item For a character $\kappa'$ of $F^\times/F^{\times2}$,
  \[
  \sum_{a\in F^\times/F^{\times2}}\kappa'(a)\,J_G(x,f^{\delta_a})
  =[N E^\times:F^{\times2}]\cdot\begin{cases}
  J^{\mathrm{st}}_G(x,f),&\kappa'=\eta_\gamma,\\
  J^{\kappa}_G(x,f),&\kappa'=\eta_\gamma\kappa_E,\\
  0,&\text{otherwise.}\end{cases}
  \]
  In particular the average of $J_G(x,f^h)$ over $h\in G_1(F)/Z(F)SL(2,F)$
  vanishes unless $D^+(\gamma)\in F^{\times2}\cup\theta F^{\times2}$; for
  $\gamma=a(t)$, unless $t$ is a square.
\item Let $E$ be a field, and let $\tilde z$ be a lift of $-1$. Then
  \[
  J^\kappa_G(x,f_{\tilde z})=J^{\mathrm{st}}_G(\tilde zx,f),
  \qquad\text{that is,}\qquad
  J^\kappa_G(x,f)=J^{\mathrm{st}}_G(\tilde zx,f_{\tilde z^{-1}}).
  \]
\item Let $v$ be finite and non-dyadic, $f=e_v$, $x\in K^*$ and
  $|D(\gamma)|_v=1$, and let $T_\gamma\cap K$ have volume $1$. Then
  $J^{\mathrm{st}}_G(x,e_v)=J^\kappa_G(x,e_v)=J_G(x,e_v)=1$.
\end{enumerate}
\end{lemma}

\begin{proof}
(a) $\delta_a\star(h\star x)=(\delta_ah)\star x$. Write
$\delta_ah=\delta_{a\det h}h_0$ with $h_0\in SL(2,F)$; then $(\delta_ah)
\star x$ is conjugate to $\delta_{a\det h}\star x$, by
$\delta_{a\det h}h_0\delta_{a\det h}^{-1}$. So the terms of~\eqref{eq:Jst}
are permuted by $a\mapsto a\det h$. The reformulation is
$J_G(y,f^{[h]})=J_G(h\star y,f)$.

(b) $J_G(x,f^{\delta_a})=J_G(\mathrm{Ad}(\delta_a)x,f)=\eta_\gamma(a)
J_G(\delta_a\star x,f)$, since $f$ is anti-genuine. So the left side is
$\sum_a\kappa'(a)\eta_\gamma(a)J_G(\delta_a\star x,f)$, and the summand
$J_G(\delta_a\star x,f)$ depends only on $a$ modulo $N E^\times$. The sum
of $\kappa'\eta_\gamma$ over a coset of $N E^\times/F^{\times2}$ is $0$
unless $\kappa'\eta_\gamma$ is trivial on $N E^\times$, that is, unless
$\kappa'\eta_\gamma\in\{1,\kappa_E\}$. For $\kappa'=1$ the condition is
$\eta_\gamma\in\{1,\kappa_E\}$.

(c) Write $\gamma=\begin{pmatrix}x_0&y_0\theta\\y_0&x_0\end{pmatrix}$ in a
suitable basis, $x_0^2-\theta y_0^2=1$. Then
$D^+(\gamma)D^+(-\gamma)=4-(2x_0)^2=-4\theta y_0^2$, so
\begin{equation}\label{eq:eta-minus}
\eta_\gamma\,\eta_{-\gamma}=(\cdot,-\theta)_v=\kappa_E\cdot(\cdot,-1)_v .
\end{equation}
By Lemma~\ref{lem:eta}, $\mathrm{Ad}(\delta_a)\tilde z=(a,-1)_v\tilde z$.
Hence
\[
\delta_a\star(\tilde zx)=\eta_{-\gamma}(a)\,(a,-1)_v\,\tilde z\,
\mathrm{Ad}(\delta_a)x=\eta_{-\gamma}(a)(a,-1)_v\eta_\gamma(a)\cdot
\tilde z\,(\delta_a\star x)=\kappa_E(a)\,\tilde z\,(\delta_a\star x).
\]
Since $J_G(\tilde zy,f)=J_G(y,f_{\tilde z})$, $J_G(\varepsilon y,f)=
\varepsilon J_G(y,f)$ for $\varepsilon\in\mu_2$, and $D(-\gamma)=D(\gamma)$,
summing over $a$ gives $J^{\mathrm{st}}_G(\tilde zx,f)=J^\kappa_G(x,
f_{\tilde z})$.

(d) $J_G(x,e_v)=1$ is the unramified computation used for
Theorem~\ref{thm:fine}. If $E$ is a field, it is unramified, and
$\delta_\varpi\star x$ lies over $\delta_\varpi\gamma\delta_\varpi^{-1}$,
which has no $SL(2,F)$-conjugate in $K$ (proof of
Theorem~\ref{thm:coarse}(a)); so $J_G(\delta_\varpi\star x,e_v)=0$.
\end{proof}

Part~(b) explains a statement of the 1995 notes: for the double cover the
orbital integral of a genuine function at a lift of $a(t)$ does not
vanish when $t$ is not a square, but its average over the $G_1(F)$-orbit
does (Appendix~\ref{sec:appendix-1995}). Part~(c) is the form that the
relation between the endoscopic data $(1,0)$ and $(0,1)$ takes here: Li's
transfer factors for the two data differ by the translation by the
central element $-1$~\cite[\S5]{Li11}, and correspondingly the
$\kappa$-orbital integral of $f$ at $x$ is the stable orbital integral of
a translate of $f$ at $\tilde zx$. For $SL(2)$ there is no such
identity: there $J^\kappa$ is a new function on the torus $T_\gamma$, and
\cite{LL} apply Poisson summation on $T_\gamma$ to it.

By Lemma~\ref{lem:Delta}(b), $x\mapsto\Delta_\psi(x)^{-1}J^{\mathrm{st}}_G
(x,f)$ is invariant under $\mu_2$ and under $\star$, so it is a function
of the stable conjugacy class of $p(x)$ in $SL(2,F)$, that is, of the
trace. Li's transfer theorem~\cite{Li11} says that it is the stable
orbital integral of a function $f'$ on $SO(3,F)$ at the element with
eigenvalues $\epsilon,1,\epsilon^{-1}$, where $\epsilon^{\pm1}$ are the
eigenvalues of $p(x)$, up to a normalization of the transfer factor. We
do not use this.

\subsection{The elliptic regular terms}\label{sec:elliptic-stable}

Let $\gamma\in G(\Q)=SL(2,k)$ be regular elliptic, $E=k[\gamma]$ its
splitting field, $T_\gamma$ its centralizer and $\kappa_E=\prod_v\kappa_{E_v}$
the quadratic character of $k^\times\backslash\A_k^\times$ attached to
$E$. Everything in Sections~\ref{sec:Ad}--\ref{sec:kappa} has an adelic
version. Let
\[
\eta_\gamma(a)=\prod_v\bigl(a_v,D^+(\gamma)\bigr)_v,\qquad a\in\A_k^\times,
\]
a quadratic character of $\A_k^\times$ which is trivial on $k^\times$ by
the product formula~\ref{H4}, because $D^+(\gamma)=2+\operatorname{tr}
\gamma\in k^\times$. For $h\in G_1(\A_k)$ and $x$ in the orbit of
$\pm s(\gamma)$ under $\mathrm{Ad}(G_1(\A_k))$ put
$h\star x=\eta_\gamma(\det h)\,\mathrm{Ad}(h)x$; its component at $v$ is
$h_v\star x_v$ in the sense of~\eqref{eq:star}. So for $a\in k^\times$,
\begin{equation}\label{eq:rational-stable}
s(\delta_a\gamma\delta_a^{-1})=\mathrm{Ad}(\delta_a)s(\gamma)=\delta_a
\star s(\gamma)
\end{equation}
by Lemma~\ref{lem:Ad}(f): \emph{rational elements that are conjugate
under $G_1(k)$ are stably conjugate in $\GA$, at every place.} This is
what makes a stabilization possible, and it is the reason for the twist
in~\eqref{eq:star}: the untwisted local conjugates
$\mathrm{Ad}(\delta_a)\gamma_v$ are not coherent in $v$. (Compare
\cite[Lemme~4.1.7]{Li15}.)

For $f\in\mathcal H(\GA)$ and $x=a\star s(\gamma)$, $a\in\A_k^\times$
(we write $a\star$ for $\delta_a\star$), let $J_G(x,f)=\int f(y^{-1}xy)\,
dy$ over $T_x(\A)\backslash G(\A)$, with the Tamagawa measures on $G(\A)$
and on $T_x(\A)=\delta_aT_\gamma(\A)\delta_a^{-1}$; no factor $|D|^{1/2}$ is
needed, since $|D(\gamma)|_{\A_k}=1$. Put
\begin{equation}\label{eq:Jst-global}
J^{\mathrm{st}}_G(\gamma,f)=\sum_{a\in\A_k^\times/N\A_E^\times}J_G\bigl(a\star
s(\gamma),f\bigr),\qquad
J^{\kappa}_G(\gamma,f)=\sum_{a\in\A_k^\times/N\A_E^\times}\kappa_E(a)\,J_G\bigl(
a\star s(\gamma),f\bigr).
\end{equation}
Here $\A_k^\times/N\A_E^\times=\bigoplus_vk_v^\times/N E_v^\times$. By
Lemma~\ref{lem:kappa-orbital}(d) only finitely many $a$ contribute, and
for $f=\bigotimes f_v$ and local measures whose products are the Tamagawa
measures,
\[
J^{\mathrm{st}}_G(\gamma,f)=\prod_vJ^{\mathrm{st}}_G(\gamma_v,f_v),\qquad
J^{\kappa}_G(\gamma,f)=\prod_vJ^{\kappa}_G(\gamma_v,f_v),
\]
where $(\gamma_v)$ are local components of $s(\gamma)$; the products do
not depend on their choice, and almost all factors are $1$.

\begin{theorem}\label{thm:elliptic-stable}
For $f\in\mathcal H(\GA)$ let
\[
S_{\mathrm{ell}}(f)=\frac12\sum_{\gamma}\mathrm{vol}\bigl(T_\gamma(\Q)
\backslash T_\gamma(\A)\bigr)\,J^{\mathrm{st}}_G(\gamma,f),
\]
the sum over a set of representatives $\gamma$ of the $G_1(k)$-conjugacy
classes of regular elliptic elements of $SL(2,k)$. The sum is finite, and
\[
\sum_{\mathcal O\ \mathrm{elliptic}}J_{\mathcal O}(f)=S_{\mathrm{ell}}(f)
+S_{\mathrm{ell}}(f_-).
\]
More precisely, for each $\gamma$ the classes $\mathcal O$ contained in
the $G_1(k)$-conjugacy class of $\gamma$ contribute
\begin{equation}\label{eq:pre-stab}
\sum_{\mathcal O\subset G_1(k)\cdot\gamma}J_{\mathcal O}(f)=\frac12\,
\mathrm{vol}\bigl(T_\gamma(\Q)\backslash T_\gamma(\A)\bigr)\Bigl(
J^{\mathrm{st}}_G(\gamma,f)+J^{\kappa}_G(\gamma,f)\Bigr),
\qquad J^\kappa_G(\gamma,f)=J^{\mathrm{st}}_G(-\gamma,f_-).
\end{equation}
The distribution $S_{\mathrm{ell}}$ is stable: for $f=\bigotimes_vf_v$ it
is a finite sum of products of the stable local distributions
$J^{\mathrm{st}}_G(\gamma_v,\cdot)$ of Lemma~\ref{lem:kappa-orbital}(a).
\end{theorem}

\begin{proof}
The $SL(2,k)$-conjugacy classes in the $G_1(k)$-class of $\gamma$ are the
classes $\mathcal O_a$ of $\delta_a\gamma\delta_a^{-1}$, $a\in k^\times/
N E^\times$. By Proposition~\ref{prop:elliptic}
and~\eqref{eq:rational-stable}, $J_{\mathcal O_a}(f)=\mathrm{vol}\cdot
J_G(a\star s(\gamma),f)$, with the same volume for all $a$, the Tamagawa
measure being compatible with $\mathrm{Ad}(\delta_a)$. Let
$\Phi(a)=J_G(a\star s(\gamma),f)$ for $a\in\A_k^\times$. It depends only
on the image of $a$ in $\A_k^\times/N\A_E^\times$, and it vanishes outside
a finite set. By the Hasse norm theorem, $k^\times/N E^\times$ injects
into $\A_k^\times/N\A_E^\times$, and by class field theory its image is the
kernel of $\kappa_E$, of index $2$~\cite{CF}. Therefore
\[
\sum_{a\in k^\times/N E^\times}\Phi(a)=\frac12\sum_{a\in\A_k^\times/N\A_E^\times}
\bigl(1+\kappa_E(a)\bigr)\Phi(a)=\frac12\Bigl(J^{\mathrm{st}}_G(\gamma,f)+
J^\kappa_G(\gamma,f)\Bigr).
\]
For the second formula of~\eqref{eq:pre-stab} we may take
$f=\bigotimes f_v$. Choose local components $\gamma_v$ of $s(\gamma)$ and
$\tilde z_v$ of $s(-1)$. Then $(\tilde z_v\gamma_v)$ are local components
of $s(-1)s(\gamma)=s(-\gamma)$, and $\bigotimes_v(f_v)_{\tilde z_v^{-1}}=
f_{s(-1)^{-1}}=f_-$. By Lemma~\ref{lem:kappa-orbital}(c) at the places
where $E_v$ is a field (at the others both sides are $J_G$),
\[
J^\kappa_G(\gamma,f)=\prod_vJ^{\mathrm{st}}_G\bigl(\tilde z_v\gamma_v,
(f_v)_{\tilde z_v^{-1}}\bigr)=J^{\mathrm{st}}_G(-\gamma,f_-).
\]
Finally $\gamma\mapsto-\gamma$ is a bijection of the set of
$G_1(k)$-classes of regular elliptic elements, with $T_{-\gamma}=
T_\gamma$. Summing~\eqref{eq:pre-stab} over $\gamma$ gives the first
formula; the sums are finite by Theorem~\ref{thm:coarse}(a). Stability is
Lemma~\ref{lem:kappa-orbital}(a) at each place.
\end{proof}

For $SL(2)$ the first formula of~\eqref{eq:pre-stab} is the
pre-stabilization of~\cite[\S5]{LL}. The second has no analog there, and
it is the content of the theorem: the unstable half of the elliptic terms
of $f$ is the stable half of the elliptic terms of $f_-$. In the notation
of~\cite[Th\'eor\`eme~5.2.6]{Li15}, where the elliptic part of the trace
formula is $\sum_{(n',n'')}\iota(\tilde G,H)\,ST^H_{\mathrm{ell}}(f^H)$ with
$\iota=\frac12$ for both pairs when $n=1$, there are likewise two
summands with coefficient $\frac12$, those of $(1,0)$ and $(0,1)$. This is
a comparison of the decomposition and of the coefficients only. Li's terms
are stable distributions on $SO(3)$ applied to transfers of $f$, with
transfer factors; ours are distributions on the cover. Matching them
requires the transfer and its normalization, which are not carried out
here (Remark~\ref{rem:sec7-status}).

\begin{example}\label{ex:order4-stable}
Let $k=\Q$ and $\gamma=w$, of order $4$, with $E=\Q(i)$, $D^+(w)=2$ and
$\eta_w=\prod_v(\cdot,2)_v$. The $G_1(\Q)$-conjugacy class of $w$ is the
union of the $SL(2,\Q)$-classes of the $\delta_aw\delta_a^{-1}$,
$a\in\Q^\times/N E^\times$; there are infinitely many, and for a given $f$
finitely many contribute. The class $a=-1$ is that of
$w^{-1}=-w$. So $-w$ lies in the $G_1(\Q)$-class of $w$,
and~\eqref{eq:pre-stab} reads
\[
\sum_{\mathcal O\subset G_1(\Q)\cdot w}J_{\mathcal O}(f)=\tfrac12\,
\mathrm{vol}\bigl(T_w(\Q)\backslash T_w(\A)\bigr)\bigl(
J^{\mathrm{st}}_G(w,f)+J^{\mathrm{st}}_G(w,f_-)\bigr).
\]
At the real place the stable class of a lift of $w$ consists of two
conjugacy classes, exchanged by $\mathrm{Ad}(\delta_{-1})$
(Example~\ref{ex:Delta-real}); at a prime $p\equiv3\pmod4$ it consists of
two classes, of which one meets $K_p^*$; at $p\equiv1\pmod4$ it is one
class.
\end{example}

\subsection{The hyperbolic terms}\label{sec:hyperbolic-gl2}

The stable class of a hyperbolic element is its conjugacy class, so the
orbital integrals $J_G(x,f)$, $x\in\AS$ regular, are stable
distributions. The invariant distributions $I_A(x,f)$ of
Section~\ref{sec:IAgamma} are not orbital integrals, and their
stabilization is a substantial part of~\cite{LL} and of~\cite{Li21}. Here
we only determine how they transform under the untwisted action of
$G_1(k_S)$; this is the correct form of the variance formulas that the
1995 notes state for conjugation by $GL(2)$.

Let $S$, $\GS$, $\AS$ be as in Section~\ref{sec:invariant}, and
$G_1(k_S)=\prod_{v\in S_k}GL(2,k_v)$, acting on $\GS$ by $\mathrm{Ad}$.
Extend the height to $G_1$: for $g\in GL(2,k_v)$ put
\[
H_{1,v}(g)=-\log\|e_2g\|_v+\tfrac12\log|\det g|_v,
\]
so that $H_{1,v}=H_v$ on $SL(2,k_v)$ (Lemma~\ref{lem:height-ineq}(a)),
$H_{1,v}(zng)=H_{1,v}(g)$ for $z$ scalar and $n\in N$, and
$H_{1,v}(\diag(t_1,t_2)g)=\frac12\log|t_1/t_2|_v+H_{1,v}(g)$; let
$H_1=\sum_{v\in S_k}H_{1,v}$. For $\delta\in G_1(k_S)$ define
$u'_B(\tilde k,\delta)=-H_1(k\delta)$, $k=p(\tilde k)$, and $f_{B,\delta}$
by~\eqref{eq:fBy} with this weight; as in Lemma~\ref{lem:uprime}(b),
$f_{B,\delta}\in C_c^\infty(\AS)_-$. For $a\in k_S^\times$ let $\kappa_a$
be the character $\tilde a\mapsto\prod_{v\in S_k}(a_v,t_v)_v$ of $\AS$,
$p(\tilde a)=a(t)$; by Lemma~\ref{lem:Ad}(d), $\pi\circ
\mathrm{Ad}(\delta_a)=\pi\kappa_a$ for $\pi\in\Pi_u(\AS)$. We need one
more property of the normalizing factors of
Theorem~\ref{thm:normalization}:
\begin{enumerate}[label=(N\arabic*),start=5]
\item\label{N5} for every $v\in S_k$, $a\in k_v^\times$ and $\pi_v$, the
  quotient $r_v(\pi_v\kappa_a,\lambda)/r_v(\pi_v,\lambda)$ does not depend
  on $\lambda$.
\end{enumerate}
Since $\pi_v\kappa_a=\chi_{\psi,v}\cdot\chi_v(a,\cdot)_v$, the condition
holds whenever $r_v(\pi_v,\lambda)$ depends on $\pi_v=\chi_{\psi,v}\chi_v$
only through $\chi_v^2$, up to a constant factor. The explicit factors of
Lemma~\ref{lem:explicit-r} are of this form; so~\ref{N5} holds where
$\pi_v$ and $\pi_v\kappa_a$ are both covered by that lemma, which is
always the case at the real and complex places. At the other places
\ref{N5} is a condition on the choice of the factors, in addition
to~\ref{N1}--\ref{N4}. We have not proved that it can always be met
(Remark~\ref{rem:N5}). The statements below say where it is used, and
Lemma~\ref{lem:variance-gl2}(c) says what happens without it.

\begin{lemma}\label{lem:variance-gl2}
Let $f\in C_c^\infty(\GS)_-$, $a\in k_S^\times$, $\delta=\delta_a$, and
$\alpha=\mathrm{Ad}(\delta)$.
\begin{enumerate}[label=(\alph*)]
\item For $x\in\GS$ with Iwasawa decompositions $x=n\tilde a\tilde k$,
  $wx=n'\tilde a'\tilde k'$,
  \[
  v_A\bigl(\alpha^{-1}(x)\bigr)=v_A(x)-H_1(k\delta)-H_1(k'\delta).
  \]
\item For $\gamma\in\AS$ regular,
  $J_A(\gamma,f^\delta)=J_A(\alpha\gamma,f)+f_{B,\delta}(\alpha\gamma)+
  f_{B,\delta}(w\,\alpha\gamma\,w^{-1})$, and $J_G(\gamma,f^\delta)=
  J_G(\alpha\gamma,f)$.
\item For $\pi\in\Pi_u(\AS)$ and $\lambda\in i\R$, with $\pi'=\pi\circ
  \alpha^{-1}=\pi\kappa_a$: $\operatorname{tr}I(\pi_\lambda,f^\delta)=
  \operatorname{tr}I(\pi'_\lambda,f)$, and
  \[
  \varphi_A(f^\delta,\pi_\lambda)=\varphi_A(f,\pi'_\lambda)+
  \widehat{f_{B,\delta}}(\pi'_\lambda)+\widehat{f_{B,\delta}}\bigl(
  (\pi')^w_{-\lambda}\bigr)-d_a(\pi,\lambda)\,\operatorname{tr}
  I(\pi'_\lambda,f),
  \]
  where $d_a(\pi,\lambda)=\sum_{v\in S_k}\bigl(\frac{r_v'}{r_v}(\pi_v
  \kappa_a,\lambda)-\frac{r_v'}{r_v}(\pi_v,\lambda)\bigr)$. The
  condition~\ref{N5} says that $d_a=0$.
\item The products $r_v(\pi_v,\lambda)r_v(\pi_v^w,-\lambda)$ and
  $r_v(\pi_v\kappa_a,\lambda)r_v((\pi_v\kappa_a)^w,-\lambda)$ are equal up
  to a sign that does not depend on $\lambda$. So~\ref{N5} is compatible
  with~\ref{N2}.
\end{enumerate}
\end{lemma}

\begin{proof}
(a) The image of $\alpha^{-1}(x)$ in $G(\Q_S)$ is $\delta^{-1}g\delta$,
$g=p(x)$. By the properties of $H_1$, $H(\delta^{-1}g\delta)=H_1(
\delta^{-1})+H(g)+H_1(k\delta)$. Since $w\delta^{-1}w^{-1}=\diag(1,a^{-1})$,
$H(w\delta^{-1}g\delta)=H_1(\diag(1,a^{-1}))+H(wg)+H_1(k'\delta)$, and
$H_1(\diag(1,a^{-1}))=-H_1(\delta^{-1})$. Subtract from
$v_A=-H(\cdot)-H(w\,\cdot)$.

(b) $x\mapsto\alpha(x)$ preserves the Haar measure of $G(\Q_S)$ and fixes
$A(\Q_S)$ pointwise, so it preserves the measure on the quotient.
Substituting $x=\alpha^{-1}(x')$ in~\eqref{eq:JA},
\[
J_A(\gamma,f^\delta)=|D(\gamma)|^{1/2}\int f\bigl(x'^{-1}\alpha(\gamma)x'
\bigr)\,v_A\bigl(\alpha^{-1}(x')\bigr)\,dx' .
\]
Insert~(a) and argue as in the proof of
Lemma~\ref{lem:variance-local}(b), which uses only the form of the
difference of the weights.

(c) Realize $I(\pi_\lambda)$ on the space of functions $\Phi$ on $\GS$
with $\Phi(n\tilde ax)=\pi(\tilde a)e^{(\lambda+1)H_S(\tilde a)}\Phi(x)$,
by right translation, and let $(A(\lambda)\Phi)(x)=\Phi(\alpha^{-1}x)$.
Since $\alpha$ preserves $N(\Q_S)$ and $H_S$ on $\AS$, $A(\lambda)$ maps
$I(\pi_\lambda)$ to $I(\pi'_\lambda)$, and $A(\lambda)^{-1}I(\pi'_\lambda,
y)A(\lambda)=I(\pi_\lambda,\alpha^{-1}y)$; integrating against $f$,
\begin{equation}\label{eq:A-conj}
A(\lambda)^{-1}\,I(\pi'_\lambda,f)\,A(\lambda)=I(\pi_\lambda,f^\delta),
\end{equation}
which gives the first statement. In the compact picture $A(\lambda)$
depends on $\lambda$ through $e^{(\lambda+1)H_S(\alpha^{-1}\tilde k)}$, so
$\frac{d}{d\lambda}A(\lambda)=\tilde HA(\lambda)$, where $\tilde H$ is
multiplication by $\tilde k\mapsto H_S(\alpha^{-1}\tilde k)=H_1(\delta^{-1})
+H_1(k\delta)$. Next we compare intertwining operators. By
Lemma~\ref{lem:Ad}(d), $\alpha^{-1}((w,1)^{-1})=(w,1)^{-1}\tilde b$ with
$\tilde b\in\AS$ over $a(a)$, and $\alpha^{-1}(n(u))=n(a^{-1}u)$. Two
changes of variable in the integral defining $M(\pi',\lambda)$ give
\[
M(\pi',\lambda)\,A(\lambda)=c(\lambda)\,A(-\lambda)\,M(\pi,\lambda),
\qquad c(\lambda)=\pi^w(\tilde b)\,|a|_S^{-\lambda},
\]
where $A(-\lambda)$ is formed with $\pi^w$. Hence, with $R=r^{-1}M$ and
$\rho(\lambda)=r(\pi',\lambda)\,c(\lambda)^{-1}r(\pi,\lambda)^{-1}$,
\[
A(-\lambda)\,R_S(\pi,\lambda)=\rho(\lambda)\,R_S(\pi',\lambda)\,A(\lambda),
\qquad\frac{\dot\rho}{\rho}=\frac{\dot r}{r}(\pi',\lambda)-\frac{\dot r}{r}
(\pi,\lambda)+\log|a|_S ,
\]
where a dot denotes $d/d\lambda$. (The scalar $r(\pi,\lambda)$ is the product of the
$r_v(\pi_v,\lambda)$, $v\in S_k$.) Differentiate, and multiply on the left
by $(\rho R_S(\pi',\lambda))^{-1}=A(\lambda)R_S(\pi,\lambda)^{-1}
A(-\lambda)^{-1}$. Writing $R$ for $R_S(\pi,\lambda)$ and $R^\natural$ for
$R_S(\pi',\lambda)$,
\[
(R^\natural)^{-1}\dot R^\natural=A(\lambda)\,R^{-1}\dot R\,A(\lambda)^{-1}
-\frac{\dot\rho}{\rho}-\tilde H-(R^\natural)^{-1}\tilde HR^\natural .
\]
Now $\tilde H=-\frac12\log|a|_S-u'$, where $u'$ is multiplication by
$u'_B(\tilde k,\delta)=-H_1(k\delta)$, and $\dot\rho/\rho=d_a(\pi,\lambda)+
\log|a|_S$. The scalars cancel:
\[
-\frac{\dot\rho}{\rho}-\tilde H-(R^\natural)^{-1}\tilde HR^\natural
=\underbrace{-\log|a|_S+\tfrac12\log|a|_S+\tfrac12\log|a|_S}_{=\,0}
-d_a(\pi,\lambda)+u'+(R^\natural)^{-1}u'R^\natural .
\]
Take
the trace against $I(\pi'_\lambda,f)$, use~\eqref{eq:A-conj},
\ref{N1}, and Lemma~\ref{lem:uprime}(c) for the multiplication operator
$-H_1(k\delta)=u'_B(\tilde k,\delta)$:
\[
\varphi_A(f,\pi'_\lambda)=\varphi_A(f^\delta,\pi_\lambda)-
\widehat{f_{B,\delta}}(\pi'_\lambda)-\widehat{f_{B,\delta}}\bigl((\pi')^w_{
-\lambda}\bigr)+d_a(\pi,\lambda)\operatorname{tr}I(\pi'_\lambda,f).
\]

(d) Apply the relation $M(\pi',\lambda)A(\lambda)=c(\lambda)A(-\lambda)
M(\pi,\lambda)$ to $\pi$ and to $\pi^w$ at $-\lambda$; the second constant
is $\pi(\tilde b)|a|_S^{\lambda}$. So $M((\pi')^w,-\lambda)M(\pi',\lambda)$
is the conjugate by $A(\lambda)$ of $\pi(\tilde b)\pi^w(\tilde b)\cdot
M(\pi^w,-\lambda)M(\pi,\lambda)$. By~\ref{N2} the two compositions are the
scalars $r(\pi',\lambda)r((\pi')^w,-\lambda)$ and $r(\pi,\lambda)
r(\pi^w,-\lambda)$, and $\pi(\tilde b)\pi^w(\tilde b)=\pi\bigl(\tilde b\,
w\tilde bw^{-1}\bigr)$, where $\tilde b\,w\tilde bw^{-1}$ lies over $1$, is
$\pm1$. The argument is local at each $v\in S_k$.
\end{proof}

\begin{proposition}\label{prop:IA-gl2}
Assume~\ref{N5}. Let $f\in C_c^\infty(\GS)_-$ and $h\in G_1(k_S)$.
\begin{enumerate}[label=(\alph*)]
\item $I_A(x,f^h)=I_A(\mathrm{Ad}(h)x,f)$ for $x\in\AS$ regular or
  central, whenever $\mathrm{Ad}(h)x\in\AS$.
\item For $h=\delta_a$ and $x$ over $a(t)$:
  $I_A(x,f^{\delta_a})=(a,t)_S\,I_A(x,f)$, with $(a,t)_S=\prod_{v\in S_k}
  (a_v,t_v)_v$; for a lift $\tilde z$ of $-1$,
  $I_A(\tilde z,f^{\delta_a})=(a,-1)_S\,I_A(\tilde z,f)$.
\item $\operatorname{tr}I(\pi,f^{\delta_a})=\operatorname{tr}I(\pi\kappa_a,
  f)$, and $\operatorname{tr}I(\pi,f^{[h]})=\operatorname{tr}I(\pi,f)$:
  the characters of the principal series are stable. (This part does not
  use~\ref{N5}. The function $f^{[h]}$ need not be smooth; the trace is
  then defined by the integral of orbital integrals in the proof.)
\end{enumerate}
\end{proposition}

\begin{proof}
(a) For $h\in SL(2,k_S)$ this is Theorem~\ref{thm:IA}(b). Let
$h=\delta=\delta_a$. By Lemma~\ref{lem:variance-gl2}(c) and Fourier
inversion on $\AS$, substituting $\pi\mapsto\pi'$ and using
$\pi(\tilde a)=\pi'(\alpha\tilde a)$,
\[
\Phi_A(\tilde a,f^\delta)=\Phi_A(\alpha\tilde a,f)+f_{B,\delta}(\alpha
\tilde a)+f_{B,\delta}(w\,\alpha\tilde a\,w^{-1}).
\]
Subtract from Lemma~\ref{lem:variance-gl2}(b). For central $x$ use the
limit of Lemma~\ref{lem:JA-central}(b), as for
Lemma~\ref{lem:variance-local}(c). (b) follows from~(a),
Lemma~\ref{lem:Ad}(d) and the fact that $I_A(\cdot,f)$ is anti-genuine.
(c) The first statement is Lemma~\ref{lem:variance-gl2}(c). By the proof
of Theorem~\ref{thm:IA}(a), $\operatorname{tr}I(\pi,f)=\int_{\AS}\pi(\tilde
a)J_G(\tilde a,f)\,d\tilde a$, an integral of hyperbolic orbital
integrals, and $J_G(\tilde a,f^{[h]})=J_G(h\star\tilde a,f)=J_G(\tilde
a,f)$ because $h\star\tilde a$ is conjugate to $\tilde a$.
\end{proof}

\begin{remarkrev}\label{rem:N5}
\emph{The condition~\ref{N5}.} By the proof of
Lemma~\ref{lem:variance-gl2}(c), \ref{N5} holds at $v$ if and only if
\[
R_v(\pi_v\kappa_a,\lambda)=C\,|a|_v^{-\lambda}\,A(-\lambda)\,R_v(\pi_v,
\lambda)\,A(\lambda)^{-1}
\]
with a constant $C$; the factor $|a|_v^{-\lambda}$ removes the scalar
$|a|_v^{\lambda}$ contained in $A(-\lambda)A(\lambda)^{-1}$. It is natural:
the inducing data $\pi_v$ and $\pi_v\kappa_a$ have the same $\chi_v^2$,
hence the same parameter on the side of $SO(3)$
(Section~\ref{sec:stable-conj}, \emph{The dual group}), and factors defined
through that parameter, as in Lemma~\ref{lem:explicit-r}, satisfy it.
Lemma~\ref{lem:variance-gl2}(d) removes one possible conflict
with~\ref{N2}; it is a necessary compatibility and not an existence
statement. We have not proved that factors with
\ref{N1}--\ref{N5} exist at the places where no explicit factor is
available (the places of the sets $S_{\mathrm{bad}}(\mu)$ of
Section~\ref{sec:normalized}), nor checked
whether the factors of~\cite{LiII} have property~\ref{N5}. One cannot
simply transport $r_v$ along the orbits of $k_v^\times/k_v^{\times2}$ on
the genuine characters: when $\pi_v^w=\pi_v\kappa_{a_0}$ with $a_0$ not a
square, that is, when $\chi_v^2=(a_0,\cdot)_v\ne1$, the relation~\ref{N2}
between $r_v(\pi_v,\cdot)$ and $r_v(\pi_v^w,\cdot)$ constrains the factor
of $\pi_v$ itself. So in this section \ref{N5} is a hypothesis at the
places where the factors are not explicit. What would remove it is the
following statement, which we formulate as a problem: \emph{for every
local field $F$ of characteristic $0$, including the dyadic ones, and
every $\psi$, there are functions $F_v(\chi^2,\lambda)$ and constants
$C_v(\chi)$ such that the factors $r_v(\chi_{\psi}\chi,\lambda)=C_v(\chi)
F_v(\chi^2,\lambda)$ satisfy \ref{N1}--\ref{N4}.} The literature we know
does not contain it. The local coefficients of~\cite{Sz09} are computed in
all residue characteristics, but we found no statement there that gives
\ref{N1}--\ref{N4} for factors of this form. The Plancherel measures
of~\cite{GS16} are computed under the assumption that the residue
characteristic is prime to the degree of the cover, which excludes the
dyadic places for the double cover. The paper~\cite{Sz16} drops that
assumption when $4\nmid n$; we have not worked out what it gives here.
Without~\ref{N5}, every formula
below that is stated under~\ref{N5} holds with the additional term
coming from $d_a$ in Lemma~\ref{lem:variance-gl2}(c). The distributions
$I_A(x,f)$ depend on the choice of the factors in any case
(Remark~\ref{rem:bad-places}).
\end{remarkrev}

So $I_A(x,\cdot)$ transforms under $G_1(k_S)$ by the character
$h\mapsto(\det h,t)_S=\eta_x(\det h)$, as the orbital integral
$J_G(x,\cdot)$ does. This does \emph{not} say that $I_A(x,\cdot)$ is
stable. Stability is invariance under $f\mapsto f^{[h]}$, where the sign
$\eta_y(\det h)$ varies with the point $y$ of the support of $f$, and
$I_A(x,f)$ depends on $f$ away from the class of $x$. At a real place, for instance, $\Phi_A(x,f)$ is built from the operators
$R_v^{-1}R_v'$, whose poles are the points $\lambda=a_i,b_i$ at which
$I(\pi_\lambda)$ has a single discrete series representation as a
subquotient (Example~\ref{ex:real-R}), while a stable character involves
two discrete series representations with different central
characters~\cite{Ad}. The stabilization of these terms in~\cite{Li21}
expresses them through stable distributions on the endoscopic groups for
both endoscopic data; it requires the weighted fundamental
lemma~\cite{Li12} and is not attempted here.

\subsection{The central terms and the unipotent coefficients}
\label{sec:unipotent}

The class $\mathcal O_z$ of $z=\pm1$ contributes $J_{\mathrm{unip}}(f_z)$,
which Proposition~\ref{prop:central} expresses through the zeta integral
$Z(f,s)$. Its second form in~\eqref{eq:Zf} is a sum over the rational
unipotent classes $\xi\in k^\times/k^{\times2}$; grouped according to the
image of $\xi$ in $k_S^\times/k_S^{\times2}$ it diverges at $s=1$, class by
class; this is the difficulty noted in
Remark~\ref{rem:sec4-status}. The way out, as in~\cite{LL}, is to expand
in the characters of the compact group $\A_k^\times/k^\times\A_k^{\times2}$
instead. This computes the coefficients $a^G(S,u_\xi)$ of
Theorem~\ref{thm:fine}.

Let $S$ be as in Theorem~\ref{thm:tf}: a finite set of places of $\Q$
containing $\infty$, $2$ and the places below which $\psi$ is ramified.
In this subsection we assume in addition that the ring $\OO_S$ of
elements of $k$ that are integral outside $S_k$ is a principal ideal
domain, which holds once $S$ is large enough (and always for $k=\Q$).
Let $\mathfrak X_S$ be the group of quadratic characters $\kappa$ of
$k^\times\backslash\A_k^\times$ that are unramified outside $S_k$,
$\kappa_S$ the restriction of $\kappa$ to $k_S^\times$, and
$L^S(s,\kappa)=\prod_{v\notin S_k}(1-\kappa_v(\varpi_v)q_v^{-s})^{-1}$. For
$\kappa=1$ this is $\zeta_k^S(s)$; write
\[
\zeta_k^S(s)=\frac{\rho_S}{s-1}+c_S+O(s-1).
\]
By Lemma~\ref{lem:w-action}(d) and its proof, $\mathfrak X_S=\{\kappa_b
\mid b\in\OO_S^\times/\OO_S^{\times2}\}$ with $\kappa_b(t)=\prod_v(b,t_v)_v$;
it has $2^{|S_k|}$ elements, by the unit theorem. For $f_S=\bigotimes_{v
\in S_k}f_v$, $z\in\{\pm1\}$ and $\Phi_v$ as in
Lemma~\ref{lem:JA-central}, put
\[
Z_v(\Phi_v,\kappa_v,s)=\int_{k_v^\times}\Phi_v(y)\kappa_v(y)|y|_v^s\,d^\times y,
\qquad Z_S(\Phi_S,\kappa_S,s)=\prod_{v\in S_k}Z_v(\Phi_v,\kappa_v,s),
\]
Tate's local zeta integrals~\cite{Ta}, holomorphic for $\mathrm{Re}\,s>0$;
by linearity $Z_S$ is defined for all $f_S\in\mathcal H(\GS)$. For
$\xi\in k_S^\times$ let
\begin{equation}\label{eq:JG-unip}
J_G(zu_\xi,f_S)=|\xi|_S\int_{\{\pm1\}N(\Q_S)\backslash G(\Q_S)}
f_S\bigl(x^{-1}z\,n(\xi)\,x\bigr)\,dx
=\prod_{v\in S_k}\frac{|\xi|_v}2\int_{k_v^\times}\Phi_v(\xi t^2)|t|_v^2\,
d^\times t,
\end{equation}
the second expression by Lemma~\ref{lem:JA-central}(c). The factor
$|\xi|_S$ makes $J_G(zu_\xi,f_S)$ depend only on the class of $\xi$ in
$k_S^\times/k_S^{\times2}$; it is the orbital integral for the measure on
the centralizer of $n(\xi)$ transported from that of $n(1)$.

\begin{proposition}\label{prop:unipotent}
Let $S$ be as above (in particular the ring $\OO_S$ of $S_k$-integers is
principal), $f_S\in\mathcal H(\GS)$, $f=f_S\otimes e^S$ and
$z\in\{\pm1\}$.
\begin{enumerate}[label=(\alph*)]
\item For $\mathrm{Re}\,s>1$,
  \[
  Z(f_z,s)=\frac12\sum_{\kappa\in\mathfrak X_S}L^S(s,\kappa)\,
  Z_S(\Phi_S,\kappa_S,s).
  \]
\item For every character $\kappa_S$ of $k_S^\times/k_S^{\times2}$,
  \[
  Z_S(\Phi_S,\kappa_S,1)=\sum_{\xi\in k_S^\times/k_S^{\times2}}\kappa_S(\xi)\,
  J_G(zu_\xi,f_S)=:J^\kappa_{\mathrm{unip}}(z,f_S),
  \]
  and $\frac{d}{ds}Z_S(\Phi_S,1,s)\big|_{s=1}=\rho_S^{-1}c_k\,J_A(z,f_S)$.
\item Consequently
  \begin{multline*}
  J_{\mathcal O_z}(f)=\mathrm{vol}\bigl(G(\Q)\backslash\GA\bigr)f_S(z)
  +\frac{c_k}2\,J_A(z,f_S)\\
  +\frac{c_S}2\,J^1_{\mathrm{unip}}(z,f_S)
  +\frac12\sum_{\kappa\in\mathfrak X_S,\ \kappa\ne1}L^S(1,\kappa)\,
  J^\kappa_{\mathrm{unip}}(z,f_S).
  \end{multline*}
  This is Theorem~\ref{thm:fine}(a), with
  \[
  a^G(S,u_\xi)=\frac12\Bigl(c_S+\sum_{\kappa\in\mathfrak X_S,\ \kappa\ne1}
  \kappa_S(\xi)\,L^S(1,\kappa)\Bigr),\qquad\xi\in k_S^\times/k_S^{\times2}.
  \]
\item For $a\in k_S^\times$,
  \[
  J_G(zu_\xi,f_S^{\delta_a})=(a,z)_S\,J_G(zu_{a\xi},f_S),\qquad
  J^\kappa_{\mathrm{unip}}(z,f_S^{\delta_a})=\kappa_S(a)\,(a,z)_S\,
  J^\kappa_{\mathrm{unip}}(z,f_S).
  \]
\end{enumerate}
\end{proposition}

\begin{proof}
We may replace $f$ by $f_z$ and take $z=1$; then $\Phi_f=\Phi_S\otimes
\bigotimes_{v\notin S_k}\mathbf 1_{\OO_v}$.

\emph{A local identity.} For $h$ integrable on $k_v^\times$ and
$\xi\in k_v^\times$,
\begin{equation}\label{eq:squares}
\int_{k_v^\times}h(\xi t^2)\,d^\times t=\frac12\sum_{\kappa_v}\kappa_v(\xi)
\int_{k_v^\times}h(y)\kappa_v(y)\,d^\times y,
\end{equation}
the sum over the characters of $k_v^\times/k_v^{\times2}$. Indeed
$t\mapsto t^2$ is two-to-one onto the open subgroup $k_v^{\times2}$ and
multiplies $d^\times t$ by $|2|_v$, so the left side is
$2|2|_v^{-1}\int_{\xi k_v^{\times2}}h\,d^\times y$; and the characteristic
function of $\xi k_v^{\times2}$ is $[k_v^\times:k_v^{\times2}]^{-1}
\sum_{\kappa_v}\kappa_v(\xi)\kappa_v$. The index is $4|2|_v^{-1}$ at every
place ($4|2|_v^{-1}$ for $v$ finite, $2$ for $v$ real, $1$ for $v$
complex, with the normalized absolute values).

(a) Let $\Gamma=\OO_S^\times$ and $Y_1=k_S^\times\times\prod_{v\notin S_k}
\OO_v^\times$. Since $\OO_S$ is principal, $\A_k^\times=k^\times Y_1$ and
$k^\times\cap Y_1=\Gamma$, so $\A_k^\times$ is the disjoint union of the
$cY_1$, $c\in k^\times/\Gamma$. For $\eta\in k^\times$ let
$J(\eta)=\int_{Y_1}\Phi_f(\eta t^2)|t|^{2s}d^\times t$. Then the integral
in the second form of~\eqref{eq:Zf} is $\sum_{c\in k^\times/\Gamma}J(\xi
c^2)$, and $(\xi,c)\mapsto\xi c^2$ is a bijection from
$(k^\times/k^{\times2})\times(k^\times/\Gamma)$ onto $k^\times/\Gamma^2$
(because $-1\in\Gamma$). So $Z(f,s)=\frac12\sum_{\eta\in k^\times/\Gamma^2}
J(\eta)$. Now $J(\eta)=0$ unless $\eta\in\OO_S$, and
by~\eqref{eq:squares} at each $v\in S_k$, applied to $h=\Phi_v|\cdot|_v^s$,
\[
J(\eta)=|\eta|_S^{-s}\,2^{-|S_k|}\sum_{\kappa_S}\kappa_S(\eta)\,
Z_S(\Phi_S,\kappa_S,s)\qquad(\eta\in\OO_S\smallsetminus\{0\}).
\]
Summing over $\eta$ modulo $\Gamma^2$: the sum of $\kappa_S$ over
$\Gamma/\Gamma^2$ is $2^{|S_k|}$ if $\kappa_S$ is trivial on $\Gamma$ and
$0$ otherwise. The characters $\kappa_S$ trivial on $\Gamma$ are the
restrictions of the $\kappa\in\mathfrak X_S$, and for such $\kappa$ and
$\eta\in\OO_S$, $\kappa_S(\eta)=\prod_{v\notin S_k}\kappa_v(\varpi_v)^{
\ord_v\eta}$ and $|\eta|_S^{-1}=\prod_{v\notin S_k}q_v^{\ord_v\eta}$, by the
product formula. The sum over $\eta\in(\OO_S\smallsetminus\{0\})/\Gamma$
runs over the nonzero ideals of $\OO_S$ and gives $L^S(s,\kappa)$.

(b) By~\eqref{eq:squares} with $h=\Phi_v|\cdot|_v$ and~\eqref{eq:JG-unip},
$J_{G,v}(u_\xi,f_v)=\frac14\sum_{\kappa_v}\kappa_v(\xi)Z_v(\Phi_v,\kappa_v,
1)$; invert on $k_v^\times/k_v^{\times2}$ and multiply over $v\in S_k$. The
constant is $\prod_{v\in S_k}4\,[k_v^\times:k_v^{\times2}]^{-1}=\prod_{v\in
S_k}|2|_v=1$, because $S_k$ contains all places at which $|2|_v\ne1$. For
the derivative, let $d^\times t=c'_v\,dt/|t|_v$ at $v\in S_k$, $dt$ the
self-dual measure, and $C'=\prod_{v\in S_k}c'_v$. Then
$Z_S(\Phi_S,1,1)=C'\prod_v\hat\Phi_v(0)$ and
\[
\frac{d}{ds}Z_S(\Phi_S,1,s)\Big|_{s=1}=C'\sum_{v\in S_k}\Bigl(\int_{k_v}
\Phi_v(y)\log|y|_v\,dy\Bigr)\prod_{w\ne v}\hat\Phi_w(0)=C'\,J_A(1,f_S),
\]
by Lemma~\ref{lem:JA-central}(b), since $\sum_{v\in S_k}\log|2|_v=0$.
Comparing the residue of~(a) at $s=1$, which is $\frac12\rho_SC'
\prod_v\hat\Phi_v(0)$, with the residue $\frac12c_k\hat\Phi_f(0)$ of
Proposition~\ref{prop:central} gives $\rho_SC'=c_k$.

(c) By~(a), the finite part of $Z(f,s)$ at $s=1$ is $\frac12\bigl(c_S
Z_S(\Phi_S,1,1)+\rho_S\frac{d}{ds}Z_S(\Phi_S,1,s)|_{s=1}\bigr)+\frac12
\sum_{\kappa\ne1}L^S(1,\kappa)Z_S(\Phi_S,\kappa_S,1)$; for $\kappa\ne1$,
$L^S(s,\kappa)$ is holomorphic at $s=1$. Insert~(b) and
Proposition~\ref{prop:central}. The coefficient of $J_G(zu_\xi,f_S)$ is
read off from the definition of $J^\kappa_{\mathrm{unip}}$.

(d) $x\mapsto\delta_ax\delta_a^{-1}$ preserves the Haar measure of
$G(\Q_S)$ and multiplies that of $N(\Q_S)$ by $|a|_S$; so
$\int h(\delta_ax\delta_a^{-1})\,dx=|a|_S\int h(x)\,dx$ for the invariant
measure on the quotient by $\{\pm1\}N(\Q_S)$. With Lemma~\ref{lem:Ad}(d),
this gives $\int f_S^{\delta_a}(
x^{-1}zn(\xi)x)\,dx=(a,z)_S\,|a|_S\int f_S(x^{-1}zn(a\xi)x)\,dx$. Multiply
by $|\xi|_S$, and sum against $\kappa_S(\xi)$.
\end{proof}

Part~(c) proves the fine expansion at the central classes directly from
the coarse one, for the sets $S$ considered, and so removes the appeal
to~\cite{A12,LiI} in Theorem~\ref{thm:fine}(a). The terms with
$\kappa\ne1$ are the analog of the terms of~\cite[\S5]{LL} that involve
$L(1,\kappa)$; in the 1995 notes they were written as a sum over the
rational unipotent classes, which diverges
(Appendix~\ref{sec:appendix-1995}).

\begin{example}\label{ex:unip-coeff}
Let $k=\Q$ and $S=\{\infty,2\}$. Then $\OO_S=\Z[\frac12]$,
$\mathfrak X_S=\{\kappa_b\mid b=\pm1,\pm2\}$, and
$\zeta^S(s)=(1-2^{-s})\zeta(s)$, so $\rho_S=\frac12$ and
$c_S=\frac12(\gamma_E+\log2)$, with Euler's constant $\gamma_E$. The
characters $\kappa_{-1}$, $\kappa_2$, $\kappa_{-2}$ are the Dirichlet
characters of conductor $4$, $8$, $8$, and
\[
L^S(1,\kappa_{-1})=\frac\pi4,\qquad L^S(1,\kappa_2)=\frac{\log(1+\sqrt2)}
{\sqrt2},\qquad L^S(1,\kappa_{-2})=\frac\pi{2\sqrt2}.
\]
For the $16$ classes $\xi\in\Q_S^\times/\Q_S^{\times2}$,
\[
a^G(S,u_\xi)=\frac12\Bigl(\frac{\gamma_E+\log2}2+\sum_{b\in\{-1,2,-2\}}
(b,\xi_\infty)_\infty(b,\xi_2)_2\,L^S(1,\kappa_b)\Bigr);
\]
for instance $a^G(S,u_\xi)\approx1.57726$ for $\xi=1$, and
$a^G(S,u_\xi)\approx-0.31886$ for the class with $\xi_\infty<0$ and $\xi_2=1$.
The coefficients are not all positive, and they depend on $\xi$ only through
the pairings with $b\in\{\pm1,\pm2\}$, the same parameters as those of the
ramified classes $\chi_{\psi_b}$ in Example~\ref{ex:ramified-Q}. In
particular $a^G(S,u_{\xi\epsilon})=a^G(S,u_\xi)$ for $\epsilon\in\OO_S^\times$,
since the $\kappa\in\mathfrak X_S$ are trivial on $\OO_S^\times$: the
classes of the rational elements $\pm1$, $\pm2$ all have the coefficient of
$\xi=1$, and the $16$ classes carry four distinct values (the invariance
holds for every $k$ and $S$ as above). With $\epsilon_b=(b,\xi_\infty)_\infty
(b,\xi_2)_2$ they are
\begin{center}
\begin{tabular}{ccc}
$(\epsilon_{-1},\epsilon_2,\epsilon_{-2})$ & a class $\xi=(\xi_\infty,\xi_2)$ &
$a^G(S,u_\xi)$\\\hline
$(+,+,+)$ & $(1,1)$ & $\phantom{-}1.57726$\\
$(-,-,+)$ & $(1,3)$ & $\phantom{-}0.16864$\\
$(+,-,-)$ & $(1,5)$ & $-0.15668$\\
$(-,+,-)$ & $(-1,1)$ & $-0.31886$
\end{tabular}
\end{center}
Each value is taken on four classes, and the four values add up to
$2c_S=\gamma_E+\log2$. The
script checks~(a) numerically at $s=3$ for a test function with
$\Phi_\infty$ not even and $\Phi_2$ the characteristic function of
$6+16\Z_2$.
\end{example}

By part~(d), $J^\kappa_{\mathrm{unip}}(1,\cdot)$ transforms under
$G_1(k_S)$ by $\kappa_S(\det h)$. Near the identity $\eta_x=1$ (because
$D^+$ is close to $4$), so there the twisted and the untwisted actions
agree: $J^1_{\mathrm{unip}}(1,\cdot)$, the sum over all unipotent classes,
is the stable unipotent distribution, and the
$J^\kappa_{\mathrm{unip}}(1,\cdot)$ with $\kappa\ne1$ are the unstable
ones. For $SL(2)$ the latter are the limits, as $\gamma\to1$ in the torus
$T_E$ with $\kappa_E=\kappa$, of the $\kappa$-orbital integrals, up to the
factor $L(1,\kappa)$~\cite[\S\S2,~5]{LL}. On the cover the same should
hold near the identity, where the cover is canonically split on the
topologically unipotent set, and by Lemma~\ref{lem:kappa-orbital}(c) the
$\kappa$-orbital integrals near $1$ are the stable orbital integrals of
$f_-$ near $-1$. So the terms $L^S(1,\kappa)J^\kappa_{\mathrm{unip}}(1,
f_S)$, $\kappa\ne1$, should be counted with $S(f_-)$, and the
corresponding terms at $z=-1$ with $S(f)$. We have not proved these limit
formulas for the cover (Remark~\ref{rem:sec7-status}).

\subsection{The spectral side}\label{sec:spectral-gl2}

On the spectral side the group $G_1(k)$ acts on the space of automorphic
forms, and the pieces of the spectral decomposition are permuted. We
record this, and then compute the one family of discrete terms whose
stable structure is visible without endoscopy: the theta functions.

For $h\in G_1(k)$ and a function $\varphi$ on $G(\Q)\backslash\GA$ put
$(U_h\varphi)(x)=\varphi(\mathrm{Ad}(h)^{-1}x)$; this is well defined by
Lemma~\ref{lem:Ad}(f). For $a\in k^\times$ let $\kappa_a$ also denote the
character $\tilde a\mapsto\prod_v(a,t_v)_v$ of $\overline{A(\A)}$.

\begin{lemma}\label{lem:spectrum-gl2}
Let $h\in G_1(k)$ and $a=\det h$.
\begin{enumerate}[label=(\alph*)]
\item $U_h$ is a unitary operator on $L^2_{\mathrm{gen}}$ with
  $U_hR(y)U_h^{-1}=R(\mathrm{Ad}(h)y)$ and $U_hR(f)U_h^{-1}=R(f^{h^{-1}})$.
  It preserves $L^2_{\mathrm{cusp}}$, $L^2_{\mathrm{res}}$ and
  $L^2_{\mathrm{cont}}$.
\item $U_h$ maps $L^2_\chi$ onto $L^2_{\chi^h}$, where $(G,\pi)^h=
  (G,\pi\circ\mathrm{Ad}(h)^{-1})$ and $(A,\mu)^h=(A,\mu\kappa_a)$. In
  particular $m(\pi\circ\mathrm{Ad}(h))=m(\pi)$, $\chi_{\psi_b}\kappa_a=
  \chi_{\psi_{ab}}$, and $U_hV_b=V_{ab}$.
\item Let $f=f_S\otimes e^S$ and $h_v\in Z(k_v)GL(2,\OO_v)$ for all
  $v\notin S_k$. Then $\operatorname{tr}(R(f^h)|_{L^2_{\chi,\mathrm{disc}}})=
  \operatorname{tr}(R(f)|_{L^2_{\chi^h,\mathrm{disc}}})$ for every $\chi$, and
  if~\ref{N5} holds, $I_\chi(f^h)=I_{\chi^h}(f)$. Hence
  $\sum_\chi I_\chi(f^h)=\sum_\chi I_\chi(f)$.
\end{enumerate}
\end{lemma}

\begin{proof}
(a) $\mathrm{Ad}(h)$ preserves $G(\Q)$ and the Tamagawa measure, which
gives unitarity and the two formulas. It preserves $N(\A)$ and its Haar
measure up to the factor $|a|_{\A_k}=1$, so $(U_h\varphi)_B=U_h(\varphi_B)$
for the constant term along $B$; hence cusp forms go to cusp forms, and
the orthogonal complement, the discrete and the continuous parts of it,
are preserved. (b) By Lemma~\ref{lem:Ad}(d), $\mathrm{Ad}(h)^{-1}$ acts
on $\overline{A(\A)}$ by the character $\kappa_a$ (for $h=\delta_a$, and
conjugation by $SL(2,k)$ does not change the class of a cuspidal datum).
So $U_h$ maps $\mathcal H_B(\mu)$ to $\mathcal H_B(\mu\kappa_a)$ and
Eisenstein series, their wave packets and their residues to the
corresponding objects for $\mu\kappa_a$. The formula for
$\chi_{\psi_b}\kappa_a$ is Lemma~\ref{lem:weil}(d). (c) By
Lemma~\ref{lem:Ad}(e), $f^h=f_S^{h_S}\otimes e^S$. The first statement
follows from~(a) and~(b): $\operatorname{tr}(R(f^h)|_{L})=
\operatorname{tr}(R(f)|_{U_hL})$ for an invariant subspace $L$. This gives
$I_\chi(f^h)=I_{\chi^h}(f)$ for $\chi\in\mathcal X_G$ and for the residual
terms. For $\mu\in\Pi_{\mathrm{aut}}(X_A)$, $\operatorname{tr}I(\mu_{it},
f^h)=\operatorname{tr}I((\mu\kappa_a)_{it},f)$ by
Lemma~\ref{lem:variance-gl2}(c), and $\frac{m'}m(\mu\kappa_a,\lambda)=
\frac{m'}m(\mu,\lambda)$: in~\eqref{eq:global-m} the partial $L$-functions
depend only on $\chi^2$, and the local factors change by constants,
by~\ref{N5}. So the terms of Theorem~\ref{thm:itf}(e),(f) for $\chi$ and
$f^h$ are those for $\chi^h$ and $f$.
\end{proof}

The geometric counterpart of~(c) is Theorem~\ref{thm:stabler}(c) below.

Now fix $b\in k^\times$ and put $\psi'=\psi_{-b}$, as in
Theorem~\ref{thm:residual}(c). For a pattern of signs
$\varepsilon=(\varepsilon_v)_v$, with $\varepsilon_v=+$ for almost all
$v$, let $r^\varepsilon_{\psi'}=\bigotimes'_vr^{\varepsilon_v}_{\psi'_v}$, an
irreducible genuine unitary representation of $\GA$ on the functions in
$\mathcal S(\A_k)$ of parity $\varepsilon_v$ at each $v$. Let
\[
\Theta_b=\bigl\{\theta^{\psi'}_\phi\ \big|\ \phi\in\mathcal S(\A_k)\bigr\}^{-}
\subset L^2_{\mathrm{gen}},
\]
the closure of the space of all theta functions~\eqref{eq:theta}, even or
not; and for $f=\bigotimes f_v$ let
\[
\Theta^{\mathrm{st}}_b(f)=\prod_v\bigl(\operatorname{tr}r^+_{\psi'_v}(f_v)-
\operatorname{tr}r^-_{\psi'_v}(f_v)\bigr)=\prod_v\int f_v(x)\,
\frac{\Delta_{\psi'_v}(x)}{|D^+(x)|_v^{1/2}}\,dx,
\]
the second expression by Proposition~\ref{prop:weil-char}. Almost all
factors are $1$. Both $\Theta_b$ and $\Theta^{\mathrm{st}}_b$ depend only
on the class of $b$ in $k^\times/k^{\times2}$: the substitution
$\phi(x)\mapsto\phi(cx)$, $c\in k^\times$, carries $\omega_{\psi'}$ to
$\omega_{\psi'_{c^{-2}}}$, preserves parities, and does not change
$\sum_{\xi\in k}$.

\begin{proposition}\label{prop:theta-packet}
\begin{enumerate}[label=(\alph*)]
\item $\phi\mapsto\theta^{\psi'}_\phi$ is zero on $r^\varepsilon_{\psi'}$ if
  $\prod_v\varepsilon_v=-1$ and injective if $\prod_v\varepsilon_v=+1$. So
  \[
  \Theta_b\cong\bigoplus_{\prod_v\varepsilon_v=+1}r^\varepsilon_{\psi'},
  \]
  a subspace of $L^2_{\mathrm{disc}}$; the summand with all
  $\varepsilon_v=+$ is $V_b$, and the others are cuspidal.
\item For $f=\bigotimes_vf_v\in\mathcal H(\GA)$,
  \[
  \operatorname{tr}\bigl(R(f)\big|_{\Theta_b}\bigr)=\tfrac12\,
  \Theta^{\mathrm{st}}_b(f)+\tfrac12\,\Theta^{\mathrm{st}}_b(f_-).
  \]
\item $\Theta^{\mathrm{st}}_b$ is a stable distribution:
  $\Theta^{\mathrm{st}}_b(f^{[h]})=\Theta^{\mathrm{st}}_b(f)$ for
  $h\in G_1(\A_k)$. Moreover $U_h\Theta_b=\Theta_{ab}$ and
  $\Theta^{\mathrm{st}}_b(f^h)=\Theta^{\mathrm{st}}_{ab}(f)$ for
  $h\in G_1(k)$, $a=\det h$.
\end{enumerate}
\end{proposition}

\begin{proof}
(a) $\theta_{\omega(y)\phi}(x)=\theta_\phi(xy)$, so $\phi\mapsto
\theta_\phi$ intertwines $\omega_{\psi'}$ with $R$. The parity operator
$P=\bigotimes_vP_v$ satisfies $\theta_{P\phi}=\theta_\phi$ (replace $\xi$
by $-\xi$ in~\eqref{eq:theta}; $P$ commutes with $\omega$). If
$\prod\varepsilon_v=-1$, then $P\phi=-\phi$ on $r^\varepsilon$ and
$\theta_\phi=0$. If $\prod\varepsilon_v=+1$, let $S_0$ be a finite set of places
containing the archimedean ones and those with $\varepsilon_v=-$, and
$\xi_0=1$. Put $\phi_v=\mathbf 1_{\OO_v}$ for $v\notin S_0$, and
$\phi_v=h_v+\varepsilon_vP_vh_v$ for $v\in S_0$, where $h_v$ is supported
in a neighborhood $U_v$ of $\xi_0$ with $U_v\cap(-U_v)=\emptyset$ and
$h_v(\xi_0)=1$. Then $\phi=\bigotimes\phi_v\in r^\varepsilon$. Since $k$ is
discrete in $\A_k$, for $U_v$ small enough the only $\xi\in k$ that are
integral outside $S_0$ and lie in $U_v\cup(-U_v)$ for every $v\in S_0$ are
$\pm\xi_0$, and $\theta_\phi(1)=\phi(\xi_0)+\phi(-\xi_0)=1+\prod_{v\in S_0}
\varepsilon_v=2$. So the map is not zero on the irreducible
$r^\varepsilon$, hence injective. The constant term of $\theta_\phi$ is
$x\mapsto(\omega(x)\phi)(0)$ (Section~\ref{sec:eisenstein}), which
vanishes when some $\phi_v$ is odd; so those $\theta_\phi$ are cuspidal,
and the even ones span $V_b$ by Theorem~\ref{thm:residual}(c). That theta
functions are square integrable, and that $\phi\mapsto\theta_\phi$ is
continuous for the $L^2$ norms on each $r^\varepsilon$, is
standard~\cite{We64,Ge2,GPS80}: on a Siegel set $\theta_\phi(n\tilde a
\tilde k)=O(e^{H(\tilde a)/2})$. The representations $r^\varepsilon$ are
pairwise inequivalent, because $r^+_{\psi'_v}$ and $r^-_{\psi'_v}$ have
different central characters; so the sum is direct and orthogonal.

(b) By~(a), $\operatorname{tr}(R(f)|_{\Theta_b})=\sum_{\prod\varepsilon_v=1}
\prod_v\operatorname{tr}r^{\varepsilon_v}_{\psi'_v}(f_v)$, a finite sum
because $\operatorname{tr}r^-_{\psi'_v}(e_v)=0$ at the unramified places.
With $t^\pm_v=\operatorname{tr}r^\pm_{\psi'_v}(f_v)$ this is
$\frac12\prod_v(t^+_v-t^-_v)+\frac12\prod_v(t^+_v+t^-_v)$. The first
product is $\Theta^{\mathrm{st}}_b(f)$, and the second is
$\operatorname{tr}\omega_{\psi'}(f)$. By Lemma~\ref{lem:weil-rep},
$\omega_{\psi'_v}(-1,1)=\gamma'_v(1)^2P_v$, and $\prod_v\gamma'_v(1)=1$
(Lemma~\ref{lem:weil}(b)); so $\omega_{\psi'}(s(-1))=P$. Therefore
$\operatorname{tr}\omega_{\psi'}(f)=\operatorname{tr}\bigl(\omega_{\psi'}(
s(-1))\,\omega_{\psi'}(f)\,P\bigr)=\operatorname{tr}\bigl(\omega_{\psi'}
(f_-)P\bigr)=\Theta^{\mathrm{st}}_b(f_-)$, since $\omega(z)\omega(f)=
\omega(f_{z^{-1}})$.

(c) Stability is Proposition~\ref{prop:weil-char} at each place. By
Lemma~\ref{lem:weil-kernel}, $\omega_{\psi'}\circ\mathrm{Ad}(\delta_a)=
\omega_{\psi'_a}$ and $\psi'_a=\psi_{-ab}$; for $a\in k^\times$ the
identity map of $\mathcal S(\A_k)$ therefore carries $\theta^{\psi'}_\phi
\circ\mathrm{Ad}(\delta_a)$ to $\theta^{\psi'_a}_\phi$, which gives
both statements (replace $a$ by $a^{-1}$, which is in the same square
class).
\end{proof}

Part~(b) has the same two-term form as Theorem~\ref{thm:elliptic-stable}: a
stable distribution evaluated at $f$ and at $f_-$, each with coefficient
$\frac12$. The multiplicity of $r^\varepsilon_{\psi'}$ in $\Theta_b$ is
$\frac12(1+\prod_v\varepsilon_v)$. Waldspurger~\cite{Wa80,Wa91} proved
that this is its multiplicity in $L^2_{\mathrm{disc}}$, and that the rest
of the discrete spectrum has a similar structure: it is a
multiplicity-free sum over the cuspidal automorphic representations
$\sigma$ of $PGL(2,\A_k)$, each contributing the representations
$\bigotimes_v\pi_v^{\varepsilon_v}$ of a packet with $\prod_v\varepsilon_v=
\varepsilon(\frac12,\sigma)$, where the local packet is $\{\pi_v^+\}$ or
$\{\pi_v^+,\pi_v^-\}$. Here the packets are those of the theta
correspondence for the additive character $\psi'$: $\pi_v^+$ is the lift
of $\sigma_v$ from $PGL(2,k_v)=SO(2,1)$, and $\pi_v^-$, when it exists, is
the lift of the Jacquet--Langlands transfer of $\sigma_v$ from the
anisotropic $SO(3)$; the labels depend on $\psi'$. See~\cite{GS12} for the
local correspondence and its epsilon dichotomy, \cite{GI} for the double
cover of $Sp(2n)$, and~\cite{Li24}. In these terms $\Theta_b$
is the packet attached to the character $\kappa_b$ of $PGL(2)$. One
difference should be noted. For the theta functions the stable
combination is the difference $\operatorname{tr}r^+-\operatorname{tr}r^-$
(Proposition~\ref{prop:weil-char}). For the packets of cuspidal $\sigma$
one expects it to be the sum $\operatorname{tr}\pi^+_v+\operatorname{tr}
\pi^-_v$, with the difference unstable. We do not know a reference that
states this for all packets in terms of the action $\star$ used here, and
we do not prove it; it is the local character identity of the endoscopic
transfer~\cite{Li19}. Two cases can be checked directly. At a finite
place, the packet of the Steinberg representation is
$\{\mathrm{sp},r^-_{\psi'}\}$~\cite[\S5]{GS12},
where $\mathrm{sp}$ is the other constituent of the principal series
that has $r^+_{\psi'}$ as a quotient (it is the lift of the trivial
representation of the anisotropic $SO(3)$, and $r^-_{\psi'}$ is the lift
of the Steinberg representation); the sum of their characters is the
character of that principal series minus $\Theta^+_{\psi'}-\Theta^-_{\psi'}$,
a difference of two stable characters
(Propositions~\ref{prop:IA-gl2}(c) and~\ref{prop:weil-char}). At a real
place a packet consists of the discrete series representations with
weights $m,m+2,\dots$ and $-m,-m-2,\dots$, $m\in\frac12+\Z$, $m>1$
(see~\cite[\S15]{Ad} for the stable discrete series characters of
$Mp(2,\R)$ and the central characters of the two members); on
$\tilde k_\theta$ the sum of their characters is $-\sin((m-1)\theta)/
\sin\theta$, which is even in $\theta$, and the difference is
$i\cos((m-1)\theta)/\sin\theta$, which is odd, while $\tilde k_\theta$ and
$\tilde k_{-\theta}$ are stably conjugate (Example~\ref{ex:Delta-real}).
In both kinds of packet the two members have opposite central characters,
so the unstable combination for $f$ is, up to a constant, the stable one
for $f_-$; this is how the pattern $S(f)+S(f_-)$ arises on the spectral
side. We do not derive any of this from the trace formula. The
multiplicity statements are quoted from~\cite{Wa80,Wa91} as recounted
in~\cite{GS12}; we have not checked them against the original papers.

By Proposition~\ref{prop:IA-gl2}(c) the terms of
Corollary~\ref{cor:itf-final} that come from the continuous spectrum,
$-\frac14\operatorname{tr}I(\chi_{\psi_b,0},f)$ and $-\frac1{4\pi}\int
\frac{m'}m(\mu,it)\operatorname{tr}I(\mu_{it},f)\,dt$, are stable
distributions. They are unchanged when $f$ is replaced by $f_-$, because
$I(\mu_\lambda,s(-1))=\mu(s(-1))=1$ for $\mu\in\Pi_{\mathrm{aut}}(X_A)$.

\subsection{The stable(r) form}\label{sec:stable-form}

We collect the results. For $f=f_S\otimes e^S$ as in
Theorem~\ref{thm:itf}, with $S$ as in Section~\ref{sec:unipotent}, write
$I(f)=\sum_{\mathcal O}I_{\mathcal O}(f)=\sum_\chi I_\chi(f)$.

Not everything in the theorem depends on~\ref{N5}. The decomposition of
the regular elliptic terms, the unipotent coefficients, the trace on the
theta functions and the stability of the principal series characters are
proved without it. It is used for the equivariance of the normalized
weighted distributions $I_A$ (Proposition~\ref{prop:IA-gl2}(a),~(b)), and
through it in the statement of part~(c) about $I_A$ and in the invariance
of the individual $I_{\mathcal O}$ and $I_\chi$ under $G_1(k)$. Parts~(a)
and~(b) as displayed do not use it.

\begin{theorem}\label{thm:stabler}
Let $S$, $f_S$ and $f=f_S\otimes e^S$ be as above. In part~(c)
assume~\ref{N5}.
\begin{enumerate}[label=(\alph*)]
\item (Geometric side.)
  \begin{multline*}
  I(f)=S_{\mathrm{ell}}(f)+S_{\mathrm{ell}}(f_-)
  +c_k\sum_{\{t,t^{-1}\}}I_A\bigl(a(t)_S,f_S\bigr)\\
  +\sum_{z=\pm1}\Bigl(\mathrm{vol}\bigl(G(\Q)\backslash\GA\bigr)f_S(z)+
  \frac{c_k}2\,I_A(z,f_S)+\frac12\sum_{\kappa\in\mathfrak X_S}
  \ell_S(\kappa)\,J^\kappa_{\mathrm{unip}}(z,f_S)\Bigr),
  \end{multline*}
  where $S_{\mathrm{ell}}$ is the stable distribution of
  Theorem~\ref{thm:elliptic-stable}, $t$ runs over $\OO_S^\times
  \smallsetminus\{\pm1\}$ modulo inversion, $\ell_S(1)=c_S$ and
  $\ell_S(\kappa)=L^S(1,\kappa)$ for $\kappa\ne1$.
\item (Spectral side.) $I(f)$ is the right side of
  Corollary~\ref{cor:itf-final}. The part of the discrete sum carried by
  the theta functions is $\frac12\sum_b\bigl(\Theta^{\mathrm{st}}_b(f)+
  \Theta^{\mathrm{st}}_b(f_-)\bigr)$, the sum over the classes $b\in k^\times/
  k^{\times2}$ with $\Theta_b$ pairwise orthogonal; the terms of the
  continuous spectrum are stable.
\item (Equivariance.) The individually indexed distributions occurring
  in~(a) transform under $G_1(k_S)$, acting by $f_S\mapsto f_S^h$, by the
  following characters of $\det h=a$: $\eta_\gamma$ for $J^{\mathrm{st}}_G(\gamma,\cdot)$ and
  $\eta_\gamma\kappa_E$ for $J^\kappa_G(\gamma,\cdot)$; $(\cdot,t)_S$ for
  $I_A(a(t)_S,\cdot)$; $(\cdot,z)_S$ for $f_S(z)$ and $I_A(z,\cdot)$;
  $\kappa_S\cdot(\cdot,z)_S$ for $J^\kappa_{\mathrm{unip}}(z,\cdot)$. The
  characters differ from one index to the next, so the sums displayed
  in~(a) are in general not eigendistributions; in particular
  $S_{\mathrm{ell}}$ is stable for the twisted action and is not an
  eigendistribution for the ordinary one. Both
  sides are invariant under $h\in G_1(k)\cap\prod_{v\notin S_k}
  Z(k_v)GL(2,\OO_v)$.
\end{enumerate}
\end{theorem}

\begin{proof}
(a) Theorem~\ref{thm:itf}(a)--(c), with Theorem~\ref{thm:elliptic-stable}
for the elliptic classes and Proposition~\ref{prop:unipotent}(c) for the
central ones; in~(c) of Theorem~\ref{thm:itf} the class of $a(t)$
contributes only if $t$ is a unit outside $S_k$. (b)
Proposition~\ref{prop:theta-packet} and the remark after it; the spaces
$\Theta_b$ for inequivalent $b$ are orthogonal because their irreducible
constituents are pairwise inequivalent (the constituents of
$\omega_{\psi'_v}$ and $\omega_{\psi''_v}$ are inequivalent when
$\psi''\ne\psi'_{c^2}$, by their restrictions to $N$). (c)
Lemma~\ref{lem:kappa-orbital}(b), Proposition~\ref{prop:IA-gl2}(b),
Lemma~\ref{lem:Ad}(d) and Proposition~\ref{prop:unipotent}(d). For the
last statement, such an $h$ acts on $f=f_S\otimes e^S$ through its
component in $G_1(k_S)$, by Lemma~\ref{lem:Ad}(e), and
Lemma~\ref{lem:spectrum-gl2}(a) applies.
\end{proof}

\paragraph{What the theorem does not say.} A stable trace formula for the
cover would be an identity
\[
I(f)=S(f)+S(f_-)
\]
with $S$ a stable distribution, and then an identification of $S(f)$
with $\frac12$ of the stable trace formula of $SO(3)$ evaluated at a
transfer of $f$; this is Li's theorem~\cite{Li21}. Parts~(a) and~(b)
exhibit this form for the regular elliptic terms and for the theta
functions. What is missing, in increasing order of difficulty, is:
\begin{enumerate}[label=(\roman*)]
\item the limit formulas that attach the terms $J^\kappa_{\mathrm{unip}}
  (z,\cdot)$, $\kappa\ne1$, to the singular elliptic points of the tori,
  and so distribute them between $S(f)$ and $S(f_-)$
  (Section~\ref{sec:unipotent});
\item the stabilization of the weighted terms $I_A(a(t),\cdot)$ and
  $I_A(z,\cdot)$, which are equivariant by
  Proposition~\ref{prop:IA-gl2} but not stable;
\item the transfer $f\mapsto f'$ to $SO(3)$ and the corresponding
  fundamental lemma~\cite{Li11,Li12};
\item on the spectral side, the description of the cuspidal spectrum by
  packets~\cite{Wa80,Wa91,GI}.
\end{enumerate}
The equivariance in~(c) is the property that the 1995 notes called
``$\kappa$-stable''. It is weaker than stability, and the part of $I(f)$
that is invariant under all of $G_1(k_S)$, which the 1995 notes took as
the stable part, is small: by Lemma~\ref{lem:kappa-orbital}(b) and
Proposition~\ref{prop:IA-gl2}(b), a regular elliptic or hyperbolic
$\gamma$ contributes to it only if, at each place of $S_k$, $D^+(\gamma)$
is a square or $\theta$ times a square; for instance if $\gamma$ is a
square in its torus.

\subsection{Status of the results}\label{sec:status7}

\begin{remarkrev}\label{rem:sec7-status}
\emph{What rests on citation.} Proved here, given the earlier sections:
Lemmas~\ref{lem:Ad} and~\ref{lem:eta} (with the cocycle~\eqref{eq:sigma-gl2}
on $GL(2)$ from~\cite{K2}; part~(c) of Lemma~\ref{lem:Ad} is also checked
directly); Lemma~\ref{lem:Delta}; Lemma~\ref{lem:weil-kernel};
Lemma~\ref{lem:kappa-orbital}; Theorem~\ref{thm:elliptic-stable}, with
the Hasse norm theorem and the index of the global norm group
from~\cite{CF}; Lemma~\ref{lem:variance-gl2};
Proposition~\ref{prop:IA-gl2}, parts~(a) and~(b) under~\ref{N5};
Proposition~\ref{prop:unipotent}; Proposition~\ref{prop:theta-packet}(b),
(c) given~(a); and Proposition~\ref{prop:weil-char} on the regular set,
with Lemma~\shellref. Taken from the literature: at the
archimedean places, the standard fact that an operator with a Schwartz
kernel is of trace class with trace the integral over the diagonal; and
at all places, that the character of the Weil representation is a locally
integrable function on the whole group~\cite{Mak,Th08}, which is used to call $\Theta^+_\psi-\Theta^-_\psi$ a
stable character and in the second formula for $\Theta^{\mathrm{st}}_b$ in
Section~\ref{sec:spectral-gl2}. Formula~\eqref{eq:weil-char-full} agrees
with~\cite[Theorem~1C]{Th08} up to the sign left open there, so its sign
rests on our computation alone; we have not compared it with the formulas
of~\cite{Mak,Th13} in detail. Also from the
literature: the convergence and square-integrability of theta functions
in Proposition~\ref{prop:theta-packet}(a), Tate's local zeta
integrals~\cite{Ta}, and the statements about Waldspurger's packets after
Proposition~\ref{prop:theta-packet}, which are not used.

\emph{What is not done.} The condition~\ref{N5} is a hypothesis at the
places where the normalizing factors are not explicit: we have not shown
that factors satisfying it exist there, nor that those of~\cite{LiII} do
(Remark~\ref{rem:N5}). The statements in the text about the relation
with the endoscopy of~\cite{Li11,Li15,Li21} are comparisons of shape: we
have not identified $J^{\mathrm{st}}_G$ with Li's stable orbital integral
beyond Proposition~\ref{prop:weil-char}, nor $\Delta_\psi$ with his
transfer factor, nor checked the normalizations in
$\iota(\tilde G,H)=\frac12$. The limit formulas of
Section~\ref{sec:unipotent} at the singular elliptic points, the
stabilization of the weighted terms, the transfer to $SO(3)$ and the
description of the cuspidal spectrum are not treated. Sections~\ref{sec:background}--\ref{sec:invariant}
and this section were refereed independently (see the Acknowledgements). An unpublished
stabilization of the trace formula for covers of $SL(2)$ by Hiraga and
Ikeda is mentioned in~\cite{Li15}; see also~\cite{HI13}.
\end{remarkrev}

\begin{remarkrev}\label{rem:verified-sec7}
The script \texttt{checks/sec7\_n2\_check.py} verifies, in exact
arithmetic at $v=\infty,2,3,5,7$ on random elements of $SL(2,\Q)$ and
$GL(2,\Q)$: Lemma~\ref{lem:Ad}(c),(d), both as a homomorphism for
$\beta_v$ and as conjugation in the cover of $GL(2)$;
Lemma~\ref{lem:eta} on the tori $T_\theta(\Q)$ and on the split torus,
and Example~\ref{ex:eta}(d); the conjugation invariance of $\Delta_\psi$,
Lemma~\ref{lem:Delta}(b),(c), and Example~\ref{ex:Delta-real}; the
identity between Weil indices behind~\eqref{eq:weil-kernel}; the
consistency of~\eqref{eq:weil-char} with~\eqref{eq:weil-char-full} under
translation by $(-1,1)$, and the agreement of~\eqref{eq:weil-char-full}
with~\cite[Theorem~1C]{Th08} up to sign;
identity~\eqref{eq:eta-minus}; and Lemma~\ref{lem:variance-gl2}(a).
Numerically: Example~\ref{ex:Delta-finite} for $q=3,5,7,11$;
Proposition~\ref{prop:unipotent}(a) for $k=\Q$, $S=\{\infty,2\}$ at
$s=3$, the local identity~\eqref{eq:squares} at $2$ and $3$, and the
constants of Example~\ref{ex:unip-coeff}, including the invariance of the
coefficients under $\OO_S^\times$; Lemma~\shellref at $2$, $3$ and $5$; the two facts used in the proof of
Proposition~\ref{prop:theta-packet}(b); and the parity in $\theta$ of the
sum and the difference of the two discrete series characters in the
paragraph after that proposition.
\end{remarkrev}


% SPDX-License-Identifier: CC-BY-4.0
% =====================================================================
% Section 8 of D. Joyner and Claude, "A hitch-hiker's guide to the
% metaplectic trace formula for SL(2): basic theory" (November 2026).
% New in the style-guide pass (style guide, Section 2.7). It restates
% results of Sections 2-7 with the status given there and lists only open
% questions that Sections 1-7 already name. It is new text and is listed
% for refereeing in ../revision_log.md (Stage 8).
% =====================================================================

\section{Summary of results and open questions}\label{sec:summary}

This section collects the main results of the paper, with the status
that the text gives them, and the questions that the paper leaves open.
The status of each statement in detail is in the subsections ``Status of
the results'' at the end of Sections~\ref{sec:background}--\ref{sec:stable}.

\subsection*{Results}

\begin{enumerate}[label=\arabic*.]
\item \emph{The cover} (Section~\ref{sec:background}). The double cover
  of $G(\A)$ is built from Kubota's cocycle, with its topology fixed
  intrinsically, and it splits over $G(\Q)$ (Lemmas~\ref{lem:kubota}
  and~\ref{lem:global}). The cocycle and the existence of the topological
  cover are taken from~\cite{Ku67,Mo68}. The rest is proved in the text.
\item \emph{Rank $0$} (Section~\ref{sec:rank0}). For the cover of a torus
  the trace formula is Poisson summation, twisted by a genuine
  automorphic character (Theorem~\ref{thm:tf-rank0} and
  Lemma~\ref{lem:transfer}). For the split torus Weil's character
  $\chi_\psi$ gives a canonical choice (Lemma~\ref{lem:weil}, whose
  parts~(a) and~(b) are Weil's).
\item \emph{The geometric side} (Section~\ref{sec:geometric}). The coarse
  geometric expansion is proved directly from the height inequality, with
  truncation parameter $T\ge0$ (Lemma~\ref{lem:height-ineq} and
  Theorem~\ref{thm:coarse}). Each class contributes a polynomial in $T$ of
  degree at most one, written down in
  Propositions~\ref{prop:hyperbolic}, \ref{prop:elliptic}
  and~\ref{prop:central}. The fine expansion, Theorem~\ref{thm:fine}, is
  quoted from~\cite{Ar86,Li14a}. Its part~(a) is derived again in
  Proposition~\ref{prop:unipotent} when the ring of $S_k$-integers is
  principal.
\item \emph{The spectral side} (Section~\ref{sec:spectral}). The
  characters fixed by the Weyl group are the Weil characters
  $\chi_{\psi_b}$ (Lemma~\ref{lem:w-action}). The normalizing factor of the
  intertwining operator is $L^S(2\lambda,\chi^2)/L^S(2\lambda+1,\chi^2)$
  (Lemma~\ref{lem:intertwining}). The residual spectrum is the sum of the
  even Weil representations $r^+_{\psi_{-b}}$ (Theorem~\ref{thm:residual}).
  Parts~(a) and~(b) of that theorem, and the matching of parameters
  in~(c), are proved here. The rest of~(c) is taken
  from~\cite{Ge76,GPS80,Wa80}. The trace formula is
  Theorem~\ref{thm:tf}. The spectral decomposition, the Eisenstein series
  and the polynomial behavior of the spectral terms are taken
  from~\cite{La76,MW95,Ar80,Li13}.
\item \emph{The invariant form} (Section~\ref{sec:invariant}). The
  invariant distributions are built by Fourier inversion on the abelian
  cover of $A$ (Theorem~\ref{thm:IA}). The invariant trace formula is
  Theorem~\ref{thm:itf}, with Corollary~\ref{cor:itf-final}, and
  $M(w,0)=-1$ for every character fixed by the Weyl group
  (Proposition~\ref{prop:Mw0}). These are proved here, given the
  existence of normalizing factors, Theorem~\ref{thm:normalization}, which
  is taken from~\cite{Ar89,Li12a}.
\item \emph{The stable(r) form} (Section~\ref{sec:stable}). Conjugation
  by $GL(2)$ can exchange the two lifts of an element
  (Lemma~\ref{lem:eta}). Stable conjugacy is conjugation by $GL(2)$
  twisted by the sign $(\det h,\det(1+\gamma))_v$, and its classes are
  described by the function $\Delta_\psi$ (Lemma~\ref{lem:Delta}). The
  character of $r^+_\psi-r^-_\psi$ is computed on the regular set
  (Proposition~\ref{prop:weil-char}). The regular elliptic terms take the
  form $S_{\mathrm{ell}}(f)+S_{\mathrm{ell}}(f_-)$
  (Theorem~\ref{thm:elliptic-stable}). The unipotent coefficients are
  $a^G(S,u_\xi)=\frac12\bigl(c_S+\sum_{\kappa\ne1}\kappa_S(\xi)
  L^S(1,\kappa)\bigr)$ (Proposition~\ref{prop:unipotent}). The theta
  functions contribute to the spectral side in the same form as the
  elliptic terms (Proposition~\ref{prop:theta-packet}). The results are
  collected in Theorem~\ref{thm:stabler}. Its statements about the
  weighted terms use the property~\ref{N5} of the normalizing factors.
  Factors with this property exist at every place
  (Proposition~\normref{}(a),~(b)). Their explicit form at the dyadic and
  ramified places, Proposition~\normref{}(c), is taken from
  Szpruch~\cite{Sz09,Sz16}.
\end{enumerate}

\subsection*{Open questions}

The following questions are named as open in Sections~\ref{sec:intro}
and~\ref{sec:stable}. They are not answered here.

\begin{enumerate}[label=\arabic*.]
\item The stabilization of the weighted terms $I_A(a(t),\cdot)$ and
  $I_A(z,\cdot)$, which are equivariant under $GL(2)$ but not stable
  (Proposition~\ref{prop:IA-gl2} and Section~\ref{sec:stable-form}).
\item The transfer $f\mapsto f'$ to $SO(3)$ and the corresponding
  fundamental lemma~\cite{Li11,Li12b}, and with them the identification
  of $\Delta_\psi$ with Li's transfer factor and of $J^{\mathrm{st}}_G$
  with his stable orbital integral (Remark~\ref{rem:sec7-status}).
\item The description of the cuspidal spectrum by
  packets~\cite{Wa80,Wa91,GI18}. The character
  identity~\eqref{eq:packet-lift} used for the packets is quoted, not
  proved.
\item Whether the normalizing factors of~\cite{Li12a} have the
  property~\ref{N5}, and an independent derivation of Szpruch's formula
  for the local coefficients at the dyadic places
  (Remark~\ref{rem:N5}).
\item Whether $T\ge0$ suffices in Theorem~\ref{thm:Jchi-poly}, as it does
  on the geometric side (Remark~\ref{rem:sec5-status}).
\item The limit formulas of Section~\ref{sec:unipotent} at the singular
  elliptic points.
\item The $n$-fold cover for $n>2$, which the 1995 notes treated and
  this revision does not (Remarks~\ref{rem:general-n}
  and~\ref{rem:general-n-rank0}).
\end{enumerate}


% SPDX-License-Identifier: CC-BY-4.0
% =====================================================================
% Glossary for D. Joyner and Claude, "A hitch-hiker's guide to the
% metaplectic trace formula for SL(2)" (November 2026).
% Added in response to the second referee report (ChatGPT). Needs longtable.
% =====================================================================

\section*{Glossary}\label{sec:glossary}
\addcontentsline{toc}{section}{Glossary}

\subsection*{Notation}

{\small
\begin{longtable}{@{}p{0.21\textwidth}p{0.55\textwidth}p{0.18\textwidth}@{}}
Symbol & Meaning & Defined in\\\hline
\endhead
\multicolumn{3}{@{}l}{\emph{Groups}}\\
$k$, $d$, $\A_k$, $\A$ & number field, its degree $[k:\Q]$, its adeles,
  the adeles of $\Q$ & \S\ref{sec:background}\\
$G$ & $\Res_{k/\Q}SL(2)$; $G(\Q)=SL(2,k)$ & \S\ref{sec:background}\\
$B$, $N$, $A$ & upper triangular Borel subgroup, its unipotent radical,
  the diagonal torus & \S\ref{sec:local}\\
$a(t)$, $n(u)$, $w$ & $\diag(t,t^{-1})$, the upper unipotent matrix with
  entry $u$, the Weyl element $\bigl(\begin{smallmatrix}0&-1\\1&0
  \end{smallmatrix}\bigr)$ & \S\ref{sec:local}, \S\ref{sec:geometric}\\
$T_\theta$, $T_\mu$ & anisotropic torus attached to $k(\sqrt\theta)$;
  centralizer of an elliptic $\mu$ & \S\ref{sec:local}, Lemma~\ref{lem:classes}\\
$K_v$, $K_m$, $K_0$, $K_f$ & maximal compact subgroup at $v$; principal
  congruence subgroup of level $\varpi^m$; $\prod_vK_v$; an open compact
  subgroup of $G(\A_f)$ & \S\ref{sec:local}, \S\ref{sec:global}\\
$A(\A)^1$, $A_B(\R)^0$ & kernel of $H$ on $A(\A)$; the split component
  $\{a(u)\mid u>0\}$ & \S\ref{sec:torus-cover}, \S\ref{sec:coarse}\\[4pt]
\multicolumn{3}{@{}l}{\emph{The cover}}\\
$\mu_2$ & $\{\pm1\}$, the kernel of the covering & \S\ref{sec:background}\\
$(a,b)_v$ & quadratic Hilbert symbol of $k_v$ & \S\ref{sec:background}\\
$\beta_v$ & Kubota's cocycle on $SL(2,k_v)$; Borel, not continuous &
  \eqref{eq:beta}, Lemma~\ref{lem:kubota}\\
$\kappa_v$, $\kappa_\A$ & Kubota's correction on $K_v$ ($v$ finite,
  non-dyadic); its product over $v$, defined on $G(\Q)$ &
  \eqref{eq:kappa}, \eqref{eq:kappa-A}\\
$\sigma_v$, $\sigma_\A$ & $\beta_v$ times the coboundary of $\kappa_v$;
  the adelic cocycle $\prod_v\sigma_v$ & \eqref{eq:sigma-v},
  \eqref{eq:sigma-A}\\
$\overline{SL(2,k_v)}$, $\overline{G(\A)}$ & the local and the adelic
  double cover; $\overline H=p^{-1}(H)$ for a subgroup $H$ &
  Lemma~\ref{lem:kubota}, \S\ref{sec:global}\\
$(g,\zeta)$ & Borel coordinates on a cover; $g\mapsto(g,1)$ is not
  continuous & \eqref{eq:law}, Example~\ref{ex:discontinuous}\\
$K_v^*$, $K_f^*$ & the image of $K_v$ under $k\mapsto(k,\kappa_v(k))$, a
  compact open subgroup of the cover; the image of $K_f$ &
  Lemma~\ref{lem:local}(b), Lemma~\ref{lem:global}(c)\\
$\overline{G(\Q_S)}$, $\overline{A(\Q_S)}$ & the covers of
  $\prod_{v\in S_k}SL(2,k_v)$ and of its diagonal torus &
  \S\ref{sec:local-dist}, \S\ref{sec:invariant}\\[4pt]
\multicolumn{3}{@{}l}{\emph{Heights, measures, functions}}\\
$H=H_B$, $H_v$, $H_S$ & the height $\log|t|$ of the $A$-component in
  $x=nak$ & \eqref{eq:H}, Lemma~\ref{lem:NAK-global}\\
$\delta_B$ & modulus character, $e^{2H}$ & \S\ref{sec:measures-local}\\
$\mathfrak a$, $\alpha$, $\alpha^\vee$ & $\Hom(X(A)_\Q,\R)\cong\R$; the
  simple root $2/d$; the coroot $d$ & \S\ref{sec:chambers},
  Remark~\ref{rem:alpha}\\
$\tau_B=\hat\tau_B$ & characteristic function of $(0,\infty)$ &
  \eqref{eq:tau}\\
$C_c^\infty(\cdot)_-$ & smooth, compactly supported, anti-genuine
  functions on a cover & \S\ref{sec:functions}\\
$e_v$, $e^S$ & unit of the spherical Hecke algebra at $v$;
  $\bigotimes_{v\notin S}e_v$ & \eqref{eq:unit}, \S\ref{sec:local-dist}\\
$\mathcal H(\overline{G(\A)})$, $\mathcal H(\overline{G(\Q_S)})$ & Hecke
  algebras of test functions; $f=f_S\otimes e^S$ & \S\ref{sec:global},
  \S\ref{sec:local-dist}\\
$f^y$ & $x\mapsto f(yxy^{-1})$ & \S\ref{sec:invariant}\\
$c_k$ & $\mathrm{vol}(A(\Q)\backslash A(\A)^1)$ & \S\ref{sec:coarse}\\[4pt]
\multicolumn{3}{@{}l}{\emph{Rank $0$}}\\
$\psi$, $\psi_b$ & a nontrivial character of $\A_k/k$; $x\mapsto\psi(bx)$
  & \S\ref{sec:torus-cover}\\
$\gamma_v(a)$ & Weil index of $x\mapsto\psi_v(ax^2)$, an eighth root of
  unity & \eqref{eq:weil-index}\\
$\chi_\psi$, $\chi_{\psi_b}$ & Weil's genuine character of
  $\overline{A(\A)}$, $(a(t),\zeta)\mapsto\zeta\prod_v\gamma_v(t_v)/
  \gamma_v(1)$ & Lemma~\ref{lem:weil}\\
$X_T$, $\Pi_{\mathrm{aut}}(X_T)$ & $T(\Q)\backslash\overline{T(\A)^1}$;
  its genuine characters & \S\ref{sec:torus-cover}\\
$I^A_{\mathcal O}$, $I^A_{\bar\chi}$ & the two kinds of terms in the
  trace formula for the cover of $A$ & \eqref{eq:IA}\\[4pt]
\multicolumn{3}{@{}l}{\emph{Geometric side}}\\
$\mathcal C$, $\mathcal O$, $\mathcal O_z$ & the set of classes in
  $G(\Q)$; a class; the central class of $z=\pm1$ &
  Lemma~\ref{lem:classes}\\
$K_P$, $K_{P,\mathcal O}$, $J_{P,\mathcal O}$ & kernels attached to
  $P\in\{B,G\}$ & \eqref{eq:KP}, \eqref{eq:JP}\\
$k^T$, $k^T_{\mathcal O}$, $j^T_{\mathcal O}$ & truncated kernels &
  \eqref{eq:kT}\\
$J^T_{\mathcal O}(f)$, $J_{\mathcal O}(f)$ & integral of
  $k^T_{\mathcal O}$, a polynomial in $T$; its value at $T=0$ &
  Theorem~\ref{thm:coarse}\\
$v_A$, $v_{A,v}$ & the weight $-H(x)-H(wx)$ & \eqref{eq:vA},
  Lemma~\ref{lem:weight-local}\\
$\Phi_f$, $Z(f,s)$ & the function $y\mapsto\int_Kf(k^{-1}n(y)k)\,dk$ and
  its zeta integral (not to be confused with $\Phi_A$) & \eqref{eq:Phi},
  \eqref{eq:Zf}\\
$J_G(\gamma,f_S)$ & orbital integral on $\overline{G(\Q_S)}$ &
  \eqref{eq:JG}\\
$J_A(\gamma,f_S)$ & weighted orbital integral, weight $v_A$; not
  invariant & \eqref{eq:JA}\\
$J_A(z,f_S)$ & its regularized value at $z=\pm1$ &
  Lemma~\ref{lem:JA-central}\\
$u_\xi$, $a^M(S,\gamma)$ & the unipotent element $n(\xi)$; the global
  coefficients of the fine expansion & Theorem~\ref{thm:fine}\\[4pt]
\multicolumn{3}{@{}l}{\emph{Spectral side}}\\
$\mathcal X$, $\mathcal X_G$, $\mathcal X_A$ & classes of cuspidal data;
  those of $(G,\pi)$; those of $(A,\mu)$ & \S\ref{sec:decomp}\\
$\mu$, $\mu^w$, $\mu_\lambda$ & a genuine automorphic character of the
  cover of $A$; its Weyl conjugate; $\mu\,e^{\lambda H}$ &
  \S\ref{sec:decomp}\\
$\chi_b$ & the ramified class of $(A,\chi_{\psi_b})$ &
  Lemma~\ref{lem:w-action}\\
$\mathcal H_B(\mu)$, $I(\mu_\lambda)$ & the space of the principal series
  and the induced representation & \eqref{eq:HBmu}\\
$E(x,\varphi,\lambda)$ & Eisenstein series & \eqref{eq:Eis}\\
$M(w,\lambda)$, $M_v$ & global and local intertwining operators &
  \eqref{eq:Mw}, Lemma~\ref{lem:intertwining}\\
$m^S(\mu,\lambda)$ & $L^S(2\lambda,\chi^2)/L^S(2\lambda+1,\chi^2)$ &
  \eqref{eq:mS}\\
$\omega_{\psi'}$, $r^\pm_{\psi'}$ & Weil representation and its even and
  odd parts & \eqref{eq:weil-rep}\\
$V_b$ & residual representation of the class $\chi_b$, isomorphic to
  $r^+_{\psi_{-b}}$ & Theorem~\ref{thm:residual}\\
$\Lambda^T$ & truncation operator & \eqref{eq:trunc}\\
$K_{P,\chi}$, $k^T_\chi$, $J^T_\chi$, $J_\chi$ & spectral kernels, their
  truncation and its integral & \S\ref{sec:Jchi}\\
$M_B^T(\mu,\lambda)$ & $2T-M(w,\lambda)^{-1}M'(w,\lambda)$ &
  \eqref{eq:MT}\\[4pt]
\multicolumn{3}{@{}l}{\emph{Invariant form}}\\
$u'_B(k,y)$, $f_{B,y}$ & the weight $-H(ky)$; the function on the cover
  of $A$ that measures $J(f^y)-J(f)$ & \eqref{eq:uprime}, \eqref{eq:fBy}\\
$V(\pi)$, $I(\pi_\lambda)$ & local principal series of
  $\overline{G(\Q_S)}$ & \S\ref{sec:invariant}\\
$r_v(\pi_v,\lambda)$, $R_v$, $R_S$ & local normalizing factor; normalized
  operator $r_v^{-1}M_v$; $\bigotimes_{v\in S_k}R_v$ &
  Theorem~\ref{thm:normalization}\\
$S_{\mathrm{bad}}(\mu)$ & the places where $r_v$ is not explicit &
  \S\ref{sec:normalized}\\
$m(\mu,\lambda)$ & global normalizing factor, $M(w,\lambda)=m\cdot R_S$ &
  \eqref{eq:global-m}\\
$\varphi_A(f,\pi)$ & $-\operatorname{tr}(R_S^{-1}R_S'\,I(\pi,f))$, a
  function of the \emph{character} $\pi$ of $\overline{A(\Q_S)}$ &
  \eqref{eq:phiA}\\
$\Phi_A(\tilde a,f)$ & the inverse Fourier transform of
  $\varphi_A(f,\cdot)$, a function of the \emph{element} $\tilde a$ of
  $\overline{A(\Q_S)}$ & \eqref{eq:PhiA}\\
$I_A(\gamma,f)$, $I_A(z,f)$ & $J_A-\Phi_A$, the invariant distributions &
  \eqref{eq:IAgamma}\\
$I_{\mathcal O}(f)$, $I_\chi(f)$ & the terms of the invariant trace
  formula & \eqref{eq:Iochi}\\[3pt]
\multicolumn{3}{@{}l}{\emph{Section~\ref{sec:stable}}}\\
$G_1$, $\delta_a$, $\mathrm{Ad}(h)$ & $GL(2)$; $\diag(a,1)$; the
  automorphism of the cover lifting conjugation by $h\in G_1$ &
  Lemma~\ref{lem:gl2}, Lemma~\ref{lem:Ad}\\
$f^h$, $f^{[h]}$, $f_z$, $f_-$ & $f\circ\mathrm{Ad}(h)$; $f(h\star\cdot)$;
  $f(z\,\cdot)$ for central $z$; $f_{-1}$ & \S\ref{sec:stable},
  \S\ref{sec:stable-conj}\\
$D^+(\gamma)$, $\eta_\gamma$ & $\det(1+\gamma)=2+\operatorname{tr}\gamma$;
  the character $a\mapsto(a,D^+(\gamma))_v$ & \eqref{eq:eta}\\
$h\star x$ & $\eta_x(\det h)\,\mathrm{Ad}(h)x$, the action whose orbits are
  the stable classes & \eqref{eq:star}\\
$\Delta_\psi$ & genuine function with values in $\mu_4$, constant on
  stable classes; $|D^+|^{1/2}$ times the character of $r^+_\psi-r^-_\psi$ &
  Lemma~\ref{lem:Delta}, Prop.~\ref{prop:weil-char}\\
$\kappa_E$, $N E^\times$ & the quadratic character of $E/F$ and its
  kernel, the norms & \S\ref{sec:stable}\\
$J^{\mathrm{st}}_G$, $J^\kappa_G$ & stable and $\kappa$-orbital
  integrals & \eqref{eq:Jst}, \eqref{eq:Jst-global}\\
$S_{\mathrm{ell}}(f)$ & the stable elliptic regular distribution &
  Theorem~\ref{thm:elliptic-stable}\\
$H_1$, $\kappa_a$ & the height on $GL(2)$; the character
  $a(t)\mapsto(a,t)$ of the cover of $A$ & \S\ref{sec:hyperbolic-gl2}\\
$\mathfrak X_S$, $L^S(s,\kappa)$, $\rho_S$, $c_S$ & quadratic Hecke
  characters unramified outside $S$; their partial $L$-functions; the
  residue and constant term of $\zeta_k^S$ at $1$ & \S\ref{sec:unipotent}\\
$J_G(zu_\xi,f_S)$, $J^\kappa_{\mathrm{unip}}(z,f_S)$ & normalized unipotent
  orbital integral; its $\kappa$-sum over $\xi$ & \eqref{eq:JG-unip},
  Prop.~\ref{prop:unipotent}\\
$U_h$, $\Theta_b$, $\Theta^{\mathrm{st}}_b$ & the action of $G_1(k)$ on
  automorphic forms; the span of the theta functions of $\psi_{-b}$; its
  stable character & \S\ref{sec:spectral-gl2}\\
\end{longtable}
}

\subsection*{Terms}

\begin{description}
\item[Anti-genuine function.] A function $f$ on a cover with
  $f(x\zeta)=\zeta^{-1}f(x)$ for $\zeta\in\mu_2$. All test functions are
  anti-genuine. For the double cover this is the same condition as
  \emph{genuine}, $f(x\zeta)=\zeta f(x)$.
\item[Borel coordinates.] The coordinates $(g,\zeta)$ in which a cover is
  described by a cocycle. They give the group law and the Borel
  structure, but not the topology (Lemma~\ref{lem:kubota}).
\item[Class (geometric).] A set of elements of $G(\Q)$ whose semisimple
  parts are conjugate: central, hyperbolic or elliptic
  (Lemma~\ref{lem:classes}). A class is \emph{unramified} if it consists
  of semisimple elements.
\item[Cuspidal datum.] A pair $(G,\pi)$ with $\pi$ cuspidal, or a pair
  $(A,\mu)$ with $\mu$ a genuine automorphic character of the cover of
  $A$; $(A,\mu)$ and $(A,\mu^w)$ are equivalent. A class $\chi$ of data
  $(A,\mu)$ is \emph{ramified} if $\mu^w=\mu$, and then
  $\mu=\chi_{\psi_b}$.
\item[Dyadic place.] A place of $k$ above $2$. The explicit splitting
  $\kappa_v$ and the explicit normalizing factors are not available
  there.
\item[Genuine representation.] A representation of a cover in which
  $\mu_2$ acts by the identity character.
\item[$(G,S)$-equivalence.] Arthur's equivalence relation on $G(\Q)$,
  for a finite set $S$ of places: equal semisimple parts up to
  conjugacy, and unipotent parts conjugate in the centralizer over
  $\Q_S$. It indexes the terms of the fine expansion
  (Theorem~\ref{thm:fine}).
\item[Endoscopic group.] For the double cover of $SL(2)$ the elliptic
  endoscopic groups are two copies of $SO(3)$, attached to the central
  elements $\pm1$ of the dual group $SL(2,\C)$~\cite{Li11}. They are
  discussed but not used.
\item[Invariant distribution.] A linear form $I$ on test functions with
  $I(f^y)=I(f)$ for all $y$ in the group.
\item[Normalizing factor.] A scalar function $r_v(\pi_v,\lambda)$ such
  that $r_v^{-1}M_v$ has the properties \ref{N1}--\ref{N4} of
  Theorem~\ref{thm:normalization}. Its existence, an explicit formula for
  it, and its role in the global factor $m$ are three different matters
  (Section~\ref{sec:normalized}).
\item[Residual spectrum.] The part of the discrete spectrum that is not
  cuspidal. For this group it is spanned by the residues of Eisenstein
  series at $\lambda=\frac12$ and consists of even Weil representations.
\item[Stable conjugacy.] On $SL(2,F)$: conjugacy under $GL(2,F)$. On the
  cover: the relation whose classes are the orbits of $h\star x$; a stable
  class lies over a stable class of $SL(2,F)$ and contains one of the two
  lifts of each element (Section~\ref{sec:stable-conj}). Conjugation by
  $GL(2,F)$ without the twist is coarser.
\item[Stable distribution.] An invariant distribution given by a class
  function that is constant on stable classes, or a stable orbital
  integral $J^{\mathrm{st}}_G$; equivalently, one that is invariant under
  $f\mapsto f^{[h]}$.
\item[Stable(r) form.] The trace formula with its terms organized by
  stable classes, before any transfer to $SO(3)$
  (Theorem~\ref{thm:stabler}).
\item[Truncation.] Subtracting from a kernel, or from an automorphic
  form, its constant term in the cusp $\{H(\delta x)>T\}$.
\item[Weighted orbital integral.] An orbital integral with the weight
  $v_A$ inserted, \eqref{eq:JA}. It arises from truncation and is not
  invariant.
\item[Weil index.] The eighth root of unity $\gamma_v(a)$ attached to the
  character of second degree $x\mapsto\psi_v(ax^2)$,
  \eqref{eq:weil-index}.
\end{description}


\section*{Acknowledgments}
\addcontentsline{toc}{section}{Acknowledgments}

The 1995 notes and the 1997 book manuscript from which they came are by
the first author. The revision of November 2026 was carried out by the
two authors together. The first author directed the work and reviewed
it. Claude (Anthropic), an AI model, took part in three roles. As a
referee, it wrote the first referee report on the 1995 notes. As a
co-author and reviser, it drafted the revised sections, the
introduction, the glossary and the summary. It also wrote the
verification scripts in \texttt{checks/}.

ChatGPT (OpenAI), an AI model, acted as an independent referee. It
refereed the November 2026 draft of
Sections~\ref{sec:background}--\ref{sec:invariant} in one report and
Section~\ref{sec:stable} in a second report. Several definitions, proofs
and examples of the present draft were added or corrected in response
to these reports.

Review by AI systems is not an independent certification of the
results. The errors that remain are the authors'. Further details on the
AI usage, including the referee reports, the responses to them and a
list of every change, can be found in the revision log
(\texttt{revision\_log.md}) that accompanies the source.

\begin{thebibliography}{GGW18}
\bibitem[Ad98]{Ad98} J. Adams, ``Lifting of characters on orthogonal and
metaplectic groups'', Duke Math. J. \textbf{92} (1998), 129--178.
\bibitem[Ar78]{Ar78} J. Arthur, ``A trace formula for reductive groups, I: terms
associated to classes in $G(\mathbb{Q})$'', Duke Math. J. \textbf{45} (1978), 911--952.
\bibitem[Ar80]{Ar80} J. Arthur, ``A trace formula for reductive groups, II:
applications of a truncation operator'', Compositio Math. \textbf{40}
(1980), 87--121.
\bibitem[Ar81]{Ar81} J. Arthur, ``The trace formula in invariant form'', Ann.
of Math. (2) \textbf{114} (1981), 1--74.
\bibitem[Ar82]{Ar82} J. Arthur, ``On a family of distributions obtained from
Eisenstein series II: explicit formulas'', Amer. J. Math. \textbf{104}
(1982), 1289--1336.
\bibitem[Ar86]{Ar86} J. Arthur, ``On a family of distributions obtained from
orbits'', Canad. J. Math. \textbf{38} (1986), 179--214.
\bibitem[Ar88a]{Ar88a} J. Arthur, ``The local behaviour of weighted orbital
integrals'', Duke Math. J. \textbf{56} (1988), 223--293.
\bibitem[Ar88b]{Ar88b} J. Arthur, ``The invariant trace formula, I. Local theory'',
J. Amer. Math. Soc. \textbf{1} (1988), 323--383.
\bibitem[Ar88c]{Ar88c} J. Arthur, ``The invariant trace formula, II. Global
theory'', J. Amer. Math. Soc. \textbf{1} (1988), 501--554.
\bibitem[Ar89]{Ar89} J. Arthur, ``Intertwining operators and residues, I.
Weighted characters'', J. Funct. Anal. \textbf{84} (1989), 19--84.
\bibitem[Ar05]{Ar05} J. Arthur, ``An introduction to the trace formula'', in
\emph{Harmonic analysis, the trace formula, and Shimura varieties}, Clay
Math. Proc. \textbf{4}, Amer. Math. Soc., 2005, 1--263.
\bibitem[BD01]{BD01} J.-L. Brylinski and P. Deligne, ``Central extensions of
reductive groups by $K_2$'', Publ. Math. IHES \textbf{94} (2001), 5--85.
\bibitem[CF67]{CF67} J. W. S. Cassels and A. Fr\"ohlich (eds.), \emph{Algebraic
number theory}, Academic Press, 1967.
\bibitem[FK86]{FK86} Y. Z. Flicker and D. A. Kazhdan, ``Metaplectic
correspondence'', Publ. Math. IHES \textbf{64} (1986), 53--110.
\bibitem[Fl79]{Fl79} D. Flath, ``Decomposition of representations into tensor
products'', in \emph{Automorphic forms, representations and
$L$-functions}, Proc. Sympos. Pure Math. \textbf{33}, Part~1, Amer. Math.
Soc., 1979, 179--183.
\bibitem[Fl80]{Fl80} Y. Z. Flicker, ``Automorphic forms on covering groups of
$GL(2)$'', Invent. Math. \textbf{57} (1980), 119--182.
\bibitem[Ga14]{Ga14} F. Gao, ``The Gindikin--Karpelevich formula and constant
terms of Eisenstein series for Brylinski--Deligne extensions'', preprint,
arXiv:1410.7624.
\bibitem[Ge76]{Ge76} S. Gelbart, \emph{Weil's representation and the spectrum of
the metaplectic group}, Lecture Notes in Math. \textbf{530}, Springer, 1976.
\bibitem[GGW18]{GGW18} W. T. Gan, F. Gao and M. H. Weissman, ``L-groups and the
Langlands program for covering groups: a historical introduction'',
Ast\'erisque \textbf{398} (2018), 1--31.
\bibitem[GI18]{GI18} W. T. Gan and A. Ichino, ``The Shimura--Waldspurger
correspondence for $\mathrm{Mp}_{2n}$'', Ann. of Math. (2) \textbf{188}
(2018), 965--1016.
\bibitem[GJ79]{GJ79} S. Gelbart and H. Jacquet, ``Forms of $GL(2)$ from the
analytic point of view'', in \emph{Automorphic forms, representations and
$L$-functions}, Proc. Sympos. Pure Math. \textbf{33}, Part~1, Amer. Math.
Soc., 1979, 213--251.
\bibitem[Go62]{Go62} R. Godement, ``Domaines fondamentaux des groupes
arithm\'etiques'', S\'eminaire Bourbaki \textbf{8} (1962/63), Exp.~257.
\bibitem[GPS80]{GPS80} S. Gelbart and I. Piatetski-Shapiro, ``Distinguished
representations and modular forms of half-integral weight'', Invent.
Math. \textbf{59} (1980), 145--188.
\bibitem[GPS83]{GPS83} S. Gelbart and I. Piatetski-Shapiro, ``Some remarks on
metaplectic cusp forms and the correspondences of Shimura and
Waldspurger'', Israel J. Math. \textbf{44} (1983), 97--126.
\bibitem[GS12]{GS12} W. T. Gan and G. Savin, ``Representations of
metaplectic groups I: epsilon dichotomy and local Langlands
correspondence'', Compos. Math. \textbf{148} (2012), 1655--1694.
\bibitem[GS16]{GS16} D. Goldberg and D. Szpruch, ``Plancherel measures for
coverings of $p$-adic $SL_2(F)$'', Int. J. Number Theory \textbf{12} (2016),
1907--1936.
\bibitem[He83]{He83} D. A. Hejhal, \emph{The Selberg trace formula for
$PSL(2,\mathbb{R})$, Vol.~2}, Lecture Notes in Math. \textbf{1001},
Springer, 1983.
\bibitem[HI13]{HI13} K. Hiraga and T. Ikeda, ``On the Kohnen plus space for
Hilbert modular forms of half-integral weight I'', Compos. Math.
\textbf{149} (2013), 1963--2010.
\bibitem[Ho10]{Ho10} T. K. Howard, ``Lifting of characters on $p$-adic
orthogonal and metaplectic groups'', Compos. Math. \textbf{146} (2010),
795--810.
\bibitem[JL70]{JL70} H. Jacquet and R. P. Langlands, \emph{Automorphic forms on
$GL(2)$}, Lecture Notes in Math. \textbf{114}, Springer, 1970.
\bibitem[Jo88]{Jo88} D. Joyner, ``On the metaplectic analog of Kazhdan's
`endoscopic' lifting'', Israel J. Math. \textbf{66} (1988), 113--154.
\bibitem[Kn02]{Kn02} A. W. Knapp, \emph{Lie groups beyond an introduction}, 2nd ed.,
Birkh\"auser, 2002.
\bibitem[KP84]{KP84} D. Kazhdan and S. J. Patterson, ``Metaplectic forms'', Publ.
Math. IHES \textbf{59} (1984), 35--142.
\bibitem[KP86]{KP86} D. A. Kazhdan and S. J. Patterson, ``Towards a
generalized Shimura correspondence'', Adv. Math. \textbf{60} (1986),
161--234.
\bibitem[Ku67]{Ku67} T. Kubota, ``Topological covering of $SL(2)$ over a local
field'', J. Math. Soc. Japan \textbf{19} (1967), 114--121.
\bibitem[Ku69]{Ku69} T. Kubota, \emph{On automorphic functions and the reciprocity
law in a number field}, Lectures in Math. \textbf{2}, Kinokuniya, Tokyo, 1969.
\bibitem[La76]{La76} R. P. Langlands, \emph{On the functional equations satisfied
by Eisenstein series}, Lecture Notes in Math. \textbf{544}, Springer, 1976.
\bibitem[Li11]{Li11} W.-W. Li, ``Transfert d'int\'egrales orbitales pour le
groupe m\'etaplectique'', Compos. Math. \textbf{147} (2011), 524--590.
\bibitem[Li12a]{Li12a} W.-W. Li, ``La formule des traces pour les rev\^etements
de groupes r\'eductifs connexes. II. Analyse harmonique locale'', Ann. Sci.
\'Ec. Norm. Sup\'er. (4) \textbf{45} (2012), 787--859.
\bibitem[Li12b]{Li12b} W.-W. Li, ``Le lemme fondamental pond\'er\'e pour le
groupe m\'etaplectique'', Canad. J. Math. \textbf{64} (2012), 497--543.
\bibitem[Li13]{Li13} W.-W. Li, ``La formule des traces pour les
rev\^etements de groupes r\'eductifs connexes. III. Le d\'eveloppement
spectral fin'', Math. Ann. \textbf{356} (2013), 1029--1064.
\bibitem[Li14a]{Li14a} W.-W. Li, ``La formule des traces pour les rev\^etements
de groupes r\'eductifs connexes. I. Le d\'eveloppement g\'eom\'etrique fin'',
J. reine angew. Math. \textbf{686} (2014), 37--109.
\bibitem[Li14b]{Li14b} W.-W. Li, ``La formule des traces pour les rev\^etements
de groupes r\'eductifs connexes. IV. Distributions invariantes'', Ann. Inst.
Fourier \textbf{64} (2014), 2379--2448.
\bibitem[Li15]{Li15} W.-W. Li, ``La formule des traces stable pour le groupe
m\'etaplectique: les termes elliptiques'', Invent. Math. \textbf{202}
(2015), 743--838.
\bibitem[Li19]{Li19} W.-W. Li, ``Spectral transfer for metaplectic groups.
I. Local character relations'', J. Inst. Math. Jussieu \textbf{18} (2019),
25--123.
\bibitem[Li21]{Li21} W.-W. Li, ``Stabilization of the trace formula for
metaplectic groups'', preprint, arXiv:2109.06581.
\bibitem[Li24]{Li24} W.-W. Li, ``Arthur packets for metaplectic groups'',
preprint, arXiv:2410.13606.
\bibitem[LL79]{LL79} J.-P. Labesse and R. P. Langlands, ``L-indistinguishability
for $SL(2)$'', Canad. J. Math. \textbf{31} (1979), 726--785.
\bibitem[Lu20]{Lu20} C. Luo, ``Endoscopic character identities for
metaplectic groups'', J. Reine Angew. Math. \textbf{768} (2020), 1--37.
\bibitem[Ma57]{Ma57} G. W. Mackey, ``Les ensembles bor\'eliens et les
extensions des groupes'', J. Math. Pures Appl. (9) \textbf{36} (1957),
171--178.
\bibitem[Ma69]{Ma69} H. Matsumoto, ``Sur les sous-groupes arithm\'etiques des
groupes semi-simples d\'eploy\'es'', Ann. Sci. \'Ecole Norm. Sup. (4)
\textbf{2} (1969), 1--62.
\bibitem[Ma99]{Ma99} K. Maktouf, ``Le caract\`ere de la repr\'esentation
m\'etaplectique et la formule du caract\`ere pour certaines
repr\'esentations d'un groupe de Lie presque alg\'ebrique sur un corps
$p$-adique'', J. Funct. Anal. \textbf{164} (1999), 249--339.
\bibitem[Mc12]{Mc12} P. J. McNamara, ``Principal series representations of
metaplectic groups over local fields'', in \emph{Multiple Dirichlet series,
L-functions and automorphic forms}, Progr. Math. \textbf{300},
Birkh\"auser, 2012, 299--327.
\bibitem[Me01]{Me01} P. Mezo, ``Comparisons of general linear groups and their
metaplectic coverings II'', Represent. Theory \textbf{5} (2001), 524--580.
\bibitem[Me02]{Me02} P. Mezo, ``Comparisons of general linear groups and their
metaplectic coverings I'', Canad. J. Math. \textbf{54} (2002), 92--137.
\bibitem[Mo68]{Mo68} C. C. Moore, ``Group extensions of $p$-adic and adelic linear
groups'', Publ. Math. IHES \textbf{35} (1968), 157--222.
\bibitem[Mo88]{Mo88} C. Moen, ``Irreducibility of unitary principal series
for covering groups of $SL(2,k)$'', Pacific J. Math. \textbf{135} (1988),
89--110.
\bibitem[MW95]{MW95} C. M\oe glin and J.-L. Waldspurger, \emph{Spectral
decomposition and Eisenstein series}, Cambridge Tracts in Math.
\textbf{113}, Cambridge Univ. Press, 1995.
\bibitem[Ni75]{Ni75} S. Niwa, ``Modular forms of half integral weight and the
integral of certain theta-functions'', Nagoya Math. J. \textbf{56} (1975),
147--161.
\bibitem[PR94]{PR94} V. Platonov and A. Rapinchuk, \emph{Algebraic groups and
number theory}, Academic Press, 1994.
\bibitem[Re99]{Re99} D. Renard, ``Endoscopy for $Mp(2n,\mathbb{R})$'', Amer. J.
Math. \textbf{121} (1999), 1215--1244.
\bibitem[Ro83]{Ro83} A. Robert, \emph{Introduction to the representation theory of
compact and locally compact groups}, London Math. Soc. Lecture Note Ser.
\textbf{80}, Cambridge, 1983.
\bibitem[Ro90]{Ro90} J. Rogawski, \emph{Automorphic representations of unitary groups
in three variables}, Ann. of Math. Studies \textbf{123}, Princeton, 1990.
\bibitem[Sa67]{Sa67} P. J. Sally, Jr., \emph{Analytic continuation of the
irreducible unitary representations of the universal covering group of
$SL(2,\mathbb{R})$}, Mem. Amer. Math. Soc. \textbf{69} (1967).
\bibitem[Sa88]{Sa88} G. Savin, ``Local Shimura correspondence'', Math. Ann.
\textbf{280} (1988), 185--190.
\bibitem[Sc98]{Sc98} J. P. Schultz, \emph{Lifting of characters of
$\widetilde{SL}_2(F)$ and $SO_{1,2}(F)$ for $F$ a nonarchimedean local
field}, Ph.D. thesis, University of Maryland, 1998.
\bibitem[Se56]{Se56} A. Selberg, ``Harmonic analysis and discontinuous groups
in weakly symmetric Riemannian spaces with applications to Dirichlet
series'', J. Indian Math. Soc. \textbf{20} (1956), 47--87.
\bibitem[Se73]{Se73} J.-P. Serre, \emph{A course in arithmetic}, Graduate Texts in
Math. \textbf{7}, Springer, 1973.
\bibitem[Sh73]{Sh73} G. Shimura, ``On modular forms of half integral weight'',
Ann. of Math. (2) \textbf{97} (1973), 440--481.
\bibitem[Sh75]{Sh75} T. Shintani, ``On the construction of holomorphic cusp
forms of half integral weight'', Nagoya Math. J. \textbf{58} (1975),
83--126.
\bibitem[Sz09]{Sz09} D. Szpruch, ``Computation of the local coefficients for
principal series representations of the metaplectic double cover of
$SL_2(F)$'', J. Number Theory \textbf{129} (2009), 2180--2213.
\bibitem[Sz10]{Sz10} D. Szpruch, \emph{The Langlands--Shahidi method for
the metaplectic group and applications}, thesis, Tel Aviv University,
2009, arXiv:1004.3516.
\bibitem[Sz16]{Sz16} D. Szpruch, ``On Shahidi local coefficients matrix'',
preprint, arXiv:1609.08122.
\bibitem[Ta67]{Ta67} J. Tate, ``Fourier analysis in number fields and Hecke's
zeta-functions'', in~\cite{CF67}, 305--347.
\bibitem[Th08]{Th08} T. Thomas, ``The character of the Weil
representation'', J. London Math. Soc. (2) \textbf{77} (2008), 221--239.
\bibitem[Th13]{Th13} T. Thomas, ``Weil representation and transfer
factor'', Algebra Number Theory \textbf{7} (2013), 1535--1570.
\bibitem[Wa80]{Wa80} J.-L. Waldspurger, ``Correspondance de Shimura'', J. Math.
Pures Appl. (9) \textbf{59} (1980), 1--132.
\bibitem[Wa81]{Wa81} J.-L. Waldspurger, ``Sur les coefficients de Fourier
des formes modulaires de poids demi-entier'', J. Math. Pures Appl. (9)
\textbf{60} (1981), 375--484.
\bibitem[Wa91]{Wa91} J.-L. Waldspurger, ``Correspondances de Shimura et
quaternions'', Forum Math. \textbf{3} (1991), 219--307.
\bibitem[We64]{We64} A. Weil, ``Sur certains groupes d'op\'erateurs
unitaires'', Acta Math. \textbf{111} (1964), 143--211.
\bibitem[We74]{We74} A. Weil, \emph{Basic number theory}, Springer, 1974.
\bibitem[We82]{We82} A. Weil, \emph{Adeles and algebraic groups}, Birkh\"auser, 1982.
\bibitem[We18]{We18} M. H. Weissman, ``L-groups and parameters for covering
groups'', Ast\'erisque \textbf{398} (2018), 33--186.
\end{thebibliography}
\end{document}
